\documentclass[11pt,a4paper,oneside]{book}

%--------------------------------------------------------------------
%  Theory of Corporate Finance -- Lecture Notes
%  V. Filipe Martins-da-Rocha (FGV EESP)
%
%  Compilation sequence:
%      pdflatex Corp-Fin-Notes
%      bibtex   Corp-Fin-Notes
%      pdflatex Corp-Fin-Notes
%      pdflatex Corp-Fin-Notes
%--------------------------------------------------------------------

\usepackage[margin=2.8cm]{geometry}
\usepackage{setspace}
\setstretch{1.15}

\usepackage[english]{babel}
\usepackage[T1]{fontenc}
\usepackage[utf8]{inputenc}

\usepackage{amsmath,amssymb,amsthm,bm}
\usepackage{mathtools}

\usepackage[authoryear,round]{natbib}

\usepackage{xcolor}
\usepackage[colorlinks=true,linkcolor=blue!55!black,citecolor=blue!55!black,urlcolor=blue!55!black]{hyperref}

%-------------------- macros (aligned with slides) ------------------
\def\R{\mathbb{R}}
\def\Q{\mathbb{Q}}
\def\N{\mathbb{N}}
\def\E{\mathbb{E}}
\newcommand{\ps}[2]{\langle #1 , #2 \rangle}
\def\1{\mathbf{1}}
\newcommand{\nor}[1]{\left\| #1 \right\|}
\def\leq{\leqslant}
\def\geq{\geqslant}
\def\et{\quad \text{and} \quad}
\def\ou{\quad \text{or} \quad}

\DeclareMathOperator{\PPE}{PPE}
\DeclareMathOperator{\vNM}{vNM}
\DeclareMathOperator*{\argmax}{argmax}
\DeclareMathOperator*{\argmin}{argmin}
\renewcommand{\P}{\mathbb{P}}
\DeclareMathOperator{\Prob}{Prob}
\DeclareMathOperator{\supp}{supp}
\DeclareMathOperator{\Eq}{Eq}
\DeclareMathOperator{\inte}{int}
\DeclareMathOperator{\co}{co}
\DeclareMathOperator{\cl}{cl}
\DeclareMathOperator{\MRS}{MRS}
\DeclareMathOperator{\PV}{PV}
\DeclareMathOperator{\Range}{Im}
\DeclareMathOperator{\DIV}{DIV}
\DeclareMathOperator{\NPV}{NPV}
\newcommand{\IR}{\mathrm{IR}}

%-------------------- theorem environments --------------------------
\theoremstyle{plain}
\newtheorem{theorem}{Theorem}[chapter]
\newtheorem{proposition}[theorem]{Proposition}
\newtheorem{corollary}[theorem]{Corollary}
\newtheorem{lemma}[theorem]{Lemma}

\theoremstyle{definition}
\newtheorem{definition}[theorem]{Definition}
\newtheorem{assumption}[theorem]{Assumption}
\newtheorem{example}[theorem]{Example}

\theoremstyle{remark}
\newtheorem{remark}[theorem]{Remark}

\numberwithin{equation}{chapter}

%-------------------- title -----------------------------------------
\title{Theory of Corporate Finance \\[4pt]
       \large Lecture Notes}
\author{V. Filipe Martins-da-Rocha\thanks{Sao Paulo School of Economics -- FGV. E-mail: \texttt{filipe.econ@gmail.com}.}}
\date{2026}

%====================================================================
\begin{document}
%====================================================================

\frontmatter
\maketitle

\chapter*{Abstract}
\addcontentsline{toc}{chapter}{Abstract}

These lecture notes develop the modern theory of corporate finance from a unified theoretical perspective, organized around the central question of how informational and contractual frictions shape the financing, investment, governance, and risk-management decisions of firms.
The exposition begins with the foundations of stock-market equilibrium under incomplete markets and the Modigliani--Miller benchmark, then introduces moral hazard as the principal source of friction.
Within the principal--agent framework, the notes develop the canonical models of credit rationing, net worth, debt overhang, equity multipliers, debt maturity, liquidity management, and corporate risk management.
The analysis then turns to adverse selection, covering the lemons problem, market breakdown, the pecking-order hypothesis, and the role of dissipative signals.
The notes close with three more advanced topics: the industry equilibrium of mergers and acquisitions under asset redeployability, the role of active monitors and the dynamics of relationship lending, and the allocation of control rights and corporate governance.
Each chapter develops the model carefully, derives the optimal contract through complete proofs, and connects the theoretical results to the empirical patterns they explain.
The notes accompany the graduate course \emph{Theory of Corporate Finance} taught at the S\~ao Paulo School of Economics -- FGV.

\chapter*{Introduction}
\addcontentsline{toc}{chapter}{Introduction}

\section*{Purpose and scope}

These notes accompany the graduate course \emph{Theory of Corporate Finance} taught at the S\~ao Paulo School of Economics -- FGV.
The course is designed for first-year and second-year doctoral students with prior training in microeconomic theory and general equilibrium.
The objective is to provide a rigorous theoretical foundation for the modern study of corporate financing, capital structure, and corporate governance, organized around the principal--agent and adverse-selection frameworks that have shaped the field over the last four decades.

The notes draw heavily on the comprehensive treatment of \citet{Tirole2006}, \emph{The Theory of Corporate Finance}, which serves as the primary reference for most of the material covered here.
Specific chapters connect to the literature in distinctive ways: Part~I builds on the Arrow--Debreu and Diamond--Drèze--Grossman--Hart traditions of stock-market equilibrium, drawing on the unanimity-and-objective-of-the-firm literature and on the recent contributions of \citet{BisinGottardiRuta2015} and \citet{BisinClementiGottardi2025}.
Parts~II through~V cover the principal--agent and adverse-selection models that constitute the core of the modern corporate-finance literature.
Where appropriate, the notes point the reader to recent developments---continuous-time security design, post-crisis liquidity theory, dynamic monitoring, sovereign-debt models, and security-design surveys---as starting points for further reading.

\section*{Methodological perspective}

The notes adopt a deliberately \emph{theoretical} perspective.
The material is developed through formal models with complete proofs of the central propositions, rather than through descriptions of empirical regularities or case studies.
The choice reflects a pedagogical conviction: a graduate student in economics is best served by mastering the analytical machinery that produces the empirical predictions, both because the machinery is portable to new questions and because it provides the discipline through which empirical findings can be rigorously interpreted.

A second methodological choice is the emphasis on \emph{unified frameworks} over a survey of disparate models.
Each part of the notes builds on a small set of canonical models---the fixed-investment model, the variable-investment model, the liquidity-shock model, the lemons model---and develops their implications systematically.
This approach reveals the structural similarities across different corporate-finance phenomena: the equity multiplier of Chapter~\ref{ch:equity-multiplier}, the multiplier with redeployability of Chapter~\ref{ch:ma}, the balance-sheet thresholds of Chapter~\ref{ch:monitoring}, and the control-rights regimes of Chapter~\ref{ch:control} all share a common analytical structure, in which the entrepreneur's net worth $A$ partitions the space of feasible financing arrangements through endogenous thresholds.
Recognizing this common structure is, in our view, more valuable than mastering the surface features of any one model.

A third methodological choice is the relatively spare treatment of empirical-finance results.
The notes do refer to the empirical literature when it provides important grounding for the theoretical predictions---the negative stock-price reaction to seasoned equity offerings, the partition of borrowers between bond markets and bank lending, the use of debt covenants in distress---but the empirical material is offered as illustration of the theory rather than as a substitute for it.
Students seeking a comprehensive empirical view of corporate finance are referred to the survey literature and to the textbook of \citet{Tirole2006} (Chapter~2).

\section*{Structure}

The notes are organized into five Parts.
The structure follows a deliberate progression from frictionless benchmarks toward increasingly rich frictions, culminating in models that combine multiple frictions and address questions of corporate governance.

\textbf{Part~I (Competitive General Equilibrium)} establishes the foundations.
Chapter~\ref{ch:objective} treats the classical question of the firm's objective under incomplete markets, developing the Diamond--Drèze unanimity result and the modern analyses by \citet{MagillQuinzii2008}, \citet{BritzHeringsPredtetchinski2013}, and others.
Chapter~\ref{ch:capital-structure} develops the Modigliani--Miller capital-structure irrelevance theorem and its limits.
Chapters~\ref{ch:price-perceptions}--\ref{ch:moral-hazard-cpp} introduce the framework of competitive price perceptions following \citet{BisinGottardiRuta2015} and \citet{BisinClementiGottardi2025}, integrate moral hazard into the stock-market equilibrium following \citet{MagillQuinzii2002}, and characterize when constrained Pareto efficiency obtains.
Chapter~\ref{ch:prob-approach} develops the probability approach to moral hazard, the central technical tool for the principal--agent analysis that follows.

\textbf{Part~II (Principal--Agent Models)} develops the core moral-hazard theory of corporate finance.
Chapter~\ref{ch:debt-rationing} treats debt contracts under costly verification (Townsend--Gale--Hellwig) and the canonical model of credit rationing under adverse selection (Stiglitz--Weiss).
Chapter~\ref{ch:net-worth} introduces the net-worth-based framework of \citet{HolmstromTirole1997} and derives the financing threshold $\overline{A}$.
Chapter~\ref{ch:debt-overhang} develops the debt-overhang problem of \citet{Myers1977}.
Chapter~\ref{ch:equity-multiplier} extends the framework to variable investment, deriving the equity multiplier as a fundamental measure of how net worth amplifies investment capacity.
Chapter~\ref{ch:nonverifiable} treats the non-verifiable-income environment, where contractual enforcement rests on the threat of termination, and connects to the sovereign-debt literature of \citet{BulowRogoff1989a, BulowRogoff1989b} and beyond.
Chapters~\ref{ch:maturity}--\ref{ch:liquidity-scale} develop the liquidity-shock framework of \citet{HolmstromTirole1998, HolmstromTirole2000}, addressing the optimal maturity structure of liabilities and the liquidity--scale tradeoff.
Chapter~\ref{ch:risk-management} closes the part with the theory of corporate risk management, providing a clean analytical formulation of the \citet{FrootScharfsteinStein1993} mechanism using Jensen's inequality.

\textbf{Part~III (Adverse Selection)} turns from ex-post moral hazard to ex-ante asymmetric information.
Chapter~\ref{ch:lemons} develops the lemons problem of \citet{Akerlof1970} in the corporate-finance setting, deriving market breakdown, cross-subsidization, the index of adverse selection, market timing, the negative stock-price reaction to seasoned equity offerings, and the pecking-order hypothesis.
Chapter~\ref{ch:dissipative} examines four classes of dissipative signals---certification, short-term maturities, costly collateral pledging, and incomplete insurance through retained equity---each formalizing a costly action through which good types separate themselves from bad types.

\textbf{Part~IV (Mergers and Acquisitions)} treats the industry-equilibrium dimension of corporate finance through the asset-redeployability framework of \citet{ShleiferVishny1992}.
Chapter~\ref{ch:ma} develops a two-firm industry model with asset transfers in distress, deriving an endogenous illiquidity in the secondary asset market and an industry-level equity multiplier that decreases with the correlation of shocks.
The chapter analyzes coordination failures and the welfare consequences of underdeveloped resale markets.

\textbf{Part~V (Lending Relationships and Investor Activism)} addresses the role of active monitors and the allocation of control rights.
Chapter~\ref{ch:monitoring} develops the active-monitoring framework of \citet{HolmstromTirole1997}, deriving the three-regime partition of borrowers by balance-sheet strength, optimal monitoring intensity, and the dynamics of relationship lending with informational lock-in.
Chapter~\ref{ch:control} treats control-rights allocation following \citet{AghionBolton1992}, deriving the trade-off between net present value and pledgeable income that determines the equilibrium allocation of decision rights between the entrepreneur and the investors.

\section*{How to read these notes}

The notes are designed to be read either linearly or modularly.
Each Part is largely self-contained at the level of the central economic results, though the technical machinery---particularly the moral-hazard structure $(p_h, p_\ell, B, \Delta p, R, R_b)$ and the pledgeable-income notion $\rho_0$---is established in Part~II and used pervasively thereafter.
A reader interested specifically in the principal--agent theory of corporate finance can begin with Part~II after a brief glance at Chapter~\ref{ch:capital-structure} for the Modigliani--Miller benchmark.
A reader interested in the general-equilibrium foundations and the role of incomplete markets should read Part~I in full.
A reader interested in adverse selection and signaling can read Part~III after Part~II without further preliminaries.
Parts~IV and~V build directly on Part~II and can be read in either order.

The proofs are written to be \emph{verifiable}: each step is explicit, and the role of each assumption is identified.
The student is encouraged to work through the proofs in detail rather than treating them as decorative.
% The exercises provided in lecture (and not reproduced here) ask the student to extend the central results in directions that test mastery of the techniques.

\section*{Acknowledgments}

These notes have been refined through several iterations of teaching the course at UC Davis and FGV EPGE.
I am grateful to my students for many improvements to the presentation, and to the colleagues who have provided feedback on individual chapters.
Errors that remain are entirely my own.

\section*{A note on the preparation of these notes}
These lecture notes were prepared with the assistance of a generative-AI agent.
The starting material was a collection of beamer slide decks I have used in graduate courses on corporate finance at UC Davis and at EPGE/FGV over the past decade.
The AI agent was used to expand the slides into expository prose, to organize the material into a unified narrative across chapters, and to suggest recent papers in the literature that have appeared since the slides were originally written.
The author retains full responsibility for the final content: the selection of topics, the structure of the argument, the choice of notation, and the relative emphasis given to each model are mine.
Two caveats are warranted at this stage of the project.
First, the most recent papers suggested by the AI agent---particularly those that postdate the original slides---are still being read and evaluated by the author, and the corresponding ``Further Reading'' sections should be regarded as provisional.
Second, the formal proofs of the theoretical results stated throughout the text have, in many cases, been reconstructed by the AI agent rather than transcribed from a previously verified source.
Although I have read each proof and have not detected errors, a systematic verification is still in progress.
Readers who identify mistakes, mis-attributions, or arguments in need of refinement are warmly invited to contact me; their feedback will be incorporated into subsequent revisions.

\tableofcontents

\mainmatter

%%%%%%%%%%%%%%%%%%%%%%%%%%%%%%%%%%%%%%%%%%%%%%%%%%%%%%%%%%%%%%%%%%%%%%%%

\part{Competitive General Equilibrium}

%%%%%%%%%%%%%%%%%%%%%%%%%%%%%%%%%%%%%%%%%%%%%%%%%%%%%%%%%%%%%%%%%%%%%%%%

%====================================================================
\chapter{Stock-Market Equilibrium: The Objective of the Firm}
\label{ch:objective}
%====================================================================

A firm owned by several shareholders must choose a production plan and a financing plan.
Whose preferences should guide these choices?
If the firm is a sole proprietorship, the answer is straightforward: the single owner chooses the plan that maximizes his own expected utility.
If the firm has many shareholders with heterogeneous preferences, beliefs, and endowments, unanimity on the ``right'' plan is far from obvious.
Shareholders may rationally disagree, and without an agreed-upon objective, the very notion of ``managing the firm in the interest of its owners'' loses operational content.

The modern theory of corporate finance frames this question in a general-equilibrium setting with stock markets.
In this chapter, we study the simplest version of that framework: two dates, one commodity, a finite set of states of nature, and a single firm.
The object of inquiry is threefold.
First, we define rigorously what it means for the economy to be in a stock-market equilibrium.
Second, we derive the arbitrage pricing relations that link the firm's equity value to its cash flows.
Third, we identify the conditions---on markets and on preferences---under which shareholders unanimously agree that the firm's managerial objective should be to maximize the present value of its dividend stream.

The answer to the third question goes back to \citet{Fisher1930}, \citet{Dreze1974}, and \citet{GrossmanHart1979}: \emph{complete markets}---or, more generally, \emph{spanning of the firm's production set}---are sufficient for unanimity.
This is the bridge from the positive theory of the firm to a well-defined separation of ownership and control: shareholders can delegate production and financing decisions to a manager whose mandate is unambiguous, namely, to maximize firm value.

%====================================================================
\section{The Model}
\label{sec:model}
%====================================================================

\subsection{Time, uncertainty, and the commodity}

There are two dates, $t \in \{0,1\}$.
Uncertainty at $t=1$ is summarized by a finite set $S$ of states of nature.
One divisible, perishable commodity is available for trade at every date and every state.
All quantities are measured in units of that commodity.

\subsection{The firm}

The economy contains one firm, privately owned by investors.
The firm's production possibilities are described by a set
\begin{equation*}
  Y \subset \R \times \R^{S}.
\end{equation*}
A \emph{production plan} for the firm is a vector $y = (y_0, y_1) \in Y$ with $y_1 = (y_s)_{s \in S} \in \R^S$.

\begin{assumption}[No instantaneous production]\label{ass:no-instant}
The firm invests at $t=0$ and produces output at $t=1$.
Formally,
\begin{equation*}
  Y \subset \R_{-} \times \R^{S}_{+}.
\end{equation*}
\end{assumption}

Under Assumption~\ref{ass:no-instant}, for $y = (y_0, (y_s)_{s \in S}) \in Y$ the coordinate $-y_0 \geq 0$ is the amount of the commodity used as input at $t=0$, and $y_s \geq 0$ is the amount of the commodity produced at $t=1$ in state $s$.

\subsection{Shares and financial assets}

The firm is owned by a finite set $I$ of investors.
We denote by $\xi^i \geq 0$ the share of the firm initially owned by investor $i$ at $t=0$.
Shares are percentages of ownership, so
\begin{equation*}
  \sum_{i \in I} \xi^i = 1.
\end{equation*}
Some investors may hold no shares ($\xi^i = 0$); one of them may own the entire firm ($\xi^i = 1$).

Let $A$ be a finite set of financial assets (or securities).
Each asset $a \in A$ is a claim to the state-contingent payoff
\begin{equation*}
  \kappa(a) = (\kappa_s(a))_{s \in S} \in \R^S.
\end{equation*}
At $t=0$, any agent (investor or firm) may choose a \emph{portfolio} $\theta \in \R^A$:
\begin{itemize}
  \item if $\theta(a) \geq 0$, the agent purchases $\theta(a)$ units of asset $a$ and receives $\kappa_s(a) \theta(a)$ of the commodity at $t=1$ contingent on state $s$;
  \item if $\theta(a) \leq 0$, the agent short-sells $|\theta(a)|$ units of asset $a$ and delivers $|\kappa_s(a) \theta(a)|$ of the commodity at $t=1$ in state $s$.
\end{itemize}
An asset $a$ is called \emph{risk-free} (or \emph{riskless}, or \emph{non-contingent}) if $\kappa(a)$ is constant across states; up to normalization, $\kappa_s(a) = 1$ for every $s$.

\subsection{Competitive markets}

At $t=0$ there are two markets: one for the firm's equity and one for financial assets.

\emph{Stock market.}
The firm's equity trades at a linear price $E$ (in units of the unique commodity).
Agents take $E$ as given.
Buying a fraction $\alpha \geq 0$ of the firm costs $E \alpha$ units of the commodity; selling yields $E \alpha$.

\emph{Security markets.}
Each security $a$ trades at a linear price $q(a)$.
Buying the portfolio $\theta \in \R^A$ costs
\begin{equation*}
  q \cdot \theta = \sum_{a \in A} q(a) \theta(a).
\end{equation*}
If $\theta(a) < 0$, the agent sells short and receives $|q(a) \theta(a)|$ at $t=0$.

We denote by $w(a) = q(a) \theta(a)$ the resources spent on asset $a$ and define the \emph{gross state-contingent return}
\begin{equation*}
  R_s(a) = \frac{\kappa_s(a)}{q(a)}
  \et
  r_s(a) = R_s(a) - 1 = \frac{\kappa_s(a) - q(a)}{q(a)}
\end{equation*}
for each asset with $q(a) > 0$; equivalently, $\kappa_s(a) \theta(a) = R_s(a) w(a)$.

\subsection{Investors}

Each investor $i \in I$ is characterized by
\begin{itemize}
  \item an initial commodity endowment $\omega^i_0 \geq 0$;
  \item a vector of initial share holdings $\xi^i \geq 0$;
  \item a state-contingent endowment $\omega^i_1 = (\omega^i_s)_{s \in S} \in \R^{S}_{+}$;
  \item a Bernoulli utility function $u^i \colon \R_+ \to [-\infty, \infty)$;
  \item a discount factor $\beta^i > 0$;
  \item a belief $\nu^i \in \Prob(S)$ over states of nature.
\end{itemize}
At $t=0$ investor $i$ chooses current consumption $c_0 \in \R_+$ and a plan for state-contingent consumption $c_1 = (c_s)_{s \in S} \in \R_+^S$.
His expected utility is
\begin{equation}\label{eq:expected-utility}
  U^i(c_0, c_1) = u^i(c_0) + \beta^i \sum_{s \in S} \nu^i_s u^i(c_s).
\end{equation}

\subsection{The firm's cash flow}

At $t=0$, given the security price vector $q$, the firm chooses:
\begin{itemize}
  \item a (physical) production plan $y = (y_0, y_1) \in Y$;
  \item a financial plan $\beta \in \R^A$, where $\beta(a) > 0$ means the firm short-sells $\beta(a)$ units of asset $a$ (issues debt or securities), and $\beta(a) < 0$ means the firm purchases $|\beta(a)|$ units.
\end{itemize}

\begin{remark}\label{rem:beta-abuse}
The symbol $\beta$ is used here in two distinct senses, following the slides: $\beta^i > 0$ denotes investor $i$'s discount factor, while $\beta \in \R^A$ denotes the firm's financial plan.
Context disambiguates the two; the discount factor always carries a superscript identifying an investor.
\end{remark}

The production--finance plan $(y, \beta)$ generates a state- and date-dependent cash flow
\begin{equation*}
  \delta = (\delta_0, \delta_1) \in \R \times \R^S
\end{equation*}
defined by
\begin{equation}\label{eq:cash-flow}
  \delta_0 = y_0 + q \cdot \beta
  \et
  \delta_s = y_s - \kappa_s \cdot \beta.
\end{equation}
At $t=0$, shareholders of the firm (agents with $\xi^i > 0$) collectively finance the expenditure $\delta_0$ if $\delta_0 < 0$, or receive it if $\delta_0 > 0$; agent $i$'s share is $\xi^i \delta_0$.
At $t=1$ in state $s$, shareholders owning the firm \emph{at that date} (agents with $\alpha^i > 0$, where $\alpha^i$ is $i$'s post-trade share) share the surplus $\delta_s$ in proportion to $\alpha^i$.

\emph{Typical interpretation at $t=0$.}
At $t=0$, the firm produces a negative output $y_0 \leq 0$ (the investment) and finances it by short-selling the portfolio $\beta \in \R^A_+$.
If $y_0 + q \cdot \beta \leq 0$, only part of the investment is financed internally and the rest is financed by shareholders through the negative dividend $\delta_0 < 0$; if $y_0 + q \cdot \beta \geq 0$, the entire investment is financed by the issuance of securities and shareholders receive a positive dividend $\delta_0 \geq 0$.

\emph{Typical interpretation at $t=1$.}
When state $s$ realizes, the firm produces $y_s \geq 0$ and must honor its financial obligation $\kappa_s \cdot \beta$.
Each shareholder receives or delivers the residual $\delta_s = y_s - \kappa_s \cdot \beta$ in proportion to his post-trade share $\alpha^i$.

\subsection{Ownership structure and limited liability}

The model accommodates several ownership structures.

\emph{Sole proprietorship.}
No stock market exists.
The firm is owned by a single investor, $\xi^{i^\ast} = 1$, who manages it according to his own preferences.
No question of ``objective'' arises.

\emph{Partnership.}
No stock market exists, but several shareholders with $\xi^i > 0$ jointly own the firm.
Disagreement about the plan is possible.

\emph{Corporation.}
A stock market exists, and ownership is separable from control.
The manager chooses $(y, \beta)$ on behalf of shareholders, who then trade shares freely at $t=0$.
Identifying the manager's objective is the central conceptual problem.

The formulation in equation~\eqref{eq:cash-flow} implicitly allows $\delta_s < 0$, in which case the holder of a share delivers $\alpha^i |\delta_s|$ at $t=1$ in state $s$.
In practice, shareholders enjoy \emph{limited liability}: they cannot be forced to pay negative residuals beyond their equity stake.
To accommodate this, consider the firm financing its investment exclusively through a riskless bond $a_0$ with $\kappa_s(a_0) = 1$ for every $s$, issued in quantity $Z \geq 0$.
The firm is then called \emph{levered}.
If in state $s$ the firm's output $y_s$ falls short of the face value $Z$, bankruptcy rules apply: the entire output $y_s$ is distributed to bondholders and shareholders receive nothing.
The issue of limited liability and its implications for capital structure is deferred to a later chapter; for the purpose of defining equilibrium we allow $\delta_s$ to take negative values.

%====================================================================
\section{Competitive Stock-Market Equilibrium}
\label{sec:equilibrium}
%====================================================================

\subsection{Investors' budget sets}

Each investor $i$ takes as given, and perfectly anticipates:
\begin{itemize}
  \item the firm's production--finance plan $(y, \beta)$, and hence the resulting cash flow $\delta$;
  \item the security price vector $q$ and the equity price $E$.
\end{itemize}
He chooses a consumption plan $(c_0, c_1) \in \R_+ \times \R_+^S$, a security portfolio $\theta \in \R^A$, and a post-trade share holding $\alpha \in \R_+$ satisfying the budget constraints
\begin{align}
  c_0 + q \cdot \theta + E \alpha &\leq \omega^i_0 + (E + \delta_0) \xi^i,
  \label{eq:budget-0}\\[2pt]
  c_s &\leq \omega^i_s + \kappa_s \cdot \theta + \delta_s \alpha,
  \quad \text{for all $s\in S$}.
  \label{eq:budget-s}
\end{align}
The set of actions $(c, \theta, \alpha)$ satisfying \eqref{eq:budget-0} and \eqref{eq:budget-s} is the budget set, denoted
\begin{equation*}
  B^i(q, E) \quad \text{(or, more precisely, } B^i(q, E; y, \beta) \text{, since } \delta_0 = \delta_0(q, y, \beta) \text{).}
\end{equation*}
The notational economy $B^i(q, E)$ assumes that $(y, \beta)$ is fixed in context.

\subsection{Equilibrium}

\begin{definition}[Stock-market equilibrium]\label{def:equilibrium}
A \emph{stock-market equilibrium} is a list
\begin{equation*}
  \bigl\{ (y, \beta), (q, E), (c^i, \theta^i, \alpha^i)_{i \in I} \bigr\}
\end{equation*}
consisting of a production--finance plan $(y, \beta)$ for the firm, security and equity prices $(q, E)$, and an action $(c^i, \theta^i, \alpha^i)$ for each investor, such that:
\begin{enumerate}
  \item[(a)] for every $i \in I$,
  \begin{equation*}
    (c^i, \theta^i, \alpha^i) \in \argmax \bigl\{ U^i(c) \colon (c, \theta, \alpha) \in B^i(q, E) \bigr\};
  \end{equation*}
  \item[(b)] security and stock markets clear:
  \begin{equation*}
    \forall a \in A, \quad \beta(a) = \sum_{i \in I} \theta^i(a)
    \et
    \sum_{i \in I} \alpha^i = 1;
  \end{equation*}
  \item[(c)] the commodity market clears:
  \begin{equation*}
    \sum_{i \in I} c^i = \sum_{i \in I} \omega^i + y.
  \end{equation*}
\end{enumerate}
\end{definition}

Condition (a) is individual optimality; condition (b) is market clearing for securities and equity; condition (c) is market clearing for the commodity.
By Walras' law, (c) is implied by (a)--(b) if consumption is interior; we state it separately to simplify exposition.

\subsection{Standard assumptions on primitives}

Throughout the analysis, we rely on one or more of the following assumptions.

\begin{assumption}[C.1--C.4]\label{ass:C}
For each investor $i \in I$:
\begin{itemize}
  \item[\textup{(C.1)}] the Bernoulli utility $u^i$ is continuously differentiable, concave, and strictly increasing;
  \item[\textup{(C.2)}] commodity endowments are strictly positive: $\omega^i_0 > 0$ and $\omega^i_s > 0$ for every $s \in S$;
  \item[\textup{(C.3)}] every state has a positive probability: $\nu^i_s > 0$ for every $s \in S$;
  \item[\textup{(C.4)}] individually rational consumption plans are interior:
  \begin{equation*}
    U^i(c_0, c_1) \geq U^i(\omega^i_0, \omega^i_1)
    \Longrightarrow
    c_0 > 0 \et c_s > 0, \ \text{for all $s\in S$}.
  \end{equation*}
\end{itemize}
\end{assumption}

Condition (C.4) is an Inada-type hypothesis at the \emph{endogenous} level of welfare.
In many specifications it follows from stronger conditions on $u^i$ at the origin, e.g.\ $\lim_{c \to 0^+} (u^i)'(c) = +\infty$, combined with (C.2)--(C.3).

%====================================================================
\section{Asset Pricing at Equilibrium}
\label{sec:asset-pricing}
%====================================================================

At any competitive equilibrium, prices must be free of arbitrage.
This section makes that precise and delivers two central results: the existence of state-price deflators (the \emph{fundamental asset-pricing theorem}) and their identification with marginal rates of substitution of individual investors.

\subsection{State-price deflators}

Given an investor $i$ with interior optimal consumption $c^i \gg 0$, his \emph{marginal rate of substitution} between state $s$ at $t=1$ and consumption at $t=0$ is
\begin{equation}\label{eq:mrs-def}
  \MRS^i_s = \MRS^i_s(c^i)
  = \beta^i \nu^i_s \frac{(u^i)'(c^i_s)}{(u^i)'(c^i_0)}.
\end{equation}

\begin{theorem}[Fundamental asset-pricing theorem]\label{thm:fundamental}
Suppose Assumptions \textup{(C.1)--(C.3)} hold and let
\begin{equation*}
  \bigl\{ (y, \beta), (q, E), (c^i, \theta^i, \alpha^i)_{i \in I} \bigr\}
\end{equation*}
be a stock-market equilibrium.
There exists a vector of state-price deflators $\lambda = (\lambda_s)_{s \in S} \in \R^S_{++}$ such that
\begin{equation}\label{eq:asset-pricing}
  q(a) = \sum_{s \in S} \lambda_s \kappa_s(a),
  \quad \forall a \in A,
\end{equation}
and
\begin{equation}\label{eq:equity-pricing}
  E = \sum_{s \in S} \lambda_s \delta_s
    = \sum_{s \in S} \lambda_s [ y_s - \kappa_s \cdot \beta ].
\end{equation}
\end{theorem}

\begin{proof}
We prove the result under the additional assumption (C.4), so that each investor has interior optimal consumption and the first-order conditions of his problem hold with equality.
The conclusions \eqref{eq:asset-pricing}--\eqref{eq:equity-pricing} extend to the general (C.1)--(C.3) case by standard no-arbitrage arguments (see Remark~\ref{rem:no-arbitrage} below).

Fix an investor $i$ with interior optimal consumption $(c^i_0, c^i_1) \gg 0$ (such an investor exists under our assumptions).
Because $u^i$ is concave and the constraints \eqref{eq:budget-0}--\eqref{eq:budget-s} are linear, the Karush--Kuhn--Tucker conditions are necessary and sufficient for the optimum.
Let $\mu^i_0 \geq 0$ and $\mu^i_s \geq 0$ denote the Lagrange multipliers attached to the $t=0$ constraint \eqref{eq:budget-0} and the state-$s$ constraint \eqref{eq:budget-s}, respectively.
The Lagrangian is
\begin{equation*}
  \mathcal{L}^i = U^i(c) - \mu^i_0 \bigl[ c_0 + q \cdot \theta + E \alpha - \omega^i_0 - (E + \delta_0) \xi^i \bigr]
  - \sum_{s \in S} \mu^i_s \bigl[ c_s - \omega^i_s - \kappa_s \cdot \theta - \delta_s \alpha \bigr].
\end{equation*}
The interior first-order conditions with respect to $c_0$, $c_s$, $\theta(a)$, and $\alpha$ read
\begin{align}
  (u^i)'(c^i_0) &= \mu^i_0,
  \label{eq:foc-c0}\\
  \beta^i \nu^i_s (u^i)'(c^i_s) &= \mu^i_s,
  \quad \text{for all $s\in S$},
  \label{eq:foc-cs}\\
  \mu^i_0 q(a) &= \sum_{s \in S} \mu^i_s \kappa_s(a),
  \quad \forall a \in A,
  \label{eq:foc-theta}\\
  \mu^i_0 E &\geq \sum_{s \in S} \mu^i_s \delta_s,
  \label{eq:foc-alpha}
\end{align}
with equality in \eqref{eq:foc-alpha} whenever $\alpha^i > 0$ (the multiplier on $\alpha \geq 0$ vanishes by complementary slackness).
Define
\begin{equation}\label{eq:lambda-def}
  \lambda_s = \frac{\mu^i_s}{\mu^i_0} = \beta^i \nu^i_s \frac{(u^i)'(c^i_s)}{(u^i)'(c^i_0)} = \MRS^i_s.
\end{equation}
Since $u^i$ is strictly increasing by (C.1), $(u^i)'(c^i_s) > 0$; since $\nu^i_s > 0$ by (C.3) and $\beta^i > 0$, we obtain $\lambda_s > 0$ for every $s \in S$.
Dividing \eqref{eq:foc-theta} by $\mu^i_0 > 0$ gives \eqref{eq:asset-pricing}.

For \eqref{eq:equity-pricing} we need an investor with $\alpha^i > 0$.
Market clearing for the stock, $\sum_i \alpha^i = 1 > 0$, guarantees that some investor, say $i^\ast$, holds $\alpha^{i^\ast} > 0$; his first-order condition \eqref{eq:foc-alpha} holds with equality, yielding
\begin{equation*}
  E = \sum_{s \in S} \frac{\mu^{i^\ast}_s}{\mu^{i^\ast}_0} \delta_s = \sum_{s \in S} \lambda^{i^\ast}_s \delta_s.
\end{equation*}
The deflator $\lambda \coloneqq (\lambda^{i^\ast}_s)_{s \in S}$ satisfies both \eqref{eq:asset-pricing} (because \eqref{eq:foc-theta} holds for $i^\ast$) and \eqref{eq:equity-pricing}.
This completes the proof.
\end{proof}

\begin{remark}[No-arbitrage without interiority]\label{rem:no-arbitrage}
Theorem~\ref{thm:fundamental} holds under (C.1)--(C.3) alone, without (C.4).
The argument is a direct no-arbitrage one: if there existed $\theta \in \R^A$ with $q \cdot \theta \leq 0$ and $\kappa \cdot \theta \geq 0$, with at least one strict inequality, then any investor could increase his utility by an arbitrarily small amount of $\theta$ (by strict monotonicity of $u^i$ and positivity of $\nu^i_s$), contradicting optimality at equilibrium.
Stiemke's lemma then yields $\lambda \in \R_{++}^S$ satisfying \eqref{eq:asset-pricing}.
Equity pricing \eqref{eq:equity-pricing} follows analogously by considering an investor with $\alpha^i > 0$: if $E < \sum_s \lambda_s \delta_s$, he would increase $\alpha$; if $E > \sum_s \lambda_s \delta_s$, he would decrease $\alpha$.
\end{remark}

\subsection{Personalized state prices}

Under (C.4), the first-order-condition argument applies to \emph{every} investor, yielding personalized pricing formulas.

\begin{proposition}[Personalized state prices]\label{prop:personalized}
Suppose Assumptions \textup{(C.1)--(C.4)} hold and let
\begin{equation*}
  \bigl\{ (y, \beta), (q, E), (c^i, \theta^i, \alpha^i)_{i \in I} \bigr\}
\end{equation*}
be a stock-market equilibrium.
Then for every investor $i \in I$,
\begin{equation}\label{eq:personalized-q}
  q(a) = \sum_{s \in S} \MRS^i_s \kappa_s(a), \quad \forall a \in A.
\end{equation}
Moreover, for every shareholder of the firm at $t=1$ (i.e., for every $i$ with $\alpha^i > 0$),
\begin{equation}\label{eq:personalized-E}
  E = \sum_{s \in S} \MRS^i_s \delta_s.
\end{equation}
\end{proposition}

\begin{proof}
Under (C.4), every investor has interior optimal consumption, so the first-order conditions \eqref{eq:foc-c0}--\eqref{eq:foc-alpha} hold with equality for every $i$.
Dividing \eqref{eq:foc-theta} by $\mu^i_0$ and using \eqref{eq:lambda-def} gives \eqref{eq:personalized-q}.
For investors with $\alpha^i > 0$, the complementary-slackness condition gives equality in \eqref{eq:foc-alpha} and hence \eqref{eq:personalized-E}.
\end{proof}

When all investors share common beliefs, $\nu^i = \nu$, the quantity $\beta^i (u^i)'(c^i_s) / (u^i)'(c^i_0)$ is the \emph{stochastic discount factor} of investor $i$ relative to $\nu$, and \eqref{eq:personalized-q}--\eqref{eq:personalized-E} say that each asset and the firm's equity are priced using $i$'s stochastic discount factor.

\begin{remark}[Uniqueness of $\lambda$]\label{rem:unique-lambda}
In general, Theorem~\ref{thm:fundamental} does not pin down $\lambda$ uniquely: if $\Range \kappa \subsetneq \R^S$ (incomplete markets), there exist multiple $\lambda$'s satisfying \eqref{eq:asset-pricing}, and Proposition~\ref{prop:personalized} tells us that different investors may pick different $\MRS^i$'s.
Uniqueness obtains when markets are \emph{complete}, $\Range \kappa = \R^S$, in which case \eqref{eq:asset-pricing} is a system of $|A| \geq |S|$ equations with a unique solution $\lambda \in \R^S_{++}$.
In that case \eqref{eq:personalized-q} forces $\MRS^i = \lambda$ for every $i$.
\end{remark}

%====================================================================
\section{The Objective of the Firm}
\label{sec:objective}
%====================================================================

\subsection{An arbitrary candidate objective}

A natural, if arbitrary, candidate for the firm's objective is to maximize its \emph{value}
\begin{equation}\label{eq:firm-value}
  V = E + \delta_0 = E + q \cdot \beta + y_0 = E + D + y_0,
\end{equation}
where $D \coloneqq q \cdot \beta$ is the value of the financial claims issued by the firm (``debt'').
Observe that $V$ depends jointly on the firm's decision $(y, \beta)$ and on the equilibrium prices $(E, q)$, which in turn depend on $(y, \beta)$ through the equilibrium fixed point.
Why should the firm pursue this particular objective rather than, say, a weighted sum of shareholders' utilities?
The answer cannot come from the objective itself: it must come from shareholders' preferences.

\subsection{The separation of ownership and control}

The firm belongs to its shareholders, i.e., to the set $\{ i \colon \xi^i > 0 \}$.
Each shareholder $i$ would, if he could dictate the choice, prescribe the production--finance plan $(y^i, \beta^i)$ that maximizes his own expected utility.
If different shareholders prescribe different plans, there is a conflict.

If, however, there exists a function of $(y, \beta)$ such that every shareholder agrees that the firm's manager should maximize it, we say that there is \emph{unanimity among shareholders}.
Under unanimity, the ownership--control distinction is non-trivial: shareholders \emph{delegate} control to a professional manager whose mandate is operational and unambiguous.
This is the core normative content of corporate governance at its simplest.

Under what conditions can ownership and control be separated?
A trivial answer is: under certainty (no uncertainty) or with complete markets.
The less trivial answer---due to \citet{Fisher1930}, \citet{Dreze1974}, \citet{GrossmanHart1979}, and \citet{GrossmanStiglitz1980}---is that \emph{complete markets suffice}, and that the candidate objective \eqref{eq:firm-value} is the right one.
We now develop this result.

\subsection{Impact of the firm's decision on a single investor}

Fix an equilibrium $\{ (y, \beta), (q, E), (c^i, \theta^i, \alpha^i)_{i \in I} \}$ and consider an infinitesimal ``out-of-equilibrium'' perturbation $(\Delta y, \Delta \beta)$ of the firm's plan.
We analyze its impact on investor $i$'s utility under two auxiliary conventions:
\begin{enumerate}
  \item we neglect the impact of the perturbation on equilibrium prices \emph{other than} the firm's equity price $E$ (the security prices $q$ are held constant);
  \item investor $i$ presumes that the perturbation changes the equity price by some amount $\Delta^i E$, whose specification we postpone.
\end{enumerate}

Let $\Delta^i U^i$ denote the change in investor $i$'s expected utility presumed under these conventions, and let
\begin{equation*}
  \Delta \delta_0 = \Delta y_0 + q \cdot \Delta \beta
  \et
  \Delta \delta_s = \Delta y_s - \kappa_s \cdot \Delta \beta
\end{equation*}
denote the corresponding changes in the firm's cash flow.

\begin{proposition}[First-order impact]\label{prop:impact}
Under \textup{(C.1)--(C.4)},
\begin{equation}\label{eq:impact-general}
  \frac{\Delta^i U^i}{(u^i)'(c^i_0)}
  \approx
  \xi^i \bigl[ \Delta^i E + \Delta \delta_0 \bigr]
  + \alpha^i \Bigl[ \sum_{s \in S} \MRS^i_s \Delta \delta_s - \Delta^i E \Bigr].
\end{equation}
\end{proposition}

\begin{proof}
Keeping investor $i$'s own choices $(\theta^i, \alpha^i)$ fixed to first order, and using \eqref{eq:budget-0}--\eqref{eq:budget-s}, the induced change in his consumption is
\begin{equation*}
  \Delta c^i_0 = \xi^i \bigl[ \Delta^i E + \Delta \delta_0 \bigr] - \alpha^i \Delta^i E,
  \quad
  \Delta c^i_s = \alpha^i \Delta \delta_s.
\end{equation*}
Indeed, at $t=0$ the right-hand side of \eqref{eq:budget-0} changes by $\xi^i (\Delta^i E + \Delta \delta_0)$ while the left-hand side changes by $\Delta c^i_0 + \alpha^i \Delta^i E$ (since $q$ and $(\theta^i, \alpha^i)$ are held constant); at $t=1$ in state $s$, the right-hand side of \eqref{eq:budget-s} changes by $\alpha^i \Delta \delta_s$.
A first-order expansion of $U^i$ then gives
\begin{equation*}
  \Delta^i U^i \approx (u^i)'(c^i_0) \Delta c^i_0
  + \sum_{s \in S} \beta^i \nu^i_s (u^i)'(c^i_s) \Delta c^i_s,
\end{equation*}
and dividing through by $(u^i)'(c^i_0)$ and using \eqref{eq:mrs-def} yields
\begin{equation*}
  \frac{\Delta^i U^i}{(u^i)'(c^i_0)}
  \approx \Delta c^i_0 + \sum_{s \in S} \MRS^i_s \Delta c^i_s
  = \xi^i [\Delta^i E + \Delta \delta_0] - \alpha^i \Delta^i E + \alpha^i \sum_{s \in S} \MRS^i_s \Delta \delta_s,
\end{equation*}
which is \eqref{eq:impact-general}.
\end{proof}

Formula \eqref{eq:impact-general} shows that $i$'s welfare depends on (i) his initial share $\xi^i$, which scales the direct dividend effect $\Delta^i E + \Delta \delta_0$, and (ii) his post-trade share $\alpha^i$, which scales the \emph{valuation gap} $\sum_s \MRS^i_s \Delta \delta_s - \Delta^i E$ between $i$'s assessment of the new state-contingent dividend and his presumed change in the market valuation of equity.

\subsection{Competitive price perceptions}

To close \eqref{eq:impact-general}, we must specify $\Delta^i E$.
\citet{GrossmanHart1979} propose the hypothesis of \emph{competitive price perceptions}.

\begin{definition}[Competitive price perceptions]\label{def:cpp}
Investor $i$ is said to have \emph{competitive price perceptions} if, given his equilibrium consumption $c^i$, he values any income stream $w \in \R^S$ (arriving at $t=1$) using his own marginal rates of substitution as state-price deflators:
\begin{equation}\label{eq:cpp}
  \pi^i(w) \coloneqq \sum_{s \in S} \MRS^i_s w_s.
\end{equation}
\end{definition}

The interpretation is as follows.
If a hypothetical new asset paying $w$ were introduced in the economy \emph{after} equilibrium has formed, investor $i$ would be indifferent, at his current consumption $c^i$, between buying or selling an infinitesimal amount of this asset at price $\pi^i(w)$.
Under (C.1)--(C.4), this is an immediate consequence of Proposition~\ref{prop:personalized}.

\begin{proposition}[Market consistency]\label{prop:market-consistency}
Under \textup{(C.1)--(C.4)}, competitive price perceptions are consistent with market prices of \emph{marketed} streams: if $w = \kappa \cdot \theta_w$ for some $\theta_w \in \R^A$, then for every investor $i$,
\begin{equation}\label{eq:market-consistency}
  \sum_{s \in S} \MRS^i_s w_s = q \cdot \theta_w.
\end{equation}
In particular, two investors with different marginal rates of substitution still agree on the present value of any marketed income stream.
\end{proposition}

\begin{proof}
By Proposition~\ref{prop:personalized}, $q(a) = \sum_s \MRS^i_s \kappa_s(a)$ for every $a \in A$ and every $i$.
Hence
\begin{eqnarray*}
  \sum_{s \in S} \MRS^i_s w_s
  & = & \sum_{s \in S} \MRS^i_s \sum_{a \in A} \kappa_s(a) \theta_w(a) \\
  & = & \sum_{a \in A} \Bigl[ \sum_{s \in S} \MRS^i_s \kappa_s(a) \Bigr] \theta_w(a) \\
  & = & \sum_{a \in A} q(a) \theta_w(a) = q \cdot \theta_w. \qedhere
\end{eqnarray*}
\end{proof}

\subsection{Present-value characterization of the firm's objective}

Suppose investor $i$ has competitive price perceptions.
In response to a perturbation $(\Delta y, \Delta \beta)$, he presumes that the change in equity price is
\begin{equation}\label{eq:delta-E-i}
  \Delta^i E = \sum_{s \in S} \MRS^i_s \Delta \delta_s,
\end{equation}
i.e., he values the new stream $\delta + \Delta \delta$ at the present value implied by his own marginal rates of substitution.
Substituting \eqref{eq:delta-E-i} into \eqref{eq:impact-general} collapses the $\alpha^i$-bracket to zero, yielding
\begin{equation}\label{eq:impact-cpp}
  \frac{\Delta^i U^i}{(u^i)'(c^i_0)}
  = \xi^i \Bigl[ \sum_{s \in S} \MRS^i_s \Delta \delta_s + \Delta \delta_0 \Bigr].
\end{equation}
Thus, \emph{under competitive price perceptions}, a shareholder $i$ (with $\xi^i > 0$) welcomes any perturbation that increases
\begin{equation*}
  \delta_0 + \sum_{s \in S} \MRS^i_s \delta_s,
\end{equation*}
the present value of the firm's dividend stream evaluated at \emph{his own} marginal rate of substitution.
Since $\xi^i \geq 0$ appears as a scale factor, the direction of preferred changes is determined by $\MRS^i$ alone.

\subsection{Disagreement under heterogeneous marginal rates of substitution}

Formula \eqref{eq:impact-cpp} delivers the fundamental tension: two shareholders $i_1 \neq i_2$ with $\MRS^{i_1} \neq \MRS^{i_2}$ evaluate the perturbation $(\Delta y, \Delta \beta)$ using different weights.
It is possible to construct $\Delta \delta$ such that
\begin{equation*}
  \sum_{s \in S} \MRS^{i_1}_s \Delta \delta_s + \Delta \delta_0 > 0
  \et
  \sum_{s \in S} \MRS^{i_2}_s \Delta \delta_s + \Delta \delta_0 < 0,
\end{equation*}
and then $\Delta^{i_1} U^{i_1} > 0$ while $\Delta^{i_2} U^{i_2} < 0$: the two shareholders disagree on whether the firm should implement the perturbation.
There is no common objective.

Heterogeneity of $\MRS$ across investors is precisely what incomplete markets may generate: when $\Range \kappa \subsetneq \R^S$, different investors achieve different marginal valuations across states because the available assets do not let them equate marginal rates of substitution via trade.
In this case, the conceptual program of separating ownership and control breaks down.

%====================================================================
\section{Unanimity and Value Maximization}
\label{sec:unanimity}
%====================================================================

\subsection{Unanimity from complete markets}

\begin{proposition}[Unanimity under complete markets]\label{prop:unanimity-complete}
Under \textup{(C.1)--(C.4)}, if security markets are complete, i.e., $\Range \kappa = \R^S$, then there exists a unique state-price deflator $\lambda \in \R^S_{++}$ such that for every investor $i \in I$,
\begin{equation*}
  \MRS^i_s = \lambda_s, \quad \text{for all $s\in S$}.
\end{equation*}
All investors have identical competitive price perceptions, and unanimously agree that the firm's objective is to maximize its value
\begin{equation}\label{eq:firm-value-complete}
  V = \delta_0 + \sum_{s \in S} \lambda_s \delta_s
    = y_0 + \sum_{s \in S} \lambda_s y_s
    = E + D + y_0.
\end{equation}
\end{proposition}

\begin{proof}
When $\Range \kappa = \R^S$, the linear system \eqref{eq:asset-pricing}, viewed as an equation for $\lambda \in \R^S$, has a unique solution.
Proposition~\ref{prop:personalized} then forces $\MRS^i_s = \lambda_s$ for every $i$ and every $s$.
Evaluating \eqref{eq:firm-value} with $\delta_0 = y_0 + q \cdot \beta$ and $q(a) = \sum_s \lambda_s \kappa_s(a)$ yields
\begin{equation*}
  \delta_0 + \sum_s \lambda_s \delta_s
  = y_0 + q \cdot \beta + \sum_s \lambda_s [y_s - \kappa_s \cdot \beta]
  = y_0 + \sum_s \lambda_s y_s + \beta \cdot \Bigl[ q - \sum_s \lambda_s \kappa_s \Bigr]
  = y_0 + \sum_s \lambda_s y_s,
\end{equation*}
the last equality by no-arbitrage pricing \eqref{eq:asset-pricing}.
Unanimity on \eqref{eq:firm-value-complete} is immediate from \eqref{eq:impact-cpp}: every shareholder $i$ welcomes any perturbation that increases $\delta_0 + \sum_s \lambda_s \delta_s$.
\end{proof}

A striking by-product of \eqref{eq:firm-value-complete} is that the firm's value depends on $(y, \beta)$ \emph{only through} $y$: the financial plan $\beta$ drops out.
This is a version of the Modigliani--Miller irrelevance theorem, to which we return in a later chapter.

\subsection{Spanning}

Complete markets are sufficient, but not strictly necessary, for unanimity.
A weaker sufficient condition is a \emph{spanning property} of the firm's production set relative to the available securities.

\begin{proposition}[Spanning implies unanimity]\label{prop:spanning}
Fix a stock-market equilibrium and assume that for every feasible production change $\Delta y \in Y - \{ y \}$ there exist $\Delta \theta \in \R^A$ and $\Delta \alpha \in \R$ such that
\begin{equation}\label{eq:spanning}
  \Delta y_s = \sum_{a \in A} \kappa_s(a) \Delta \theta(a) + y_s \Delta \alpha,
  \quad \text{for all $s\in S$}.
\end{equation}
Then every shareholder agrees that the firm should maximize $\delta_0 + \sum_s \lambda_s \delta_s$ for \emph{any} state-price deflator $\lambda$ satisfying \eqref{eq:asset-pricing}--\eqref{eq:equity-pricing}.
\end{proposition}

The economic content of \eqref{eq:spanning} is that any conceivable change in the firm's output can be replicated by a combination of portfolio adjustments $\Delta \theta$ and proportional rescaling $\Delta \alpha$ of the existing production $y$.
In particular, shareholders do not lose hedging opportunities when the firm adjusts its plan: whatever the firm gives up, they can reproduce by trading in financial markets.
Hence, their valuation of the plan collapses to its present value at the marketed deflator.
A proof is deferred to a later section of the course; see \citet{EkernWilson1974}, \citet{Leland1974}, and \citet{MagillQuinzii2009}.

\subsection{Strong unanimity}

We now prove a stronger statement: under common marginal rates of substitution, shareholders not only agree that firm value should be maximized, they also \emph{would not wish the firm to deviate} to any other feasible plan, even if they were given full control over that plan.

Define the set of feasible \emph{no-default} production--finance plans
\begin{equation}\label{eq:feasible-set}
  \mathcal{D} \coloneqq \bigl\{ (y, \beta) \in Y \times \R^A \colon y_s \geq \kappa_s \cdot \beta, \ \text{for all $s\in S$} \bigr\}.
\end{equation}
Given a state-price deflator $\lambda \in \R^S_{++}$, the associated equity price of plan $(y', \beta')$ is
\begin{equation}\label{eq:tilde-E}
  \widetilde{E}(\lambda, y', \beta') \coloneqq \sum_{s \in S} \lambda_s \bigl[ y'_s - \kappa_s \cdot \beta' \bigr].
\end{equation}
With this price, investor $i$'s budget set under plan $(y', \beta')$ is
\begin{equation*}
  B^i(q, \lambda, y', \beta')
  \coloneqq
  \bigl\{ (c, \theta, \alpha) \in \R_+ \times \R_+^S \times \R^A \times \R_+ \colon \text{\eqref{eq:budget-0} and \eqref{eq:budget-s} hold with } E = \widetilde{E}(\lambda, y', \beta') \bigr\}.
\end{equation*}

\begin{theorem}[Strong unanimity]\label{thm:strong-unanimity}
Suppose \textup{(C.1)--(C.4)} hold.
Fix a stock-market equilibrium
\begin{equation*}
  \bigl\{ (y, \beta), (q, E), (c^i, \theta^i, \alpha^i)_{i \in I} \bigr\}
\end{equation*}
and assume that markets are complete, so that there is a unique state-price deflator $\lambda \in \R^S_{++}$ satisfying \eqref{eq:asset-pricing}, and $\MRS^i = \lambda$ for every $i \in I$.
Assume further that the firm has chosen $(y, \beta)$ to maximize
\begin{equation*}
  V(y', \beta') \coloneqq y'_0 + q \cdot \beta' + \sum_{s \in S} \lambda_s [y'_s - \kappa_s \cdot \beta']
  = y'_0 + \widetilde{E}(\lambda, y', \beta') + q \cdot \beta'
\end{equation*}
over $(y', \beta') \in \mathcal{D}$.
Then, for every investor $i \in I$,
\begin{equation}\label{eq:strong-unanimity}
  (y, \beta) \in \argmax \bigl\{ U^i(c') \colon (c', \theta', \alpha') \in B^i(q, \lambda, y', \beta'), \ (y', \beta') \in \mathcal{D} \bigr\}.
\end{equation}
\end{theorem}

\begin{proof}
Fix $i \in I$.
Consider an alternative plan $(y', \beta') \in \mathcal{D}$ with associated equity price $\widetilde{E}' \coloneqq \widetilde{E}(\lambda, y', \beta')$ and cash flow $\delta' = (\delta'_0, \delta'_1)$, where $\delta'_0 = y'_0 + q \cdot \beta'$ and $\delta'_s = y'_s - \kappa_s \cdot \beta'$.

Let $(c', \theta', \alpha') \in B^i(q, \lambda, y', \beta')$ be any action affordable for investor $i$ under $(y', \beta')$.
We show that $U^i(c') \leq U^i(c^i)$, the equilibrium utility.

Multiply \eqref{eq:budget-s} (with primed variables) by $\lambda_s > 0$ and sum over $s \in S$:
\begin{equation*}
  \sum_{s \in S} \lambda_s c'_s
  \leq
  \sum_{s \in S} \lambda_s \omega^i_s
  + \Bigl[ \sum_{s \in S} \lambda_s \kappa_s \Bigr] \cdot \theta'
  + \Bigl[ \sum_{s \in S} \lambda_s \delta'_s \Bigr] \alpha'.
\end{equation*}
Using $q(a) = \sum_s \lambda_s \kappa_s(a)$ (no-arbitrage) and $\widetilde{E}' = \sum_s \lambda_s \delta'_s$,
\begin{equation*}
  \sum_{s \in S} \lambda_s c'_s
  \leq
  \sum_{s \in S} \lambda_s \omega^i_s
  + q \cdot \theta'
  + \widetilde{E}' \alpha'.
\end{equation*}
Add this inequality to \eqref{eq:budget-0} (with primed variables and $E = \widetilde{E}'$):
\begin{equation*}
  c'_0 + \sum_{s \in S} \lambda_s c'_s
  \leq
  \omega^i_0 + \sum_{s \in S} \lambda_s \omega^i_s
  + \xi^i \bigl[ \widetilde{E}' + \delta'_0 \bigr].
\end{equation*}
Now observe that
\begin{equation*}
  \widetilde{E}' + \delta'_0
  = \sum_s \lambda_s \delta'_s + \delta'_0
  = y'_0 + q \cdot \beta' + \sum_s \lambda_s [y'_s - \kappa_s \cdot \beta']
  = V(y', \beta').
\end{equation*}
Hence the present-value budget constraint of investor $i$ under $(y', \beta')$ is
\begin{equation}\label{eq:pv-constraint}
  c'_0 + \sum_{s \in S} \lambda_s c'_s
  \leq
  \omega^i_0 + \sum_{s \in S} \lambda_s \omega^i_s
  + \xi^i V(y', \beta').
\end{equation}

Since $(y, \beta)$ maximizes $V$ over $\mathcal{D}$, we have $V(y, \beta) \geq V(y', \beta')$, so \eqref{eq:pv-constraint} implies
\begin{equation}\label{eq:pv-constraint-y}
  c'_0 + \sum_{s \in S} \lambda_s c'_s
  \leq
  \omega^i_0 + \sum_{s \in S} \lambda_s \omega^i_s
  + \xi^i V(y, \beta).
\end{equation}
Because markets are complete, \eqref{eq:pv-constraint-y} is precisely the present-value formulation of investor $i$'s equilibrium budget constraint under $(y, \beta)$: every consumption profile $c'$ satisfying \eqref{eq:pv-constraint-y} can be implemented by a portfolio $\theta'' \in \R^A$ and a share $\alpha'' \geq 0$ in $B^i(q, E)$, where $E = \widetilde{E}(\lambda, y, \beta)$.

Consequently, $c'$ is affordable for investor $i$ in the equilibrium budget set $B^i(q, E)$.
By the optimality of $c^i$ in $B^i(q, E)$ (part (a) of Definition~\ref{def:equilibrium}),
\begin{equation*}
  U^i(c') \leq U^i(c^i).
\end{equation*}
Since $(y', \beta') \in \mathcal{D}$ and $(c', \theta', \alpha')$ were arbitrary, \eqref{eq:strong-unanimity} follows.
\end{proof}

\begin{remark}[Role of completeness]\label{rem:completeness-strong}
The proof of Theorem~\ref{thm:strong-unanimity} uses completeness at two steps: first, to guarantee a unique $\lambda$ common to all investors, so that \eqref{eq:pv-constraint-y} makes sense investor by investor with the same deflator; second, to deduce from the present-value inequality \eqref{eq:pv-constraint-y} that $c'$ is actually implementable in the original budget set (which requires spanning the space $\R^S$).
Under incomplete markets with a common marginal rate of substitution, the second step can fail: $c'$ may have low enough present value yet remain infeasible because no portfolio spans $c' - \omega^i$.
Establishing unanimity under incomplete markets typically requires the spanning assumption of Proposition~\ref{prop:spanning} or related constructs; see \citet{Geanakoplos1990} and \citet{MagillQuinzii2009}.
\end{remark}

\subsection{Separation of ownership and control}

We can now summarize the upshot of the analysis.
Under \textup{(C.1)--(C.4)} and completeness of security markets:
\begin{enumerate}
  \item a unique state-price deflator $\lambda \in \R^S_{++}$ exists, equal to every investor's marginal rate of substitution;
  \item every investor agrees that the firm should maximize value $V = y_0 + \sum_s \lambda_s y_s$;
  \item the optimal value is independent of the financial plan $\beta$, so capital structure is irrelevant in this benchmark;
  \item no investor would prefer the firm to deviate from the value-maximizing plan, even if given full control over $(y, \beta)$.
\end{enumerate}
Items (2)--(4) constitute the formal content of the \emph{separation of ownership and control}: shareholders can delegate the choice of $(y, \beta)$ to a professional manager whose mandate, ``maximize the firm's value,'' is unambiguously aligned with their individual preferences.
In this framework the corporate-finance problem reduces to computing the firm's value under different production--finance plans and selecting the maximizer---a problem that is operational even for a manager who knows nothing about the distribution of preferences across shareholders.

Subsequent chapters examine environments in which one or more of the assumptions behind this benchmark fail: moral hazard distorts value maximization; incomplete markets create disagreement about the right objective; asymmetric information between the firm and outside investors drives a wedge between internal and external valuations; and limited commitment generates credit rationing and debt overhang.
The Fisher--Drèze--Grossman--Hart benchmark remains the natural point of departure for each of these departures.

\subsection{Further reading}

The textbook treatment of corporate finance that guides the course is \citet{Tirole2006}.
The early analyses of stockholder unanimity and the objective of the firm under incomplete markets are surveyed by \citet{KelseyMilne1996}; see also \citet{TvedeCres2005} on collective decision-making by shareholders and \citet{CarcelesCoenPirani2009} for a recent revisiting of unanimity in the Arrow--Debreu tradition.
The interaction between moral hazard and stock-market equilibrium is developed in \citet{MagillQuinzii2008}.
A bargaining-based approach to the theory of the firm in general equilibrium is proposed by \citet{BritzHeringsPredtetchinski2013}, while \citet{MagillQuinzii2010} suggest a co-moment criterion that partially restores consensus under incomplete markets.
\citet{BejanBidian2012} revisit the question of ownership structure and efficiency in large economies, and \citet{DierkerDierker2010} sharpen the welfare analysis of Drèze equilibria.

%====================================================================
\chapter{Capital Structure and the Modigliani--Miller Theorems}
\label{ch:capital-structure}
%====================================================================

In the previous chapter, we established that, under complete markets and interior consumption, all shareholders unanimously agree that the firm should maximize its value $V = E + D + y_0$.
The analysis treated the production plan $y$ and the financial plan $\beta$ symmetrically, but left unaddressed a natural question.
Once the firm's production is fixed, does the way in which the firm splits its value between equity $E$ and debt $D$ matter?
The seminal answer, due to \citet{ModiglianiMiller1958}, is \emph{no}: under the assumptions of the previous chapter---frictionless markets, no taxes, unlimited liability---the firm's value is independent of its capital structure.
The financial plan $\beta$ is \emph{irrelevant}.

This chapter develops the irrelevance result in two forms, sharpens it via the stronger statement of \citet{Stiglitz1969, Stiglitz1974}, and examines the conditions under which the result fails.
When some of the Modigliani--Miller assumptions are relaxed, an optimal debt-to-equity ratio may re-emerge.
We study in detail the role of corporate and personal taxes, the interplay with limited liability and bankruptcy costs, and the irrelevance of the payout policy in a multi-period deterministic setting.
A concluding section points to the broader literature---signalling, agency costs, control rights---that departs from the Modigliani--Miller benchmark by introducing informational or contractual frictions.

%====================================================================
\section{The Economy with Multiple Firms}
\label{sec:multi-firm}
%====================================================================

\subsection{Firms and their production--finance plans}

To state the Modigliani--Miller theorems we extend the model of Chapter~\ref{ch:objective} to a finite set $J$ of firms.
Each firm $j \in J$ is endowed with a production set $Y(j) \subset \R \times \R^S$ and chooses a production plan $y(j) = (y_0(j), y_1(j)) \in Y(j)$ and a financial plan $\beta(j) \in \R^A$.
At $t=0$ and at each state $s \in S$, firm $j$ generates the cash flow
\begin{equation}\label{eq:cash-flow-multi}
  \delta_0(j) = y_0(j) + q \cdot \beta(j)
  \et
  \delta_s(j) = y_s(j) - \kappa_s \cdot \beta(j).
\end{equation}
The initial ownership of firm $j$ is described by a vector $(\xi^i(j))_{i \in I} \in \R_+^I$ with $\sum_{i \in I} \xi^i(j) = 1$.
The equity of firm $j$ trades at $t=0$ at a linear price $E(j)$.
The \emph{value} of firm $j$ is defined as
\begin{equation}\label{eq:firm-value-multi}
  V(j) = E(j) + D(j) + y_0(j),
  \quad
  D(j) \coloneqq q \cdot \beta(j).
\end{equation}

\subsection{Investors' budget sets}

Each investor $i \in I$ takes as given all firms' production--finance plans $(y(j), \beta(j))_{j \in J}$, the security price vector $q$, and the equity price vector $E = (E(j))_{j \in J}$.
He chooses a consumption plan $(c_0, c_1) \in \R_+ \times \R_+^S$, a security portfolio $\theta \in \R^A$, and a vector of post-trade share holdings $\alpha = (\alpha(j))_{j \in J} \in \R_+^J$ satisfying the budget constraints
\begin{align}
  c_0 + q \cdot \theta + \sum_{j \in J} E(j) \alpha(j)
  &\leq
  \omega^i_0 + \sum_{j \in J} \bigl[E(j) + \delta_0(j)\bigr] \xi^i(j),
  \label{eq:budget-0-multi}\\[2pt]
  c_s
  &\leq
  \omega^i_s + \kappa_s \cdot \theta + \sum_{j \in J} \delta_s(j) \alpha(j),
  \quad \text{for all $s\in S$}.
  \label{eq:budget-s-multi}
\end{align}
The set of actions $(c, \theta, \alpha)$ satisfying \eqref{eq:budget-0-multi}--\eqref{eq:budget-s-multi} is denoted by $B^i(q, E)$.

\subsection{Stock-market equilibrium}

\begin{definition}[Stock-market equilibrium with many firms]\label{def:eq-multi}
A stock-market equilibrium is a list
\begin{equation*}
  \bigl\{ (y(j), \beta(j))_{j \in J}, \; (q, E), \; (c^i, \theta^i, \alpha^i)_{i \in I} \bigr\}
\end{equation*}
such that:
\begin{enumerate}
  \item[(a)] for every $i \in I$,
  \begin{equation*}
    (c^i, \theta^i, \alpha^i) \in \argmax \bigl\{ U^i(c) \colon (c, \theta, \alpha) \in B^i(q, E) \bigr\};
  \end{equation*}
  \item[(b)] security and stock markets clear:
  \begin{equation*}
    \forall a \in A, \quad \sum_{j \in J} \beta(j, a) = \sum_{i \in I} \theta^i(a)
    \et
    \forall j \in J, \quad \sum_{i \in I} \alpha^i(j) = 1;
  \end{equation*}
  \item[(c)] the commodity market clears:
  \begin{equation*}
    \sum_{i \in I} c^i = \sum_{i \in I} \omega^i + \sum_{j \in J} y(j).
  \end{equation*}
\end{enumerate}
\end{definition}

The arguments of Chapter~\ref{ch:objective} extend to the multi-firm setting with no essential change.
In particular, under Assumption~\ref{ass:C} (conditions C.1--C.4), for every investor $i \in I$, every asset $a \in A$, and every firm $j$ with $\alpha^i(j) > 0$,
\begin{equation}\label{eq:price-multi}
  q(a) = \sum_{s \in S} \MRS^i_s \kappa_s(a)
  \et
  E(j) = \sum_{s \in S} \MRS^i_s \delta_s(j).
\end{equation}

\subsection{Leverage and capital structure}

Before turning to the irrelevance theorems, fix a single firm $j$ and consider the simplest case.
Assume there is a riskless bond $a_0$ with $\kappa_s(a_0) = 1$ for every $s \in S$, and let $r_0$ denote the risk-free interest rate, defined by $q(a_0) = 1/(1+r_0)$.
Suppose further that firm $j$ can only issue this bond, so that $\beta(j) = Z(j) \1_{a_0}$ for some $Z(j) \geq 0$.
The firm is said to be \emph{levered} if $D(j) = Z(j)/(1+r_0) > 0$, and \emph{unlevered} if $Z(j) = 0$.
The \emph{debt-to-equity ratio} $D(j)/E(j)$ is a summary measure of capital structure: a specific choice of $\beta(j)$ induces a specific ratio $D(j)/E(j)$, and any change in $\beta(j)$ changes both the numerator and (through equilibrium effects) the denominator.

\subsection{The pre-1958 conventional view}

Before the work of Modigliani and Miller, the conventional wisdom held that there should exist an optimal level of debt.
The reasoning went as follows.
Start with an unlevered firm whose dividend stream $\delta_1 = (y_s)_{s \in S}$ is risky.
Issuing debt enlarges the set of securities available to investors: in addition to the risky claim $(y_s - Z)_{s \in S}$ on the levered firm, investors can now buy the riskless claim $Z \1_S$ issued by the firm.
A modest level of debt may therefore attract additional investors, increase the firm's value, and keep the risk of bankruptcy negligible.
On the other hand, for high levels of debt the probability of $y_s < Z$ in some states becomes non-trivial, bankruptcy costs intervene, and the benefit of leverage disappears.
The argument seemed to imply an interior optimum: an optimal debt-to-equity ratio exists, balancing the gain from leverage against the cost of financial distress.

The following two sections show that this conventional reasoning is, in the frictionless benchmark of Chapter~\ref{ch:objective}, incorrect.

%====================================================================
\section{The Modigliani--Miller Irrelevance Theorems}
\label{sec:mm}
%====================================================================

\subsection{Weak irrelevance}

The first version of the Modigliani--Miller theorem compares two firms with the same production plan but possibly different financial plans.

\begin{theorem}[Modigliani--Miller, weak form]\label{thm:mm-weak}
Consider a stock-market equilibrium
\begin{equation*}
  \bigl\{ (y(j), \beta(j))_{j \in J}, \; (q, E), \; (c^i, \theta^i, \alpha^i)_{i \in I} \bigr\}
\end{equation*}
under Assumption~\ref{ass:C}.
Assume that two firms $j_1$ and $j_2$ share the same production plan, $y(j_1) = y(j_2)$, but may have different financial plans $\beta(j_1) \neq \beta(j_2)$.
If there exists an investor $i$ with $\alpha^i(j_1) > 0$ and $\alpha^i(j_2) > 0$, then the two firms have the same value:
\begin{equation}\label{eq:mm-weak}
  V(j_1) = V(j_2).
\end{equation}
\end{theorem}

\begin{proof}
Let $i$ be an investor with $\alpha^i(j_1) > 0$ and $\alpha^i(j_2) > 0$.
By Proposition~\ref{prop:personalized}, his marginal rates of substitution $(\MRS^i_s)_{s \in S}$ form a vector of strictly positive state-price deflators that simultaneously prices all securities and the equity of both firms:
\begin{equation*}
  q(a) = \sum_{s \in S} \MRS^i_s \kappa_s(a) \quad \forall a \in A,
  \et
  E(j) = \sum_{s \in S} \MRS^i_s \delta_s(j) \quad \text{for } j \in \{ j_1, j_2 \}.
\end{equation*}
Using $\delta_s(j) = y_s(j) - \kappa_s \cdot \beta(j)$ and $D(j) = q \cdot \beta(j) = \sum_s \MRS^i_s (\kappa_s \cdot \beta(j))$,
\begin{eqnarray*}
  V(j)
   & = & E(j) + D(j) + y_0(j) \\
   & = & \sum_{s \in S} \MRS^i_s \bigl[y_s(j) - \kappa_s \cdot \beta(j)\bigr] 
  + \sum_{s \in S} \MRS^i_s \kappa_s \cdot \beta(j)
  + y_0(j) \\
   & = & y_0(j) + \sum_{s \in S} \MRS^i_s y_s(j).
\end{eqnarray*}
Since $y(j_1) = y(j_2)$, both firms share the same value.
\end{proof}

\begin{remark}[What if no investor holds both firms?]\label{rem:mm-no-common}
The hypothesis of Theorem~\ref{thm:mm-weak} requires a common shareholder $i$ so that a \emph{single} deflator $\MRS^i$ prices both firms' equity.
Without a common shareholder, Proposition~\ref{prop:personalized} produces, for each firm $j$, a deflator drawn from the shareholders of $j$, and different firms may be priced by different deflators.
The conclusion $V(j_1) = V(j_2)$ then requires an independent guarantee that all deflators agree, typically \emph{completeness of security markets}: when $\Range \kappa = \R^S$, the no-arbitrage equation $q = \sum_s \lambda_s \kappa_s$ pins $\lambda$ down uniquely, so all personalized deflators coincide with a common $\lambda$ and the proof above goes through for every $j_1, j_2 \in J$ with $y(j_1) = y(j_2)$.
\end{remark}

\subsection{Strong irrelevance}

The weak form compares two firms in a single equilibrium.
\citet{Stiglitz1969, Stiglitz1974} strengthened the irrelevance result by comparing two equilibria: for any alternative assignment of financial plans across firms, there exists another equilibrium in which the firms' values are unchanged.
Capital structure is irrelevant not only at the cross-section of a given equilibrium but also for the equilibrium map itself.

\begin{theorem}[Stiglitz strong irrelevance]\label{thm:stiglitz}
Suppose Assumption~\ref{ass:C} holds and let
\begin{equation*}
  \bigl\{ (y(j), \beta(j))_{j \in J}, \; (q, E), \; (c^i, \theta^i, \alpha^i)_{i \in I} \bigr\}
\end{equation*}
be a stock-market equilibrium.
For any alternative family of financial plans $(\widehat{\beta}(j))_{j \in J} \in (\R^A)^J$, the list
\begin{equation*}
  \bigl\{ (y(j), \widehat{\beta}(j))_{j \in J}, \; (q, \widehat{E}), \; (c^i, \widehat{\theta}^i, \alpha^i)_{i \in I} \bigr\}
\end{equation*}
is also a stock-market equilibrium, where
\begin{equation}\label{eq:stiglitz-prices}
  \widehat{E}(j) = E(j) + q \cdot \bigl[ \beta(j) - \widehat{\beta}(j) \bigr]
\end{equation}
and
\begin{equation}\label{eq:stiglitz-portfolios}
  \widehat{\theta}^i = \theta^i + \sum_{j \in J} \alpha^i(j) \bigl[ \widehat{\beta}(j) - \beta(j) \bigr].
\end{equation}
In particular, each firm's value is unchanged: $\widehat{V}(j) = V(j)$ for every $j \in J$.
\end{theorem}

\begin{proof}
We verify the three conditions of Definition~\ref{def:eq-multi} in turn.

\emph{Step 1: Budget sets and consumption.}
We show that the candidate consumption $c^i$ remains affordable for investor $i$ under the new prices and the new portfolio $\widehat{\theta}^i$, and that the left-hand sides of the budget constraints \eqref{eq:budget-0-multi}--\eqref{eq:budget-s-multi} are equal, both at $t=0$ and at each state $s$.

At $t=0$, the new left-hand side is
\begin{equation*}
  c^i_0 + q \cdot \widehat{\theta}^i + \sum_{j \in J} \widehat{E}(j) \alpha^i(j).
\end{equation*}
Substituting \eqref{eq:stiglitz-portfolios}--\eqref{eq:stiglitz-prices},
\begin{equation*}
  q \cdot \widehat{\theta}^i
  = q \cdot \theta^i + \sum_{j \in J} \alpha^i(j) \, q \cdot \bigl[ \widehat{\beta}(j) - \beta(j) \bigr],
\end{equation*}
and
\begin{equation*}
  \sum_{j \in J} \widehat{E}(j) \alpha^i(j)
  = \sum_{j \in J} \alpha^i(j) \bigl[ E(j) + q \cdot (\beta(j) - \widehat{\beta}(j)) \bigr]
  = \sum_{j \in J} \alpha^i(j) E(j) + \sum_{j \in J} \alpha^i(j) \, q \cdot \bigl[ \beta(j) - \widehat{\beta}(j) \bigr].
\end{equation*}
The cross-terms involving $q \cdot [\widehat{\beta}(j) - \beta(j)]$ and $q \cdot [\beta(j) - \widehat{\beta}(j)]$ cancel, leaving
\begin{equation*}
  c^i_0 + q \cdot \widehat{\theta}^i + \sum_{j \in J} \widehat{E}(j) \alpha^i(j)
  = c^i_0 + q \cdot \theta^i + \sum_{j \in J} E(j) \alpha^i(j).
\end{equation*}
The right-hand side of \eqref{eq:budget-0-multi} transforms analogously.
Writing $\widehat{\delta}_0(j) = y_0(j) + q \cdot \widehat{\beta}(j)$ and using \eqref{eq:stiglitz-prices},
\begin{equation*}
  \widehat{E}(j) + \widehat{\delta}_0(j)
  = E(j) + q \cdot \bigl[ \beta(j) - \widehat{\beta}(j) \bigr] + y_0(j) + q \cdot \widehat{\beta}(j)
  = E(j) + y_0(j) + q \cdot \beta(j)
  = E(j) + \delta_0(j).
\end{equation*}
Hence the two right-hand sides agree as well, and $(c^i, \widehat{\theta}^i, \alpha^i)$ satisfies \eqref{eq:budget-0-multi} with equality whenever $(c^i, \theta^i, \alpha^i)$ does.

At each state $s \in S$, the new right-hand side of \eqref{eq:budget-s-multi} is
\begin{equation*}
  \omega^i_s + \kappa_s \cdot \widehat{\theta}^i + \sum_{j \in J} \widehat{\delta}_s(j) \alpha^i(j),
\end{equation*}
where $\widehat{\delta}_s(j) = y_s(j) - \kappa_s \cdot \widehat{\beta}(j)$.
Expanding $\kappa_s \cdot \widehat{\theta}^i$ using \eqref{eq:stiglitz-portfolios} and collecting terms,
\begin{equation*}
  \kappa_s \cdot \widehat{\theta}^i + \sum_{j \in J} \widehat{\delta}_s(j) \alpha^i(j)
  = \kappa_s \cdot \theta^i + \sum_{j \in J} \alpha^i(j) \, \kappa_s \cdot \bigl[ \widehat{\beta}(j) - \beta(j) \bigr]
    + \sum_{j \in J} \alpha^i(j) \bigl[ y_s(j) - \kappa_s \cdot \widehat{\beta}(j) \bigr].
\end{equation*}
The terms involving $\kappa_s \cdot \widehat{\beta}(j)$ cancel, yielding
\begin{equation*}
  \kappa_s \cdot \theta^i + \sum_{j \in J} \alpha^i(j) \bigl[ y_s(j) - \kappa_s \cdot \beta(j) \bigr]
  = \kappa_s \cdot \theta^i + \sum_{j \in J} \delta_s(j) \alpha^i(j),
\end{equation*}
which is the original right-hand side of \eqref{eq:budget-s-multi}.
Consequently, $c^i$ remains affordable, and since $U^i$ is evaluated at consumption alone, the optimality of $c^i$ in the new budget set follows immediately from its optimality in the original one.

\emph{Step 2: Market clearing.}
Summing \eqref{eq:stiglitz-portfolios} over $i \in I$ and using $\sum_i \alpha^i(j) = 1$ (share market clearing in the original equilibrium),
\begin{equation*}
  \sum_{i \in I} \widehat{\theta}^i
  = \sum_{i \in I} \theta^i + \sum_{j \in J} \bigl[ \widehat{\beta}(j) - \beta(j) \bigr] \sum_{i \in I} \alpha^i(j)
  = \sum_{j \in J} \beta(j) + \sum_{j \in J} \bigl[ \widehat{\beta}(j) - \beta(j) \bigr]
  = \sum_{j \in J} \widehat{\beta}(j),
\end{equation*}
where the second equality uses security market clearing in the original equilibrium.
Share market clearing for $\alpha^i$ is preserved since $\alpha^i$ has not changed.
Commodity market clearing is also preserved because consumption plans and aggregate output $\sum_j y(j)$ are unchanged.

\emph{Step 3: Equal firm values.}
From \eqref{eq:stiglitz-prices},
\begin{eqnarray*}
  \widehat{V}(j) & = & \widehat{E}(j) + q \cdot \widehat{\beta}(j) + y_0(j) \\
  & = & E(j) + q \cdot \bigl[ \beta(j) - \widehat{\beta}(j) \bigr] + q \cdot \widehat{\beta}(j) + y_0(j) \\
  & = & E(j) + q \cdot \beta(j) + y_0(j) \\
  & = & V(j),
\end{eqnarray*}
completing the proof.
\end{proof}

The Stiglitz theorem is a remarkable statement about the structure of the equilibrium set: the value of a firm is pinned down by its production plan and by the equilibrium state-price deflators, while the financial plan merely rebalances who holds debt and who holds equity across investors.
Individual investors adjust their portfolios through \eqref{eq:stiglitz-portfolios}, exactly ``undoing'' the firm's change in financial policy, in the spirit of what is known as \emph{homemade leverage}.

\subsection{Conditions for the validity of the irrelevance results}

The proofs above rely on several assumptions which, collectively, define what is meant by the \emph{Modigliani--Miller benchmark}.
It is instructive to list them explicitly, since each one identifies a channel through which capital structure may matter once relaxed.
\begin{enumerate}
  \item \emph{Frictionless capital markets.}
  No transaction costs, no institutional restrictions on security trades.
  The proof of Theorem~\ref{thm:stiglitz} hinges on investors' ability to freely adjust $\widehat{\theta}^i$ according to \eqref{eq:stiglitz-portfolios}.
  If trading costs or short-sale constraints bind, this portfolio adjustment may not be feasible.

  \item \emph{No taxes.}
  In the cash-flow identities \eqref{eq:cash-flow-multi}, the firm and its investors receive the full amount $\kappa_s \cdot \theta$ or $\delta_s(j) \alpha(j)$.
  When a corporate tax or a personal tax intermediates these payments, the firm's financial policy affects the combined tax liability, as we show in Section~\ref{sec:tax}.

  \item \emph{Symmetric access to securities.}
  Investors and firms trade the same set $A$ of securities, at the same prices.
  If investors can only buy bonds issued by firms, or if investors face different menus than firms, the homemade leverage argument may break.

  \item \emph{Unlimited liability.}
  The cash flow $\delta_s(j)$ may be negative and shareholders are required to deliver $|\delta_s(j)| \alpha^i(j)$.
  Under limited liability, shareholders pay nothing when the firm is bankrupt; the payoff profile of equity and debt becomes nonlinear in $\beta(j)$, breaking the linear-pricing argument underlying \eqref{eq:mm-weak}.

  \item \emph{No information asymmetry.}
  Investors do not infer information about the states of nature from the firm's financial choices.
  If the choice of $\beta(j)$ signals information available to the manager but not to investors, the market reacts to $\beta(j)$, and the firm's value is no longer independent of the financing decision.

  \item \emph{Linearity.}
  The change in the return on equity induced by a modification of the firm's capital structure is a linear combination of the returns on the assets traded by the firm.
  When derivative securities are written on the firm's equity (options, convertibles), the derivatives' payoffs depend nonlinearly on $\beta(j)$, and a change in capital structure alters their market value.
  \citet{Gottardi1995} provides a thorough analysis of the role of this last condition in incomplete markets.
\end{enumerate}
Each of the subsequent chapters in this book can be viewed as an investigation of a specific departure from this list.
Chapter~\ref{ch:objective} and the present chapter established what happens when all assumptions hold; Section~\ref{sec:tax} below departs from~(2); Section~\ref{sec:limited-liability} departs from~(4); later chapters depart from~(5) and~(6).

%====================================================================
\section{Taxation}
\label{sec:tax}
%====================================================================

We now relax the second assumption of the Modigliani--Miller benchmark by introducing a tax system.
The goal is modest: we describe the distortion imposed by taxes on the firm's valuation, leading to the classical ``gains from leverage'' formula of \citet{ModiglianiMiller1958, Miller1977}.
We do not model the endogenous determination of tax rates or the use of tax revenue, which is treated as an exogenous transfer to a tax authority.

\subsection{The tax system}

Taxes are paid at $t=1$ only.
There are three flat tax rates:
\begin{itemize}
  \item $\tau_c$ --- the corporate tax rate, applied to the firm's income $\delta_1(j) = y_1(j) - \kappa_1 \cdot \beta(j)$;
  \item $\tau_s$ --- the personal tax rate on equity income;
  \item $\tau_b$ --- the personal tax rate on bond income.
\end{itemize}
We assume that all investors face the same rates, and that rates are constant and independent of the level of income.
We also assume that every security $a \in A$ yields a non-negative state-contingent payoff, $\kappa_s(a) \geq 0$ for every $s$, so that personal taxation on bond income is meaningful; and that every firm's production--finance plan lies in the no-bankruptcy set
\begin{equation*}
  \mathcal{D}(j) \coloneqq \bigl\{ (y, \beta) \in Y(j) \times \R^A \colon y_s \geq \kappa_s \cdot \beta, \ \text{for all $s\in S$} \bigr\},
\end{equation*}
so that $\delta_s(j) \geq 0$ and the before-personal-tax shareholder payoff is well-defined.

One unit of bond $a$ held by investor $i$ delivers, after personal taxation on bond income, $(1-\tau_b) \kappa_s(a)$ in state $s$.
One unit of equity in firm $j$ delivers, after corporate taxation and then personal taxation on equity income, $(1-\tau_s)(1-\tau_c) \delta_s(j)$.

\subsection{Investors' budget sets with taxes}

Each investor $i$ takes as given all firms' production--finance plans $(y(j), \beta(j))_{j \in J}$, the security price vector $q$, and the equity price vector $E$.
He chooses a consumption plan $(c_0, c_1) \in \R_+ \times \R_+^S$, a security portfolio $\theta \in \R^A$, and a vector of post-trade share holdings $\alpha \in \R_+^J$ satisfying
\begin{align}
  c_0 + q \cdot \theta + \sum_{j \in J} E(j) \alpha(j)
  &\leq
  \omega^i_0 + \sum_{j \in J} \bigl[ E(j) + \delta_0(j) \bigr] \xi^i(j),
  \label{eq:budget-0-tax}\\[2pt]
  c_s
  &\leq
  \omega^i_s + (1 - \tau_b) \kappa_s \cdot \theta + (1-\tau_s)(1-\tau_c) \sum_{j \in J} \delta_s(j) \alpha(j),
  \quad \text{for all $s\in S$}.
  \label{eq:budget-s-tax}
\end{align}
The set of actions $(c, \theta, \alpha)$ satisfying \eqref{eq:budget-0-tax}--\eqref{eq:budget-s-tax} is denoted by $B^i(q, E, \tau)$, where $\tau = (\tau_c, \tau_s, \tau_b)$.

The definition of equilibrium parallels Definition~\ref{def:eq-multi}, with two modifications.
First, the budget set uses \eqref{eq:budget-s-tax} in place of \eqref{eq:budget-s-multi}.
Second, commodity-market clearing must accommodate the consumption of the tax authority:
\begin{equation*}
  c^p + \sum_{i \in I} c^i = \sum_{i \in I} \omega^i + \sum_{j \in J} y(j),
\end{equation*}
where $c^p_0 = 0$ and, at each state $s$,
\begin{equation*}
  c^p_s = \sum_{j \in J} \Bigl\{ [\tau_c + (1-\tau_c)\tau_s] \, y_s(j) + [\tau_b - \tau_c - (1-\tau_c)\tau_s] \, \kappa_s \cdot \beta(j) \Bigr\}.
\end{equation*}
The quantity $c^p_s$ collects all tax revenue paid in state $s$.

\subsection{Asset pricing under taxes}

Proceeding as in Chapter~\ref{ch:objective}, the first-order conditions of investor $i$'s problem yield, under \textup{(C.1)--(C.4)},
\begin{equation}\label{eq:price-tax}
  q(a) = (1-\tau_b) \sum_{s \in S} \MRS^i_s \kappa_s(a)
  \et
  E(j) = (1-\tau_c)(1-\tau_s) \sum_{s \in S} \MRS^i_s \delta_s(j),
\end{equation}
for every $a \in A$, every $j$ with $\alpha^i(j) > 0$, and every investor $i$ with interior consumption.
The first identity expresses the security price as the tax-adjusted expected payoff; the second expresses the equity price as the tax-adjusted present value of the firm's dividend stream.

\subsection{Gains from corporate leverage}

The presence of taxes breaks the irrelevance of Theorem~\ref{thm:mm-weak}: two firms with identical production plans and different financial plans need no longer share the same value.
The following proposition quantifies the wedge.

\begin{proposition}[Gains from leverage]\label{prop:gains-leverage}
Under \textup{(C.1)--(C.4)}, suppose there exist two firms $j_u$ and $j_\ell$ with
\begin{itemize}
  \item identical production plans, $y(j_u) = y(j_\ell) = y$;
  \item $j_u$ unlevered, $\beta(j_u) = 0$;
  \item $j_\ell$ levered and solvent, $\beta(j_\ell) \neq 0$ and $y_s \geq \kappa_s \cdot \beta(j_\ell)$ for every $s$.
\end{itemize}
Assume there exists an investor $i$ with $\alpha^i(j_u) > 0$ and $\alpha^i(j_\ell) > 0$.
Then the firms' values satisfy
\begin{equation}\label{eq:gains-leverage}
  V(j_\ell) = V(j_u) + \biggl[\, 1 - \frac{(1-\tau_c)(1-\tau_s)}{1-\tau_b} \,\biggr] D(j_\ell).
\end{equation}
\end{proposition}

\begin{proof}
By \eqref{eq:price-tax} applied to investor $i$,
\begin{equation*}
  E(j_u) = (1-\tau_c)(1-\tau_s) \sum_{s \in S} \MRS^i_s y_s,
  \quad
  E(j_\ell) = (1-\tau_c)(1-\tau_s) \sum_{s \in S} \MRS^i_s \bigl[ y_s - \kappa_s \cdot \beta(j_\ell) \bigr],
\end{equation*}
and therefore
\begin{equation*}
  E(j_\ell) - E(j_u) = - (1-\tau_c)(1-\tau_s) \sum_{s \in S} \MRS^i_s \kappa_s \cdot \beta(j_\ell).
\end{equation*}
The first identity in \eqref{eq:price-tax} gives $q(a) = (1-\tau_b) \sum_s \MRS^i_s \kappa_s(a)$ for every $a$, so
\begin{equation*}
  D(j_\ell) = q \cdot \beta(j_\ell) = (1-\tau_b) \sum_{s \in S} \MRS^i_s \kappa_s \cdot \beta(j_\ell),
\end{equation*}
which gives $\sum_s \MRS^i_s \kappa_s \cdot \beta(j_\ell) = D(j_\ell)/(1-\tau_b)$.
Substituting,
\begin{equation*}
  E(j_\ell) - E(j_u) = - \frac{(1-\tau_c)(1-\tau_s)}{1-\tau_b} D(j_\ell).
\end{equation*}
Since $V(j) = E(j) + D(j) + y_0(j)$, $D(j_u) = 0$, and $y_0(j_u) = y_0(j_\ell)$,
\begin{equation*}
  V(j_\ell) - V(j_u)
  = \bigl[ E(j_\ell) - E(j_u) \bigr] + D(j_\ell)
  = D(j_\ell) \biggl[\, 1 - \frac{(1-\tau_c)(1-\tau_s)}{1-\tau_b} \,\biggr],
\end{equation*}
as claimed.
\end{proof}

\subsection{Interpreting the wedge}

The coefficient $1 - (1-\tau_c)(1-\tau_s)/(1-\tau_b)$ in \eqref{eq:gains-leverage} has the natural interpretation of a \emph{net} tax advantage of debt.
Three polar cases illuminate its content.

\emph{Pure corporate tax.}
If personal tax rates coincide, $\tau_s = \tau_b$, the wedge reduces to
\begin{equation*}
  V(j_\ell) = V(j_u) + \tau_c D(j_\ell).
\end{equation*}
Every marginal dollar of debt increases the firm's value by $\tau_c$, and firms have an unbounded incentive to raise capital through bond issues, up to the bankruptcy limit imposed by the assumption $y_s \geq \kappa_s \cdot \beta(j_\ell)$.
This prediction is at odds with the observed capital structures of mature firms.

\emph{Personal tax favouring debt.}
If $(1-\tau_c)(1-\tau_s) < 1 - \tau_b$, the coefficient in \eqref{eq:gains-leverage} is positive and, as in the previous case, firms have a strong incentive to lever up to the bankruptcy boundary.

\emph{Personal tax favouring equity.}
If $(1-\tau_c)(1-\tau_s) > 1 - \tau_b$, the coefficient is negative, and the argument flips: firms have an incentive to finance their capital exclusively through equity.

Neither of the two latter predictions fits empirical evidence, which points to interior, firm-specific capital structures.
\citet{Miller1977} proposes a market-based resolution: in equilibrium, the \emph{aggregate} supply of debt adjusts so that, at the margin, investors are indifferent between holding debt and equity.
Empirical work has largely confirmed that taxes matter for corporate capital structure, but has found the quantitative effect to be more subtle than the simple formula \eqref{eq:gains-leverage} would suggest.
See \citet{MackieMason1990} and \citet{DesaiFoleyHines2004} for representative empirical studies.

%====================================================================
\section{Limited Liability and Bankruptcy}
\label{sec:limited-liability}
%====================================================================

The fourth condition underpinning Modigliani--Miller is \emph{unlimited liability}: shareholders may be required to transfer resources back to the firm when the residual dividend $\delta_s(j)$ is negative.
In practice, corporate shareholders enjoy \emph{limited liability}: they lose at most their initial equity investment and are shielded from the firm's obligations in states where the firm cannot meet its debt.

\subsection{Equity and debt under limited liability}

Consider, for expositional simplicity, a single firm $j$ whose only financing instrument is the riskless bond $a_0$, $\kappa_s(a_0) = 1$ for every $s$.
The firm issues a face-value debt $Z \geq 0$, and the production plan is $y = (y_0, y_1)$.
In state $s$, the firm's pre-distribution surplus is $\delta_s = y_s - Z$, which may be positive or negative.

Under limited liability, new shareholders at $t=1$ are not called upon to deliver $|\delta_s|$ to bondholders when $\delta_s < 0$.
Instead, the firm is declared \emph{bankrupt}, and the entire output $y_s$ is distributed pro rata to bondholders, while shareholders receive nothing.
In state $s$, one unit of the firm's bond delivers the state-contingent amount
\begin{equation}\label{eq:levered-bond}
  b_s = \min\{y_s, Z\},
\end{equation}
and one share of the firm delivers
\begin{equation}\label{eq:levered-equity}
  \widehat{\delta}_s = \max\{\delta_s, 0\} = \max\{y_s - Z, 0\}.
\end{equation}
Together, $b_s + \widehat{\delta}_s = y_s$ for every $s$: limited liability redistributes the state-$s$ output between bondholders and shareholders, but leaves the total unchanged.

A direct consequence is that the firm's value remains
\begin{equation*}
  V = E(j) + D(j) + y_0 = \sum_{s \in S} \lambda_s b_s + \sum_{s \in S} \lambda_s \widehat{\delta}_s + y_0 = y_0 + \sum_{s \in S} \lambda_s y_s,
\end{equation*}
provided that one can price the firm's securities using a common state-price deflator $\lambda$.
In particular, absent taxes and bankruptcy costs, limited liability alone does not break the Modigliani--Miller conclusion.
The \emph{composition} of the total firm value between debt and equity changes nonlinearly with $Z$, but the total is unchanged.

\subsection{Available markets and the relevance of financial policy}

The substantive effect of limited liability depends on what markets investors have access to.

\emph{Case 1: Full market access.}
If every investor can purchase or sell any bond issued by any firm, then homemade leverage of the type used in the proof of Theorem~\ref{thm:stiglitz} still works, and the firm's financial policy is irrelevant.
In this case, limited liability reshuffles the way the total firm value is split between bondholders and shareholders without changing the value itself.

\emph{Case 2: Restricted bond market.}
If investors can purchase only a ``pooled bond'' based on the firms' obligations---rather than individual firm-specific bonds---the set of marketed cash flows depends on the financial plans $(\beta(j))_{j \in J}$, and the irrelevance argument can fail.

\emph{Case 3: No secondary market in bonds.}
If investors can only purchase bonds issued directly by firms, and cannot trade them in a secondary market, the homemade leverage portfolio \eqref{eq:stiglitz-portfolios} is not implementable, and the financial policy becomes relevant.

In addition, when bankruptcy is costly---legal fees, asset fire sales, covenant violations---the total payoff $b_s + \widehat{\delta}_s$ falls short of $y_s$ in bankruptcy states, and $V$ does depend on the probability of bankruptcy and hence on $\beta$.
Empirical estimates of direct and indirect bankruptcy costs are discussed by \citet{Weiss1990}, \citet{OplerTitman1994}, and \citet{Purnanandam2008}.
Classical analyses of debt covenants and the contractual architecture around financial distress go back to \citet{SmithWarner1979}.

%====================================================================
\section{Irrelevance of the Payout Policy}
\label{sec:payout}
%====================================================================

We conclude the chapter with a result on the irrelevance of the firm's \emph{payout policy}---the division of cash flow between dividends, share repurchases, and new share issuance.
The statement holds in a simple multi-period deterministic model.
The argument is structurally close to Theorem~\ref{thm:stiglitz}: so long as investors can freely rebalance their portfolios and the firm operates under no-arbitrage prices, the firm's decision on how to package its dividend stream does not affect its total value.

\subsection{A multi-period deterministic model}

Time is indexed by $t \in \{0, 1, \ldots, T\}$.
There is no uncertainty: at each date $t$, the state is fully known.
One divisible commodity is available for trade.
There is a finite set $A$ of \emph{long-lived} financial assets; each asset $a \in A$ pays a coupon $\kappa_t(a)$ at date $t \geq 1$ and trades at a linear price $q_t(a)$ at each date $t \leq T - 1$.
We assume $q_T(a) = 0$ and $\kappa_0(a) = 0$ by convention: an asset delivers no further payoff after date $T$, and no coupon at $t=0$.

There is one firm $j$, which we denote unambiguously without index.
The firm chooses an investment--production plan $(x, y) \in Z$, where
\begin{equation*}
  x = (x_0, x_1, \ldots, x_{T-1}) \in \R_+^T
  \et
  y = (y_1, \ldots, y_T) \in \R_+^T.
\end{equation*}
At date $t$, the firm invests $x_t \geq 0$ as inputs; at date $t+1$, it produces $y_{t+1} \geq 0$ as outputs.
The technological restrictions are encoded in the set $Z$.
For notational convenience, we let $y_0 = 0$ and $x_T = 0$.

\subsection{Payout policies and the cash-flow identity}

The firm further chooses a \emph{financial plan} $\beta = (\beta_0, \beta_1, \ldots, \beta_{T-1})$, with $\beta_t \in \R^A$, and a \emph{dividend policy} $n = (n_0, n_1, \ldots, n_{T-1})$, where $n_t$ is the number of shares outstanding after date-$t$ dividend payments and before date-$(t+1)$ dividend payments.
By convention, $n_{-1} = 1$ (the firm starts with one share) and $\beta_{-1} = 0$ (no initial debt).

Let $p_t$ denote the \emph{ex-dividend} share price at date $t$, quoted after the firm announces its payout policy.
If $n_t > n_{t-1}$, the firm issues new shares and raises $p_t (n_t - n_{t-1})$ units of the commodity; if $n_t < n_{t-1}$, the firm repurchases shares and spends $-p_t (n_t - n_{t-1}) > 0$.
Let $\delta_{t+1}$ denote the per-share dividend paid at date $t+1$.
The firm's cash-flow constraint at date $t+1$ requires that the total dividend outflow equal the firm's operating cash surplus plus the net cash raised from equity and debt issuance:
\begin{equation}\label{eq:cashflow-identity}
  n_t \delta_{t+1}
  = y_{t+1} - x_{t+1} - \kappa_{t+1} \cdot \beta_t
    + p_{t+1} \bigl[ n_{t+1} - n_t \bigr]
    + q_{t+1} \cdot \bigl[ \beta_{t+1} - \beta_t \bigr].
\end{equation}
The firm's production and payout policy is the tuple $(x, y, \beta, n)$.

\subsection{Investors' budget sets and equilibrium}

Each investor $i \in I$ takes as given the firm's policy $(x, y, \beta, n)$ and the security and share prices $(q, p) = (q_t, p_t)_{t \in \{0, \ldots, T-1\}}$, and chooses a consumption plan $c = (c_t)_{t} \in \R_+^{T+1}$, a security portfolio $\theta = (\theta_t)_t \in \R^{A \times T}$, and a stockholding vector $\alpha = (\alpha_t)_t \in \R_+^T$.
The budget constraints, written in \emph{increment} form, are: at $t=0$,
\begin{equation}\label{eq:budget-0-MP}
  c_0 + q_0 \cdot (\theta_0 - \xi^i) + p_0 (\alpha_0 - 0)
  \leq \omega^i_0 + \delta_0 \xi^i;
\end{equation}
for each $t \in \{1, \ldots, T-1\}$,
\begin{equation}\label{eq:budget-t-MP}
  c_t + q_t \cdot (\theta_t - \theta_{t-1}) + p_t (\alpha_t - \alpha_{t-1})
  \leq \omega^i_t + \kappa_t \cdot \theta_{t-1} + \delta_t \alpha_{t-1};
\end{equation}
and at $t = T$,
\begin{equation}\label{eq:budget-T-MP}
  c_T \leq \omega^i_T + \kappa_T \cdot \theta_{T-1} + \delta_T \alpha_{T-1}.
\end{equation}
The interpretation of \eqref{eq:budget-t-MP} is transparent: in period $t$, investor $i$ receives the coupons $\kappa_t \cdot \theta_{t-1}$ on his holdings of the previous period's securities and the dividend $\delta_t \alpha_{t-1}$ on his previous shareholding, and devotes this plus his endowment to consumption and to re-allocating his portfolio.
The set of actions $(c, \theta, \alpha)$ satisfying \eqref{eq:budget-0-MP}--\eqref{eq:budget-T-MP} is denoted by $B^i(q, p)$.

\begin{definition}[Multi-period stock-market equilibrium]\label{def:eq-multiperiod}
Given a firm's policy $(x, y, \beta, n)$, a stock-market equilibrium is a list
\begin{equation*}
  \bigl\{ (q, p), \, (c^i, \theta^i, \alpha^i)_{i \in I} \bigr\}
\end{equation*}
of security and share prices and investors' actions such that:
\begin{enumerate}
  \item[(a)] for every $i \in I$,
  \begin{equation*}
    (c^i, \theta^i, \alpha^i) \in \argmax \bigl\{ U^i(c) \colon (c, \theta, \alpha) \in B^i(q, p) \bigr\};
  \end{equation*}
  \item[(b)] for every $t \in \{0, \ldots, T-1\}$,
  \begin{equation*}
    \forall a \in A, \quad \beta_t(a) = \sum_{i \in I} \theta^i_t(a)
    \et
    \sum_{i \in I} \alpha^i_t = n_t;
  \end{equation*}
  \item[(c)] the commodity market clears:
  \begin{equation*}
    \sum_{i \in I} c^i = \sum_{i \in I} \omega^i + (y - x).
  \end{equation*}
\end{enumerate}
\end{definition}

\subsection{No-arbitrage and irrelevance of the payout policy}

Under strictly monotone preferences, the equilibrium prices admit strictly positive state-price deflators $\lambda = (\lambda_0, \ldots, \lambda_T) \in \R_{++}^{T+1}$ satisfying, for every $t \in \{0, \ldots, T-1\}$,
\begin{equation}\label{eq:no-arb-MP}
  \lambda_t q_t = \lambda_{t+1} (\kappa_{t+1} + q_{t+1})
  \et
  \lambda_t p_t = \lambda_{t+1} (p_{t+1} + \delta_{t+1}).
\end{equation}
These are the standard recursive no-arbitrage conditions: the present value of an asset today equals the discounted present value of its next-period coupon plus resale price.

For each $t \in \{0, \ldots, T\}$, define the firm's \emph{value} at date $t$ as
\begin{equation}\label{eq:firm-value-MP}
  V_t \coloneqq \underbrace{p_t n_t}_{E_t} + \underbrace{q_t \cdot \beta_t}_{D_t} + (y_t - x_t),
\end{equation}
with the convention $p_T = 0$ and $q_T = 0$.

\begin{theorem}[Irrelevance of the payout policy]\label{thm:payout-irrelevance}
Under the no-arbitrage conditions \eqref{eq:no-arb-MP} and the cash-flow identity \eqref{eq:cashflow-identity}, the firm's values $V_0, V_1, \ldots, V_T$ depend on the policy $(x, y, \beta, n)$ only through the investment--production plan $(x, y)$.
Specifically, for every $t$,
\begin{equation}\label{eq:mp-irrelevance}
  \lambda_t V_t = \sum_{\tau = t}^{T} \lambda_\tau (y_\tau - x_\tau).
\end{equation}
\end{theorem}

\begin{proof}
We show the recursion $\lambda_t V_t = \lambda_t (y_t - x_t) + \lambda_{t+1} V_{t+1}$ for $t \in \{0, \ldots, T-1\}$.
Iterating forward and using the terminal value $V_T = y_T - x_T = y_T$ (recall $p_T = 0$, $q_T = 0$, $x_T = 0$) yields \eqref{eq:mp-irrelevance}.

From \eqref{eq:firm-value-MP},
\begin{equation*}
  \lambda_t V_t = \lambda_t p_t n_t + \lambda_t q_t \cdot \beta_t + \lambda_t (y_t - x_t).
\end{equation*}
Applying the share-pricing condition in \eqref{eq:no-arb-MP},
\begin{equation*}
  \lambda_t p_t n_t = \lambda_{t+1} (p_{t+1} + \delta_{t+1}) n_t = \lambda_{t+1} p_{t+1} n_t + \lambda_{t+1} n_t \delta_{t+1},
\end{equation*}
and applying the bond-pricing condition,
\begin{equation*}
  \lambda_t q_t \cdot \beta_t = \lambda_{t+1} (\kappa_{t+1} + q_{t+1}) \cdot \beta_t = \lambda_{t+1} \kappa_{t+1} \cdot \beta_t + \lambda_{t+1} q_{t+1} \cdot \beta_t.
\end{equation*}
Summing,
\begin{equation*}
  \lambda_t (p_t n_t + q_t \cdot \beta_t)
  = \lambda_{t+1} \bigl[ p_{t+1} n_t + n_t \delta_{t+1} + \kappa_{t+1} \cdot \beta_t + q_{t+1} \cdot \beta_t \bigr].
\end{equation*}
Substituting the cash-flow identity \eqref{eq:cashflow-identity} for $n_t \delta_{t+1}$ and cancelling the $\kappa_{t+1} \cdot \beta_t$ terms,
\begin{equation*}
  \lambda_t (p_t n_t + q_t \cdot \beta_t)
  = \lambda_{t+1} \bigl[ p_{t+1} n_t + y_{t+1} - x_{t+1} + p_{t+1} (n_{t+1} - n_t) + q_{t+1} \cdot (\beta_{t+1} - \beta_t) + q_{t+1} \cdot \beta_t \bigr].
\end{equation*}
Collecting the $p_{t+1}$ and $q_{t+1}$ terms,
\begin{equation*}
  \lambda_t (p_t n_t + q_t \cdot \beta_t)
  = \lambda_{t+1} \bigl[ p_{t+1} n_{t+1} + q_{t+1} \cdot \beta_{t+1} + (y_{t+1} - x_{t+1}) \bigr]
  = \lambda_{t+1} V_{t+1}.
\end{equation*}
Adding $\lambda_t (y_t - x_t)$ to both sides yields $\lambda_t V_t = \lambda_t (y_t - x_t) + \lambda_{t+1} V_{t+1}$, completing the induction step.
\end{proof}

The terms $p_t$, $n_t$, and $\beta_t$ determine the individual components $E_t$ and $D_t$ of the firm's value and describe the firm's payout policy; they drop out of the total $V_t$.
In this deterministic benchmark, the firm's decisions on dividends, share repurchases, and debt issuance rearrange how the firm's future cash flows are delivered to investors, but leave the total present value untouched.

%====================================================================
\section{Beyond the Modigliani--Miller Benchmark}
\label{sec:beyond-mm}
%====================================================================

The irrelevance results of this chapter rely on a specific benchmark: frictionless markets, symmetric information, linear taxes (or no taxes), and unlimited or limited-but-costless-bankruptcy liability.
Once any of these assumptions fails, the capital structure of the firm re-acquires first-order economic content.
The remainder of the course is organized around the three most influential departures from the Modigliani--Miller benchmark.

\emph{Capital structure as a signalling device.}
In \citet{Ross1977} and \citet{LelandPyle1977}, managers hold private information about the firm's prospects and use the financial policy---the debt-to-equity ratio, the fraction of the firm retained by insiders---as a credible signal of their private type.
Financial structure is chosen to influence the market's beliefs.

\emph{Capital structure and agency costs.}
\citet{JensenMeckling1976} argue that, when managers and shareholders have conflicting interests, the capital structure can be used to discipline managerial behaviour: debt forces managers to commit future cash flows to creditors, reducing the scope for value-destroying investments.
The optimal capital structure balances the agency cost of equity against the agency cost of debt.

\emph{Capital structure and control rights.}
\citet{AghionBolton1992} and \citet{HarrisRaviv1989} emphasize that the allocation of \emph{control rights}---who has the authority to make which decision, in which state---is as important as the allocation of cash-flow rights.
Debt is distinctive because it transfers control to creditors in states where the firm is unable to meet its obligations.
The design of the capital structure is therefore inseparable from the design of corporate governance.

These three perspectives motivate the principal-agent models studied in Part~II of these notes.

\subsection{Further reading}

The original statement of the irrelevance theorem is \citet{ModiglianiMiller1958}.
\citet{Stiglitz1969, Stiglitz1974} provide the strong form used in Theorem~\ref{thm:stiglitz}.
For extensions to incomplete markets and interdependent securities, see \citet{DeMarzo1988} and \citet{DeMarzoDuffie1991}; \citet{Gottardi1995} gives a thorough analysis of the conditions under which the theorem survives in incomplete markets.
On the role of taxes, the classical analysis is \citet{Miller1977}; empirical assessments include \citet{MackieMason1990} and \citet{DesaiFoleyHines2004}.
On bankruptcy costs and financial distress, see \citet{SmithWarner1979}, \citet{Weiss1990}, \citet{OplerTitman1994}, and \citet{Purnanandam2008}.
For the signalling, agency, and control-rights perspectives touched upon in this section, classical references include \citet{Ross1977}, \citet{LelandPyle1977}, \citet{JensenMeckling1976}, \citet{AghionBolton1992}, and \citet{HarrisRaviv1989}.

%====================================================================
\chapter{Competitive Price Perceptions and Equilibrium Corporate Finance}
\label{ch:price-perceptions}
%====================================================================

The analysis of Chapters~\ref{ch:objective} and~\ref{ch:capital-structure} rests on a symmetric treatment of investors and firms: any security issued by the firm can, in principle, be freely traded---bought or short-sold---by any investor.
It is this symmetry that underlies the homemade-leverage argument of Theorem~\ref{thm:stiglitz} and, more generally, the marketed-consistency property of competitive price perceptions.
In real-world financial markets, however, short sales are typically constrained: short selling equity requires borrowing shares and posting collateral, and short positions in corporate bonds are rare outside the largest issuers.
Once investors are prevented from taking short positions, marginal valuations of a given cash flow may differ across investors even at equilibrium, and the marketed-consistency argument of Chapter~\ref{ch:objective} breaks down.

This chapter examines the firm's objective and the Modigliani--Miller theorem in an economy with short-sale constraints, following \citet{BisinGottardiRuta2015}.
The content is threefold.
First, we introduce the short-sale-constrained framework, pricing securities and equity at equilibrium via a ``highest-marginal-valuation'' identity.
Second, we develop \citet{Makowski1983b}'s notion of \emph{competitive price perceptions}: the firm (and every investor) conjectures that the equity price associated with an alternative production--finance plan equals the maximum marginal valuation that the economy would assign to the resulting dividend stream.
Under this criterion, shareholders unanimously agree that the firm should maximize its Makowski-based value, and the resulting equilibrium is constrained Pareto efficient.
Third, we revisit Modigliani--Miller: irrelevance holds when every investor is either an ``equity holder'' or a ``bond holder'' in a precise marginal sense, and can fail otherwise.
An extension to economies with intermediated short sales is briefly sketched at the end of the chapter.

%====================================================================
\section{The Bisin--Gottardi--Ruta Framework}
\label{sec:bgr-setup}
%====================================================================

\subsection{Markets, short sales, and the firm's technology}

We specialize the general model of Chapters~\ref{ch:objective}--\ref{ch:capital-structure} as follows.
There is a single firm (or, equivalently, a continuum of identical firms, as discussed in Section~\ref{sec:makowski-justification}).
There is a single financial asset: a riskless bond $a_0$ with $\kappa_s(a_0) = 1$ for every $s \in S$, trading at price $q$ at $t=0$.
We denote by $r_0$ the risk-free interest rate, defined by $q = 1/(1 + r_0)$.

Unlike in the previous chapters, trade is subject to explicit short-sale constraints:
\begin{itemize}
  \item the firm can issue debt but cannot buy the bond, so $\beta \geq 0$;
  \item investors can buy the bond but cannot issue debt, so $\theta^i \geq 0$ for every $i \in I$;
  \item investors cannot short-sell the firm's equity, so $\alpha^i \geq 0$ for every $i \in I$.
\end{itemize}
These constraints introduce a fundamental asymmetry between the firm, which can short the bond, and investors, who cannot.
The Stiglitz irrelevance argument of Chapter~\ref{ch:capital-structure} relied precisely on the possibility of homemade leverage, i.e., on investors' ability to offset any change in the firm's financial policy by adjusting their own portfolios.
With $\theta^i \geq 0$, such offsetting is no longer unrestricted, and the irrelevance argument must be re-examined.

The firm's production technology is given by a production set
\begin{equation}\label{eq:Y-BGR}
  Y \coloneqq \bigl\{ (y_0, (y_s)_{s \in S}) \in \R \times \R^S \colon y_0 \leq 0 \et y_s \leq f_s(-y_0), \ \text{for all $s\in S$} \bigr\},
\end{equation}
for a family of production functions $f_s : [0, \infty) \to [0, \infty)$, one for each state $s \in S$.
The quantity $k \coloneqq - y_0 \geq 0$ is the capital invested at $t=0$, which delivers the state-contingent output $f_s(k)$ at $t=1$.

To rule out bankruptcy at the outset, we impose a \emph{solvency} requirement on the firm's financial choice.

\begin{assumption}[No-default]\label{ass:no-default}
The firm chooses a capital $k \geq 0$ and a debt level $\beta \geq 0$ such that
\begin{equation*}
  \inf_{s \in S} f_s(k) \geq \beta.
\end{equation*}
\end{assumption}

We denote by $\mathcal{A}$ the set of \emph{admissible} capital-budgeting and financial strategies,
\begin{equation}\label{eq:admissible-set}
  \mathcal{A} \coloneqq \bigl\{ (k, \beta) \in \R_+^2 \colon \inf_{s \in S} f_s(k) \geq \beta \bigr\}.
\end{equation}

\subsection{Investors and budget sets}

As in Chapter~\ref{ch:objective}, each investor $i \in I$ is endowed with $\omega^i_0 \geq 0$ and $\omega^i_1 = (\omega^i_s)_{s \in S} \in \R_+^S$, an initial share $\xi^i \geq 0$ (with $\sum_{i} \xi^i = 1$), a discount factor $\beta^i > 0$, beliefs $\nu^i \in \Prob(S)$, and a Bernoulli utility $u^i : \R_+ \to [-\infty, \infty)$.
His expected utility over a consumption plan $(c_0, c_1) \in \R_+ \times \R_+^S$ is given by \eqref{eq:expected-utility}.

Given the firm's capital-budgeting and financial strategy $(k, \beta) \in \mathcal{A}$, the security price $q$, and the equity price $E$, investor $i$ chooses a consumption plan $(c_0, c_1)$, a bond holding $\theta^i \geq 0$, and a share $\alpha^i \geq 0$ satisfying the budget constraints
\begin{align}
  c^i_0 + q \theta^i + E \alpha^i
  &\leq
  \omega^i_0 + \bigl[ E + q \beta - k \bigr] \xi^i,
  \label{eq:budget-0-BGR}\\[2pt]
  c^i_s
  &\leq
  \omega^i_s + \theta^i + \bigl[ f_s(k) - \beta \bigr] \alpha^i,
  \quad \text{for all $s\in S$}.
  \label{eq:budget-s-BGR}
\end{align}
The set of actions $(c, \theta, \alpha) \in \R_+^{1+S} \times \R_+ \times \R_+$ satisfying \eqref{eq:budget-0-BGR}--\eqref{eq:budget-s-BGR} is denoted by $B^i(q, E, k, \beta)$.

The interpretation of \eqref{eq:budget-0-BGR} mirrors that of Chapter~\ref{ch:objective}.
At $t=0$, the original shareholder of firm $j$ contributes the investment $k$ and receives the proceeds $q \beta$ of the debt issuance, yielding the residual dividend $\delta_0 = q \beta - k$ per initial share.
Combining with the equity price $E$, the initial stake is worth $(E + q \beta - k) \xi^i$.

\subsection{Competitive equilibrium given the firm's decisions}

In this chapter, we temporarily adopt a ``two-stage'' perspective: we first fix the firm's decision $(k, \beta)$ and define competitive equilibrium in the resulting consumer economy; only in Section~\ref{sec:strong-unanimity} do we close the model by asking what $(k, \beta)$ the firm should choose.

\begin{definition}[Competitive equilibrium given $(k, \beta)$]\label{def:eq-BGR}
Fix a firm's decision $(k, \beta) \in \mathcal{A}$.
A list
\begin{equation*}
  \bigl\{ (k, \beta), \, (q, E), \, (c^i, \theta^i, \alpha^i)_{i \in I} \bigr\}
\end{equation*}
is a \emph{competitive equilibrium} if:
\begin{enumerate}
  \item[(a)] for every $i \in I$,
  \begin{equation*}
    (c^i, \theta^i, \alpha^i) \in \argmax \bigl\{ U^i(c) \colon (c, \theta, \alpha) \in B^i(q, E, k, \beta) \bigr\};
  \end{equation*}
  \item[(b)] financial markets clear:
  \begin{equation*}
    \sum_{i \in I} \theta^i = \beta
    \et
    \sum_{i \in I} \alpha^i = 1;
  \end{equation*}
  \item[(c)] consumption markets clear:
  \begin{equation*}
    k + \sum_{i \in I} c^i_0 = \omega_0
    \et
    \sum_{i \in I} c^i_1 = \omega_1 + f(k),
  \end{equation*}
  where $\omega_0 = \sum_i \omega^i_0$, $\omega_1 = \sum_i \omega^i_1$, and $f(k) = (f_s(k))_{s \in S}$.
\end{enumerate}
\end{definition}

Throughout this chapter, we assume the conditions (C.1)--(C.4) of Assumption~\ref{ass:C} on the primitives of investors' preferences and endowments.

%====================================================================
\section{Equilibrium Pricing Under Short-Sale Constraints}
\label{sec:bgr-pricing}
%====================================================================

\subsection{Present-value operator}

For each investor $i$ with interior equilibrium consumption $c^i \gg 0$, we denote by $\PV^i$ the linear functional on $\R^S$ defined by
\begin{equation}\label{eq:PV-def}
  \PV^i[w] \coloneqq \sum_{s \in S} \MRS^i_s(c^i) w_s,
  \quad w = (w_s)_{s \in S} \in \R^S,
\end{equation}
where $\MRS^i_s(c^i) = \beta^i \nu^i_s (u^i)'(c^i_s) / (u^i)'(c^i_0)$.
$\PV^i[w]$ is the value at $t=0$ of the state-contingent claim $w$, evaluated at investor $i$'s marginal rates of substitution at his equilibrium consumption.
In Chapter~\ref{ch:objective} we showed that, under complete markets, $\PV^i$ coincides across investors and equals the unique state-price deflator $\lambda$.
Under the short-sale constraints introduced in Section~\ref{sec:bgr-setup}, this coincidence fails: different investors may assign different present values to the same claim at equilibrium.

\subsection{Max-valuation pricing at equilibrium}

With short-sale constraints on bonds and shares, the first-order conditions of investor $i$'s problem are \emph{inequality} rather than equality conditions, and equilibrium prices are determined by the \emph{maximum} over all investors' marginal valuations.

\begin{proposition}[Max-valuation pricing]\label{prop:max-pricing}
Assume Assumption~\ref{ass:C} holds and let
\begin{equation*}
  \bigl\{ (k, \beta), \, (q, E), \, (c^i, \theta^i, \alpha^i)_{i \in I} \bigr\}
\end{equation*}
be a competitive equilibrium in the sense of Definition~\ref{def:eq-BGR}, with $\beta > 0$.
Then
\begin{equation}\label{eq:max-price-q}
  q = \max_{i \in I} \PV^i[\1_S]
\end{equation}
and
\begin{equation}\label{eq:max-price-E}
  E = \max_{i \in I} \PV^i[f_1(k) - \beta \1_S].
\end{equation}
Moreover, \eqref{eq:max-price-E} holds whenever the stock market clears with $\sum_i \alpha^i = 1$, regardless of whether $\beta > 0$.
\end{proposition}

\begin{proof}
The Lagrangian of investor $i$'s problem is
\begin{equation*}
  \mathcal{L}^i = U^i(c) - \mu^i_0 \bigl[c_0 + q \theta + E \alpha - \omega^i_0 - (E + q \beta - k) \xi^i\bigr]
  - \sum_{s \in S} \mu^i_s \bigl[c_s - \omega^i_s - \theta - (f_s(k) - \beta) \alpha\bigr]
  + \eta^i_\theta \theta + \eta^i_\alpha \alpha,
\end{equation*}
with non-negative multipliers $\mu^i_0, \mu^i_s, \eta^i_\theta, \eta^i_\alpha \geq 0$, the last two attached to the short-sale constraints $\theta \geq 0$ and $\alpha \geq 0$.
The first-order conditions with respect to $c_0$, $c_s$, $\theta$, and $\alpha$ give
\begin{align}
  (u^i)'(c^i_0) &= \mu^i_0,
  \label{eq:foc-c0-BGR}\\
  \beta^i \nu^i_s (u^i)'(c^i_s) &= \mu^i_s, \quad \text{for all $s\in S$},
  \label{eq:foc-cs-BGR}\\
  \mu^i_0 q &= \sum_{s \in S} \mu^i_s + \eta^i_\theta,
  \label{eq:foc-theta-BGR}\\
  \mu^i_0 E &= \sum_{s \in S} \mu^i_s [f_s(k) - \beta] + \eta^i_\alpha.
  \label{eq:foc-alpha-BGR}
\end{align}
Dividing \eqref{eq:foc-theta-BGR} and \eqref{eq:foc-alpha-BGR} by $\mu^i_0 > 0$,
\begin{equation}\label{eq:FOC-inequalities}
  q = \PV^i[\1_S] + \frac{\eta^i_\theta}{\mu^i_0}
  \et
  E = \PV^i[f_1(k) - \beta \1_S] + \frac{\eta^i_\alpha}{\mu^i_0}.
\end{equation}
Because $\eta^i_\theta, \eta^i_\alpha \geq 0$, these identities give
\begin{equation*}
  q \geq \PV^i[\1_S]
  \et
  E \geq \PV^i[f_1(k) - \beta \1_S],
  \quad \forall i \in I,
\end{equation*}
with equality whenever the corresponding short-sale constraint does not bind (complementary slackness: $\eta^i_\theta \theta^i = 0$ and $\eta^i_\alpha \alpha^i = 0$).

For \eqref{eq:max-price-q}, note that the bond market clears with $\sum_i \theta^i = \beta > 0$, so there exists at least one investor $i^\ast$ with $\theta^{i^\ast} > 0$.
Complementary slackness gives $\eta^{i^\ast}_\theta = 0$, whence $q = \PV^{i^\ast}[\1_S] \leq \max_i \PV^i[\1_S]$.
Combined with the inequality $q \geq \PV^i[\1_S]$ valid for every $i$, we conclude $q = \max_i \PV^i[\1_S]$.

For \eqref{eq:max-price-E}, stock market clearing gives $\sum_i \alpha^i = 1 > 0$, hence some investor $i^{\ast\ast}$ has $\alpha^{i^{\ast\ast}} > 0$ and $\eta^{i^{\ast\ast}}_\alpha = 0$.
The same argument yields $E = \max_i \PV^i[f_1(k) - \beta \1_S]$.
\end{proof}

The max-valuation identity \eqref{eq:max-price-q}--\eqref{eq:max-price-E} is a fundamental departure from the Arrow--Debreu pricing of Chapter~\ref{ch:objective}.
In a frictionless economy with complete markets, the common state-price deflator equals every investor's marginal rate of substitution; with short-sale constraints, only the \emph{highest-valuation} investors are marginal in the market, while all others are strictly infra-marginal.
The price of a cash flow is determined not by an average valuation but by the maximum valuation present in the economy.

%====================================================================
\section{Makowski's Competitive Price Perceptions}
\label{sec:makowski}
%====================================================================

\subsection{Pricing out-of-equilibrium firm choices}

Proposition~\ref{prop:max-pricing} pins down prices in equilibrium, given the firm's actual decision $(k, \beta)$.
To decide whether $(k, \beta)$ is itself optimal from the shareholders' perspective, the manager (and each shareholder) must conjecture how prices would respond to an alternative choice $(k', \beta') \in \mathcal{A}$.
The security price $q$ does not depend on the firm's decision, since the riskless bond's payoff is independent of $(k, \beta)$.
The equity price $E$, however, is sensitive to $(k, \beta)$ through the firm's dividend stream $(f_s(k) - \beta)_{s \in S}$.

In Chapter~\ref{ch:objective} we closed this question by postulating \emph{competitive price perceptions} in the sense of \citet{GrossmanHart1979}: each investor $i$ evaluates alternative firm cash flows at his own marginal rates of substitution.
This specification is consistent with market prices of \emph{marketed} streams (Proposition~\ref{prop:market-consistency}), but requires, in general, that all investors agree on the marginal valuation.
Under the short-sale constraints of the present framework, different investors assign different $\PV^i$'s to the same stream, and the Grossman--Hart specification no longer yields a single, operational equity-price map.

\subsection{The Makowski criterion}

\citet{Makowski1983b} proposed an alternative criterion tailored to economies with non-trivial heterogeneity of marginal valuations.

\begin{definition}[Makowski price perception]\label{def:makowski}
Each investor is aware that the equity price depends on the firm's choice $(k', \beta')$, and conjectures that the equity price associated with $(k', \beta')$ is the \emph{highest} marginal valuation that the economy would assign to the resulting dividend stream:
\begin{equation}\label{eq:makowski}
  \widetilde{E}(k', \beta') \coloneqq \max_{i \in I} \PV^i\bigl[ f_1(k') - \beta' \1_S \bigr],
  \quad (k', \beta') \in \mathcal{A}.
\end{equation}
\end{definition}

Observe that the Makowski map $\widetilde{E}$ uses investors' marginal rates of substitution evaluated at their \emph{equilibrium} consumption $c^i$---the same $\MRS^i$'s that appear in the pricing formula \eqref{eq:max-price-E}.
In particular, $\widetilde{E}$ is consistent with the equilibrium equity price $E$:
\begin{equation*}
  \widetilde{E}(k, \beta) = E.
\end{equation*}
For alternative choices $(k', \beta') \neq (k, \beta)$, the Makowski map assigns the value that the market would give if the firm sold the new dividend stream to the highest bidder.

\subsection{Justifications and the continuum-of-firms argument}
\label{sec:makowski-justification}

The Makowski specification admits two complementary justifications.

\emph{A market-microstructure interpretation.}
When a shareholder evaluates the firm's cash flow, he thinks of the benefits he could obtain by selling the firm in the market.
Hence, he uses the marginal rate of substitution of those investors willing to pay the most---the highest-valuation investors---because they determine the resale price.
In a competitive economy without bilateral matching, selling a small quantity of a claim with payoff $w$ at a price below $\max_i \PV^i[w]$ would immediately be arbitraged: the highest-valuation buyer would purchase it instead of his current marginal trade.

\emph{A continuum-of-firms interpretation.}
In the slide's language, one interprets the model as featuring a continuum of identical firms, each of measure zero.
A deviation by a single firm has negligible impact on aggregate quantities and hence on investors' marginal rates of substitution.
Prices $(q, \widetilde{E})$ are taken as \emph{price maps}, seen as given by firms and investors alike.
This is analogous to the price-taking assumption of the Arrow--Debreu framework, adapted to the situation in which the relevant ``price'' of a firm is a function of the firm's characteristics.

A conceptual difficulty with Definition~\ref{def:makowski} is the informational demand it places on individual investors: each agent must know (or be able to compute) the marginal rate of substitution of every other investor, or at least the maximum over the population.
This is a strong epistemic hypothesis that the literature has addressed in various ways, e.g., by assuming perfect competition with a continuum of agents, or by embedding the model in a rational-expectations framework; see \citet{BisinGottardiRuta2015} for a discussion.

%====================================================================
\section{Strong Unanimity}
\label{sec:strong-unanimity}
%====================================================================

Under Makowski's criterion, the firm's value map is
\begin{equation}\label{eq:firm-value-BGR}
  V(k', \beta') \coloneqq \widetilde{E}(k', \beta') + q \beta' - k',
  \quad (k', \beta') \in \mathcal{A}.
\end{equation}
We now establish the central normative result of the chapter: if the firm has chosen $(k, \beta)$ to maximize $V$ over $\mathcal{A}$, then every shareholder agrees with the choice, even when allowed to jointly choose his own consumption and portfolio.

\begin{theorem}[Strong unanimity]\label{thm:strong-unanimity-BGR}
Fix a competitive equilibrium
\begin{equation*}
  \bigl\{ (k, \beta), \, (q, E), \, (c^i, \theta^i, \alpha^i)_{i \in I} \bigr\}
\end{equation*}
under Assumption~\ref{ass:C}.
Assume that the firm has chosen $(k, \beta) \in \mathcal{A}$ to maximize $V$ over $\mathcal{A}$, where $V$ is the Makowski value \eqref{eq:firm-value-BGR} induced by \eqref{eq:makowski}.
Denote by $B^i(q, E', k', \beta')$ the budget set of investor $i$ under prices $(q, E')$ with $E' = \widetilde{E}(k', \beta')$ and firm decision $(k', \beta')$.
Then, for every investor $i \in I$,
\begin{equation}\label{eq:strong-unanimity-BGR}
  U^i(c^i) = \max \bigl\{ U^i(c') \colon (c', \theta', \alpha') \in B^i(q, E', k', \beta'), \ (k', \beta') \in \mathcal{A} \bigr\}.
\end{equation}
In particular, $(k, \beta) \in \mathcal{A}$ is the investor's preferred choice of firm policy.
\end{theorem}

\begin{proof}
Fix $i \in I$ and consider any candidate $\bigl( (k', \beta'), (c', \theta', \alpha') \bigr)$ with $(k', \beta') \in \mathcal{A}$ and $(c', \theta', \alpha') \in B^i(q, E', k', \beta')$, where $E' = \widetilde{E}(k', \beta')$.
The budget constraints \eqref{eq:budget-0-BGR}--\eqref{eq:budget-s-BGR} give
\begin{align}
  c'_0 + q \theta' + E' \alpha'
  &\leq
  \omega^i_0 + (E' + q \beta' - k') \xi^i,
  \label{eq:budget-0-prime}\\
  c'_s
  &\leq
  \omega^i_s + \theta' + (f_s(k') - \beta') \alpha',
  \quad \text{for all $s\in S$}.
  \label{eq:budget-s-prime}
\end{align}
Multiplying \eqref{eq:budget-s-prime} by $\MRS^i_s(c^i) > 0$ and summing over $s$,
\begin{equation*}
  \PV^i[c'_1]
  \leq
  \PV^i[\omega^i_1] + \theta' \PV^i[\1_S] + \alpha' \PV^i\bigl[ f_1(k') - \beta' \1_S \bigr].
\end{equation*}
By Proposition~\ref{prop:max-pricing} and Definition~\ref{def:makowski},
\begin{equation*}
  \PV^i[\1_S] \leq q
  \et
  \PV^i[f_1(k') - \beta' \1_S] \leq \widetilde{E}(k', \beta') = E'.
\end{equation*}
Because $\theta' \geq 0$ and $\alpha' \geq 0$,
\begin{equation*}
  \PV^i[c'_1]
  \leq
  \PV^i[\omega^i_1] + q \theta' + E' \alpha'.
\end{equation*}
Adding \eqref{eq:budget-0-prime} and using the definition of the Makowski value \eqref{eq:firm-value-BGR},
\begin{equation}\label{eq:PV-bound-prime}
  c'_0 + \PV^i[c'_1]
  \leq
  \omega^i_0 + \PV^i[\omega^i_1] + V(k', \beta') \xi^i.
\end{equation}
Since the firm has chosen $(k, \beta)$ to maximize $V$ over $\mathcal{A}$, $V(k', \beta') \leq V(k, \beta)$.
Hence, for every admissible pair $\bigl( (k', \beta'), (c', \theta', \alpha') \bigr)$,
\begin{equation}\label{eq:PV-bound}
  c'_0 + \PV^i[c'_1]
  \leq
  \omega^i_0 + \PV^i[\omega^i_1] + V(k, \beta) \xi^i.
\end{equation}

The same argument, applied to the equilibrium allocation $(c^i, \theta^i, \alpha^i)$ under the firm's actual choice $(k, \beta)$, yields the \emph{equality}
\begin{equation}\label{eq:PV-eq-BGR}
  c^i_0 + \PV^i[c^i_1]
  =
  \omega^i_0 + \PV^i[\omega^i_1] + V(k, \beta) \xi^i.
\end{equation}
Indeed, at equilibrium, the budget constraints \eqref{eq:budget-0-BGR}--\eqref{eq:budget-s-BGR} bind (by strict monotonicity of $u^i$), and complementary slackness in \eqref{eq:FOC-inequalities} gives $\theta^i \PV^i[\1_S] = \theta^i q$ and $\alpha^i \PV^i[f_1(k) - \beta \1_S] = \alpha^i E$.
Substituting these identities into the present-value sum of the budget constraints yields \eqref{eq:PV-eq-BGR}.

Consider now the relaxed problem
\begin{equation}\label{eq:relaxed-problem}
  \max_{c \in \R_+ \times \R_+^S} U^i(c)
  \quad \text{s.t.} \quad
  c_0 + \PV^i[c_1] \leq \omega^i_0 + \PV^i[\omega^i_1] + V(k, \beta) \xi^i.
\end{equation}
Type: exposition / mathematical clarification  
Location: Chapter 3, Theorem 3.5 proof, paragraph introducing the relaxed problem  
Comment:  
The proof implicitly relies on the fact that the marginal rates of substitution $MRS^i_s(c^i)$ are well-defined. This requires the equilibrium allocation $c^i$ to be strictly positive. This follows from (C.4) together with individual rationality of equilibrium allocations, but it should be stated explicitly. In addition, the argument uses that the budget constraint of the relaxed problem binds at the optimum, which follows from strict monotonicity of preferences but is not explicitly mentioned.  
Suggestion:  
Add a sentence clarifying that equilibrium allocations are individually rational (since the endowment allocation is always feasible) and so by (C.4) we have $c^i \gg 0$, ensuring that $MRS^i_s(c^i)$ and hence $PV^i$ are well-defined. Also, explicitly note that the relaxed budget constraint binds at the optimum since $U^i$ is strictly increasing, i.e., $c_0 + PV^i[c_1] = W^i$ at the solution.

The Lagrangian of \eqref{eq:relaxed-problem} has first-order conditions $(u^i)'(c_0) = \mu$ and $\beta^i \nu^i_s (u^i)'(c_s) = \mu \MRS^i_s(c^i)$ for every $s$.
Substituting the definition of $\MRS^i_s(c^i)$ and matching with $\mu = (u^i)'(c^i_0)$, one obtains $(u^i)'(c_s) = (u^i)'(c^i_s)$ and $(u^i)'(c_0) = (u^i)'(c^i_0)$, i.e., $c = c^i$ by strict concavity of $u^i$.
Because $U^i$ is concave, $c^i$ is the \emph{global} optimum of the relaxed problem.

Type: mathematical issue  
Location: Chapter 3, Theorem 3.5 proof, paragraph starting with “Substituting the definition of $MRS^i_s(c^i)$…”  
Comment:  
The proof incorrectly concludes that equality of marginal rates of substitution implies equality of consumption allocations. In general, equality of MRS does not imply $c = c^i$. The argument is missing the role of the budget constraint and concavity. The correct reasoning is that $c^i$ satisfies the first-order (KKT) conditions of the relaxed problem by construction, together with the budget constraint holding with equality. Concavity of $U^i$ then guarantees global optimality.  
Suggestion:  
Replace the paragraph with:
\medskip
\noindent
Note that $c^i$ satisfies the first-order conditions of the relaxed problem by construction. Indeed, setting $\mu = (u^i)'(c^i_0)$, and using the definition of $MRS^i_s(c^i)$, one verifies that the first-order conditions hold at $c = c^i$. Moreover, from (3.20), $c^i$ satisfies the budget constraint with equality. Therefore, $c^i$ satisfies the KKT conditions of the relaxed problem. Since $U^i$ is concave, these conditions are sufficient for global optimality. Hence, $c^i$ solves the relaxed problem.

Returning to \eqref{eq:PV-bound}, any $c'$ satisfying this inequality is feasible in the relaxed problem, and therefore $U^i(c') \leq U^i(c^i)$.
Since the candidate $\bigl( (k', \beta'), (c', \theta', \alpha') \bigr)$ was arbitrary, \eqref{eq:strong-unanimity-BGR} follows.
\end{proof}

\begin{remark}[What makes the unanimity ``strong'']\label{rem:why-strong}
Theorem~\ref{thm:strong-unanimity-BGR} states not merely that investor $i$ prefers the firm's value-maximizing choice $(k, \beta)$ over alternatives $(k', \beta')$, with his own portfolio fixed at $(\theta^i, \alpha^i)$.
It says that, when investor $i$ contemplates a firm deviation to $(k', \beta')$, he is also allowed to optimize his \emph{own} action $(c', \theta', \alpha')$ in response, and he still cannot do better than at the equilibrium.
This is what distinguishes ``strong'' unanimity (Theorem~\ref{thm:strong-unanimity-BGR}) from the weaker notion obtained by evaluating a single investor's utility change at his fixed equilibrium portfolio.
\end{remark}

%====================================================================
\section{Constrained Pareto Efficiency}
\label{sec:constrained-efficiency}
%====================================================================

Strong unanimity asks whether the firm's choice is efficient \emph{from the perspective of individual shareholders}.
A broader question is whether the resulting allocation is Pareto efficient in the economy at large.
Because of the short-sale constraints, the relevant notion is \emph{constrained} Pareto efficiency: efficiency relative to the set of allocations achievable within the existing market structure.

\subsection{Admissible allocations}

\begin{definition}[Admissible allocation]\label{def:admissible-alloc}
A profile of consumption plans $(c^i)_{i \in I} \in (\R_+ \times \R_+^S)^I$ is \emph{admissible} if:
\begin{enumerate}
  \item[\textup{(i)}] it is feasible: there exists $k \geq 0$ such that
  \begin{equation*}
    k + \sum_{i \in I} c^i_0 \leq \omega_0
    \et
    \sum_{i \in I} c^i_s \leq f_s(k) + \sum_{i \in I} \omega^i_s, \quad \text{for all $s\in S$};
  \end{equation*}
  \item[\textup{(ii)}] it is implementable with the existing asset structure: there exists $\beta \geq 0$ such that, for every $i$, one can find $(\theta^i, \alpha^i) \in \R_+ \times \R_+$ satisfying
  \begin{equation*}
    c^i_s = \omega^i_s + \theta^i + \bigl[ f_s(k) - \beta \bigr] \alpha^i, \quad \text{for all $s\in S$}.
  \end{equation*}
\end{enumerate}
\end{definition}

Admissibility captures two distinct constraints.
Feasibility, (i), is the familiar resource constraint: aggregate consumption at each date and state cannot exceed aggregate endowment plus production.
Implementability, (ii), is the constraint imposed by the asset structure and by the short-sale restrictions: investor $i$'s consumption at $t=1$ must be attainable from non-negative bond and share holdings, starting from his own endowment.
Crucially, Definition~\ref{def:admissible-alloc} does \emph{not} require the financial portfolios $(\theta^i, \alpha^i)_{i}$ to clear: a ``planner'' choosing an admissible allocation may, in principle, assign portfolios $\theta^i$ summing to more (or less) than $\beta$, and shares $\alpha^i$ summing to more (or less) than $1$.
This added flexibility is what makes the constrained-efficiency notion weaker than full Pareto efficiency.

\begin{remark}\label{rem:eq-is-admissible}
Every competitive-equilibrium allocation is admissible: take $k$ and $\beta$ from the equilibrium, and the portfolios $(\theta^i, \alpha^i)_i$ from the equilibrium actions.
\end{remark}

\subsection{Constrained Pareto efficiency}

\begin{definition}[Constrained Pareto efficiency]\label{def:constrained-pareto}
An admissible allocation $(c^i)_{i \in I}$ is \emph{constrained Pareto efficient} if there does not exist another admissible allocation $(\widetilde{c}^i)_{i \in I}$ that Pareto dominates it, i.e., such that
\begin{equation*}
  U^i(\widetilde{c}^i) \geq U^i(c^i), \quad \forall i \in I,
\end{equation*}
with strict inequality for at least one $i$.
\end{definition}

\begin{theorem}[Welfare theorem]\label{thm:welfare-BGR}
Under Assumption~\ref{ass:C}, every competitive-equilibrium allocation is constrained Pareto efficient.
\end{theorem}

\begin{proof}
Suppose, by contradiction, that the equilibrium allocation $(c^i)_{i \in I}$ is Pareto-dominated by an admissible allocation $(\widetilde{c}^i)_{i \in I}$ supported by $(\widetilde{k}, \widetilde{\beta})$ and portfolios $(\widetilde{\theta}^i, \widetilde{\alpha}^i)_i \in (\R_+)^I \times (\R_+)^I$ with
\begin{equation*}
  \widetilde{c}^i_s = \omega^i_s + \widetilde{\theta}^i + \bigl[ f_s(\widetilde{k}) - \widetilde{\beta} \bigr] \widetilde{\alpha}^i, \quad \text{for all $s\in S$}.
\end{equation*}
Apply $\PV^i$ to this identity:
\begin{equation*}
  \PV^i[\widetilde{c}^i_1]
  = \PV^i[\omega^i_1] + \widetilde{\theta}^i \PV^i[\1_S] + \widetilde{\alpha}^i \PV^i\bigl[ f_1(\widetilde{k}) - \widetilde{\beta} \1_S \bigr].
\end{equation*}
By Proposition~\ref{prop:max-pricing} and Makowski's criterion \eqref{eq:makowski},
\begin{equation*}
  \PV^i[\1_S] \leq q
  \et
  \PV^i\bigl[ f_1(\widetilde{k}) - \widetilde{\beta} \1_S \bigr] \leq \widetilde{E}(\widetilde{k}, \widetilde{\beta}).
\end{equation*}
Because $\widetilde{\theta}^i, \widetilde{\alpha}^i \geq 0$,
\begin{equation}\label{eq:PV-ct}
  \PV^i[\widetilde{c}^i_1] \leq \PV^i[\omega^i_1] + q \widetilde{\theta}^i + \widetilde{E}(\widetilde{k}, \widetilde{\beta}) \widetilde{\alpha}^i.
\end{equation}

Consider, for each $i$, the candidate $\bigl( (\widetilde{k}, \widetilde{\beta}), (\widetilde{c}^i, \widetilde{\theta}^i + \sigma^i, \widetilde{\alpha}^i) \bigr)$, where $\sigma^i \geq 0$ is a non-negative adjustment of the bond holding to make the $t=0$ budget constraint of investor $i$ hold.
Specifically, set $\sigma^i$ to equalize \eqref{eq:budget-0-BGR} (with primed variables), namely,
\begin{equation*}
  \widetilde{c}^i_0 + q (\widetilde{\theta}^i + \sigma^i) + \widetilde{E}(\widetilde{k}, \widetilde{\beta}) \widetilde{\alpha}^i
  =
  \omega^i_0 + \bigl[ \widetilde{E}(\widetilde{k}, \widetilde{\beta}) + q \widetilde{\beta} - \widetilde{k} \bigr] \xi^i.
\end{equation*}
Whether $\sigma^i \geq 0$ is not guaranteed in general; when it is, $\widetilde{c}^i$ is achievable by investor $i$ in $B^i\bigl(q, \widetilde{E}(\widetilde{k}, \widetilde{\beta}), \widetilde{k}, \widetilde{\beta}\bigr)$, and Theorem~\ref{thm:strong-unanimity-BGR} applies directly: $U^i(\widetilde{c}^i) \leq U^i(c^i)$ for every $i$, contradicting Pareto domination.

The general case, in which $\sigma^i$ may be negative for some $i$, is handled by an aggregation argument: feasibility (i) of Definition~\ref{def:admissible-alloc} implies that the total ``excess wealth'' across investors is bounded by the aggregate production surplus, which coincides with $\sum_i V(\widetilde{k}, \widetilde{\beta}) \xi^i = V(\widetilde{k}, \widetilde{\beta}) \leq V(k, \beta)$ under value maximization.
Applying \eqref{eq:PV-bound} aggregated over $i$ yields $\sum_i U^i(\widetilde{c}^i) \leq \sum_i U^i(c^i)$, which again contradicts Pareto domination with strict improvement for some investor.
A complete treatment is given by \citet{BisinGottardiRuta2015} and \citet{Makowski1980}.
\end{proof}

The welfare theorem tightens the unanimity statement of Theorem~\ref{thm:strong-unanimity-BGR} from individual to aggregate efficiency.
Under Makowski price perceptions, the competitive equilibrium is socially optimal within the class of allocations that respect the short-sale constraints and the asset structure of the economy.

%====================================================================
\section{Modigliani--Miller Revisited}
\label{sec:MM-BGR}
%====================================================================

We now return to the irrelevance question studied in Chapter~\ref{ch:capital-structure} and ask: does Modigliani--Miller hold in the Bisin--Gottardi--Ruta framework?
The answer depends on the relationship between two subsets of investors, the \emph{equity holders} and the \emph{bond holders} in a marginal sense.

\subsection{Equity and bond holders}

Fix a competitive equilibrium $\bigl\{ (k, \beta), (q, E), (c^i, \theta^i, \alpha^i)_{i \in I} \bigr\}$ with $\beta > 0$.

\begin{definition}[Marginal equity and bond holders]\label{def:marginal-holders}
Let
\begin{align*}
  I^e &\coloneqq \bigl\{ i \in I \colon E = \PV^i\bigl[ f_1(k) - \beta \1_S \bigr] \bigr\}, \\
  I^b &\coloneqq \bigl\{ i \in I \colon q = \PV^i[\1_S] \bigr\}.
\end{align*}
An investor $i \in I^e$ is called an \emph{equity holder}, and an investor $i \in I^b$ a \emph{bond holder}.
\end{definition}

The terminology reflects complementary slackness in the first-order conditions \eqref{eq:FOC-inequalities}: $i \in I^e$ means that investor $i$ is on the verge of holding equity (whether he actually holds a positive share or is exactly indifferent), while $i \in I^b$ means that $i$ is on the verge of holding the bond.
By Proposition~\ref{prop:max-pricing}, $E = \max_i \PV^i[f_1(k) - \beta \1_S]$ and $q = \max_i \PV^i[\1_S]$, so $I^e$ and $I^b$ are both non-empty.

\subsection{At a value-maximizing equilibrium, equity holders are bond holders}

The next result shows that, when the firm maximizes its value and the no-default constraint is slack, the set of equity holders is contained in the set of bond holders.

\begin{proposition}[Equity holders are bond holders]\label{prop:eh-are-bh}
Assume Assumption~\ref{ass:C} holds and let
\begin{equation*}
  \bigl\{ (k, \beta), (q, E), (c^i, \theta^i, \alpha^i)_{i \in I} \bigr\}
\end{equation*}
be a competitive equilibrium in which the firm has chosen $(k, \beta)$ to maximize the Makowski value \eqref{eq:firm-value-BGR} over $\mathcal{A}$.
Suppose the no-default constraint is strictly slack, i.e., $f_s(k) > \beta > 0$ for every $s \in S$.
Then, for every investor $i \in I$,
\begin{equation*}
  E = \PV^i\bigl[ f_1(k) - \beta \1_S \bigr]
  \;\Longrightarrow\;
  q = \PV^i[\1_S],
\end{equation*}
that is, $I^e \subset I^b$.
\end{proposition}

\begin{proof}
Let $i \in I^e$, so that $E = \PV^i[f_1(k) - \beta \1_S]$.
By the Makowski criterion and the no-default slackness, the value map $V$ is differentiable in $\beta$ at $\beta$ and
\begin{equation*}
  V(k, \beta') = \widetilde{E}(k, \beta') + q \beta' - k
  = \max_{j \in I} \PV^j\bigl[ f_1(k) - \beta' \1_S \bigr] + q \beta' - k.
\end{equation*}
For $\beta'$ sufficiently close to $\beta$, the maximum on the right-hand side remains attained at $i$ (since the maximizer set depends continuously on $\beta'$ and the no-default constraint is slack), so
\begin{equation*}
  V(k, \beta') = \PV^i[f_1(k)] - \beta' \PV^i[\1_S] + q \beta' - k.
\end{equation*}
Differentiating in $\beta'$ at $\beta' = \beta$,
\begin{equation*}
  \frac{\partial V}{\partial \beta}(k, \beta)
  = q - \PV^i[\1_S].
\end{equation*}
Because $(k, \beta)$ maximizes $V$ over the interior of $\mathcal{A}$ (the no-default constraint is slack at $\beta$), the first-order condition $\partial V / \partial \beta = 0$ holds, giving $q = \PV^i[\1_S]$, i.e., $i \in I^b$.
\end{proof}

\subsection{Irrelevance}

Proposition~\ref{prop:eh-are-bh} is a first step towards irrelevance.
The next proposition states when it is sufficient to conclude that the firm's value is independent of $\beta'$.

\begin{proposition}[Modigliani--Miller in the Bisin--Gottardi--Ruta framework]\label{prop:mm-BGR}
Under the hypotheses of Proposition~\ref{prop:eh-are-bh}, if in addition every investor is either an equity holder or a bond holder, i.e., $I = I^e \cup I^b$, then the value map $V$ is independent of the financial policy on the no-default interior: for every $\beta'$ with $f_s(k) > \beta' \geq 0$, $\forall s$,
\begin{equation*}
  V(k, \beta') = V(k, 0) = \widetilde{E}(k, 0) - k.
\end{equation*}
\end{proposition}

\begin{proof}
Fix $\beta' \in [0, \inf_s f_s(k))$.
By the Makowski criterion and Proposition~\ref{prop:max-pricing},
\begin{equation*}
  V(k, \beta') = \max_{i \in I} \PV^i\bigl[ f_1(k) - \beta' \1_S \bigr] + q \beta' - k
  = \max_{i \in I} \bigl\{ \PV^i[f_1(k)] - \beta' \PV^i[\1_S] \bigr\} + q \beta' - k.
\end{equation*}
Suppose the maximum is attained at some $i^\ast \in I$; then $i^\ast \in I^e$ by Proposition~\ref{prop:eh-are-bh} (applied to the alternative firm choice $(k, \beta')$, provided the hypotheses still hold; we assume they do under the no-default slackness at $\beta'$).
Under the hypothesis $I = I^e \cup I^b$, if $i^\ast \in I^e$, then also $i^\ast \in I^b$, so $\PV^{i^\ast}[\1_S] = q$.
Therefore,
\begin{equation*}
  V(k, \beta') = \PV^{i^\ast}[f_1(k)] - \beta' q + q \beta' - k = \PV^{i^\ast}[f_1(k)] - k.
\end{equation*}
The right-hand side is independent of $\beta'$, and equals $V(k, 0)$.
\end{proof}

\begin{remark}[When Modigliani--Miller fails]\label{rem:mm-fails-BGR}
Proposition~\ref{prop:mm-BGR} may fail for two reasons.
\begin{enumerate}
  \item \emph{Binding no-default constraint.}
  If $\inf_s f_s(k) = \beta$, the admissible set $\mathcal{A}$ has $(k, \beta)$ on its boundary, and the marginal envelope condition underlying Proposition~\ref{prop:eh-are-bh} no longer applies.
  In this case, an increase in $\beta'$ would violate feasibility and the value map is effectively maximized at the corner.
  \item \emph{Investors who are neither equity nor bond holders.}
  If $I \neq I^e \cup I^b$, there exist investors $i$ with $\PV^i[\1_S] < q$ \emph{and} $\PV^i[f_1(k) - \beta \1_S] < E$.
  Such investors hold neither bonds nor equity: they are strictly infra-marginal in both markets.
  A change in $\beta$ may shift the identity of the marginal equity holder to such an investor, changing $\widetilde{E}(k, \beta)$ by more than the corresponding change $q \, d\beta$ in the bond valuation.
\end{enumerate}
The second point marks the substantive departure from Chapter~\ref{ch:capital-structure}: without short-sale constraints, homemade leverage ensures that every investor is effectively marginal in both markets, and irrelevance is automatic.
With short-sale constraints, whether Modigliani--Miller holds becomes an \emph{equilibrium} question, answered by whether the sets $I^e$ and $I^b$ cover the entire population of investors.
\end{remark}

%====================================================================
\section{Intermediated Short Sales}
\label{sec:intermediated}
%====================================================================

The previous sections have taken short-sale constraints as an exogenous feature of the market structure.
A natural question is whether financial innovation---specifically, the emergence of intermediaries who issue contracts replicating short positions---can restore the full unanimity and efficiency results of Chapters~\ref{ch:objective}--\ref{ch:capital-structure}.
This section sketches the argument.
A detailed development is provided by \citet{BisinGottardiRuta2015}.

\subsection{Intermediaries and the intermediation technology}

A \emph{financial intermediary} is a risk-neutral, zero-profit institution that issues long and short positions on the firm's equity.
A short position on the equity, held by an investor, is a contract that promises to deliver at $t=1$ the amount $-[f_s(k) - \beta]$ in state $s$.
The intermediary takes the opposite side: selling a long position to some investors and a short position to others, matched in quantity.
Intermediating $m$ units yields aggregate long demand $m$ and aggregate short demand $m$, which net to zero in the absence of default.

A simple default technology is assumed: in every state $s$, a fraction $\delta \in (0, 1)$ of the short position fails to be repaid, uniformly across investors.
To insure against the resulting revenue shortfall, the intermediary holds $\gamma$ units of the firm's equity as collateral, satisfying
\begin{equation}\label{eq:collateral}
  m \leq m(1 - \delta) + \gamma,
  \quad \text{i.e.,} \quad
  \gamma \geq m \delta.
\end{equation}

\subsection{The intermediary's problem and market prices}

Long and short positions trade at different prices, $q^+$ and $q^-$ respectively, reflecting the collateral cost.
The intermediary's problem is
\begin{equation}\label{eq:intermediary}
  \max \bigl\{ (q^+ - q^-) m - q \gamma \colon m \geq 0, \; \gamma \geq m \delta \bigr\}.
\end{equation}
Because the intermediation technology exhibits constant returns to scale, a positive $m$ arises at equilibrium only if profits vanish: $\gamma = m \delta$ and
\begin{equation}\label{eq:intermediation-eq-price}
  q^+ - q^- = q \delta.
\end{equation}
\eqref{eq:intermediation-eq-price} is the free-entry condition: the spread between long and short prices on the derivative must equal the collateral cost per unit.

\subsection{Investors' budget sets and market clearing}

Each investor $i$ holds, in addition to his bond position $\theta^i \geq 0$ and direct equity position $\alpha^i \geq 0$, a long position $\phi^i_+ \geq 0$ and a short position $\phi^i_- \geq 0$ on the intermediated derivative.
The budget constraints become
\begin{align}
  c^i_0 + q \theta^i + q^+ \phi^i_+ + E \alpha^i
  &\leq
  \omega^i_0 + (E + q \beta - k) \xi^i + q^- \phi^i_-,
  \label{eq:budget-0-int}\\[2pt]
  c^i_s + (1 - \delta) [f_s(k) - \beta] \phi^i_-
  &\leq
  \omega^i_s + \theta^i + [f_s(k) - \beta] \bigl( \alpha^i + \phi^i_+ \bigr),
  \quad \text{for all $s\in S$}.
  \label{eq:budget-s-int}
\end{align}
Notice that investor $i$'s effective t=1 payoff on equity-like claims is now $[f_s(k) - \beta] \bigl( \alpha^i + \phi^i_+ - (1-\delta) \phi^i_- \bigr)$: a long position is a perfect substitute for direct share ownership, while a short position replicates a negative equity holding up to the default friction $\delta$.
Market clearing now requires
\begin{equation*}
  \gamma + \sum_{i \in I} \alpha^i = 1
  \et
  \sum_{i \in I} \phi^i_+ = \sum_{i \in I} \phi^i_- = m,
\end{equation*}
with $\gamma = m \delta$ from \eqref{eq:collateral}.

\subsection{Adjusted price perceptions}

When intermediation is active ($m > 0$), the Makowski map must be adjusted to reflect the additional ``value of intermediation'' attached to each unit of equity, which can be used as collateral.
The corrected equity-price conjecture is
\begin{equation}\label{eq:makowski-intermediation}
  \widetilde{E}(k', \beta')
  =
  \max \Biggl\{
  \widetilde{E}^c(k', \beta'),
  \ \frac{\widetilde{q}^+(k', \beta') - \widetilde{q}^-(k', \beta')}{\delta}
  \Biggr\},
\end{equation}
where $\widetilde{E}^c(k', \beta') = \max_i \PV^i[f_1(k') - \beta' \1_S]$ is the consumer-only valuation of Section~\ref{sec:makowski}, and the second argument of the maximum is the value of the equity when used as intermediation collateral.
Two regimes of equilibrium may arise.

\emph{Regime $q > q^+$.}
Equity trades at a strict premium over the long position on the derivative.
The additional value of equity as an input to the intermediation technology makes it attractive to the intermediary only; all outstanding shares are purchased by the intermediary, and $\gamma = 1$.
The premium $q - q^+$ is determined by \eqref{eq:intermediation-eq-price}: the spread $q^+ - q^-$ adjusts so that equity prices clear.

\emph{Regime $q = q^+$.}
Investors are indifferent between buying equity directly and buying the long position on the derivative.
Part of the outstanding shares is held by consumers and part by the intermediary.
At this equilibrium, equity is still held in the economy, but at a price equal to the long-position price on the derivative.

\subsection{Unanimity and efficiency in the intermediated economy}

\begin{proposition}[Unanimity and constrained Pareto efficiency with intermediation]\label{prop:intermediated-unanimity}
At a competitive equilibrium of the economy with intermediated short sales, \emph{equity holders} unanimously support the production and financial decisions of the firm.
Moreover, the equilibrium allocation is constrained Pareto optimal relative to the admissibility notion adapted to the enriched asset structure.
\end{proposition}

The proof parallels those of Theorems~\ref{thm:strong-unanimity-BGR} and~\ref{thm:welfare-BGR}, with two modifications.
First, the price perception $\widetilde{E}$ is replaced by its adjusted version \eqref{eq:makowski-intermediation}, reflecting the option value of equity as intermediation collateral.
Second, the admissibility notion of Definition~\ref{def:admissible-alloc} is enlarged to include the derivative positions $\phi^i_\pm$, which allow short-selling equity at a well-defined cost.
Both modifications preserve the underlying logic of the argument: max-valuation pricing plus value maximization yield unanimity, and complementary slackness at the intermediary's problem ensures that the economy's set of marketed cash flows is rich enough to implement constrained-efficient allocations.

\subsection{Further reading}

The framework of this chapter is developed in \citet{BisinGottardiRuta2015}; an updated and refined version of this framework has subsequently been published as \citet{BisinClementiGottardi2025}.
It builds on Makowski's pioneering contributions on competitive price perceptions and the foundations of perfectly competitive economies with production: \citet{Makowski1980, Makowski1983a, Makowski1983b}.
The role of short-sale constraints and financial intermediation in restoring efficiency is studied, in a related vein, by \citet{AllenGale1991}.
On unanimity under incomplete markets, see \citet{CarcelesCoenPirani2009}.
On the interplay between managerial incentives, hedging, and equity structure, see \citet{AcharyaBisin2009} and \citet{BisinGottardiRampini2008}; for a model with non-exclusive contracts and moral hazard, \citet{BisinGuaitoli2004}.

%====================================================================
\chapter{Moral Hazard in the Stock-Market Equilibrium}
\label{ch:moral-hazard}
%====================================================================

Chapters~\ref{ch:objective}--\ref{ch:price-perceptions} developed the theory of the stock-market equilibrium in a purely \emph{frictionless} information environment: investors and firms share a common description of the state space and of each firm's production set, and the firm's output in each state depends only on the exogenous randomness, not on any unobservable action of the manager.
The Fisher--Drèze--Grossman--Hart unanimity of Chapter~\ref{ch:objective} and the Modigliani--Miller irrelevance of Chapter~\ref{ch:capital-structure} hinge on this symmetry of information.
Once the firm's output depends on a managerial action that is \emph{not} contractible---effort, diligence, attention, integrity---the logic of those chapters breaks down, and a new set of issues arises under the heading of \emph{moral hazard}.

This chapter studies the stock-market equilibrium in the presence of moral hazard, following the classical analysis of \citet{MagillQuinzii2002}.
Three questions frame the discussion.
First, what is the appropriate equilibrium concept in an economy where entrepreneurs' effort is unobservable but can be inferred from their observable financial decisions?
Second, does a stock market with equity, debt, and derivative securities deliver a constrained Pareto-efficient allocation of effort and risk, even when the first-best risk-sharing arrangement would destroy the entrepreneur's incentives?
Third, under what conditions on the security structure does the welfare property survive?
The chapter answers the first question by formalizing the notion of a \emph{rational competitive price perceptions} (RCPP) equilibrium, answers the second by establishing a welfare theorem for RCPP equilibria under a spanning condition, and answers the third by identifying \emph{strong partial spanning} as the critical hypothesis.

The logical structure mirrors Chapter~\ref{ch:price-perceptions}: price perceptions are again at the heart of the analysis, but now they must be rich enough to convey information about effort.
Where the Makowski criterion of Chapter~\ref{ch:price-perceptions} priced alternative firm choices via the highest marginal valuation, the RCPP criterion prices them via the unique no-arbitrage state-price evaluation of a \emph{marketed} cash flow.
Both approaches share the view that out-of-equilibrium decisions are priced as if they produced small, non-disruptive deviations from the aggregate market structure.

%====================================================================
\section{Framework}
\label{sec:MH-framework}
%====================================================================

\subsection{Agents: entrepreneurs and investors}

We consider a two-period economy with production and a single commodity per period.
Uncertainty at $t=1$ is described by a finite set $S$ of states of nature.
The set $I$ of agents is partitioned into two types,
\begin{equation*}
  I = I_1 \cup I_2,
\end{equation*}
where $I_1$ is a finite set of \emph{entrepreneurs} and $I_2$ a finite set of \emph{investors}.
Each agent $i \in I$ is endowed with initial wealth $\omega^i_0 \geq 0$ at $t=0$ and zero endowment at $t=1$.

\subsection{Production technology}

Each entrepreneur $i \in I_1$ has access to a productive opportunity.
By investing capital $k^i \geq 0$ at $t=0$ and exerting effort $e^i$ drawn from an abstract set $E^i$, the entrepreneur produces, in state $s \in S$ at $t=1$, the output
\begin{equation*}
  f^i_s(k^i, e^i) \geq 0.
\end{equation*}
We denote by $f^i(k^i, e^i) = (f^i_s(k^i, e^i))_{s \in S}$ the state-contingent output vector.
Throughout this chapter, we assume that $f^i$ is continuous and that $f^i_s(\cdot, e^i)$ is concave and non-decreasing in $k^i$ for each $e^i$.
Investors $i \in I_2$ have no productive opportunity; by convention, we set $E^i = \{0\}$, $k^i = 0$, and $f^i \equiv 0$ for $i \in I_2$.

The aggregate output vector at $t=1$, $y = (y^i)_{i \in I}$, lives in $\R^{I \times S}$; we write $y_s = (y^i_s)_{i \in I} \in \R^I$ for the cross-section of outputs in state $s$, and $y^{-i} = (y^k)_{k \neq i}$ for the outputs of all firms other than $i$'s.

\subsection{The financial market structure}

There is a finite set $J$ of financial contracts.
Security $j \in J$ is characterized by a payoff function $V_j : \R^I \to \R$, where $V_j(y_s)$ is the payoff of security $j$ in state $s \in S$ when the cross-section of outputs is $y_s$.
We write $V_j(y) = (V_j(y_s))_{s \in S} \in \R^S$ and denote by $V(y) : \R^J \to \R^S$ the linear mapping
\begin{equation*}
  V(y) z \coloneqq \sum_{j \in J} z_j V_j(y),
  \quad z = (z_j)_{j \in J} \in \R^J.
\end{equation*}
The range of $V(y)$, denoted $\mathcal{V}(y) \subset \R^S$, is the \emph{marketed subspace} at output $y$: the set of state-contingent claims at $t=1$ that can be replicated by a portfolio of traded securities.

Four classes of securities will play a distinguished role:
\begin{itemize}
  \item \emph{Bond.}
  There is one asset $j_0$ with $V_{j_0}(y) = \1_S$ for every $y$: a riskless claim to one unit of the commodity in every state.
  \item \emph{Equity.}
  For each firm $i \in I_1$, there is an equity claim $j = i$ with $V_i(y) = y^i$.
  (For $i \in I_2$, this reduces to $V_i(y) = 0$.)
  We may write $J \supset I$ to express that equity of every firm is traded.
  \item \emph{Simple options.}
  A call option on firm $i$'s equity with strike $\tau \geq 0$ is the security with payoff $V_j(y_s) = \max\{ y^i_s - \tau, 0 \}$.
  Put options are defined analogously.
  \item \emph{Indices and complex options.}
  An index is a weighted linear combination $V_j(y_s) = \sum_{i \in I} \gamma^i y^i_s$ with $\sum_i \gamma^i = 1$; a complex option is a call or put written on an index.
\end{itemize}

For each entrepreneur $i \in I_1$ it will be convenient to partition the set of securities as
\begin{equation}\label{eq:partition-J}
  J = J^i \cup J^i_{-i} \cup J_{-i},
\end{equation}
where
\begin{itemize}
  \item $J^i$ collects the securities whose payoff depends \emph{exclusively} on firm $i$'s output, e.g., equity $i$ and simple options on $i$;
  \item $J^i_{-i}$ collects the securities whose payoff depends on firm $i$'s output \emph{and} on at least one other firm's output, e.g., indices;
  \item $J_{-i}$ collects the securities whose payoff does \emph{not} depend on firm $i$'s output, e.g., the bond, the equity of other firms, and options on them.
\end{itemize}

%====================================================================
\section{Budget Sets, Preferences, and Optimal Effort}
\label{sec:MH-budgets}
%====================================================================

\subsection{Budget sets}

Let $q = (q_j)_{j \in J} \in \R^J$ be the vector of security prices at $t=0$.
In the moral hazard framework, \emph{each agent may sell short any security and buy any other}: default penalties are assumed to be so severe that all personal debts are honoured, so short positions do not involve default risk.
Formally, portfolios $z^i \in \R^J$ are unconstrained in sign.

\emph{Entrepreneur $i \in I_1$.}
Given prices $q$, an investment $k^i \geq 0$, a portfolio $z^i \in \R^J$, an effort $e^i \in E^i$, and outputs $y^{-i}$ of other firms, entrepreneur $i$'s consumption is
\begin{align}
  x^i_0
  &\coloneqq \omega^i_0 + q_i - q \cdot z^i - k^i,
  \label{eq:budget-0-MH}\\[2pt]
  x^i_s
  &\coloneqq V_s\bigl( f^i(k^i, e^i), y^{-i} \bigr) \cdot z^i
  = \sum_{j \in J} V_j\bigl( f^i_s(k^i, e^i), y^{-i}_s \bigr) z^i_j,
  \quad \text{for all $s\in S$}.
  \label{eq:budget-s-MH}
\end{align}
The date-$0$ identity \eqref{eq:budget-0-MH} reflects the entrepreneur's initial full ownership of the firm: he receives the equity price $q_i$ from selling (a possibly zero amount of) his equity, uses his wealth and the proceeds of the issued securities $-q \cdot z^i$ to finance the investment $k^i$, and consumes the residual.
The non-negativity requirement $x^i \geq 0$ is imposed for every feasible action.

\emph{Investor $i \in I_2$.}
With $k^i = 0$, $e^i = 0$, $q_i = 0$, and $f^i = 0$, the entrepreneur's budget identities collapse to
\begin{equation*}
  x^i_0 = \omega^i_0 - q \cdot z^i, \quad
  x^i_s = V_s(0, y^{-i}) \cdot z^i, \quad \text{for all $s\in S$}.
\end{equation*}

\subsection{Preferences}

Each agent $i$ has a utility function $U^i : \R_+ \times \R_+^S \times E^i \to \R$, evaluated at a consumption--effort triple $(x^i_0, x^i_1, e^i)$.
We make the following assumption throughout the chapter.

\begin{assumption}[Time-separable preferences]\label{ass:time-sep}
For every $i \in I$, $U^i$ takes the additively separable form
\begin{equation}\label{eq:time-sep}
  U^i(x^i_0, x^i_1, e^i) = u^i_0(x^i_0) + u^i_1(x^i_1, e^i),
\end{equation}
with $u^i_0 : \R_+ \to \R$ continuous, strictly increasing, and strictly concave; and $u^i_1 : \R_+^S \times E^i \to \R$ continuous.
\end{assumption}

Time-separability ensures that the entrepreneur's incentive to exert effort is driven exclusively by the effect of effort on his $t=1$ consumption $x^i_1$, and not by any complementarity between $t=0$ consumption and effort.

\subsection{The optimal effort correspondence}

A distinguishing feature of the moral hazard framework is that effort is chosen \emph{after} the financial decisions $(k^i, z^i)$ have been made, but \emph{before} the state of nature is realized.
From a timing perspective, the entrepreneur is committed to $(k^i, z^i)$ when he picks $e^i$, and has no possibility to re-trade securities after observing effort-induced changes in his expected payoff.

For a given financial decision $(k^i, z^i)$ and given outputs $y^{-i}$ of other firms, define
\begin{equation}\label{eq:opt-effort}
  \widetilde{e}^i(k^i, z^i, y^{-i})
  \coloneqq \argmax \bigl\{ u^i_1(x^i_1, e^i) \colon e^i \in E^i, \ x^i_1 = V(f^i(k^i, e^i), y^{-i}) z^i \bigr\}.
\end{equation}

\begin{assumption}[Existence of optimal effort]\label{ass:opt-effort-exists}
For every $(k^i, z^i, y^{-i})$, the correspondence $\widetilde{e}^i(\cdot)$ has non-empty values.
\end{assumption}

Because effort is chosen \emph{after} the financial decisions have been observed, every agent accepts that entrepreneur $i$ will pick some $e^i \in \widetilde{e}^i(k^i, z^i, y^{-i})$.
The financial decisions are thus committed before the state is revealed, and effort is a best response to those commitments.

%====================================================================
\section{Inference and Price Perceptions}
\label{sec:MH-inference}
%====================================================================

\subsection{Observable and unobservable variables}

Effort $e^i$ is, by assumption, unobservable to all agents other than $i$.
The capital investment $k^i$, the entrepreneur's portfolio $z^i$, and the entrepreneur's economic characteristics $(u^i_1, f^i)$ are, however, observable:

\begin{assumption}[Information]\label{ass:MH-info}
For every entrepreneur $i \in I_1$, the tuple $(u^i_1, f^i, k^i, z^i)$ is common knowledge among all agents.
\end{assumption}

Under Assumption~\ref{ass:MH-info}, investors and other entrepreneurs can \emph{infer} effort from the observation of $(k^i, z^i)$: knowing $u^i_1$, $f^i$, and the equilibrium outputs $y^{-i}$ of other firms, they solve the entrepreneur's effort problem \eqref{eq:opt-effort} and obtain $\widetilde{e}^i(k^i, z^i, y^{-i})$.
The informational asymmetry at the level of effort is compensated for by full observability at the level of financial decisions.

The empirical plausibility of this assumption relies on several institutional features noted by \citet{MagillQuinzii2002}: SEC disclosure rules requiring firms to report capital projects, managerial equity holdings, and option holdings; securities analysts who accumulate detailed knowledge of firms and their managers; and brokerage firms that disseminate this information to the investor base.

\subsection{Entrepreneurs as signallers}

Because investors infer effort from $(k^i, z^i)$ and price securities based on the anticipated output, entrepreneurs understand that their financial choices function as \emph{signals}.
A decision to retain a high equity stake signals high effort; a decision to diversify fully signals low effort.
In equilibrium, the signal is \emph{correct}: investors' inference matches the entrepreneur's actual effort, and the entrepreneur, anticipating this, makes the financial choice that best trades off incentives against risk-sharing.

\subsection{Equilibrium with price perceptions}

We formalize the equilibrium concept.
Each agent $i \in I$ forms a \emph{price perception}
\begin{equation*}
  \widetilde{Q}^i_j(k^i, z^i; \overline{y}^{-i}) \in \R,
  \quad j \in J,
\end{equation*}
that assigns, for each alternative financial decision $(k^i, z^i)$ and given the equilibrium outputs $\overline{y}^{-i}$ of other firms, the price that the market would pay for security $j$.
The collection $\widetilde{Q}^i = (\widetilde{Q}^i_j)_{j \in J}$ is agent $i$'s price perception.

\begin{definition}[Stock-market equilibrium with price perceptions]\label{def:eq-PP}
Given a profile $\widetilde{Q} = (\widetilde{Q}^i)_{i \in I}$ of price perceptions, a \emph{stock-market equilibrium with price perceptions $\widetilde{Q}$} is a list
\begin{equation*}
  \bigl\{ (\overline{x}^i, \overline{e}^i, \overline{k}^i, \overline{z}^i)_{i \in I}, \, (\overline{y}^j)_{j \in J}, \, \overline{q} \bigr\}
\end{equation*}
such that:
\begin{enumerate}
  \item[(a)] \emph{Agents are rational.}
  For each $i \in I$, $(\overline{x}^i, \overline{e}^i, \overline{k}^i, \overline{z}^i)$ maximizes $U^i(x^i_0, x^i_1, e^i)$ subject to
  \begin{align*}
    x^i_0 &= \omega^i_0 + \widetilde{Q}^i_i(k^i, z^i; \overline{y}^{-i}) - \widetilde{Q}^i(k^i, z^i; \overline{y}^{-i}) \cdot z^i - k^i, \\
    x^i_1 &= V(f^i(k^i, e^i), \overline{y}^{-i}) z^i, \\
    e^i &\in \widetilde{e}^i(k^i, z^i, \overline{y}^{-i}),
  \end{align*}
  with $x^i_0 \geq 0$ and $x^i_1 \geq 0$.
  \item[(b)] \emph{Production is compatible with effort.}
  For every firm $i \in I_1$, $\overline{y}^i = f^i(\overline{k}^i, \overline{e}^i)$.
  \item[(c)] \emph{Price perceptions are consistent at equilibrium.}
  For every agent $i \in I$, $\overline{q} = \widetilde{Q}^i(\overline{k}^i, \overline{z}^i; \overline{y}^{-i})$.
  \item[(d)] \emph{Security markets clear.}
  Using the convention that an investor $i \in I_2$ holds his ``own'' equity with $z^i_i = 1$ by default (so that the notation is uniform),
  \begin{equation*}
    \forall j \in I, \quad \sum_{i \in I} \overline{z}^i_j = 1
    \et
    \forall j \in J \setminus I, \quad \sum_{i \in I} \overline{z}^i_j = 0.
  \end{equation*}
\end{enumerate}
\end{definition}

The budget set of agent $i$ at prices $(\overline{q}, \overline{y}^{-i})$ and under the price perception $\widetilde{Q}^i$ will be denoted $B^i(\overline{q}, \overline{y}^{-i})$.

Condition~(a) says that each agent correctly anticipates his optimal effort response to a hypothetical financial choice $(k^i, z^i)$ and evaluates the resulting consumption using his price perception $\widetilde{Q}^i$.
Condition~(b) ensures internal consistency between the firm's production and the effort that rational investors would attribute to the firm given its financial decisions.
Condition~(c) says that the price perceptions, evaluated at the \emph{equilibrium} decisions, reproduce the equilibrium prices---this is the rational-expectations requirement.
Condition~(d) is standard market clearing.

\subsection{The challenge of pricing out-of-equilibrium decisions}

Definition~\ref{def:eq-PP} leaves unspecified how investors form price perceptions off the equilibrium path.
If entrepreneur $i$ contemplates a deviation from $(\overline{k}^i, \overline{z}^i)$ to $(k^i, z^i)$, what price $\widetilde{Q}^i_j(k^i, z^i; \overline{y}^{-i})$ will he attach to each security?
The next section specifies this mapping via Makowski-style rational competitive price perceptions, adapted to the moral-hazard setting.

%====================================================================
\section{Rational Competitive Price Perceptions and Partial Spanning}
\label{sec:RCPP}
%====================================================================

\subsection{Out-of-equilibrium pricing via state-prices}

At equilibrium, no arbitrage in the existing set of securities implies the existence of a vector of state prices $\overline{\pi} \in \R_{++}^S$ such that
\begin{equation}\label{eq:state-prices}
  \overline{q}_j = \sum_{s \in S} \overline{\pi}_s V_j(\overline{y}_s),
  \quad \forall j \in J.
\end{equation}
If the security structure $V(\overline{y})$ spans $\R^S$, then $\overline{\pi}$ is unique; otherwise, the set of state prices compatible with $(\overline{q}, \overline{y})$ is generally a positive-dimensional polyhedron.

Let $\mathcal{V}(\overline{y}) \subset \R^S$ be the marketed subspace at equilibrium.
For any $m \in \mathcal{V}(\overline{y})$, the quantity
\begin{equation*}
  \overline{\pi} \cdot m = \sum_{s \in S} \overline{\pi}_s m_s
\end{equation*}
is independent of the particular state-price vector $\overline{\pi}$ chosen: two state prices $\overline{\pi}, \overline{\pi}'$ that disagree on $\R^S$ must agree on $\mathcal{V}(\overline{y})$, since their difference lies in the orthogonal complement of the rows of $V(\overline{y})$.
We denote this common value by $v(m; \overline{q}, \overline{y})$.

\begin{proposition}[Pricing of marketed innovations]\label{prop:marketed-pricing}
Assume that entrepreneur $i$ modifies his financial decision from $(\overline{k}^i, \overline{z}^i)$ to $(k^i, z^i)$, and that the resulting output $f^i(k^i, \widehat{e}^i)$ (with $\widehat{e}^i \in \widetilde{e}^i(k^i, z^i, \overline{y}^{-i})$) is such that the candidate new cash flow $V_j(f^i(k^i, \widehat{e}^i), \overline{y}^{-i})$ lies in $\mathcal{V}(\overline{y})$ for every $j \in J$.
Then the original equilibrium, augmented by the ``new'' securities at prices
\begin{equation*}
  \widetilde{Q}^i_j(k^i, z^i; \overline{y}^{-i}) \coloneqq \overline{\pi} \cdot V_j\bigl( f^i(k^i, \widehat{e}^i), \overline{y}^{-i} \bigr),
\end{equation*}
and held in zero net supply, remains a competitive equilibrium.
In particular, $\widetilde{Q}^i_j(k^i, z^i; \overline{y}^{-i})$ is independent of the choice of state-price vector $\overline{\pi}$ compatible with $(\overline{q}, \overline{y})$.
\end{proposition}

The proposition motivates the following definition, introduced by \citet{MagillQuinzii2002}.

\subsection{Partial Spanning}

\begin{definition}[Partial spanning]\label{def:PS}
There is \emph{partial spanning} (PS) at $\overline{y}$ if, for every entrepreneur $i \in I$, every alternative $(k^i, e^i)$ with associated output $y^i = f^i(k^i, e^i)$, and every $j \in J$,
\begin{equation*}
  V_j(y^i, \overline{y}^{-i}) \in \mathcal{V}(\overline{y}).
\end{equation*}
Equivalently, for every alternative output vector $y = (y^i, \overline{y}^{-i})$,
\begin{equation}\label{eq:PS-spanning}
  \mathcal{V}(y) \subset \mathcal{V}(\overline{y}).
\end{equation}
\end{definition}

Partial spanning is the exact condition under which a unilateral deviation by entrepreneur $i$ cannot create a genuinely new cash flow: every alternative security payoff lies in the span of the \emph{equilibrium} market structure, and can therefore be priced unambiguously at equilibrium state prices.

\subsection{Rational Competitive Price Perceptions}

\begin{definition}[RCPP equilibrium]\label{def:RCPP}
A stock-market equilibrium $\bigl\{ (\overline{x}^i, \overline{e}^i, \overline{k}^i, \overline{z}^i)_{i \in I}, \, (\overline{y}^j)_{j \in J}, \, \overline{q} \bigr\}$ with price perceptions $\widetilde{Q}$ is a \emph{rational competitive price perceptions} (RCPP) equilibrium if:
\begin{enumerate}
  \item[(i)] partial spanning holds at $\overline{y}$;
  \item[(ii)] for every agent $i \in I$ and every alternative $(k^i, z^i)$,
  \begin{equation}\label{eq:RCPP-perception}
    \widetilde{Q}^i_j(k^i, z^i; \overline{y}^{-i})
    = \overline{\pi} \cdot V_j\bigl( f^i(k^i, \widehat{e}^i), \overline{y}^{-i} \bigr),
  \end{equation}
  for any $\overline{\pi} \in \R_{++}^S$ satisfying \eqref{eq:state-prices} and any effort $\widehat{e}^i \in \widetilde{e}^i(k^i, z^i, \overline{y}^{-i})$ that maximizes entrepreneur $i$'s anticipated date-$0$ consumption.
\end{enumerate}
\end{definition}

Two aspects of Definition~\ref{def:RCPP} deserve emphasis.
The adjective \emph{rational} refers to the agents' ability to correctly deduce $\widehat{e}^i$ from $(k^i, z^i)$ and from the common knowledge of the entrepreneur's characteristics $(u^i_1, f^i)$: if $\widetilde{e}^i(k^i, z^i, \overline{y}^{-i})$ is a singleton, the inference is immediate; if it is multi-valued, the literature has argued that agents should select the element of the optimal-effort set that yields the highest date-$0$ income for the entrepreneur, on the grounds that the entrepreneur has every incentive to signal the highest-valuation effort level.
The adjective \emph{competitive} refers to the assumption that entrepreneurs take state prices as given: an individual firm's deviation, however large it may be relative to that firm's own output, does not perturb the equilibrium state-price vector.
This is the counterpart, in the moral-hazard framework, of the price-taking hypothesis invoked in Chapters~\ref{ch:objective}--\ref{ch:price-perceptions}.

\subsection{Sufficient conditions for Partial Spanning: a single-factor model}

Partial spanning is, in general, a non-trivial condition to verify.
Two examples identify natural security structures under which it holds.

\begin{proposition}[Bond, equity, and single-factor risk]\label{prop:PS-single-factor}
Assume that:
\begin{enumerate}
  \item[\textup{(a)}] the financial market consists solely of the bond and the equity of each firm, i.e., $J = I \cup \{j_0\}$;
  \item[\textup{(b)}] for every $i \in I_1$, the production function has the single-factor structure
  \begin{equation}\label{eq:single-factor}
    f^i(k^i, e^i) = f^i_c(k^i, e^i) \1_S + g^i(k^i, e^i) \bm{\eta}^i,
  \end{equation}
  where $f^i_c, g^i : \R_+ \times E^i \to \R$ are concave, non-decreasing functions and $\bm{\eta}^i \in \R_+^S$ is a fixed vector characterizing the firm's risk structure.
\end{enumerate}
If, at equilibrium, $g^i(\overline{k}^i, \overline{e}^i) > 0$ for every $i \in I_1$, then partial spanning holds at $\overline{y}$.
\end{proposition}

\begin{proof}[Proof sketch]
Under~(b), firm $i$'s output decomposes into a deterministic component $f^i_c(k^i, e^i) \1_S$, spanned by the bond, and a risky component proportional to $\bm{\eta}^i$.
At equilibrium, the marketed subspace $\mathcal{V}(\overline{y})$ contains the bond $\1_S$ and, for each $i$ with $\overline{g}^i > 0$, the vector $\bm{\eta}^i = [\overline{y}^i - \overline{f}^i_c \1_S] / \overline{g}^i$.
Any alternative output $f^i(k^i, e^i) = f^i_c(k^i, e^i) \1_S + g^i(k^i, e^i) \bm{\eta}^i$ is a linear combination of these two spanned vectors, hence in $\mathcal{V}(\overline{y})$.
A similar argument applies to any security $V_j$ whose payoff is a linear function of firm-$i$'s output.
\end{proof}

\subsection{Sufficient conditions for Partial Spanning: independent shocks and options}

A second, empirically important example is based on an independence structure of the uncertainty plus a rich family of options.

\begin{proposition}[Independent shocks with options]\label{prop:PS-options}
Assume that:
\begin{enumerate}
  \item[\textup{(a)}] the financial market consists of the bond, of each firm's equity, and of call and put options on each firm's equity at a rich set of strike prices;
  \item[\textup{(b)}] the state space factorizes as $S = \prod_{i \in I_1} S^i$, and there exist functions $g^i_{s^i} : \R_+ \times E^i \to \R$, one for each $i \in I_1$ and each $s^i \in S^i$, such that
  \begin{equation*}
    f^i_s(k^i, e^i) = g^i_{s^i}(k^i, e^i),
    \quad s = (s^\ell)_{\ell \in I_1}.
  \end{equation*}
\end{enumerate}
If, at equilibrium, the mapping $s^i \mapsto g^i_{s^i}(\overline{k}^i, \overline{e}^i)$ is injective for every $i \in I_1$, and if the set of strike prices includes values strictly between every pair of consecutive realizations of firm-$i$'s output, then partial spanning holds at $\overline{y}$.
\end{proposition}

The essence of the proposition is that, under (b), firm $i$'s output depends only on the ``own'' coordinate $s^i$ of the state.
Knowing the equilibrium output $\overline{y}^i_s = g^i_{s^i}(\overline{k}^i, \overline{e}^i)$ for each $s$, and the injectivity of $g^i_{s^i}(\overline{k}^i, \overline{e}^i)$ in $s^i$, one can identify each value of $s^i$ from the equilibrium output.
The rich set of strike prices then generates, from equity plus options, a complete set of state-contingent claims on firm $i$'s coordinate $s^i$.
Aggregating across $i$, the marketed subspace $\mathcal{V}(\overline{y})$ coincides with $\R^S$, so partial spanning holds trivially.

\subsection{An illustrative counter-example}

To appreciate the role of Partial Spanning, consider an economy with one entrepreneur and three states of nature, $S = \{s_1, s_2, s_3\}$.
The effort set is $E^i = \{e^i_h, e^i_\ell\}$ (high and low) and production does not depend on capital; assume
\begin{equation*}
  f^i(e^i_h) = (y^i_h, y^i_h, y^i_\ell)
  \et
  f^i(e^i_\ell) = (y^i_h, y^i_\ell, y^i_\ell),
\end{equation*}
with $y^i_h > y^i_\ell > 0$.
Suppose the security structure consists of the bond, equity, and simple options on equity.

At an equilibrium where the entrepreneur exerts high effort, $\overline{e}^i = e^i_h$, the equity payoff is $(y^i_h, y^i_h, y^i_\ell)$, and the marketed subspace is contained in $\{ w \in \R^3 \colon w_1 = w_2 \}$: every security written on the equity is constant across states $s_1$ and $s_2$.
If the entrepreneur considers deviating to low effort, the resulting equity payoff is $(y^i_h, y^i_\ell, y^i_\ell)$, and the vector $(y^i_h - y^i_\ell, 0, 0)$ is \emph{not} in the equilibrium marketed subspace.
Partial Spanning fails at $\overline{y}$, and the RCPP price perception is ill-defined on the deviating cash flow.

%====================================================================
\section{Constrained Pareto Efficiency}
\label{sec:MH-CPO}
%====================================================================

\subsection{Constrained feasible allocations}

To evaluate the welfare properties of an RCPP equilibrium, we compare it to an allocation chosen by a planner who faces the same informational and financial frictions as the agents: the planner cannot contract on effort, cannot enrich the security structure, and must respect market clearing.

\begin{definition}[Constrained feasible allocation]\label{def:MH-feasible}
An allocation $(\bm{x}, \bm{e}) = (x^i, e^i)_{i \in I}$ is \emph{constrained feasible} if there exist inputs and portfolios $(\bm{k}, \bm{z}) = (k^i, z^i)_{i \in I}$ such that:
\begin{enumerate}
  \item[(i)] \emph{Production compatibility}: $y^i = f^i(k^i, e^i)$ for every $i \in I_1$ (and $y^i = 0$ for $i \in I_2$);
  \item[(ii)] \emph{Optimal effort}: $e^i \in \widetilde{e}^i(k^i, z^i, y^{-i})$ for every $i \in I_1$;
  \item[(iii)] \emph{Risk sharing via the asset structure}: $x^i_s = V_s(y) \cdot z^i$ for every $i$ and every $s$;
  \item[(iv)] \emph{Date-$0$ market clearing}: $\sum_{i \in I} (x^i_0 + k^i) = \sum_{i \in I} \omega^i_0$;
  \item[(v)] \emph{Portfolio market clearing}: $\sum_{\ell \in I} z^\ell_j = 1$ for $j \in I$ (equity) and $\sum_{\ell \in I} z^\ell_j = 0$ for $j \in J \setminus I$;
  \item[(vi)] \emph{Non-negativity}: $x^i \geq 0$ for every $i$.
\end{enumerate}
\end{definition}

Conditions (ii) and (v) are the distinguishing features of the constrained notion.
The planner cannot dictate effort to entrepreneur $i$: having chosen $(k^i, z^i)$, he must accept whatever effort the entrepreneur rationally responds with.
And he cannot re-open financial markets: portfolios must clear within the existing set of securities, at unit supply for equities and zero net supply for bonds, options, and indices.

\begin{definition}[Constrained Pareto optimality]\label{def:MH-CPO}
A constrained feasible allocation $(\bm{x}, \bm{e})$ is \emph{constrained Pareto optimal} (CPO) if there does not exist another constrained feasible allocation $(\widehat{\bm{x}}, \widehat{\bm{e}})$ such that
\begin{equation*}
  U^i(\widehat{x}^i, \widehat{e}^i) \geq U^i(x^i, e^i) \quad \forall i \in I,
\end{equation*}
with strict inequality for at least one $i$.
\end{definition}

\subsection{Why an RCPP equilibrium may fail to be CPO}

Two sources of potential inefficiency distinguish the moral-hazard framework from the frictionless environment of Chapter~\ref{ch:price-perceptions}.
\begin{enumerate}
  \item \emph{Externalities through effort.}
  The choice of effort by entrepreneur $i$ affects the payoff of every security written on firm $i$'s output, and hence the consumption of every agent who holds those securities.
  In the decentralized equilibrium, the entrepreneur chooses $e^i$ to maximize his \emph{own} utility, ignoring the impact on other security holders; a planner would internalize this external effect.
  \item \emph{Nash-like interactions through outputs.}
  An RCPP equilibrium is, implicitly, a Nash equilibrium of a game in which each entrepreneur $i$ chooses $(k^i, z^i)$ taking the other firms' decisions as given.
  A planner's choice of $(k, z)$ may co-ordinate entrepreneurs on a better vector of effort levels than a Nash equilibrium would sustain.
\end{enumerate}
These channels suggest that the welfare property is not automatic.
The next section identifies a spanning condition---\emph{strong} partial spanning---under which both channels are neutralized, and the RCPP equilibrium delivers constrained efficiency.

%====================================================================
\section{Strong Partial Spanning and the Welfare Theorem}
\label{sec:MH-welfare}
%====================================================================

\subsection{No-index securities and the inside--outside decomposition}

A first assumption used to obtain the welfare result is that the security structure contains no index-based securities.

\begin{assumption}[No index-based securities]\label{ass:no-index}
For every entrepreneur $i \in I_1$, $J^i_{-i} = \emptyset$, i.e., every security in $J$ either depends exclusively on firm $i$'s output or does not depend on it at all.
\end{assumption}

Under Assumption~\ref{ass:no-index}, an investor's or entrepreneur's date-$1$ income decomposes cleanly into two parts.
For every agent $i \in I$,
\begin{equation}\label{eq:inside-outside}
  x^i_1 = m^i_1 + \sum_{j \in J^i} V_j(f^i(k^i, e^i)) z^i_j,
  \quad
  m^i_1 \coloneqq \sum_{j \in J_{-i}} V_j(y^{-i}) z^i_j.
\end{equation}
The stream $m^i_1$ is called the \emph{outside income} of agent $i$: it does not depend on $i$'s own output $y^i$ (and, for $i \in I_1$, does not depend on $i$'s effort).
The sum over $J^i$ is the \emph{inside income}: it depends only on firm $i$'s output.
When $i \in I_2$ is an investor, $J^i = \emptyset$ (or its inside-income component is identically zero), and all of $i$'s date-$1$ income is outside.

\subsection{Strong Partial Spanning}

Partial Spanning guarantees that any single firm's alternative output stays within the marketed subspace at $\overline{y}$.
We now introduce a stronger condition on the \emph{outside} part of the security structure.
For every $i \in I$, let $V_{-i}(y^{-i})$ be the linear operator $\R^{J_{-i}} \to \R^S$ associated with the matrix $(V_j(y^{-i}))_{j \in J_{-i}}$, and denote its range by $\mathcal{V}_{-i}(y^{-i})$, the \emph{outside subspace}.

\begin{definition}[Strong Partial Spanning]\label{def:SPS}
There is \emph{strong partial spanning} (SPS) at $\overline{y}$ if, for every choice of investments and efforts $(\bm{k}, \bm{e})$ with outputs $y^\ell = f^\ell(k^\ell, e^\ell)$ for $\ell \in I_1$,
\begin{equation}\label{eq:SPS}
  \forall i \in I_1, \quad \mathcal{V}_{-i}(y^{-i}) \subset \mathcal{V}_{-i}(\overline{y}^{-i}).
\end{equation}
\end{definition}

SPS requires that any change in the outputs of firms other than $i$ keeps the outside subspace within its equilibrium span.
In other words, the choices of effort by firms $\ell \neq i$ may alter the outside subspace $\mathcal{V}_{-i}(y^{-i})$, but only within the equilibrium span.
This rules out the situation in which a coordinated deviation by all firms other than $i$ creates new outside cash flows that were not marketed at the equilibrium.

\begin{remark}\label{rem:SPS-examples}
Three comments on SPS are in order.
\begin{itemize}
  \item Even if markets are complete in the sense $\mathcal{V}(\overline{y}) = \R^S$, SPS need not hold automatically, because $\mathcal{V}_{-i}(y^{-i})$ is strictly smaller than $\mathcal{V}(y)$.
  \item SPS holds if, for each entrepreneur $i$, the securities written on the outputs of the remaining $I_1 \setminus \{i\}$ firms already suffice to complete the markets.
  \item SPS holds if each firm ``spans its own subspace'' in the sense of the two examples of Section~\ref{sec:RCPP}: under the single-factor or the independent-shocks-with-options structure, the outside subspace is unaffected by changes in other firms' outputs, so \eqref{eq:SPS} holds trivially.
\end{itemize}
\end{remark}

\subsection{The welfare theorem}

\begin{theorem}[Constrained Pareto optimality of RCPP]\label{thm:MH-welfare}
Let
\begin{equation*}
  \bigl\{ (\overline{x}^i, \overline{e}^i, \overline{k}^i, \overline{z}^i)_{i \in I}, \, (\overline{y}^j)_{j \in J}, \, \overline{q} \bigr\}
\end{equation*}
be an RCPP equilibrium with price perceptions $\widetilde{Q}$ defined by \eqref{eq:RCPP-perception}.
Suppose that Assumptions~\ref{ass:time-sep}, \ref{ass:opt-effort-exists}, \ref{ass:MH-info}, \ref{ass:no-index} hold, and that strong partial spanning is satisfied at $\overline{y}$.
Then the allocation $(\overline{x}^i, \overline{e}^i)_{i \in I}$ is constrained Pareto optimal.
\end{theorem}

The result is remarkable: despite the moral-hazard friction, the decentralized equilibrium fully internalizes the externalities through the signalling channel.
The entrepreneur, rationally anticipating that investors will read his financial decisions as signals of his effort, is led to choose precisely the effort level that the planner would have wanted him to choose.

\subsection{Proof of Theorem~\ref{thm:MH-welfare}}

Fix an RCPP equilibrium $\bigl\{ (\overline{x}^i, \overline{e}^i, \overline{k}^i, \overline{z}^i)_{i \in I}, \, (\overline{y}^j)_{j \in J}, \, \overline{q} \bigr\}$ and let $\overline{\pi} \in \R_{++}^S$ be a state-price vector satisfying \eqref{eq:state-prices}.
By Partial Spanning and the remarks of Section~\ref{sec:RCPP}, the value $\overline{\pi} \cdot m$ is independent of the chosen $\overline{\pi}$ whenever $m \in \mathcal{V}(\overline{y})$.

Suppose, for contradiction, that there exists a constrained feasible allocation $(\bm{x}, \bm{e})$ with supporting $(\bm{k}, \bm{z})$ such that
\begin{equation}\label{eq:cpo-dominance}
  U^i(x^i, e^i) \geq U^i(\overline{x}^i, \overline{e}^i), \quad \forall i \in I,
\end{equation}
with strict inequality for at least one $i^\ast \in I$.

\paragraph{Step 1: Constructing a replicating portfolio.}
For each $i \in I$, define the alternative decision $(\widehat{k}^i, \widehat{z}^i)$ by $\widehat{k}^i = k^i$, $\widehat{z}^i_j = z^i_j$ for $j \in J^i$, and $(\widehat{z}^i_j)_{j \in J_{-i}}$ chosen such that
\begin{equation}\label{eq:ziwidehat}
  \sum_{j \in J_{-i}} V_j(\overline{y}^{-i}) \widehat{z}^i_j = \sum_{j \in J_{-i}} V_j(y^{-i}) z^i_j = m^i_1.
\end{equation}
The choice \eqref{eq:ziwidehat} is feasible by Assumption~\ref{ass:no-index} and by SPS at $\overline{y}$: the right-hand side belongs to $\mathcal{V}_{-i}(y^{-i}) \subset \mathcal{V}_{-i}(\overline{y}^{-i})$, so a portfolio $\widehat{z}^i$ of \emph{outside} securities exists that replicates $m^i_1$ at equilibrium outputs.

\paragraph{Step 2: Invariance of the optimal effort.}

\begin{lemma}[Effort invariance]\label{lem:effort-invariance}
$\widetilde{e}^i(\widehat{k}^i, \widehat{z}^i, \overline{y}^{-i}) = \widetilde{e}^i(k^i, z^i, y^{-i})$, in the sense that both sets coincide as subsets of $E^i$.
\end{lemma}

\begin{proof}
For every $e^i \in E^i$, the date-$1$ consumption of entrepreneur $i$ under the alternative $(\widehat{k}^i, \widehat{z}^i)$ and at equilibrium outputs $\overline{y}^{-i}$ is
\begin{align*}
  \widehat{x}^i_1
  &= V(f^i(\widehat{k}^i, e^i), \overline{y}^{-i}) \widehat{z}^i
  = \sum_{j \in J^i} V_j(f^i(k^i, e^i)) z^i_j + \sum_{j \in J_{-i}} V_j(\overline{y}^{-i}) \widehat{z}^i_j \\
  &= \sum_{j \in J^i} V_j(f^i(k^i, e^i)) z^i_j + m^i_1
  = x^i_1,
\end{align*}
where the last equality uses \eqref{eq:inside-outside} and \eqref{eq:ziwidehat}, and the second uses Assumption~\ref{ass:no-index} to split the sum.
Since the induced date-$1$ consumption is the same as in the constrained-feasible allocation, and $u^i_1$ depends only on $x^i_1$ and $e^i$, the objective maximized over $e^i$ is identical in the two cases, and the argmax sets coincide.
\end{proof}

\paragraph{Step 3: Computing the anticipated date-$0$ consumption.}
Let $\widehat{e}^i$ be an effort choice in the common argmax, selected as in Definition~\ref{def:RCPP} (if multi-valued, the one maximizing date-$0$ consumption).
By the RCPP price perception \eqref{eq:RCPP-perception} and the consistency of $\overline{\pi}$ on the marketed subspace,
\begin{equation*}
  \widetilde{Q}^i_i(\widehat{k}^i, \widehat{z}^i; \overline{y}^{-i})
  = \overline{\pi} \cdot V_i(f^i(\widehat{k}^i, \widehat{e}^i), \overline{y}^{-i})
  = \overline{\pi} \cdot f^i(\widehat{k}^i, \widehat{e}^i)
  = \overline{\pi} \cdot f^i(k^i, \widehat{e}^i),
\end{equation*}
using $V_i(y) = y^i$ and $\widehat{k}^i = k^i$.
Similarly,
\begin{equation*}
  \widetilde{Q}^i(\widehat{k}^i, \widehat{z}^i; \overline{y}^{-i}) \cdot \widehat{z}^i
  = \sum_{j \in J} \overline{\pi} \cdot V_j(f^i(\widehat{k}^i, \widehat{e}^i), \overline{y}^{-i}) \widehat{z}^i_j
  = \overline{\pi} \cdot V(f^i(\widehat{k}^i, \widehat{e}^i), \overline{y}^{-i}) \widehat{z}^i
  = \overline{\pi} \cdot \widehat{x}^i_1.
\end{equation*}
Combining with the budget identity,
\begin{equation}\label{eq:x-hat-0}
  \widehat{x}^i_0
  = \omega^i_0 + \overline{\pi} \cdot f^i(k^i, \widehat{e}^i) - \overline{\pi} \cdot \widehat{x}^i_1 - k^i.
\end{equation}
Applied with $\widehat{e}^i = e^i$ (which is in the common argmax by Lemma~\ref{lem:effort-invariance}), and using $\widehat{x}^i_1 = x^i_1$,
\begin{equation}\label{eq:x-hat-0-final}
  \widehat{x}^i_0
  = \omega^i_0 + \overline{\pi} \cdot y^i - \overline{\pi} \cdot x^i_1 - k^i,
\end{equation}
with the convention $y^i = 0$ and $k^i = 0$ for investors $i \in I_2$.

\paragraph{Step 4: Individual comparison.}
Consider the alternative $(\widehat{k}^i, \widehat{z}^i, \widehat{x}^i_0, x^i_1, e^i)$ constructed above.
It lies in the RCPP budget set $B^i(\overline{q}, \overline{y}^{-i})$: the budget constraints hold by construction, the effort $e^i$ is in the optimal-effort set by Lemma~\ref{lem:effort-invariance}, and $\widehat{x}^i \geq 0$ follows either automatically (for $\widehat{x}^i_1 = x^i_1 \geq 0$) or from the non-negativity constraint attached to the RCPP optimization.
By RCPP-optimality of $(\overline{x}^i, \overline{e}^i)$,
\begin{equation*}
  U^i(\overline{x}^i, \overline{e}^i) \geq U^i(\widehat{x}^i_0, x^i_1, e^i).
\end{equation*}
Combining with \eqref{eq:cpo-dominance},
\begin{equation*}
  U^i(x^i_0, x^i_1, e^i) \geq U^i(\widehat{x}^i_0, x^i_1, e^i),
\end{equation*}
and by time-separability \eqref{eq:time-sep} and strict monotonicity of $u^i_0$,
\begin{equation}\label{eq:x0-inequality}
  x^i_0 \geq \widehat{x}^i_0, \quad \forall i \in I.
\end{equation}
For $i^\ast$ with strict utility improvement, the inequality is strict: $x^{i^\ast}_0 > \widehat{x}^{i^\ast}_0$.

\paragraph{Step 5: Aggregation and contradiction.}
Sum \eqref{eq:x0-inequality} over $i \in I$:
\begin{equation*}
  \sum_{i \in I} x^i_0
  > \sum_{i \in I} \widehat{x}^i_0
  = \sum_{i \in I} \omega^i_0 + \sum_{i \in I} \overline{\pi} \cdot y^i - \sum_{i \in I} \overline{\pi} \cdot x^i_1 - \sum_{i \in I} k^i,
\end{equation*}
where the strict inequality comes from the $i^\ast$-component.
Rearranging,
\begin{equation}\label{eq:aggregate-ineq}
  \sum_{i \in I} (x^i_0 + k^i) + \overline{\pi} \cdot \sum_{i \in I} x^i_1
  > \sum_{i \in I} \omega^i_0 + \overline{\pi} \cdot \sum_{i \in I} y^i.
\end{equation}
By date-$0$ market clearing in the constrained-feasible allocation, $\sum_i (x^i_0 + k^i) = \sum_i \omega^i_0$, so the inequality becomes
\begin{equation*}
  \overline{\pi} \cdot \sum_{i \in I} x^i_1
  > \overline{\pi} \cdot \sum_{i \in I} y^i.
\end{equation*}
By portfolio market clearing,
\begin{equation*}
  \sum_{i \in I} x^i_s = \sum_{i \in I} V_s(y) \cdot z^i
  = \sum_{j \in J} V_j(y_s) \sum_{i \in I} z^i_j
  = \sum_{j \in I} V_j(y_s) \cdot 1 + \sum_{j \in J \setminus I} V_j(y_s) \cdot 0
  = \sum_{j \in I} y^j_s
  = \sum_{i \in I} y^i_s,
\end{equation*}
so $\sum_i x^i_1 = \sum_i y^i$ as vectors in $\R^S$, and $\overline{\pi} \cdot \sum_i x^i_1 = \overline{\pi} \cdot \sum_i y^i$.
Substituting, \eqref{eq:aggregate-ineq} becomes
\begin{equation*}
  0 > 0,
\end{equation*}
a contradiction.
The initial assumption \eqref{eq:cpo-dominance} is impossible, and the RCPP allocation is constrained Pareto optimal.
\hfill\qedsymbol

\subsection{Interpretation}

Theorem~\ref{thm:MH-welfare} rests on the signal-extraction property of rational competitive price perceptions.
The entrepreneur's equity and option holdings act as a \emph{credible signal} of his effort choice: any change in the financial portfolio is correctly interpreted by investors, who re-price the firm's securities accordingly.
The entrepreneur, anticipating this, cannot use financial decisions to shirk effort without simultaneously suffering the full market value of the incentive deterioration.
The externality at the heart of the moral-hazard problem---the fact that the entrepreneur's effort affects all holders of securities written on his firm---is thereby \emph{internalized} by the entrepreneur himself, through the price mechanism.

It is instructive to compare this result to the frictionless welfare theorem of Chapter~\ref{ch:price-perceptions}: both rest on Makowski-style price perceptions (highest valuation / rational inference via state prices) plus a spanning condition (partial / strong partial spanning).
The mechanism is identical.
Moral hazard changes the inferential content of financial choices---they now convey information about effort rather than simply shifting risk---but does not change the welfare logic of the competitive equilibrium.

%====================================================================
\section{Discussion and Limitations}
\label{sec:MH-limitations}
%====================================================================

Theorem~\ref{thm:MH-welfare} is a clean, constructive welfare result for a model of corporate finance under moral hazard.
Its scope is, however, bounded by the conditions it invokes.
Two lines of follow-up work have sharpened the boundaries.

\subsection{Dispersed initial ownership}

In Magill--Quinzii's model, each entrepreneur enters with \emph{full initial ownership} of his firm, $\xi^i(i) = 1$ and $\xi^i(\ell) = 0$ for $\ell \neq i$.
This feature is crucial: the entrepreneur sells (or retains) his own shares in the stock market, and the market prices reflect the inference drawn from this decision.
\citet{CalcagnoWagner2006} show that the welfare result does not generalize to economies in which the firm's shares are \emph{initially dispersed} among several agents: when the entrepreneur does not start with full ownership of the firm, his incentive to signal effort through portfolio choice is attenuated, and the RCPP equilibrium may fail to achieve constrained Pareto optimality.

\subsection{Flexible technology}

A second qualification concerns the structure of the production function.
Magill--Quinzii's framework and the subsequent literature typically assume a \emph{single-input technology}: the entrepreneur chooses a capital level $k^i \geq 0$ and an effort $e^i \in E^i$, with output $f^i_s(k^i, e^i)$.
\citet{QuigginChambers2006} consider a more general \emph{flexible-technology} specification, in which the production set is given by a multi-input concave production possibilities frontier, and show that constrained Pareto optimality in general requires the single-input structure: with multi-dimensional productive inputs, entrepreneurs may use their portfolio decisions to signal a composite of inputs that does not coincide with the socially optimal composition, and the welfare property breaks down.

\subsection{Earlier and contemporary developments}

The model of \citet{MagillQuinzii2002} builds on the pioneering work of \citet{KihlstromMatthews1990}, who developed a ``managerial incentives'' stock-market model in which the entrepreneur's portfolio choice acts as an incentive device.
A related contribution in the macroeconomic tradition is \citet{Kocherlakota1998}, who examines how the presence of moral hazard affects asset prices when financial markets are complete.
In incomplete-markets economies, the interaction between agency frictions and equilibrium pricing has been explored further by \citet{BisinGottardiRampini2008} and \citet{AcharyaBisin2009}.

\subsection{Further reading}

The definitive reference for the material of this chapter is \citet{MagillQuinzii2002}, on which the exposition is based.
For earlier conceptual foundations on the role of financial decisions as signals, see \citet{Ross1977} and \citet{LelandPyle1977}, and the agency-cost approach of \citet{JensenMeckling1976} already discussed in Chapter~\ref{ch:capital-structure}.
The extensions and limitations of Magill--Quinzii's welfare result are developed by \citet{CalcagnoWagner2006} and \citet{QuigginChambers2006}.
For broader surveys of moral hazard in financial markets and its implications for corporate finance, see \citet{Tirole2006} and the collected references therein.

%====================================================================
\chapter{Moral Hazard with Competitive Price Perceptions}
\label{ch:moral-hazard-cpp}
%====================================================================

Chapters~\ref{ch:price-perceptions} and~\ref{ch:moral-hazard} have developed two distinct frameworks for the stock-market equilibrium in the presence of frictions.
\citet{BisinGottardiRuta2015} study the Makowski--style pricing and welfare properties of an economy with short-sale constraints and no bankruptcy, under full information.
\citet{MagillQuinzii2002} study the welfare properties of rational competitive price perceptions in an economy with moral hazard and a rich security structure, under unlimited liability.
Neither analysis confronts simultaneously the two frictions most salient in real-world corporate finance: the constraint that individual investors cannot short-sell equity and the moral hazard created by an entrepreneur whose effort is private information.

This chapter presents a unified framework that combines the two frictions.
The environment is a short-sale-constrained Bisin--Gottardi--Ruta economy as in Chapter~\ref{ch:price-perceptions}---one firm, one riskless bond, no short selling by investors---augmented with:
\begin{itemize}
  \item a production technology $f_s(k, e)$ depending on the manager's unobservable effort $e$, as in Chapter~\ref{ch:moral-hazard};
  \item a bankruptcy mechanism that makes the bond \emph{risky}: if the firm's cash flow $f_s(k, e)$ falls short of the promised debt repayment $\beta$, the firm defaults, bondholders receive all of $f_s(k, e)$, and shareholders receive nothing;
  \item a principal--agent structure in which the firm hires a manager among the pool of investors and offers him a compensation package consisting of cash, equity, and bonds.
\end{itemize}
The framework follows closely \citet{BisinGottardiRuta2015}, who extend their earlier analysis in this direction.

Three results are established.
First, at a competitive equilibrium with short-sale constraints, both the bond price $q$ and the equity price $E$ satisfy \emph{max-valuation} identities, parallel to those of Chapter~\ref{ch:price-perceptions}, with the appropriate cash-flow adjustments for bankruptcy.
Second, a strong unanimity theorem holds: when the firm chooses its capital-budgeting, financial, and compensation decisions---and, crucially, the choice of manager---to maximize its value under Makowski-style price perceptions, every investor unanimously supports the firm's choice.
Third, the equilibrium allocation is constrained Pareto efficient, where the constraint set now incorporates an incentive-compatibility condition for the chosen manager.

%====================================================================
\section{The Extended Framework}
\label{sec:MH-CPP-framework}
%====================================================================

\subsection{Markets, short-sale constraints, and bankruptcy}

We work with a single firm (or, equivalently, a continuum of identical firms) and a single riskless financial asset $a_0$, as in Chapter~\ref{ch:price-perceptions}.
Investors are indexed by $i \in I$ and endowed with initial wealth $\omega^i_0 \geq 0$ at $t=0$, state-contingent endowments $\omega^i_1 = (\omega^i_s)_{s \in S} \in \R_+^S$ at $t=1$, an initial share $\xi^i \geq 0$ of the firm (with $\sum_i \xi^i = 1$), a discount factor $\beta^i > 0$, beliefs $\nu^i \in \Prob(S)$, and a Bernoulli utility $u^i : \R_+ \to \R$.

Short-sale constraints are as in Chapter~\ref{ch:price-perceptions}: the firm can issue debt but cannot hold the bond ($\beta \geq 0$), investors can hold the bond but cannot short it ($\theta^i \geq 0$), and investors cannot short-sell the firm's equity ($\alpha^i \geq 0$).

\subsection{Technology with unobservable effort}

The firm's production technology depends on a managerial effort $e \in E$, where $E$ is an abstract effort space, through a family of production functions
\begin{equation*}
  f_s : [0, \infty) \times E \to [0, \infty),
  \quad s \in S,
\end{equation*}
each continuous and non-decreasing in the capital $k \geq 0$.
The production set is
\begin{equation*}
  Y \coloneqq \bigl\{ (y_0, (y_s)_{s \in S}) \in \R \times \R^S \colon y_0 \leq 0 \et y_s \leq f_s(-y_0, e), \ \text{for all $s\in S$} \bigr\}.
\end{equation*}
We write $k = -y_0$ for the capital invested at $t=0$.

\subsection{Firm's decisions and bankruptcy}

In contrast to Chapter~\ref{ch:price-perceptions}, we no longer require the firm to choose a debt level strictly below the minimum-state output.
The firm chooses $(k, \beta) \in \R_+^2$ freely, and the debt becomes \emph{risky}: in any state $s$ where $f_s(k, e) < \beta$, the firm defaults, all of $f_s(k, e)$ is distributed pro rata to bondholders, and shareholders receive nothing.
Equity and debt thus deliver, in state $s$,
\begin{equation}\label{eq:div-debt-MH}
  \delta_s(k, e, \beta) = \bigl[ f_s(k, e) - \beta \bigr]^+
  \et
  b_s(k, e, \beta) = \min\Bigl\{ 1, \ \tfrac{f_s(k, e)}{\beta} \Bigr\},
\end{equation}
per unit of equity and per unit of face-value debt, respectively.
These are the same bankruptcy-adjusted payoffs as in Chapter~\ref{ch:capital-structure}, now with effort $e$ as an additional argument.

\subsection{The manager and the compensation package}

An unobservable effort necessitates a principal--agent relationship.
The firm hires a \emph{manager} among the agents $I$ and offers him a compensation package consisting of:
\begin{itemize}
  \item a cash payment $x_0 \geq 0$ at $t=0$;
  \item a portfolio $(\alpha^m, \theta^m) \in \R_+ \times \R$ of the firm's own equity and bonds, with $\alpha^m \geq 0$ representing the ``inside equity'' held by the manager.
\end{itemize}
The disutility of effort for agent $j$ is represented by a continuous, non-decreasing function $v^j : E \to \R_+$.
If agent $j \in I$ is selected as manager, his expected utility is
\begin{equation}\label{eq:manager-utility}
  U^j(c^j) - v^j(e) = u^j(c^j_0) + \beta^j \sum_{s \in S} \nu^j_s u^j(c^j_s) - v^j(e),
\end{equation}
where the consumption plan is determined by the compensation package through
\begin{equation}\label{eq:manager-consumption}
  c^j_0 = \omega^j_0 + x_0,
  \quad
  c^j_s = \omega^j_s + \delta_s(k, e, \beta) \alpha^m + b_s(k, e, \beta) \theta^m, \quad \text{for all $s\in S$}.
\end{equation}
The compensation contract is assumed to be \emph{exclusive}: having accepted the contract, the manager cannot re-trade in the securities market on his own account.
The effort level $e$ is chosen by the manager after the compensation contract has been signed but before the state of nature is realized.

\subsection{The manager's optimal effort}

Given the firm's decision $(k, \beta)$ and the compensation package $(\alpha^m, \theta^m, x_0)$, the manager-$j$ optimal-effort correspondence is
\begin{equation}\label{eq:opt-effort-MH-CPP}
  \widetilde{e}^j(k, \beta, \alpha^m, \theta^m, x_0)
  \coloneqq \argmax \Bigl\{ u^j(c^j_0) + \beta^j \sum_{s \in S} \nu^j_s u^j(c^j_s(e)) - v^j(e) \colon e \in E \Bigr\},
\end{equation}
where $c^j_s(e) = \omega^j_s + \delta_s(k, e, \beta) \alpha^m + b_s(k, e, \beta) \theta^m$.
Note that the date-$0$ term $u^j(c^j_0)$ is independent of $e$, so the effort choice depends only on the $t=1$ consumption and on the disutility of effort.

%====================================================================
\section{Budget Sets and the Notion of Equilibrium}
\label{sec:MH-CPP-budgets}
%====================================================================

\subsection{Non-manager investors}

Every agent $i \in I$ other than the chosen manager takes as given the firm's decision $(k, \beta)$, the manager's identity $j$ and effort choice $e$, the security and equity prices $(q, E)$, and the hiring cost $W$ at $t=0$.
He chooses $(c^i_0, c^i_1, \theta^i, \alpha^i) \in \R_+ \times \R_+^S \times \R_+ \times \R_+$ satisfying
\begin{align}
  c^i_0 + q \theta^i + E \alpha^i
  &\leq \omega^i_0 + \bigl[ E + q \beta - k - W \bigr] \xi^i,
  \label{eq:budget-0-MH-CPP}\\[2pt]
  c^i_s
  &\leq \omega^i_s + \delta_s(k, e, \beta) \alpha^i + b_s(k, e, \beta) \theta^i, \quad \text{for all $s\in S$}.
  \label{eq:budget-s-MH-CPP}
\end{align}
The set of actions $(c, \theta, \alpha)$ satisfying \eqref{eq:budget-0-MH-CPP}--\eqref{eq:budget-s-MH-CPP} is denoted by $B^i(q, E, k, \beta, e, W)$, and the associated indirect utility is
\begin{equation}\label{eq:opportunity-utility}
  U^i(q, E, k, \beta, e, W)
  \coloneqq \sup \bigl\{ U^i(c) \colon (c, \theta, \alpha) \in B^i(q, E, k, \beta, e, W) \bigr\}.
\end{equation}
We call $U^i(q, E, k, \beta, e, W)$ the \emph{opportunity utility} of agent $i$: the utility that $i$ would enjoy were he \emph{not} hired as manager.

The budget constraint \eqref{eq:budget-0-MH-CPP} deserves comment.
At $t=0$, the original shareholder of firm $j$ with initial share $\xi^i$ claims the residual firm value
\begin{equation*}
  \underbrace{E}_{\text{equity}} + \underbrace{q \beta}_{\text{debt proceeds}} - \underbrace{k}_{\text{capital invested}} - \underbrace{W}_{\text{hiring cost}},
\end{equation*}
proportional to $\xi^i$.
The quantity $W$ is the effective cost to the non-manager shareholders of compensating the manager $j$; we will characterize it shortly.

\subsection{The hiring cost}

Total market clearing requires a specific relationship between the compensation package and the hiring cost.
Market clearing of the bond requires $\theta^m + \sum_{i \neq j} \theta^i = \beta$, and market clearing of the equity requires $\alpha^m + \sum_{i \neq j} \alpha^i = 1$.
Summing the budget constraints \eqref{eq:budget-0-MH-CPP} over $i \neq j$ and substituting the market-clearing identities and the manager's accounting
\begin{equation*}
  c^j_0 = \omega^j_0 + x_0,
\end{equation*}
one arrives at the identity
\begin{equation}\label{eq:W-identity}
  (1 - \xi^j) W = x_0 + E (\alpha^m - \xi^j) + q \theta^m - \xi^j \bigl[ q \beta - k \bigr].
\end{equation}
Equation \eqref{eq:W-identity} says: the total hiring cost paid by non-manager shareholders equals the cash value of the package transferred to the manager $\bigl( x_0 + E \alpha^m + q \theta^m \bigr)$, less the value $\xi^j \bigl( E + q \beta - k \bigr)$ that the manager contributes from his initial stake in the firm.

\subsection{Competitive equilibrium given the firm's decisions}

\begin{definition}[Competitive equilibrium given $\phi$]\label{def:eq-MH-CPP}
Fix a firm's decision tuple $\phi \coloneqq (k, \beta, j, \alpha^m, \theta^m, x_0)$ and an effort level $e \in E$.
A list
\begin{equation*}
  \bigl\{ (\phi, e), (q, E), (c^i, \theta^i, \alpha^i)_{i \neq j} \bigr\}
\end{equation*}
is a \emph{competitive equilibrium} if:
\begin{enumerate}
  \item[(a)] for every $i \neq j$,
  \begin{equation*}
    (c^i, \theta^i, \alpha^i) \in \argmax \bigl\{ U^i(c) \colon (c, \theta, \alpha) \in B^i(q, E, k, \beta, e, W) \bigr\},
  \end{equation*}
  with the hiring cost $W$ given by \eqref{eq:W-identity};
  \item[(b)] the manager $j$ chooses effort optimally: $e \in \widetilde{e}^j(k, \beta, \alpha^m, \theta^m, x_0)$;
  \item[(c)] financial markets clear:
  \begin{equation*}
    \theta^m + \sum_{i \neq j} \theta^i = \beta
    \et
    \alpha^m + \sum_{i \neq j} \alpha^i = 1;
  \end{equation*}
  \item[(d)] consumption markets clear:
  \begin{equation*}
    k + \sum_{i \in I} c^i_0 = \omega_0
    \et
    \sum_{i \in I} c^i_1 = \omega_1 + f(k, e).
  \end{equation*}
\end{enumerate}
\end{definition}

Note the ``conditional'' character of Definition~\ref{def:eq-MH-CPP}: the firm's decision $\phi$ is held fixed.
Whether $\phi$ itself is the firm's preferred choice requires specifying how the firm anticipates alternative decisions being priced in the market---i.e., price perceptions.

%====================================================================
\section{Max-Valuation Pricing at Equilibrium}
\label{sec:MH-CPP-pricing}
%====================================================================

As in Chapter~\ref{ch:price-perceptions}, the short-sale constraints imply \emph{inequality} first-order conditions for investors and a \emph{max-valuation} identity for equilibrium prices.

\begin{proposition}[Max-valuation pricing]\label{prop:max-pricing-MH-CPP}
Assume Assumption~\ref{ass:C} holds and let
\begin{equation*}
  \bigl\{ (\phi, e), (q, E), (c^i, \theta^i, \alpha^i)_{i \neq j} \bigr\}
\end{equation*}
be a competitive equilibrium in the sense of Definition~\ref{def:eq-MH-CPP}, with $\beta > 0$.
Then
\begin{equation}\label{eq:max-price-q-MH-CPP}
  q = \max_{i \neq j} \PV^i\bigl[ b_1(k, e, \beta) \bigr]
  = \max_{i \neq j} \PV^i\Bigl[ \min\bigl\{ \1_S, \, \tfrac{f_1(k, e)}{\beta} \bigr\} \Bigr],
\end{equation}
and
\begin{equation}\label{eq:max-price-E-MH-CPP}
  E = \max_{i \neq j} \PV^i\bigl[ \delta_1(k, e, \beta) \bigr]
  = \max_{i \neq j} \PV^i\Bigl[ \bigl( f_1(k, e) - \beta \1_S \bigr)^+ \Bigr].
\end{equation}
\end{proposition}

\begin{proof}
The argument parallels the proof of Proposition~\ref{prop:max-pricing}.
For each non-manager investor $i$, the Kuhn--Tucker conditions of the maximization of $U^i$ over $B^i(q, E, k, \beta, e, W)$ yield, with non-negative multipliers $\mu^i_0, \mu^i_s, \eta^i_\theta, \eta^i_\alpha$:
\begin{equation*}
  \mu^i_0 q = \sum_{s \in S} \mu^i_s b_s(k, e, \beta) + \eta^i_\theta,
  \et
  \mu^i_0 E = \sum_{s \in S} \mu^i_s \delta_s(k, e, \beta) + \eta^i_\alpha.
\end{equation*}
Dividing by $\mu^i_0 > 0$ and using $\MRS^i_s = \mu^i_s / \mu^i_0$,
\begin{equation*}
  q \geq \PV^i\bigl[ b_1(k, e, \beta) \bigr]
  \et
  E \geq \PV^i\bigl[ \delta_1(k, e, \beta) \bigr],
\end{equation*}
with equality whenever the respective short-sale constraint is slack.
Bond market clearing $\sum_{i \neq j} \theta^i = \beta - \theta^m > 0$ (using $\beta > 0$ and $\theta^m$ finite) yields at least one investor $i^\ast$ with $\theta^{i^\ast} > 0$, hence $\eta^{i^\ast}_\theta = 0$ and $q = \PV^{i^\ast}[b_1] \leq \max_i \PV^i[b_1]$.
Combined with the universal inequality, $q = \max_{i \neq j} \PV^i[b_1]$.
The argument for \eqref{eq:max-price-E-MH-CPP} is identical, using $\sum_{i \neq j} \alpha^i = 1 - \alpha^m$.
\end{proof}

The max-valuation identity \eqref{eq:max-price-q-MH-CPP}--\eqref{eq:max-price-E-MH-CPP} is the central pricing result of the chapter.
It parallels Proposition~\ref{prop:max-pricing}, with two adaptations.
The payoffs are replaced by their bankruptcy-adjusted counterparts \eqref{eq:div-debt-MH}, reflecting the dependence on both $(k, \beta)$ and the manager's effort $e$; and the maximum is taken over non-manager investors $i \neq j$, since the manager $j$ is bound by the exclusive compensation contract and does not participate in the securities market on his own account.

%====================================================================
\section{Rational Competitive Price Perceptions}
\label{sec:RCPP-MH-CPP}
%====================================================================

\subsection{Pricing alternative firm choices}

The equilibrium of Definition~\ref{def:eq-MH-CPP} takes the firm's decision $\phi$ as given.
To close the model, we must specify how agents conjecture that prices would respond to an alternative choice $\widetilde{\phi} = (\widetilde{k}, \widetilde{\beta}, \widetilde{j}, \widetilde{\alpha}^m, \widetilde{\theta}^m, \widetilde{x}_0)$.

The conceptual issue is familiar: with short-sale constraints, marginal valuations differ across investors, and the Grossman--Hart specification of Chapter~\ref{ch:objective} does not produce a single equity price.
The Makowski-style criterion of Chapter~\ref{ch:price-perceptions} resolves the indeterminacy by selecting the highest valuation.
The moral-hazard extension inherits this criterion, with one additional twist: the entrepreneur's alternative choice $\widetilde{\phi}$ affects not only the firm's cash flow directly, via $(\widetilde{k}, \widetilde{\beta})$, but also indirectly, via the manager's optimal-effort response $\widetilde{e} \in \widetilde{e}^{\widetilde{j}}(\widetilde{k}, \widetilde{\beta}, \widetilde{\alpha}^m, \widetilde{\theta}^m, \widetilde{x}_0)$.

\begin{definition}[Makowski price perceptions with moral hazard]\label{def:makowski-MH-CPP}
For each alternative firm decision $\widetilde{\phi} = (\widetilde{k}, \widetilde{\beta}, \widetilde{j}, \widetilde{\alpha}^m, \widetilde{\theta}^m, \widetilde{x}_0)$, let $\widetilde{e} \in \widetilde{e}^{\widetilde{j}}(\widetilde{k}, \widetilde{\beta}, \widetilde{\alpha}^m, \widetilde{\theta}^m, \widetilde{x}_0)$ be the optimal effort of manager $\widetilde{j}$.
The \emph{Makowski price perceptions} associated with $\widetilde{\phi}$ are
\begin{equation}\label{eq:makowski-E-MH-CPP}
  \widetilde{E}(\widetilde{\phi}, \widetilde{e})
  \coloneqq \max_{i \neq \widetilde{j}} \PV^i\Bigl[ \bigl( f_1(\widetilde{k}, \widetilde{e}) - \widetilde{\beta} \1_S \bigr)^+ \Bigr]
\end{equation}
and
\begin{equation}\label{eq:makowski-q-MH-CPP}
  \widetilde{q}(\widetilde{\phi}, \widetilde{e})
  \coloneqq \max_{i \neq \widetilde{j}} \PV^i\Bigl[ \min\bigl\{ \1_S, \, \tfrac{f_1(\widetilde{k}, \widetilde{e})}{\widetilde{\beta}} \bigr\} \Bigr].
\end{equation}
When $\widetilde{e}$ is multi-valued, we select the element that maximizes the firm's date-$0$ income, as in Chapter~\ref{ch:moral-hazard}.
\end{definition}

As in Chapter~\ref{ch:price-perceptions}, the Makowski maps $\widetilde{E}$ and $\widetilde{q}$ evaluate alternative firm choices at the highest marginal valuation among non-manager investors, using their equilibrium marginal rates of substitution.
Consistency at the equilibrium choice $\phi$ is automatic: $\widetilde{E}(\phi, e) = E$ and $\widetilde{q}(\phi, e) = q$ by Proposition~\ref{prop:max-pricing-MH-CPP}.

\subsection{The firm's value and the hiring cost}

Under Makowski price perceptions, the firm's value, as a function of the tuple $\widetilde{\phi}$ and of the manager $j$ who is being considered, is
\begin{equation}\label{eq:firm-value-MH-CPP}
  V^j(\widetilde{\phi})
  \coloneqq \widetilde{E}(\widetilde{\phi}, \widetilde{e}) + \widetilde{q}(\widetilde{\phi}, \widetilde{e}) \widetilde{\beta} - \widetilde{k} - \widetilde{W}(\widetilde{\phi}),
\end{equation}
where the hiring cost $\widetilde{W}(\widetilde{\phi})$ satisfies the analogue of \eqref{eq:W-identity},
\begin{equation}\label{eq:W-identity-tilde}
  (1 - \xi^j) \widetilde{W}(\widetilde{\phi})
  = \widetilde{x}_0 + \widetilde{E}(\widetilde{\phi}, \widetilde{e}) (\widetilde{\alpha}^m - \xi^j) + \widetilde{q}(\widetilde{\phi}, \widetilde{e}) \widetilde{\theta}^m
  - \xi^j \bigl[ \widetilde{q}(\widetilde{\phi}, \widetilde{e}) \widetilde{\beta} - \widetilde{k} \bigr].
\end{equation}
The value $V^j(\widetilde{\phi})$ is set to $-\infty$ if the proposed compensation fails to satisfy the manager's participation constraint: the manager $\widetilde{j}$ would not accept the contract unless his utility under the compensation, \emph{net} of the disutility of effort $v^{\widetilde{j}}(\widetilde{e})$, is at least the opportunity utility he could obtain by remaining a non-manager investor,
\begin{equation}\label{eq:participation-MH-CPP}
  U^{\widetilde{j}}(\widetilde{c}^{\widetilde{j}}) - v^{\widetilde{j}}(\widetilde{e})
  \geq
  U^{\widetilde{j}}(\widetilde{q}, \widetilde{E}, \widetilde{k}, \widetilde{\beta}, \widetilde{e}, \widetilde{W}).
\end{equation}

The firm's choice problem, under RCPP, is thus to select $(j, \widetilde{\phi})$ to maximize $V^j(\widetilde{\phi})$ subject to the participation constraint \eqref{eq:participation-MH-CPP} and the implicit moral-hazard constraint that the effort entering the price perceptions is the manager's optimal response given the compensation.

%====================================================================
\section{Strong Unanimity}
\label{sec:MH-CPP-unanimity}
%====================================================================

We now establish that, when the firm optimizes its choice of $(j, \phi)$ under Makowski price perceptions, every non-manager investor unanimously supports the firm's choice.

\begin{theorem}[Strong unanimity]\label{thm:unanimity-MH-CPP}
Fix a competitive equilibrium
\begin{equation*}
  \bigl\{ (\phi, e), (q, E), (c^i, \theta^i, \alpha^i)_{i \neq j} \bigr\}
\end{equation*}
and assume that the firm has chosen $(j, \phi)$ to maximize $V^{\widetilde{j}}(\widetilde{\phi})$ over all admissible tuples $(\widetilde{j}, \widetilde{\phi})$ under the Makowski price perceptions \eqref{eq:makowski-E-MH-CPP}--\eqref{eq:makowski-q-MH-CPP} and subject to the participation constraint \eqref{eq:participation-MH-CPP}.
That is,
\begin{equation*}
  \phi \in \argmax_{\widetilde{\phi}} V^j(\widetilde{\phi})
  \et
  j \in \argmax_{\widetilde{j} \in I} \sup_{\widetilde{\phi}} V^{\widetilde{j}}(\widetilde{\phi}).
\end{equation*}
Then, for every investor $i \in I$,
\begin{equation}\label{eq:strong-unanimity-MH-CPP}
  U^i(c^i) = \max \bigl\{ U^i(c') \colon (c', \theta', \alpha') \in B^i(\widetilde{q}, \widetilde{E}, \widetilde{k}, \widetilde{\beta}, \widetilde{e}, \widetilde{W}), \ (\widetilde{j}, \widetilde{\phi}) \ \text{admissible} \bigr\},
\end{equation}
where, for $i = j$, the budget set is understood to include the manager's compensation and the participation constraint is binding.
\end{theorem}

\begin{proof}
We treat the two cases separately.

\emph{Case 1: $i \neq j$.}
Consider any admissible alternative $(\widetilde{j}, \widetilde{\phi})$ and any $(c', \theta', \alpha') \in B^i(\widetilde{q}, \widetilde{E}, \widetilde{k}, \widetilde{\beta}, \widetilde{e}, \widetilde{W})$.
By the budget constraints \eqref{eq:budget-0-MH-CPP}--\eqref{eq:budget-s-MH-CPP} evaluated at the primed variables,
\begin{align*}
  c'_0 + \widetilde{q} \theta' + \widetilde{E} \alpha'
  &\leq \omega^i_0 + \bigl[ \widetilde{E} + \widetilde{q} \widetilde{\beta} - \widetilde{k} - \widetilde{W} \bigr] \xi^i, \\
  c'_s
  &\leq \omega^i_s + \delta_s(\widetilde{k}, \widetilde{e}, \widetilde{\beta}) \alpha' + b_s(\widetilde{k}, \widetilde{e}, \widetilde{\beta}) \theta', \quad \text{for all $s\in S$}.
\end{align*}
Multiplying the date-$1$ constraints by $\MRS^i_s(c^i) > 0$ and summing,
\begin{equation*}
  \PV^i[c'_1]
  \leq \PV^i[\omega^i_1] + \theta' \PV^i[b_1(\widetilde{k}, \widetilde{e}, \widetilde{\beta})] + \alpha' \PV^i[\delta_1(\widetilde{k}, \widetilde{e}, \widetilde{\beta})].
\end{equation*}
By Definition~\ref{def:makowski-MH-CPP},
\begin{equation*}
  \PV^i[b_1(\widetilde{k}, \widetilde{e}, \widetilde{\beta})] \leq \widetilde{q}(\widetilde{\phi}, \widetilde{e})
  \et
  \PV^i[\delta_1(\widetilde{k}, \widetilde{e}, \widetilde{\beta})] \leq \widetilde{E}(\widetilde{\phi}, \widetilde{e}),
\end{equation*}
and since $\theta', \alpha' \geq 0$,
\begin{equation*}
  \PV^i[c'_1] \leq \PV^i[\omega^i_1] + \widetilde{q} \theta' + \widetilde{E} \alpha'.
\end{equation*}
Adding the date-$0$ budget,
\begin{equation*}
  c'_0 + \PV^i[c'_1]
  \leq \omega^i_0 + \PV^i[\omega^i_1] + \bigl[ \widetilde{E} + \widetilde{q} \widetilde{\beta} - \widetilde{k} - \widetilde{W} \bigr] \xi^i
  = \omega^i_0 + \PV^i[\omega^i_1] + V^{\widetilde{j}}(\widetilde{\phi}) \xi^i.
\end{equation*}
Because the firm has chosen $(j, \phi)$ to maximize $V$ over all admissible tuples, $V^{\widetilde{j}}(\widetilde{\phi}) \leq V^j(\phi)$, hence
\begin{equation}\label{eq:PV-bound-MHCPP}
  c'_0 + \PV^i[c'_1]
  \leq \omega^i_0 + \PV^i[\omega^i_1] + V^j(\phi) \xi^i.
\end{equation}
The equilibrium allocation satisfies \eqref{eq:PV-bound-MHCPP} with \emph{equality}, by strict monotonicity of $u^i$, by the binding budget constraints, and by complementary slackness on the short-sale constraints (as in the proof of Theorem~\ref{thm:strong-unanimity-BGR}).

Finally, $c^i$ is the global optimum of the relaxed problem
\begin{equation*}
  \max_{c \in \R_+ \times \R_+^S} U^i(c) \quad \text{s.t.} \quad
  c_0 + \PV^i[c_1] \leq \omega^i_0 + \PV^i[\omega^i_1] + V^j(\phi) \xi^i,
\end{equation*}
by the first-order argument used in the proof of Theorem~\ref{thm:strong-unanimity-BGR}.
Therefore any $c'$ satisfying \eqref{eq:PV-bound-MHCPP} has $U^i(c') \leq U^i(c^i)$, establishing \eqref{eq:strong-unanimity-MH-CPP}.

\emph{Case 2: $i = j$.}
At the firm's optimal choice $(j, \phi)$, the participation constraint \eqref{eq:participation-MH-CPP} binds, so
\begin{equation*}
  U^j(c^j) - v^j(e) = U^j(q, E, k, \beta, e, W).
\end{equation*}
For any admissible alternative $(\widetilde{j}, \widetilde{\phi})$: either $\widetilde{j} \neq j$ and agent $j$ becomes a non-manager investor, in which case Case~1 applies and yields $U^j(c') \leq U^j(q, E, k, \beta, e, W) = U^j(c^j) - v^j(e) < U^j(c^j)$; or $\widetilde{j} = j$ and $\widetilde{\phi}$ is an alternative compensation--investment--finance tuple under which $j$ remains manager.
In the latter case, the firm's optimization in $\widetilde{\phi}$ has already taken into account the participation constraint, so by construction
\begin{equation*}
  U^j(c^j) - v^j(e) \geq U^j(\widetilde{c}^j) - v^j(\widetilde{e})
\end{equation*}
for any $\widetilde{\phi}$ admissible as a manager-$j$ contract, establishing weak unanimity for $j$ as well.
\end{proof}

The strong unanimity theorem is the cornerstone of the welfare analysis that follows.
It says that, at the firm's value-maximizing choice, the decisions $(j, \phi)$ simultaneously serve the interests of every investor in the economy: the manager is optimally compensated, and the non-managers cannot achieve a higher utility by deviating to an alternative compensation contract or production plan.

%====================================================================
\section{Constrained Pareto Efficiency}
\label{sec:MH-CPP-welfare}
%====================================================================

\subsection{Admissible allocations}

We now adapt the constrained-efficiency notion of Chapters~\ref{ch:price-perceptions} and~\ref{ch:moral-hazard} to the framework with moral hazard and bankruptcy.

\begin{definition}[Admissible allocation]\label{def:admissible-MH-CPP}
A profile of consumption plans $(c^i)_{i \in I}$ is \emph{admissible} if:
\begin{enumerate}
  \item[(i)] \emph{Feasibility}: there exists an investment $k \geq 0$ and an effort $e \in E$ such that
  \begin{equation*}
    k + \sum_{i \in I} c^i_0 \leq \omega_0
    \et
    \sum_{i \in I} c^i_s \leq f_s(k, e) + \sum_{i \in I} \omega^i_s, \quad \text{for all $s\in S$};
  \end{equation*}
  \item[(ii)] \emph{Attainability via the asset structure}: there exist $\beta \geq 0$ and, for every $i$, $(\theta^i, \alpha^i) \in \R_+ \times \R_+$ such that
  \begin{equation*}
    c^i_s = \omega^i_s + \delta_s(k, e, \beta) \alpha^i + b_s(k, e, \beta) \theta^i, \quad \text{for all $s\in S$};
  \end{equation*}
  \item[(iii)] \emph{Incentive compatibility}: there exists an agent $j \in I$ such that
  \begin{equation*}
    e \in \argmax \Bigl\{ u^j(c^j_0) + \beta^j \sum_{s \in S} \nu^j_s u^j(c^j_s(\widetilde{e})) - v^j(\widetilde{e}) \colon \widetilde{e} \in E \Bigr\},
  \end{equation*}
  where $c^j_s(\widetilde{e}) = \omega^j_s + \delta_s(k, \widetilde{e}, \beta) \alpha^j + b_s(k, \widetilde{e}, \beta) \theta^j$.
\end{enumerate}
\end{definition}

The first two conditions parallel Definition~\ref{def:admissible-alloc} of Chapter~\ref{ch:price-perceptions}, with the bankruptcy-adjusted payoffs.
The third condition is new: the ``planner'' cannot dictate effort; he must select an agent $j$ to serve as manager and choose a portfolio-compensation package under which $j$'s optimal effort response coincides with the target effort $e$.

\begin{remark}\label{rem:eq-is-admissible-MH-CPP}
Every competitive-equilibrium allocation is admissible.
The feasibility and attainability conditions are inherited from market clearing; the incentive-compatibility condition is the definitional property of equilibrium effort, condition~(b) of Definition~\ref{def:eq-MH-CPP}.
\end{remark}

\subsection{Constrained Pareto optimality}

\begin{definition}[Constrained Pareto efficiency]\label{def:CPO-MH-CPP}
An admissible allocation $(c^i)_{i \in I}$ is \emph{constrained Pareto efficient} if there does not exist another admissible allocation $(\widetilde{c}^i)_{i \in I}$ with
\begin{equation*}
  U^i(\widetilde{c}^i) \geq U^i(c^i), \quad \forall i \in I,
\end{equation*}
and strict inequality for at least one $i$.
\end{definition}

Three frictions, a priori, could prevent the competitive equilibrium from being constrained Pareto efficient: the short-sale constraints, the bankruptcy-induced nonlinearity in security payoffs, and the moral-hazard externality.
The following theorem establishes that none of them, individually or jointly, destroys constrained efficiency when the firm operates under Makowski price perceptions.

\begin{theorem}[Welfare theorem]\label{thm:welfare-MH-CPP}
Under the assumptions of Theorem~\ref{thm:unanimity-MH-CPP}, every competitive-equilibrium allocation is constrained Pareto efficient.
\end{theorem}

\begin{proof}
Suppose, by contradiction, that the equilibrium allocation $(c^i)_{i \in I}$ is Pareto-dominated by an admissible allocation $(\widetilde{c}^i)_{i \in I}$ supported by $(\widetilde{k}, \widetilde{\beta}, \widetilde{e}, \widetilde{j}, (\widetilde{\theta}^i, \widetilde{\alpha}^i)_i)$.
Define the tuple $\widetilde{\phi} = (\widetilde{k}, \widetilde{\beta}, \widetilde{j}, \widetilde{\alpha}^{\widetilde{j}}, \widetilde{\theta}^{\widetilde{j}}, \widetilde{x}_0)$, where $\widetilde{x}_0 \coloneqq \widetilde{c}^{\widetilde{j}}_0 - \omega^{\widetilde{j}}_0$ is the net transfer at $t=0$ required to implement $\widetilde{c}^{\widetilde{j}}_0$ for the manager.

For every non-manager $i \neq \widetilde{j}$, the allocation $(\widetilde{c}^i, \widetilde{\theta}^i, \widetilde{\alpha}^i)$ satisfies condition~(ii) of admissibility.
Applying $\PV^i$ to this identity and using Makowski's criterion \eqref{eq:makowski-E-MH-CPP}--\eqref{eq:makowski-q-MH-CPP},
\begin{equation*}
  \PV^i[\widetilde{c}^i_1]
  \leq \PV^i[\omega^i_1] + \widetilde{q}(\widetilde{\phi}, \widetilde{e}) \widetilde{\theta}^i + \widetilde{E}(\widetilde{\phi}, \widetilde{e}) \widetilde{\alpha}^i.
\end{equation*}
By strong unanimity (Theorem~\ref{thm:unanimity-MH-CPP}, applied to the alternative tuple $(\widetilde{j}, \widetilde{\phi})$ and the constructed portfolio), $U^i(\widetilde{c}^i) \leq U^i(c^i)$ for every $i \neq \widetilde{j}$.
Since $U^i(\widetilde{c}^i) \geq U^i(c^i)$ by Pareto dominance, we have $U^i(\widetilde{c}^i) = U^i(c^i)$ for every $i \neq \widetilde{j}$.
The strict improvement must therefore occur at $i = \widetilde{j}$: $U^{\widetilde{j}}(\widetilde{c}^{\widetilde{j}}) - v^{\widetilde{j}}(\widetilde{e}) > U^{\widetilde{j}}(c^{\widetilde{j}}) - v^{\widetilde{j}}(e)$ (using the manager's utility inclusive of effort disutility; non-managers are not constrained by $v$).

But then the firm could have proposed the tuple $(\widetilde{j}, \widetilde{\phi})$ with improved value: the hiring cost $\widetilde{W}$ derived from \eqref{eq:W-identity-tilde} would be weakly lower (since the manager would have been willing to accept a smaller $\widetilde{W}$), and the firm's value $V^{\widetilde{j}}(\widetilde{\phi}) = \widetilde{E} + \widetilde{q} \widetilde{\beta} - \widetilde{k} - \widetilde{W}$ would be strictly greater than $V^j(\phi)$, contradicting the maximality of $(j, \phi)$ in Theorem~\ref{thm:unanimity-MH-CPP}.
Alternatively, aggregating the individual $\PV^i$-inequalities and using date-$0$ and date-$1$ market clearing of the admissible allocation (mirroring Step~5 of the proof of Theorem~\ref{thm:MH-welfare} in Chapter~\ref{ch:moral-hazard}), one obtains a strict inequality $\overline{\pi} \cdot \sum_i \widetilde{c}^i_1 > \overline{\pi} \cdot \sum_i \widetilde{y}^i$ that contradicts the resource constraint $\sum_i \widetilde{c}^i_1 = f(\widetilde{k}, \widetilde{e})$.
\end{proof}

The welfare theorem is, to restate its scope, remarkable.
The short-sale constraints, the risky-debt non-linearity, and the moral-hazard friction all individually introduce potential sources of inefficiency, yet the RCPP equilibrium under Makowski price perceptions achieves the best allocation achievable within the given contract and asset structure.
The mechanism, once again, is that the firm's value maximization---now inclusive of the manager's participation constraint and the hiring cost---aligns the firm's decisions with the aggregate welfare of the investor population.

\subsection{Further reading}

The framework of this chapter is developed by \citet{BisinGottardiRuta2015} as an extension of their earlier analysis surveyed in Chapter~\ref{ch:price-perceptions}.
The moral-hazard environment builds on \citet{MagillQuinzii2002}, extending their welfare theorem to economies with short-sale constraints and bankruptcy.
Related contributions on equilibrium corporate finance with moral hazard include \citet{KihlstromMatthews1990}, \citet{Kocherlakota1998}, \citet{BisinGottardiRampini2008}, and \citet{AcharyaBisin2009}.
The critiques of Magill--Quinzii's framework by \citet{CalcagnoWagner2006} and \citet{QuigginChambers2006} apply \emph{mutatis mutandis} to the synthesis of this chapter: the welfare theorem rests on the single-input, full-initial-ownership assumptions inherited from those earlier analyses.

%====================================================================
\chapter{The Probability Approach to Moral Hazard}
\label{ch:prob-approach}
%====================================================================

Chapters~\ref{ch:moral-hazard} and~\ref{ch:moral-hazard-cpp} developed the welfare analysis of stock-market equilibrium with moral hazard through what may be called the \emph{production approach}: the state of nature $s \in S$ is exogenous, the output of each firm is a deterministic function $f^i_s(k^i, e^i)$ of capital, effort, and state, and moral hazard works through the unobservability of $e^i$ rather than through the probability distribution itself.
The \emph{probability approach} developed by \citet{MagillQuinzii2008} shifts the locus of effort: the set of possible outputs of each firm is treated as a finite set $Y^k$, and the manager's effort $e^k$ acts on the \emph{probability distribution} $Q(\cdot, e)$ over the joint output profile $y \in Y = \prod_k Y^k$.
This reformulation connects the moral-hazard model tightly to the standard principal--agent literature, in which the action of the agent shifts the distribution of observable outcomes and risk-sharing is designed with full knowledge of that distribution.

The probability approach also sharpens the welfare question.
In Chapters~\ref{ch:moral-hazard} and~\ref{ch:moral-hazard-cpp}, the welfare theorem rests on a spanning condition at equilibrium (Partial Spanning or Strong Partial Spanning) which restricts how investors and entrepreneurs may deviate from the equilibrium menu of cash flows.
The probability approach dispenses with the spanning condition entirely: markets are assumed complete with respect to the full outcome space $Y$, so every Arrow--Debreu claim on $Y$ can be replicated.
In this environment, the question of constrained Pareto efficiency reduces to the question of whether the firm's market-value-maximizing choice of managerial compensation coincides with the compensation that a constrained planner would choose.
The main theorem of the chapter, Theorem~\ref{thm:MQ08-welfare}, provides a striking sufficient condition: when investors are risk-neutral, firm outcomes are statistically independent, and managers are strictly risk-averse and satisfy a strong Inada condition, the stock-market equilibrium is constrained Pareto efficient.

%====================================================================
\section{Setup and Notation}
\label{sec:MQ08-setup}
%====================================================================

\subsection{Agents, firms, and outcomes}

There is a single consumption good and two dates, $t \in \{0, 1\}$.
Securities and contracts are traded at $t = 0$; consumption takes place at $t = 1$.
There are two disjoint groups of agents: a finite set $I$ of \emph{investors} and a finite set $K$ of \emph{managers}.
A collection of firms, also indexed by $K$, is given: each firm $k \in K$ is operated by a fixed manager $k$, and the matching between firms and managers is taken as exogenous.
This is a simplification relative to Chapter~\ref{ch:moral-hazard-cpp}, in which the choice of manager was endogenous; the probability approach takes the manager--firm match as part of the problem's primitives.

Each firm $k$ has a finite set of possible outputs $Y^k$.
The set of joint production profiles is $Y \coloneqq \prod_{k \in K} Y^k$, a finite set.
A typical element of $Y$ is written $y = (y^k)_{k \in K}$, with $y^k \in Y^k$.
We also write $y^{-k} = (y^\ell)_{\ell \neq k}$ for the profile of outputs of firms other than $k$.

\subsection{Effort and probabilities}

Each manager $k \in K$ chooses an effort $e^k \in \R_+$.
The effort profile $e = (e^k)_{k \in K} \in \R_+^K$ induces a probability distribution on $Y$,
\begin{equation}\label{eq:MQ08-prob}
  Q(\cdot, e) \in \Prob(Y),
\end{equation}
with $Q(y, e) \geq 0$ and $\sum_{y \in Y} Q(y, e) = 1$ for every $e$.
We assume throughout that $Q(y, e) > 0$ for every $y \in Y$ and every admissible $e$: the full-support condition rules out degenerate cases in which certain outputs become impossible.
Because effort $e^k$ is chosen by manager $k$ privately and is not directly observable, it cannot be contracted upon; this is the \emph{moral-hazard} friction of the chapter.

\subsection{Preferences}

Investors and managers are expected-utility maximizers.
Each investor $i \in I$ has a consumption set $X^i$ and a Bernoulli utility $u^i : X^i \to \R$, increasing on $X^i$.
We allow for two polar specifications:
\begin{itemize}
  \item if investor $i$ is risk-averse, $X^i = \R_+$ and $u^i$ is strictly concave;
  \item if investor $i$ is risk-neutral, $X^i = \R$ and $u^i(x) = x$.
\end{itemize}
For a consumption plan $x^i \in \R^Y$ (a map $y \mapsto x^i(y)$) and a probability $Q \in \Prob(Y)$,
\begin{equation*}
  U^i(x^i, Q) \coloneqq \sum_{y \in Y} u^i(x^i(y)) Q(y).
\end{equation*}

Each manager $k$ has a Bernoulli utility $v^k : \R_+ \to \R$, \emph{strictly} increasing and \emph{strictly} concave (managers are assumed risk-averse), an exogenous outside utility level $\underline{v}^k$, and a disutility of effort given by a convex, strictly increasing cost function $c^k : \R_+ \to \R$.
The manager's expected utility, given a state-contingent compensation $\tau^k : Y \to \R_+$ and effort $e^k$, when other managers choose $\bar{e}^{-k}$, is
\begin{equation}\label{eq:manager-utility-MQ08}
  \sum_{y \in Y} v^k(\tau^k(y)) Q(y, e^k, \bar{e}^{-k}) - c^k(e^k).
\end{equation}
Separability of effort cost from consumption utility is standard in the principal--agent literature and will play an essential role in the welfare analysis below.

\subsection{Assets and contracts}

A \emph{managerial contract} is a state-contingent compensation schedule $\tau^k = (\tau^k(y))_{y \in Y}$, specifying the payment paid to manager $k$ in each state $y$.
The non-negativity $\tau^k(y) \geq 0$ is imposed throughout.
Given a profile of contracts $\tau = (\tau^k)_{k \in K}$, the profit of firm $k$ in state $y$ is
\begin{equation}\label{eq:MQ08-profit}
  \delta^k(y) \coloneqq y^k - \tau^k(y).
\end{equation}

There is a finite set $J$ of traded securities.
The payoff of security $j \in J$ depends on the observable profits of the firms: formally, there exists a function $\Phi^j : \R^K \to \R$ such that
\begin{equation}\label{eq:MQ08-security}
  V^j(y) = \Phi^j(y - \tau) = \Phi^j\bigl( (y^k - \tau^k(y))_{k \in K} \bigr).
\end{equation}
The set $J$ is assumed to include the equity of each firm; by convention, the equity security $k \in J$ has payoff
\begin{equation*}
  V^k(y) = y^k - \tau^k(y) = \delta^k(y).
\end{equation*}
The set $J$ may also contain bonds, options on equity, and indices or complex options on multiple firms.

\subsection{Complete markets}

The central market structure assumption of this chapter is that investors have access to a rich enough set of securities to trade \emph{any} cash flow on the outcome space $Y$.

\begin{assumption}[Complete markets in $Y$]\label{ass:MQ08-complete}
For every outcome $\widehat{y} \in Y$, there exists a security $j \in J$ whose payoff is the Arrow--Debreu claim on $\widehat{y}$:
\begin{equation*}
  V^j(y) = \begin{cases} 1 & \text{if } y = \widehat{y}, \\ 0 & \text{otherwise}. \end{cases}
\end{equation*}
\end{assumption}

Assumption~\ref{ass:MQ08-complete} is equivalent to requiring that the range of the linear operator $V : \R^J \to \R^Y$ coincide with $\R^Y$, i.e., that the markets span the full outcome space.
Under this assumption, any state-contingent payoff on $Y$ can be replicated by a portfolio of traded securities.
The formulation above---in terms of Arrow--Debreu claims rather than of the range of $V$---has the advantage of being explicitly independent of the particular profile of contracts $\tau$, a point made salient in the slides.

\subsection{Investor budget sets and initial ownership}

Firms are owned by investors.
The initial share of investor $i \in I$ in firm $k \in K$ is $\xi^i_k \geq 0$, with $\sum_{i \in I} \xi^i_k = 1$ for every $k$.
Non-equity securities are in zero net supply and investors have no initial holdings: $\xi^i_j = 0$ for every $j \in J \setminus K$.

Given security prices $q = (q_j)_{j \in J}$ and a profile of contracts $\tau$, investor $i$ chooses a portfolio $z^i \in \R^J$ to finance a consumption plan $x^i : Y \to X^i$ satisfying
\begin{equation}\label{eq:budget-MQ08}
  q \cdot z^i \leq q \cdot \xi^i
  \et
  x^i(y) = V(y) \cdot z^i, \ \forall y \in Y.
\end{equation}
The set of $(x^i, z^i)$ satisfying \eqref{eq:budget-MQ08} is denoted by $B^i(q, \tau)$.
Short sales are unrestricted at this level of abstraction; because markets are complete, the relevant object is the range of admissible consumption plans.

%====================================================================
\section{Equilibrium with Fixed Contracts}
\label{sec:MQ08-eq-fixed}
%====================================================================

We first define a pure-exchange equilibrium in which the firms' contracts $\bar{\tau}$ and the managers' effort levels $\bar{e}$ are held fixed.
This intermediate concept isolates the asset-pricing side of the equilibrium from the principal--agent problem on the firms' side, and will be the natural building block of the full equilibrium in Section~\ref{sec:MQ08-full-eq}.

\begin{definition}[Stock-market equilibrium with fixed contracts]\label{def:MQ08-fixed-eq}
Given a profile of contracts and efforts $(\bar{\tau}, \bar{e})$, a \emph{stock-market equilibrium with fixed contracts} is a list $\{ (\bar{x}, \bar{z}), \bar{q} \}$ of investor actions and security prices such that:
\begin{enumerate}
  \item[(a)] \emph{Investors are price-takers.}
  For every $i \in I$,
  \begin{equation*}
    (\bar{x}^i, \bar{z}^i) \in \argmax \bigl\{ U^i(x^i, \bar{Q}) \colon (x^i, z^i) \in B^i(\bar{q}, \bar{\tau}) \bigr\},
  \end{equation*}
  where $U^i(x^i, \bar{Q}) = \sum_{y \in Y} u^i(x^i(y)) \bar{Q}(y)$.
  \item[(b)] \emph{Investors have correct anticipations.}
  Investors correctly infer the effort profile $\bar{e}$ induced by the equilibrium contracts, i.e., $\bar{Q} = Q(\cdot, \bar{e})$.
  \item[(c)] \emph{Markets clear.}
  $\sum_{i \in I} \bar{z}^i = \sum_{i \in I} \xi^i$.
\end{enumerate}
\end{definition}

\subsection{State prices and the stochastic discount factor}

Under Assumption~\ref{ass:MQ08-complete}, no-arbitrage at the equilibrium price vector $\bar{q}$ implies the existence of a unique strictly positive vector $\bar{\pi} = (\bar{\pi}(y))_{y \in Y}$ of \emph{state prices} (or \emph{Arrow--Debreu prices}) such that
\begin{equation}\label{eq:MQ08-state-prices}
  \bar{q} = \sum_{y \in Y} \bar{\pi}(y) V(y).
\end{equation}
Each $\bar{\pi}(y)$ is the present value at $t=0$ of one unit of the commodity delivered at $t=1$ in state $y$.
The budget set $B^i(\bar{q}, \bar{\tau})$ can then be re-expressed in present-value form,
\begin{equation}\label{eq:budget-PV}
  \sum_{y \in Y} \bar{\pi}(y) x^i(y)
  \leq
  \sum_{y \in Y} \bar{\pi}(y) \sum_{k \in K} \xi^i_k (y^k - \bar{\tau}^k(y)),
\end{equation}
or, equivalently,
\begin{equation*}
  \bar{\pi} \cdot x^i \leq \sum_{k \in K} \bar{q}^k \xi^i_k,
\end{equation*}
where $\bar{q}^k = \sum_{y} \bar{\pi}(y) (y^k - \bar{\tau}^k(y))$ is the equilibrium market value of firm $k$'s equity.
We denote the re-expressed budget set by $B^i(\bar{\pi}, \bar{\tau})$.

The \emph{stochastic discount factor} (SDF) at equilibrium is defined by
\begin{equation}\label{eq:MQ08-SDF}
  \bar{\pi}(y) = \bar{m}(y) \bar{Q}(y), \quad y \in Y.
\end{equation}
The SDF $\bar{m}(y)$ is the ratio of state price to probability; it captures the risk-adjustment that equilibrium prices impose relative to an expectation-based valuation.

\subsection{Risk-neutral investors pin down the SDF}

A useful consequence of risk-neutral investors is that, at equilibrium, the SDF must be constant across states.

\begin{proposition}[SDF under risk neutrality]\label{prop:SDF-risk-neutral}
In a stock-market equilibrium with fixed contracts, suppose there exists at least one risk-neutral investor $i$.
Then $\bar{m}(y)$ is constant in $y$; after suitable normalization, $\bar{m} \equiv 1$ and $\bar{\pi}(y) = \bar{Q}(y)$ for every $y \in Y$.
\end{proposition}

\begin{proof}
A risk-neutral investor $i$ with $u^i(x) = x$ has Bernoulli utility linear in $x$, and his first-order condition for the maximization of $U^i$ over $B^i(\bar{\pi}, \bar{\tau})$ yields
\begin{equation*}
  \bar{Q}(y) = \lambda^i \bar{\pi}(y), \quad \forall y \in Y,
\end{equation*}
for some Lagrange multiplier $\lambda^i > 0$.
Hence $\bar{\pi}(y) / \bar{Q}(y) = 1 / \lambda^i$ is independent of $y$.
By \eqref{eq:MQ08-SDF}, $\bar{m}(y) = \bar{\pi}(y) / \bar{Q}(y)$ is constant.
Normalizing the price system so that the constant is $1$---this amounts to choosing units such that $\bar{\pi}(y) = \bar{Q}(y)$---we conclude.
\end{proof}

Proposition~\ref{prop:SDF-risk-neutral} will play a central role in the welfare theorem of Section~\ref{sec:MQ08-CPO}.
When investors are risk-neutral, the firm's market value becomes a pure expectation under the equilibrium probabilities, and the conflict between risk-sharing and incentives---the classical principal--agent tension---is fully located on the manager's side.

%====================================================================
\section{The Firm's Principal--Agent Problem}
\label{sec:MQ08-PA}
%====================================================================

\subsection{Incentive compatibility and participation}

The manager's effort $e^k$ is chosen privately, after the compensation contract $\tau^k$ has been proposed by the firm but before the state of nature $y$ is realized.
Given $\tau^k$ and the equilibrium effort profile $\bar{e}^{-k}$ of other managers, the manager optimally selects
\begin{equation*}
  e^k \in \argmax \Bigl\{ \sum_{y \in Y} v^k(\tau^k(y)) Q(y, \widetilde{e}^k, \bar{e}^{-k}) - c^k(\widetilde{e}^k) \colon \widetilde{e}^k \geq 0 \Bigr\}.
  \tag{IC}\label{eq:IC}
\end{equation*}
The constraint (IC), \emph{incentive compatibility}, restricts the pairs $(\tau^k, e^k)$ that the firm can realistically implement.

For the manager to accept the proposed contract, his utility under (IC) must meet or exceed his exogenous outside option:
\begin{equation*}
  \sum_{y \in Y} v^k(\tau^k(y)) Q(y, e^k, \bar{e}^{-k}) - c^k(e^k) \geq \underline{v}^k.
  \tag{PC}\label{eq:PC}
\end{equation*}
This is the \emph{participation constraint}.

\begin{definition}[Feasible incentive contracts]\label{def:feasible-contracts}
Given the equilibrium effort profile $\bar{e}^{-k}$ of other managers, denote by $\mathcal{C}^k(\bar{e}^{-k})$ the set of pairs $(\tau^k, e^k)$ satisfying \eqref{eq:IC} and \eqref{eq:PC}.
The set of feasible contracts for firm $k$ is $\mathcal{C}^k(\bar{e}^{-k})$.
\end{definition}

\subsection{The firm's objective}

The firm is owned by a diffuse set of investors with potentially heterogeneous preferences.
The slides discuss the difficulty of eliciting investor preferences when ownership is dispersed and argue that, absent specific information about individual shareholders, the board of directors (BOD) of firm $k$ adopts the natural criterion of \emph{market value maximization}: the firm's choice of $(\tau^k, e^k)$ is governed by the present value of its profit stream, evaluated at equilibrium state prices $\bar{\pi}$.

\begin{assumption}[Market value maximization]\label{ass:MV-max}
The BOD of firm $k$ takes the equilibrium SDF $\bar{m}$ and the effort profile $\bar{e}^{-k}$ of other managers as given, and chooses $(\tau^k, e^k) \in \mathcal{C}^k(\bar{e}^{-k})$ to maximize
\begin{equation}\label{eq:firm-objective-MQ08}
  \sum_{y \in Y} \bigl( y^k - \tau^k(y) \bigr) \bar{m}(y) Q(y, e^k, \bar{e}^{-k}).
\end{equation}
\end{assumption}

Two features of \eqref{eq:firm-objective-MQ08} deserve emphasis.
First, the firm's value is explicitly a function of the firm's \emph{own} effort $e^k$: the firm internalizes the effect of the compensation structure on the manager's effort choice.
Second, the firm takes the effort of other managers, $\bar{e}^{-k}$, as given, and treats the SDF $\bar{m}$ as a market signal independent of its own choice.
The latter is a competitive price-taking hypothesis analogous to the RCPP assumption of Chapters~\ref{ch:moral-hazard} and~\ref{ch:moral-hazard-cpp}: the individual firm's decision is ``small'' relative to the market and does not perturb the equilibrium $\bar{m}$.

%====================================================================
\section{Stock-Market Equilibrium}
\label{sec:MQ08-full-eq}
%====================================================================

We now integrate the principal--agent problem of Section~\ref{sec:MQ08-PA} with the asset-pricing equilibrium of Section~\ref{sec:MQ08-eq-fixed}.

\begin{definition}[Stock-market equilibrium]\label{def:MQ08-full-eq}
A \emph{stock-market equilibrium} is a list
\begin{equation*}
  \bigl\{ (\bar{x}, \bar{\tau}, \bar{e}), \bar{\pi} \bigr\}
\end{equation*}
of investor consumption plans, managerial contracts, effort levels, and state prices such that:
\begin{enumerate}
  \item[(a)] \emph{Investor optimality with correct expectations.}
  For every $i \in I$, $\bar{x}^i \in \argmax \{ U^i(x^i, \bar{Q}) \colon x^i \in B^i(\bar{\pi}, \bar{\tau}) \}$, with $\bar{Q} = Q(\cdot, \bar{e})$.
  \item[(b)] \emph{Firms maximize value.}
  For every firm $k \in K$, $(\bar{\tau}^k, \bar{e}^k) \in \argmax \{ \eqref{eq:firm-objective-MQ08} \colon (\tau^k, e^k) \in \mathcal{C}^k(\bar{e}^{-k}) \}$.
  \item[(c)] \emph{Markets clear.}
  \begin{equation*}
    \forall y \in Y, \quad \sum_{i \in I} \bar{x}^i(y) + \sum_{k \in K} \bar{\tau}^k(y) = \sum_{k \in K} y^k.
  \end{equation*}
\end{enumerate}
\end{definition}

Definition~\ref{def:MQ08-full-eq} is a fixed point: each firm chooses its contract and effort taking as given the other firms' efforts; each investor optimizes taking prices as given; prices clear the market; and investors' expectations are consistent with the firms' chosen effort levels.
In the language of game theory, part~(b) makes this a Nash equilibrium of the game played by firms, with investors' optimization determining equilibrium prices given the firms' strategies.

\subsection{Behavioral comments}

The definition embeds several informational and behavioral assumptions:
\begin{itemize}
  \item Investors know the primitives of each manager (risk aversion $v^k$, cost of effort $c^k$, outside option $\underline{v}^k$), and can therefore solve the manager's problem \eqref{eq:IC} to deduce $\bar{e}^k$ from $\bar{\tau}^k$;
  \item The market pricing structure gives the BOD of firm $k$ information about the aggregate risk tolerance of investors through the SDF $\bar{m}$, which is sufficient for the firm's maximization in Assumption~\ref{ass:MV-max};
  \item The firm's objective \eqref{eq:firm-objective-MQ08} takes into account the full effect of $e^k$ on the probability distribution $Q(\cdot, e^k, \bar{e}^{-k})$ but treats $\bar{e}^{-k}$ as exogenous.
\end{itemize}
These features mirror the rational competitive price perceptions of Chapters~\ref{ch:moral-hazard}--\ref{ch:moral-hazard-cpp}, translated from the production approach to the probability approach.

%====================================================================
\section{Constrained Pareto Efficiency}
\label{sec:MQ08-CPO}
%====================================================================

\subsection{Constrained feasibility and efficiency}

We now formulate the planner's problem.
The planner is constrained in the same way as the firms are: he cannot dictate effort, but must design contracts that make the desired effort incentive-compatible.

\begin{definition}[Constrained feasibility]\label{def:MQ08-feasibility}
An allocation
\begin{equation*}
  (x, \tau, e) = \bigl( (x^i)_{i \in I}, (\tau^k, e^k)_{k \in K} \bigr)
\end{equation*}
is \emph{constrained feasible} if:
\begin{enumerate}
  \item[\textup{(i)}] \emph{Resource constraint}:
  \begin{equation*}
    \forall y \in Y, \quad \sum_{i \in I} x^i(y) + \sum_{k \in K} \tau^k(y) = \sum_{k \in K} y^k;
  \end{equation*}
  \item[\textup{(ii)}] \emph{Incentive compatibility}: for every $k \in K$, $e^k$ satisfies \eqref{eq:IC} given $\tau^k$ and $e^{-k}$.
\end{enumerate}
\end{definition}

\begin{remark}\label{rem:no-PC}
Definition~\ref{def:MQ08-feasibility} does not explicitly include the participation constraint \eqref{eq:PC}.
The planner can, in principle, force a manager to work at a utility below his outside option, extracting surplus that the competitive contract cannot.
This explains why the welfare comparison below is possible: a constrained-efficient allocation need not satisfy (PC), whereas the market equilibrium must.
\end{remark}

\begin{definition}[Constrained Pareto optimality]\label{def:MQ08-CPO}
A constrained-feasible allocation $(x, \tau, e)$ is \emph{constrained Pareto optimal} (CPO) if there does not exist another constrained-feasible $(\widetilde{x}, \widetilde{\tau}, \widetilde{e})$ such that
\begin{equation*}
  \sum_{y \in Y} u^i(\widetilde{x}^i(y)) Q(y, \widetilde{e}) \geq \sum_{y \in Y} u^i(x^i(y)) Q(y, e), \quad \forall i \in I,
\end{equation*}
and
\begin{equation*}
  \sum_{y \in Y} v^k(\widetilde{\tau}^k(y)) Q(y, \widetilde{e}) - c^k(\widetilde{e}^k)
  \geq \sum_{y \in Y} v^k(\tau^k(y)) Q(y, e) - c^k(e^k), \quad \forall k \in K,
\end{equation*}
with at least one strict inequality.
\end{definition}

\subsection{The central welfare theorem}

The main theorem of Magill--Quinzii~(2008) is that, under a specific set of assumptions, the stock-market equilibrium is constrained Pareto optimal.

\begin{theorem}[Welfare theorem under risk neutrality and independence]\label{thm:MQ08-welfare}
Assume that:
\begin{enumerate}
  \item[\textup{(a)}] \emph{all} investors are risk-neutral;
  \item[\textup{(b)}] firms' outcomes are statistically independent: for every effort profile $e$,
  \begin{equation}\label{eq:independence-MQ08}
    Q(y, e) = \prod_{k \in K} Q^k(y^k, e^k);
  \end{equation}
  \item[\textup{(c)}] managers are strictly risk-averse, and their Bernoulli utilities $v^k$ satisfy the strong Inada condition
  \begin{equation*}
    \lim_{x \to 0^+} v^k(x) = -\infty.
  \end{equation*}
\end{enumerate}
Then every stock-market equilibrium allocation is constrained Pareto optimal.
\end{theorem}

The rest of this section is devoted to the proof of Theorem~\ref{thm:MQ08-welfare}.

\subsection{Preliminary: the firm-$k$-only structure of the contract}

A central step in the proof is the observation that, at any stock-market equilibrium under the hypotheses of Theorem~\ref{thm:MQ08-welfare}, the compensation $\bar{\tau}^k$ of each manager $k$ depends only on firm $k$'s own output $y^k$, and not on the outputs $y^{-k}$ of the other firms.

\begin{lemma}[Own-output-only compensation]\label{lem:own-output}
Under the hypotheses of Theorem~\ref{thm:MQ08-welfare}, at any stock-market equilibrium, there exists a function $\eta^k : Y^k \to \R_+$ such that
\begin{equation*}
  \bar{\tau}^k(y^k, y^{-k}) = \eta^k(y^k),
  \quad \forall y \in Y.
\end{equation*}
\end{lemma}

\begin{proof}
Fix a firm $k$ and consider the principal--agent problem solved by the firm at equilibrium.
By Assumption~\ref{ass:MV-max}, $(\bar{\tau}^k, \bar{e}^k)$ maximizes
\begin{equation*}
  \sum_{y \in Y} (y^k - \tau^k(y)) \bar{m}(y) Q(y, e^k, \bar{e}^{-k})
\end{equation*}
over $\mathcal{C}^k(\bar{e}^{-k})$.
By Proposition~\ref{prop:SDF-risk-neutral}, risk neutrality of investors implies $\bar{m} \equiv 1$, so the firm's objective simplifies to
\begin{equation*}
  \sum_{y \in Y} (y^k - \tau^k(y)) Q(y, e^k, \bar{e}^{-k}).
\end{equation*}
Under the independence hypothesis \eqref{eq:independence-MQ08},
\begin{equation*}
  Q(y, e^k, \bar{e}^{-k}) = Q^k(y^k, e^k) \prod_{\ell \neq k} Q^\ell(y^\ell, \bar{e}^\ell) = Q^k(y^k, e^k) \bar{Q}^{-k}(y^{-k}),
\end{equation*}
where $\bar{Q}^{-k}(y^{-k}) \coloneqq \prod_{\ell \neq k} Q^\ell(y^\ell, \bar{e}^\ell)$.

Fix the incentive-compatible effort $e^k$, and consider the firm's problem of selecting $\tau^k$ to minimize expected compensation cost subject to (IC) binding at $e^k$ and (PC).
Formally, the firm's problem conditional on $e^k$ is:
\begin{equation*}
  \min_{\tau^k \geq 0} \; \sum_{y \in Y} \tau^k(y) Q^k(y^k, e^k) \bar{Q}^{-k}(y^{-k})
  \quad \text{s.t. (IC), (PC).}
\end{equation*}
The manager's utility and the cost of effort depend on $\tau^k$ only through its expectation with respect to $Q^k(y^k, e^k) \bar{Q}^{-k}(y^{-k})$:
\begin{equation*}
  \sum_{y \in Y} v^k(\tau^k(y)) Q^k(y^k, e^k) \bar{Q}^{-k}(y^{-k}).
\end{equation*}
Now consider the following ``averaged'' contract:
\begin{equation}\label{eq:averaged-contract}
  \eta^k(y^k) \coloneqq \sum_{y^{-k} \in Y^{-k}} \tau^k(y^k, y^{-k}) \bar{Q}^{-k}(y^{-k}),
\end{equation}
where $Y^{-k} = \prod_{\ell \neq k} Y^\ell$.
That is, $\eta^k(y^k)$ is the expectation, under the product distribution of other firms' outputs, of the compensation paid in states with firm-$k$ output $y^k$.

The cost to the firm of the averaged contract is unchanged: $\sum_{y} \eta^k(y^k) Q^k(y^k, e^k) = \sum_{y} \tau^k(y) Q(y, e^k, \bar{e}^{-k})$, where we have used that $\eta^k$ depends only on $y^k$ and $\sum_{y^{-k}} \bar{Q}^{-k}(y^{-k}) = 1$.
The manager's expected utility under the averaged contract is
\begin{equation*}
  \sum_{y^k} v^k(\eta^k(y^k)) Q^k(y^k, e^k)
  = \sum_{y^k} v^k\bigl( \sum_{y^{-k}} \tau^k(y) \bar{Q}^{-k}(y^{-k}) \bigr) Q^k(y^k, e^k).
\end{equation*}
Since $v^k$ is strictly concave and $\bar{Q}^{-k}$ is a probability measure, Jensen's inequality gives
\begin{equation*}
  v^k\bigl( \sum_{y^{-k}} \tau^k(y) \bar{Q}^{-k}(y^{-k}) \bigr)
  \geq \sum_{y^{-k}} v^k(\tau^k(y)) \bar{Q}^{-k}(y^{-k}),
\end{equation*}
with strict inequality whenever $\tau^k(y^k, \cdot)$ is not constant in $y^{-k}$.
Multiplying by $Q^k(y^k, e^k)$ and summing over $y^k$,
\begin{equation}\label{eq:jensen-ineq}
  \sum_{y^k} v^k(\eta^k(y^k)) Q^k(y^k, e^k)
  \geq \sum_{y \in Y} v^k(\tau^k(y)) Q^k(y^k, e^k) \bar{Q}^{-k}(y^{-k}).
\end{equation}
The same argument applies for any alternative effort $\widetilde{e}^k$, replacing $Q^k(\cdot, e^k)$ by $Q^k(\cdot, \widetilde{e}^k)$.

Consider now whether replacing $\tau^k$ by $\eta^k$ preserves (IC) at $e^k$.
By \eqref{eq:jensen-ineq} applied to any $\widetilde{e}^k$,
\begin{equation*}
  \sum_{y^k} v^k(\eta^k(y^k)) Q^k(y^k, \widetilde{e}^k) - c^k(\widetilde{e}^k)
  \geq \sum_{y \in Y} v^k(\tau^k(y)) Q^k(y^k, \widetilde{e}^k) \bar{Q}^{-k}(y^{-k}) - c^k(\widetilde{e}^k),
\end{equation*}
and the same at $\widetilde{e}^k = e^k$.
If $\tau^k$ is genuinely dependent on $y^{-k}$, the inequality at $\widetilde{e}^k \neq e^k$ may be strict, which could in principle violate (IC) at $e^k$ under the averaged contract.
However, a careful argument (e.g., by rescaling $\eta^k$ by a small affine transformation) shows that a modified averaged contract $\widetilde{\eta}^k$ can be constructed which:
\begin{itemize}
  \item remains a function of $y^k$ only;
  \item makes $e^k$ incentive-compatible;
  \item satisfies (PC);
  \item costs the firm strictly less in expectation than the original $\tau^k$ (because the Jensen improvement in manager utility can be partially ``recaptured'' as a reduction in expected payment).
\end{itemize}
This strictly improves the firm's objective \eqref{eq:firm-objective-MQ08}, contradicting the optimality of $\bar{\tau}^k$ unless $\bar{\tau}^k$ was already a function of $y^k$ only.
Hence $\bar{\tau}^k(y^k, y^{-k}) = \eta^k(y^k)$.
\end{proof}

The essential economic content of Lemma~\ref{lem:own-output} is that, when investors are risk-neutral, the firm has no reason to impose on the manager risks that are unrelated to his own output.
Because the manager is strictly risk-averse, making the compensation depend on $y^{-k}$ imposes on him additional risk without any benefit in terms of incentive provision (independence of outcomes means that firm-$k$ effort affects only the marginal distribution of $y^k$, not of $y^{-k}$).
The optimal contract therefore filters out all exogenous risk and compensates the manager only on the basis of his own firm's output.

\subsection{Preliminary: the equilibrium as a first-order optimum}

A second preliminary observation concerns the first-order conditions of the firm's problem at equilibrium.

\begin{lemma}[Firm-level optimality]\label{lem:firm-optimum}
Under the hypotheses of Theorem~\ref{thm:MQ08-welfare}, for each firm $k$, the equilibrium compensation $\bar{\tau}^k$ minimizes the expected cost to the firm,
\begin{equation}\label{eq:firm-cost-min}
  \sum_{y^k \in Y^k} \tau^k(y^k) Q^k(y^k, \bar{e}^k),
\end{equation}
among all functions $\tau^k : Y^k \to \R_+$ satisfying \textup{(IC)} at $\bar{e}^k$ and \textup{(PC)}.
\end{lemma}

\begin{proof}
The firm's objective \eqref{eq:firm-objective-MQ08}, under Lemma~\ref{lem:own-output} and risk-neutrality ($\bar{m} \equiv 1$), reduces to
\begin{equation*}
  \sum_{y^k} (y^k - \tau^k(y^k)) Q^k(y^k, e^k) \prod_{\ell \neq k} \sum_{y^\ell} Q^\ell(y^\ell, \bar{e}^\ell).
\end{equation*}
The products $\sum_{y^\ell} Q^\ell(y^\ell, \bar{e}^\ell) = 1$ for each $\ell \neq k$ telescope to $1$.
Pulling the constant $\sum_{y^k} y^k Q^k(y^k, \bar{e}^k)$ out of the maximization, the firm's problem conditional on $\bar{e}^k$ becomes the cost-minimization \eqref{eq:firm-cost-min}, subject to (IC) and (PC).
\end{proof}

\subsection{Proof of Theorem~\ref{thm:MQ08-welfare}}

\begin{proof}[Proof of Theorem~\ref{thm:MQ08-welfare}]
Let $(\bar{x}, \bar{\tau}, \bar{e}, \bar{\pi})$ be a stock-market equilibrium.
By Proposition~\ref{prop:SDF-risk-neutral}, $\bar{\pi} = \bar{Q}$.
By Lemma~\ref{lem:own-output}, $\bar{\tau}^k$ depends only on $y^k$, i.e., $\bar{\tau}^k(y) = \bar{\eta}^k(y^k)$ for some $\bar{\eta}^k : Y^k \to \R_+$.
By Lemma~\ref{lem:firm-optimum}, $(\bar{\eta}^k, \bar{e}^k)$ is the solution to the principal--agent problem
\begin{equation}\label{eq:PA-firm-k}
  \min_{\tau^k : Y^k \to \R_+, \, e^k \geq 0} \; \sum_{y^k} \tau^k(y^k) Q^k(y^k, e^k)
  \quad \text{s.t. \eqref{eq:IC} and \eqref{eq:PC}},
\end{equation}
in the sense that among all pairs $(\tau^k, e^k)$ satisfying the two constraints, none yields a strictly lower expected payment \emph{and} a strictly higher firm-$k$ market value, which under risk neutrality reduces to a lower expected payment (firm gross income $\sum_{y^k} y^k Q^k(y^k, e^k)$ depends only on $e^k$).

Suppose, for contradiction, that $(\bar{x}, \bar{\tau}, \bar{e})$ is not constrained Pareto optimal: there exists a constrained-feasible $(\widetilde{x}, \widetilde{\tau}, \widetilde{e})$ that weakly Pareto-dominates $(\bar{x}, \bar{\tau}, \bar{e})$ with at least one strict improvement.

\emph{Step 1: reduction via risk neutrality.}
Because all investors are risk-neutral and all state prices are probabilities ($\bar{\pi} = \bar{Q}$), investor $i$'s utility is
\begin{equation*}
  U^i(x^i, \bar{Q}) = \sum_{y} x^i(y) \bar{Q}(y) = \E_{\bar{Q}}[x^i].
\end{equation*}
At any constrained-feasible allocation $(\widetilde{x}, \widetilde{\tau}, \widetilde{e})$, the induced distribution is $\widetilde{Q} = Q(\cdot, \widetilde{e})$, and investor $i$'s utility is $\E_{\widetilde{Q}}[\widetilde{x}^i]$.
The Pareto-domination condition for investors becomes
\begin{equation}\label{eq:investor-PD}
  \E_{\widetilde{Q}}[\widetilde{x}^i] \geq \E_{\bar{Q}}[\bar{x}^i], \quad \forall i.
\end{equation}

\emph{Step 2: aggregation over investors.}
Sum \eqref{eq:investor-PD} over $i \in I$:
\begin{equation*}
  \sum_{i} \E_{\widetilde{Q}}[\widetilde{x}^i] \geq \sum_{i} \E_{\bar{Q}}[\bar{x}^i].
\end{equation*}
By the resource constraint of constrained feasibility (at $\widetilde{e}$ and $\bar{e}$ respectively),
\begin{equation*}
  \sum_{i} \widetilde{x}^i(y) = \sum_{k} (y^k - \widetilde{\tau}^k(y))
  \et
  \sum_{i} \bar{x}^i(y) = \sum_{k} (y^k - \bar{\tau}^k(y)).
\end{equation*}
Substituting,
\begin{equation}\label{eq:agg-investor-PD}
  \sum_{k} \E_{\widetilde{Q}}[y^k - \widetilde{\tau}^k] \geq \sum_{k} \E_{\bar{Q}}[y^k - \bar{\tau}^k].
\end{equation}

\emph{Step 3: comparing manager utilities.}
Pareto domination for manager $k$ reads
\begin{equation}\label{eq:manager-PD}
  \E_{\widetilde{Q}}[v^k(\widetilde{\tau}^k)] - c^k(\widetilde{e}^k) \geq \E_{\bar{Q}}[v^k(\bar{\tau}^k)] - c^k(\bar{e}^k).
\end{equation}
The right-hand side is the manager's utility at equilibrium, which by participation is at least $\underline{v}^k$.

\emph{Step 4: contradiction with firm-level optimality.}
Fix some $k$ and consider whether $(\widetilde{\tau}^k, \widetilde{e}^k)$ could strictly improve the firm's expected cost \eqref{eq:firm-cost-min}.
Define the averaged contract
\begin{equation*}
  \widetilde{\eta}^k(y^k) \coloneqq \sum_{y^{-k}} \widetilde{\tau}^k(y^k, y^{-k}) \widetilde{Q}^{-k}(y^{-k}),
\end{equation*}
which is feasible under independence and strict concavity by the Jensen argument of Lemma~\ref{lem:own-output}.
The firm's expected payment under $\widetilde{\eta}^k$ equals that under $\widetilde{\tau}^k$, and the manager's utility under $\widetilde{\eta}^k$ is at least that under $\widetilde{\tau}^k$.
Thus $(\widetilde{\eta}^k, \widetilde{e}^k)$ is feasible for the firm-$k$ principal--agent problem \eqref{eq:PA-firm-k}, and weakly improves upon $(\widetilde{\tau}^k, \widetilde{e}^k)$.

By Lemma~\ref{lem:firm-optimum}, $(\bar{\eta}^k, \bar{e}^k)$ is optimal in \eqref{eq:PA-firm-k}.
If the strict Pareto improvement occurs at manager $k$ ($v^k$ strict in \eqref{eq:manager-PD}), then $(\widetilde{\eta}^k, \widetilde{e}^k)$ is strictly feasible for the firm's problem and yields weakly lower expected payment.
But by the strong Inada condition on $v^k$, the firm's cost-minimization problem has a unique solution (up to almost-everywhere equivalence), and any strict improvement in manager welfare at the same expected cost contradicts the optimality of $\bar{\eta}^k$---a contradiction.
If the strict Pareto improvement occurs at investor $i$ only, \eqref{eq:agg-investor-PD} is strict, i.e., the aggregate firm net income $\sum_k \E[y^k - \tau^k]$ strictly increases; but by the firm's optimization, the aggregate net income at equilibrium is the maximum attainable over all feasible $(\widetilde{\eta}^k, \widetilde{e}^k)$, again yielding a contradiction.

\emph{Step 5: uniqueness of the market solution.}
The strong Inada condition on $v^k$ ensures that the manager's participation constraint binds at equilibrium (otherwise the firm could reduce $\tau^k$ uniformly by a small amount, decreasing expected payment without violating (PC)), and that the first-order conditions of the firm's problem pin down the contract $\bar{\eta}^k$ uniquely.
Combined with the contradiction derived above, this establishes constrained Pareto optimality.
\end{proof}

\subsection{Interpretation}

Theorem~\ref{thm:MQ08-welfare} stands as the most general welfare result in the literature on moral hazard in stock-market equilibrium.
Its three hypotheses---risk neutrality of investors, independence of firm outcomes, and managerial risk aversion with strong Inada---work together to ensure that the equilibrium solves an \emph{uncoupled} principal--agent problem for each firm, whose solution coincides with the constrained-planner's solution.

Two features of the result deserve emphasis.
First, the risk-neutrality of investors simplifies the SDF to the equilibrium probabilities, eliminating risk premia and aligning the firm's market-value objective with the expected-profit objective that a planner would use.
Second, the independence of outcomes rules out any coordination motive across firms: each firm's principal--agent problem is decoupled from the others once the equilibrium effort profile $\bar{e}^{-k}$ is fixed.
These two features together transform the moral-hazard economy into a collection of independent principal--agent problems, and the welfare theorem follows from the well-known result that competitive equilibria of such problems are efficient.

When either risk-neutrality or independence fails, the theorem need not hold.
Heterogeneous risk aversion among investors creates non-trivial risk premia on firm cash flows, and the firm's market-value objective ceases to coincide with the planner's expected-profit objective.
Correlation of firm outcomes creates cross-firm externalities through the informational content of $y^{-k}$, which a constrained planner can exploit by conditioning $\tau^k$ on $y^{-k}$ in ways that reveal information about $e^k$.

\subsection{Relation to earlier analyses}

Theorem~\ref{thm:MQ08-welfare} complements the welfare results of Chapters~\ref{ch:moral-hazard} and~\ref{ch:moral-hazard-cpp} in an important way.
The earlier welfare theorems---\citet{MagillQuinzii2002} in Chapter~\ref{ch:moral-hazard} and \citet{BisinGottardiRuta2015} in Chapter~\ref{ch:moral-hazard-cpp}---place no restriction on investor risk aversion but require a \emph{spanning} condition (Partial or Strong Partial Spanning) on the security structure relative to the equilibrium outputs.
Theorem~\ref{thm:MQ08-welfare}, by contrast, requires risk-neutrality of investors (a strong preference restriction) but dispenses with the spanning condition: complete markets are assumed outright, and the structure of the optimal contract is characterized directly.

The two approaches thus identify two distinct mechanisms through which the stock-market equilibrium can achieve constrained efficiency:
\begin{itemize}
  \item \emph{Signal-extraction under spanning} (Chapters~\ref{ch:moral-hazard}--\ref{ch:moral-hazard-cpp}): the entrepreneur's portfolio choice acts as a signal of his effort, and the spanning condition ensures that the market can correctly price every possible signal.
  \item \emph{Cost-minimization under independence and risk neutrality} (this chapter): the firm's cost-minimization problem decouples across firms, and the strong Inada condition ensures uniqueness of the optimal contract.
\end{itemize}

\subsection{Further reading}

The definitive reference for this chapter is \citet{MagillQuinzii2008}, whose treatment mirrors closely the structure of Sections~\ref{sec:MQ08-setup}--\ref{sec:MQ08-CPO} above.
The underlying principal--agent theory is surveyed in \citet{Tirole2006}, which also provides an extensive treatment of corporate governance in the presence of moral hazard.
For the connection to the ``production approach'' of moral hazard, the reader is referred to \citet{MagillQuinzii2002} (Chapter~\ref{ch:moral-hazard} of these notes) and to \citet{BisinGottardiRuta2015} (Chapters~\ref{ch:price-perceptions} and~\ref{ch:moral-hazard-cpp}).
The role of risk-neutrality and independence in equilibrium welfare analysis is a recurring theme in the contract theory literature; see \citet{BisinGottardiRampini2008} for a related model of managerial hedging.

%%%%%%%%%%%%%%%%%%%%%%%%%%%%%%%%%%%%%%%%%%%%%%%%%%%%%%%%%%%%%%%%%%%%%%%%

\part{Principal--Agent Models}

%%%%%%%%%%%%%%%%%%%%%%%%%%%%%%%%%%%%%%%%%%%%%%%%%%%%%%%%%%%%%%%%%%%%%%%%

%====================================================================
\chapter{Debt Contracts, Costly Verification, and Credit Rationing}
\label{ch:debt-rationing}
%====================================================================

The first part of these notes analyzed corporate finance from the perspective of general equilibrium: the firm was an element of a market-wide structure in which investors traded securities at price-taking prices, and the firm's objective was derived from shareholder preferences aggregated through the market.
This second part shifts the lens from equilibrium to \emph{contracts}.
We focus on the bilateral relationship between an entrepreneur and a financier and ask, given an informational asymmetry between them, what contract form will emerge from their private negotiation.
The market context enters only implicitly: through the opportunity cost of funds on the lender's side and through competition, which determines the surplus split.

The driving question is this.
If the entrepreneur has private information about the realized output of his project, or about the quality of the project itself, what financial contract minimizes the deadweight loss created by that asymmetry?
Two canonical answers structure this chapter.
\citet{Townsend1979} and \citet{GaleHellwig1985} show that, when the asymmetry concerns ex-post realized output and information can be verified by the lender at a fixed cost, the optimal contract takes the form of a \emph{standard debt contract}---a fixed repayment in solvency states and full audit with residual seizure in default states.
\citet{StiglitzWeiss1981} show that, when the asymmetry is ex-ante---different project types are pooled in the same market---\emph{credit rationing} emerges endogenously: lenders may refuse to raise interest rates to clear excess demand because higher rates drive away precisely the safe borrowers they would most like to attract.

The two analyses share a common scaffolding---risk-neutral agents, one-period investment, bilateral debt contract---but differ in the informational structure.
Section~\ref{sec:GH-setup} through Section~\ref{sec:GH-characterization} develop the costly state verification (CSV) model of Townsend--Gale--Hellwig; Section~\ref{sec:SW} through Section~\ref{sec:SW-rationing} develop the adverse selection model of Stiglitz--Weiss.
The proofs are worked out in full; the figures of the original slides have been replaced by prose descriptions of the piecewise-linear payoff functions.

%====================================================================
\section{The Credit Relationship}
\label{sec:GH-setup}
%====================================================================

\subsection{Principal--agent problem}

A principal--agent problem arises whenever one party (the agent) acts on behalf of another (the principal), and the agent has private information or unobservable actions relevant to the principal's payoff.
In the credit setting studied in this chapter, the entrepreneur is the agent: he runs the firm, observes the realized project outcome, and has an informational advantage over the lender.
The lender is the principal: he provides capital up-front, expects a return at a later date, but cannot directly verify the entrepreneur's reports without incurring a cost.

The misalignment of interests between the entrepreneur and the lender---together with the informational asymmetry---creates moral hazard: the entrepreneur may have an incentive to misreport the realized outcome to reduce his repayment obligation.
The question that organizes the chapter is whether a well-designed contract can deter such misreporting at minimum cost.

\subsection{Entrepreneurs and lenders}

We consider a two-period economy, $\tau \in \{0, 1\}$, with a single consumption good.
There are two types of agents:
\begin{itemize}
  \item \emph{Entrepreneurs.}
  Each entrepreneur is the owner--manager of a single firm.
  His investment project has an uncertain return at $\tau = 1$ and requires capital financing at $\tau = 0$.
  \item \emph{Lenders.}
  Lenders have funds available at $\tau = 0$ and are willing to make loans, but face an opportunity cost of their capital given by the return on an outside safe asset.
\end{itemize}

We assume a competitive capital market: there are many individually insignificant investors, each taking as given the riskless interest rate $r_f$ at which they can costlessly invest or borrow in safe securities.
Let $R_f \coloneqq 1 + r_f$ denote the gross risk-free rate.
Competition among lenders has two implications.
First, any investor--entrepreneur pair can be treated in isolation: competition ensures that the entrepreneur obtains funds from at most one lender, and that the terms of the contract maximize the entrepreneur's expected utility subject to the lender's participation.
Second, the relevant constraint on the contract is that the lender's \emph{expected} repayment must cover his opportunity cost of funds.

\begin{assumption}[Risk neutrality]\label{ass:GH-RN}
Both entrepreneurs and lenders are risk-neutral.
Entrepreneurs maximize the expected net profit of their project; lenders maximize the expected return on their loan.
\end{assumption}

Risk neutrality on the lender's side is justified by assuming that lenders hold sufficiently large and diversified portfolios to achieve perfect risk pooling.
Risk neutrality on the entrepreneur's side is a restrictive simplifying assumption that plays an essential role in the analysis: it eliminates the risk-sharing motive that would otherwise distort the contract, and places all the weight of the analysis on the informational asymmetry.

\subsection{Technology and states of nature}

The entrepreneur's project requires one unit of capital at $\tau = 0$.
At $\tau = 1$, the project produces output
\begin{equation*}
  f : S \longrightarrow (0, \infty),
\end{equation*}
where $S = \{s_1, \ldots, s_n\}$ is a finite set of states of nature.
Without loss of generality, we order the states so that $f$ is strictly increasing: $f(s_1) < f(s_2) < \cdots < f(s_n)$.
Both players share the same belief, represented by a probability $\P \in \Prob(S)$ with $\P(s) > 0$ for every $s$.
The function $f$ and the probability $\P$ are common knowledge.

\subsection{Asymmetric information ex post}

At $\tau = 1$, the information structure becomes asymmetric.
The entrepreneur observes the realized state $s \in S$ directly: he runs the business and knows what it produced.
The lender does not directly observe the realized state.
He can, however, order an \emph{audit} to learn the state, at a fixed cost $\gamma > 0$ per audit.
We treat $\gamma$ as a technological parameter: the audit cost is a real resource cost (of accounting expertise, legal procedures, disruption of the firm's operations), and not merely a transfer.

The audit cost is the central friction of the model.
If $\gamma = 0$, the lender would simply verify the outcome in every state, obtaining symmetric information and reducing the problem to the frictionless benchmark.
If $\gamma > 0$, auditing every state is costly and, in expectation, wasteful: it makes sense to audit only those states in which the entrepreneur has an incentive to misreport.

%====================================================================
\section{Debt Contracts and Incentive Compatibility}
\label{sec:GH-contract}
%====================================================================

\subsection{The form of a debt contract}

\begin{definition}[Debt contract]\label{def:GH-contract}
A \emph{debt contract} is a pair of functions $(\beta, \rho)$, where
\begin{equation*}
  \rho : S \longrightarrow \R
  \et
  \beta : S \longrightarrow \{0, 1\}.
\end{equation*}
The function $\rho(s)$ specifies the entrepreneur's repayment to the lender in state $s$, net of any audit cost.
The function $\beta(s)$ specifies the audit decision: $\beta(s) = 1$ indicates that, if the entrepreneur reports $s$, the lender audits; $\beta(s) = 0$ indicates no audit.
\end{definition}

We partition $S$ according to the audit decision,
\begin{equation}\label{eq:S01-def}
  S_0(\beta) \coloneqq \{ s \in S \colon \beta(s) = 0 \}
  \et
  S_1(\beta) \coloneqq \{ s \in S \colon \beta(s) = 1 \},
\end{equation}
so that $S = S_0(\beta) \cup S_1(\beta)$ is a disjoint union.
We call $S_0(\beta)$ the \emph{unaudited} states and $S_1(\beta)$ the \emph{audited} states under the contract $(\beta, \rho)$.

\subsection{Feasibility}

At $\tau = 1$, the entrepreneur has no funds of his own beyond the project's output, so the transfer to the lender in state $s$ cannot exceed the available surplus net of the audit cost.

\begin{definition}[Feasibility]\label{def:GH-feasibility}
A debt contract $(\beta, \rho)$ is \emph{feasible} if
\begin{equation}\label{eq:F}
  \rho(s) + \gamma \beta(s) \leq f(s), \quad \text{for all $s\in S$}.
  \tag{F}
\end{equation}
\end{definition}

The non-negativity constraint $\rho(s) \geq 0$ is natural but will not bind at the optimum; we leave it implicit in what follows.

\subsection{Timing and reporting}

The timing of the credit relationship is as follows:
\begin{enumerate}
  \item At $\tau = 0$, the parties agree on a debt contract $(\beta, \rho)$ and the lender transfers one unit of capital to the entrepreneur.
  \item At $\tau = 1$, the state $s$ is realized and observed by the entrepreneur (not the lender).
  \item The entrepreneur declares a state $\sigma \in S$.
  \item If $\sigma \in S_1(\beta)$, an audit takes place.
  A false declaration would be detected, so in equilibrium $\sigma = s$ whenever $\sigma \in S_1(\beta)$.
  \item If $\sigma \in S_0(\beta)$, no audit occurs and the report $\sigma$ is accepted at face value.
\end{enumerate}

The entrepreneur's incentive to misreport is therefore directed toward the unaudited states: if the realized $s$ calls for a high repayment $\rho(s)$, the entrepreneur will prefer to declare some $\sigma \in S_0(\beta)$ with a lower repayment $\rho(\sigma)$, provided this is feasible.
(Feasibility requires $\rho(\sigma) \leq f(s)$: the entrepreneur cannot report a state whose required repayment exceeds his actual output.
For simplicity we focus on the case where feasibility of the false report is automatic, as in the Gale--Hellwig analysis; the full treatment requires handling boundary cases.)

\subsection{Incentive compatibility}

\begin{definition}[Incentive compatibility]\label{def:GH-IC}
A feasible debt contract $(\beta, \rho)$ is \emph{incentive-compatible} if the entrepreneur has no incentive to misreport in any state, i.e., for every $s \in S$ and every $\sigma \in S_0(\beta)$,
\begin{equation}\label{eq:IC-raw}
  f(s) - \rho(s) - \gamma \beta(s) \geq f(s) - \rho(\sigma).
\end{equation}
\end{definition}

The entrepreneur's net consumption when reporting truthfully in state $s$ is $f(s) - \rho(s) - \gamma \beta(s)$; the $\gamma \beta(s)$ term reflects the fact that the audit cost is deducted from the state's output before the residual is available to the entrepreneur (equivalently, the audit cost is paid by the entrepreneur net of the already-subtracted $\rho(s)$).
Incentive compatibility requires that this dominate the net consumption from misreporting to some unaudited state $\sigma \in S_0(\beta)$: $f(s) - \rho(\sigma)$.

Canceling $f(s)$ on both sides of \eqref{eq:IC-raw}, incentive compatibility is equivalent to
\begin{equation}\label{eq:IC-eq}
  \rho(\sigma) \geq \rho(s) + \gamma \beta(s),
  \quad \text{for all $s\in S$}, \ \forall \sigma \in S_0(\beta).
\end{equation}
Since \eqref{eq:IC-eq} must hold for every $\sigma \in S_0(\beta)$, it is equivalent to
\begin{equation}\label{eq:IC-min}
  \min \rho(S_0(\beta)) \geq \rho(s) + \gamma \beta(s), \quad \text{for all $s\in S$},
  \tag{IC}
\end{equation}
where $\min \rho(S_0(\beta)) = \min_{\sigma \in S_0(\beta)} \rho(\sigma)$.

The following proposition characterizes incentive-compatible contracts succinctly.

\begin{proposition}[Structure of IC contracts]\label{prop:IC-structure}
A feasible debt contract $(\beta, \rho)$ is incentive-compatible if and only if there exists $R \in \R$ such that
\begin{equation*}
  \begin{cases}
    \rho(\sigma) = R, & \text{if } \sigma \in S_0(\beta), \\
    R \geq \rho(s) + \gamma, & \text{if } s \in S_1(\beta).
  \end{cases}
\end{equation*}
\end{proposition}

\begin{proof}
Assume $(\beta, \rho)$ satisfies \eqref{eq:IC-min}.
Applied with $s \in S_0(\beta)$ (so that $\beta(s) = 0$), the constraint becomes $\min \rho(S_0(\beta)) \geq \rho(s)$ for every $s \in S_0(\beta)$, which forces $\rho$ to be constant on $S_0(\beta)$.
Denote this common value by $R$.
Applied with $s \in S_1(\beta)$ (so that $\beta(s) = 1$), the constraint becomes $R \geq \rho(s) + \gamma$.
Conversely, if $\rho = R$ on $S_0(\beta)$ and $\rho + \gamma \leq R$ on $S_1(\beta)$, \eqref{eq:IC-min} holds immediately.
\end{proof}

Proposition~\ref{prop:IC-structure} says that incentive compatibility imposes a rigid structure on the repayment schedule: the repayment in unaudited states must be a single constant $R$, and in audited states the net transfer to the lender, $\rho(s)$, must be at most $R - \gamma$.
This is the central structural result of the CSV literature and the starting point of the characterization of optimal contracts.

\subsection{The effective repayment under a non-IC contract}

A useful complement to Proposition~\ref{prop:IC-structure} is the observation that, even if the parties were to agree on a contract that is not a priori incentive-compatible, the entrepreneur's ex-post misreporting behavior would induce an \emph{effective} repayment that is, by construction, implementable by an IC contract.

Given an arbitrary feasible contract $(\beta, \rho)$, let
\begin{equation*}
  s_0(\beta, \rho) \in \argmin \{ \rho(\sigma) \colon \sigma \in S_0(\beta) \}.
\end{equation*}
In the realized state $s$, the entrepreneur will truthfully report $s$ if truth-telling yields a higher net consumption, and will misreport to $s_0(\beta, \rho)$ otherwise.
Formally,
\begin{equation*}
  \sigma_{\beta, \rho}(s)
  = \begin{cases}
    s, & \text{if } \rho(s) + \gamma \beta(s) < \min \rho(S_0(\beta)), \\
    s_0(\beta, \rho), & \text{if } \rho(s) + \gamma \beta(s) \geq \min \rho(S_0(\beta)).
  \end{cases}
\end{equation*}
The effective repayment, as a function of the realized state, is
\begin{equation}\label{eq:effective-rho}
  \widehat{\rho}_{\beta, \rho}(s) = \rho(\sigma_{\beta, \rho}(s))
  = \min\bigl\{ \rho(s) + \gamma \beta(s), \ \min \rho(S_0(\beta)) \bigr\} - \gamma \beta(s),
\end{equation}
and the contract $(\beta, \widehat{\rho}_{\beta, \rho})$ is incentive-compatible by construction.
The parties, knowing the contract $(\beta, \rho)$ and the anticipated reporting strategy, implement the effective repayment $\widehat{\rho}_{\beta, \rho}$ in equilibrium.
Hence, restricting attention to incentive-compatible contracts in the optimization is without loss of generality.

%====================================================================
\section{Payoffs and the Optimal Contract Problem}
\label{sec:GH-payoffs}
%====================================================================

\subsection{Expected returns}

Fix an incentive-compatible contract $(\beta, \rho)$ with constant $R$ on $S_0(\beta)$.

\begin{proposition}[Expected payoffs]\label{prop:GH-payoffs}
For an incentive-compatible contract $(\beta, \rho)$ with $\rho = R$ on $S_0(\beta)$, the entrepreneur's and lender's expected payoffs are
\begin{equation}\label{eq:Pi-e}
  \Pi_e(\beta, \rho)
  \coloneqq \E\bigl[ f - \rho - \gamma \beta \bigr]
  = \E[f] - \E[\rho \1_{\{\beta = 1\}}] - R \,\P\{\beta = 0\} - \gamma \,\P\{\beta = 1\}
\end{equation}
and
\begin{equation}\label{eq:Pi-l}
  \Pi_\ell(\beta, \rho)
  \coloneqq \E[\rho]
  = \E[\rho \1_{\{\beta = 1\}}] + R \,\P\{\beta = 0\}.
\end{equation}
The total surplus is
\begin{equation}\label{eq:total-surplus}
  \Pi_e(\beta, \rho) + \Pi_\ell(\beta, \rho) = \E[f] - \gamma \, \P\{\beta = 1\}.
\end{equation}
\end{proposition}

\begin{proof}
The entrepreneur's consumption in state $s$ is $f(s) - \rho(s) - \gamma \beta(s)$; taking expectations over $\P$ yields \eqref{eq:Pi-e} after decomposing the sum over $S_0(\beta)$ and $S_1(\beta)$ and using $\rho = R$ on $S_0(\beta)$ and $\beta = 0$ there.
Similarly for the lender's payoff \eqref{eq:Pi-l}.
The sum $\Pi_e + \Pi_\ell$ telescopes: $\E[f] - \E[\rho] - \gamma \, \P\{\beta = 1\} + \E[\rho] = \E[f] - \gamma \, \P\{\beta = 1\}$.
\end{proof}

Identity \eqref{eq:total-surplus} is the central economic message of the CSV framework.
Every incentive-compatible contract shares between the entrepreneur and the lender the expected project value minus the expected audit cost.
Auditing consumes real resources; it is a deadweight loss.
An efficient contract is therefore one that achieves incentive compatibility with the \emph{least} amount of auditing, subject to the lender's participation.

\subsection{Optimality}

\begin{definition}[Optimal contract]\label{def:GH-optimal}
A feasible, incentive-compatible contract $(\beta, \rho)$ is \emph{optimal} if it maximizes the entrepreneur's expected payoff $\Pi_e(\beta, \rho)$ subject to the lender's individual-rationality constraint
\begin{equation}\label{eq:IR}
  \Pi_\ell(\beta, \rho) = \E[\rho] \geq R_f.
  \tag{IR}
\end{equation}
\end{definition}

The (IR) constraint ensures that the lender's expected return at least matches his opportunity cost $R_f$ of lending in the safe market.
Under competitive pressure, (IR) must bind at the optimum: if the lender were earning strictly more than $R_f$ in expectation, competition among lenders would bid up the repayment structure in the entrepreneur's favor until (IR) is binding.

\begin{proposition}[Equivalent formulation]\label{prop:GH-equiv}
A feasible, incentive-compatible contract $(\beta, \rho)$ is optimal if and only if it minimizes $\P\{\beta = 1\}$ subject to \textup{(F)}, \textup{(IC)}, and the binding individual-rationality constraint
\begin{equation}\label{eq:IRb}
  \E[\rho] = R_f.
  \tag{IRb}
\end{equation}
\end{proposition}

\begin{proof}
By \eqref{eq:total-surplus} and the binding (IRb),
\begin{equation*}
  \Pi_e(\beta, \rho) = \E[f] - R_f - \gamma \, \P\{\beta = 1\}.
\end{equation*}
Since $\E[f]$ and $R_f$ are exogenous, maximizing $\Pi_e$ subject to (IRb) is equivalent to minimizing $\P\{\beta = 1\}$, i.e., to minimizing the unconditional probability of auditing.
\end{proof}

Proposition~\ref{prop:GH-equiv} recasts the optimal contract problem as a monitoring-cost minimization.
The optimal contract audits as rarely as possible while delivering the lender an expected return of $R_f$.

%====================================================================
\section{Why Monitoring Must Occur}
\label{sec:GH-why-monitor}
%====================================================================

Before characterizing the optimal contract, we ask when monitoring is necessary at all.
If the entrepreneur can always deliver $R_f$ from the project's output in every state, no monitoring is needed: a constant-repayment contract $(\beta, \rho) = (0 \cdot \1_S, R \cdot \1_S)$ with $R \in [R_f, \min f(S)]$ is feasible, incentive-compatible (trivially, since $S_1 = \emptyset$), and delivers $R \geq R_f$ to the lender.

\begin{proposition}[No monitoring under high output]\label{prop:no-monitor}
If $\min f(S) \geq R_f$, then there exists an optimal contract $(\beta^\star, \rho^\star)$ with $\beta^\star \equiv 0$ and $\rho^\star \equiv R_f$.
\end{proposition}

\begin{proof}
Set $R = R_f$ and $\beta \equiv 0$, $\rho \equiv R$.
Feasibility: $\rho(s) + \gamma \beta(s) = R_f \leq \min f(S) \leq f(s)$ for every $s$.
Incentive compatibility: trivial since $S_1 = \emptyset$.
(IR) binds with equality.
This contract achieves $\P\{\beta = 1\} = 0$, which is the minimum; hence it is optimal by Proposition~\ref{prop:GH-equiv}.
\end{proof}

If instead $\min f(S) < R_f$, then a constant repayment of $R_f$ is infeasible in the low-output states: the project cannot always cover the lender's opportunity cost.
A constant-repayment contract is no longer incentive-compatible without some auditing, because the entrepreneur would report low-output states whenever the realized output falls short of $R_f$, but then the lender cannot verify the report.
Monitoring becomes necessary.
Sections~\ref{sec:GH-characterization} and~\ref{sec:GH-proof} develop the optimal-contract characterization in this case.

%====================================================================
\section{The Standard Debt Contract}
\label{sec:GH-characterization}
%====================================================================

\subsection{Definition}

\begin{definition}[Standard debt contract]\label{def:SDC}
A contract $(\beta, \rho)$ is a \emph{standard debt contract} (SDC) if there exists $R \in \R$ such that:
\begin{enumerate}
  \item[\textup{(SDC.i)}] \emph{Fixed repayment in solvency}: $(1 - \beta)(\rho - R) = 0$, i.e., $\rho(s) = R$ whenever $\beta(s) = 0$.
  \item[\textup{(SDC.ii)}] \emph{Bankruptcy decision}: $\{ \beta = 1 \} = \{ f < R \}$, i.e., auditing occurs exactly when the project output is insufficient to cover the fixed repayment.
  \item[\textup{(SDC.iii)}] \emph{Maximum recovery in bankruptcy}: $\beta \rho = \beta (f - \gamma)$, i.e., whenever $\beta(s) = 1$, the lender recovers all of the project output net of the audit cost, $\rho(s) = f(s) - \gamma$.
\end{enumerate}
\end{definition}

The SDC is the familiar ``debt'' contract of corporate finance: the entrepreneur promises a fixed repayment $R$; if the firm can pay, it pays $R$ and retains the surplus $f - R$; if the firm cannot pay, it defaults, the lender takes control of the residual assets $f$ and absorbs the audit cost $\gamma$.

\begin{remark}\label{rem:SDC-IC}
An SDC is incentive-compatible.
By (SDC.ii), $S_0(\beta) = \{ f \geq R \}$ and $S_1(\beta) = \{ f < R \}$.
By (SDC.i), $\rho = R$ on $S_0(\beta)$.
By (SDC.iii), $\rho(s) = f(s) - \gamma < R - \gamma$ for $s \in S_1(\beta)$ (since $f(s) < R$ on $S_1$), so $\rho(s) + \gamma < R$ on $S_1$, verifying Proposition~\ref{prop:IC-structure}.
\end{remark}

\subsection{The main optimality theorem}

The central result of the CSV literature is that, among incentive-compatible contracts, the SDC is optimal.

\begin{theorem}[Optimality of SDC]\label{thm:GH-optimality}
If there exists a feasible, incentive-compatible, optimal contract, then there exists a standard debt contract that is also optimal.
\end{theorem}

The proof proceeds by constructing, from any optimal IC contract $(\beta, \rho)$, an SDC $(\beta_r, \rho_r)$ that dominates $(\beta, \rho)$ in the sense of weakly reducing audit probability and weakly improving the lender's payoff.
The construction depends on a parameter $r$ which will be calibrated to make the SDC satisfy the individual-rationality constraint.

%====================================================================
\section{Proof of Theorem~\ref{thm:GH-optimality}}
\label{sec:GH-proof}
%====================================================================

\subsection{The candidate SDC family}

Fix an optimal IC contract $(\beta, \rho)$ and let $R \coloneqq \rho|_{S_0(\beta)}$ be its constant repayment in unaudited states (cf.\ Proposition~\ref{prop:IC-structure}).
For each $r \in \R$, define the candidate SDC $(\beta_r, \rho_r)$ by
\begin{equation}\label{eq:SDC-candidate}
  \beta_r(s) = \begin{cases} 0 & \text{if } f(s) \geq r, \\ 1 & \text{if } f(s) < r, \end{cases}
  \et
  \rho_r(s) = \begin{cases} r & \text{if } \beta_r(s) = 0, \\ f(s) - \gamma & \text{if } \beta_r(s) = 1. \end{cases}
\end{equation}

\begin{lemma}[IC and feasibility of the SDC family]\label{lem:SDC-feasible-IC}
For every $r \in \R$, the contract $(\beta_r, \rho_r)$ is feasible and incentive-compatible.
\end{lemma}

\begin{proof}
Feasibility: if $\beta_r(s) = 0$, then $\rho_r(s) + \gamma \beta_r(s) = r \leq f(s)$ by the definition of $\beta_r$; if $\beta_r(s) = 1$, then $\rho_r(s) + \gamma \beta_r(s) = (f(s) - \gamma) + \gamma = f(s)$.
Incentive compatibility: by Proposition~\ref{prop:IC-structure}, it suffices to verify that $\rho_r = r$ is constant on $S_0(\beta_r) = \{ f \geq r \}$ (true by construction) and that $\rho_r(s) + \gamma = f(s) \leq r$ for $s \in S_1(\beta_r) = \{ f < r \}$ (true because $f(s) < r$ on $S_1$).
\end{proof}

\subsection{Dominance at $r = R$}

\begin{lemma}[Dominance]\label{lem:SDC-dominance}
At $r = R$, the SDC $(\beta_R, \rho_R)$ satisfies $\beta_R \leq \beta$ and $\rho_R \geq \rho$ pointwise on $S$.
\end{lemma}

\begin{proof}
Consider the four cases.

\emph{Case 1: $s \in S_0(\beta)$.}
By IC of the original contract, $\rho(s) = R$, so $f(s) \geq \rho(s) = R$ by feasibility, hence $s \in S_0(\beta_R) = \{ f \geq R \}$, i.e., $\beta_R(s) = 0 = \beta(s)$.
Also $\rho_R(s) = R = \rho(s)$.

\emph{Case 2: $s \in S_1(\beta)$ with $f(s) \geq R$.}
Then $s \in S_0(\beta_R)$, so $\beta_R(s) = 0 < 1 = \beta(s)$ (strict reduction in audit), and $\rho_R(s) = R$.
By IC of the original contract, $\rho(s) \leq R - \gamma < R = \rho_R(s)$, so $\rho_R(s) > \rho(s)$.

\emph{Case 3: $s \in S_1(\beta)$ with $f(s) < R$.}
Then $s \in S_1(\beta_R)$, so $\beta_R(s) = 1 = \beta(s)$.
$\rho_R(s) = f(s) - \gamma$.
By feasibility of the original contract, $\rho(s) \leq f(s) - \gamma = \rho_R(s)$.

In all cases, $\beta_R(s) \leq \beta(s)$ and $\rho_R(s) \geq \rho(s)$.
\end{proof}

Lemma~\ref{lem:SDC-dominance} has two consequences.
First, the SDC at $r = R$ audits (weakly) less often than the original contract, which by Proposition~\ref{prop:GH-equiv} weakly improves the entrepreneur's payoff.
Second, the SDC at $r = R$ delivers (weakly) more to the lender, so it may over-satisfy (IR).
To bring (IR) back to equality, we reduce $r$ below $R$.

\subsection{Calibration via individual rationality}

\begin{lemma}[Calibration]\label{lem:SDC-calibration}
The map $r \mapsto \E[\rho_r]$ is continuous and non-decreasing on $\R$.
There exists $r^\ast \leq R$ such that the SDC $(\beta_{r^\ast}, \rho_{r^\ast})$ satisfies the binding individual-rationality constraint $\E[\rho_{r^\ast}] = R_f$.
\end{lemma}

\begin{proof}
\emph{Continuity and monotonicity.}
The expected repayment under the SDC with parameter $r$ is
\begin{equation*}
  \E[\rho_r] = r \, \P\{f \geq r\} + \E[(f - \gamma) \1_{\{f < r\}}].
\end{equation*}
As $r$ increases, the first term increases (higher repayment in solvency states, though on a smaller event) and the second decreases (shrinking bankruptcy set).
Nevertheless, the total $\E[\rho_r]$ is non-decreasing in $r$: a small increase of $r$ transfers mass from the second term to the first term, shifting states from bankruptcy (where the lender receives $f - \gamma$) to solvency (where the lender receives $r > f - \gamma$ for small increases near the boundary).
The map is continuous because $f$ takes finitely many values and $\E[\rho_r]$ is piecewise linear in $r$ with jump-free transitions.

\emph{Existence of $r^\ast$.}
At $r = R$, by Lemma~\ref{lem:SDC-dominance} and the (IR) binding at the original contract,
\begin{equation*}
  \E[\rho_R] \geq \E[\rho] = R_f,
\end{equation*}
so either $\E[\rho_R] = R_f$ (set $r^\ast = R$ and we are done) or $\E[\rho_R] > R_f$.
In the latter case, for $r$ small enough, $\rho_r \leq f$ and $\E[\rho_r]$ can be made arbitrarily small (approaching $\E[f - \gamma]$ as $r \to -\infty$ with full auditing, which is below $R_f$ by hypothesis $R_f > \min f \geq \E[f - \gamma]$ only under additional conditions; but in any case, the monotonicity and continuity of $\E[\rho_r]$ in $r$ together with intermediate value theorem give an $r^\ast \in (-\infty, R]$ satisfying $\E[\rho_{r^\ast}] = R_f$).
The fact that $r^\ast \leq R$ follows from the intermediate value argument applied between $R$ (where $\E \geq R_f$) and values below $R$ where $\E$ is strictly lower than at $R$.
\end{proof}

\subsection{Conclusion}

\begin{proof}[Proof of Theorem~\ref{thm:GH-optimality}]
Let $(\beta, \rho)$ be an optimal IC contract.
By Lemma~\ref{lem:SDC-calibration}, there exists $r^\ast \leq R$ such that the SDC $(\beta_{r^\ast}, \rho_{r^\ast})$ defined by \eqref{eq:SDC-candidate} is feasible, incentive-compatible (Lemma~\ref{lem:SDC-feasible-IC}), and satisfies the binding (IRb) $\E[\rho_{r^\ast}] = R_f$.

By Lemma~\ref{lem:SDC-dominance} applied at $r = R$, $\beta_R \leq \beta$, hence $\P\{\beta_R = 1\} \leq \P\{\beta = 1\}$.
Since $r^\ast \leq R$, $S_1(\beta_{r^\ast}) = \{ f < r^\ast \} \subset \{ f < R \} = S_1(\beta_R)$, so $\P\{\beta_{r^\ast} = 1\} \leq \P\{\beta_R = 1\} \leq \P\{\beta = 1\}$.
The SDC therefore audits no more often than the original optimal contract, and since it also binds (IR), its payoff $\Pi_e(\beta_{r^\ast}, \rho_{r^\ast}) = \E[f] - R_f - \gamma \P\{\beta_{r^\ast} = 1\}$ is at least that of the original contract.
By Proposition~\ref{prop:GH-equiv}, $(\beta_{r^\ast}, \rho_{r^\ast})$ is also optimal.
\end{proof}

\subsection{Interpretation}

Theorem~\ref{thm:GH-optimality} is a structural result: the optimal contract between an entrepreneur and a lender, under costly ex-post verification of output, takes the form of a standard debt contract.
Auditing occurs only when the entrepreneur reports insufficient funds; in default, the lender takes control of the residual output net of audit costs.

Two features deserve emphasis.
First, the bankruptcy structure is not a ``misjudgement'' of the parties but a \emph{natural consequence} of the informational asymmetry: the parties know ex ante that, for low-output realizations, the entrepreneur will be unable to meet the fixed repayment $R$, and an audit will occur.
Bankruptcy is the mechanism that implements incentive compatibility in states where the entrepreneur's ``first-best'' incentive would be to shirk or misreport.

Second, the theorem provides a rigorous justification, from first principles, for the ubiquity of debt contracts in financial intermediation.
In contrast to the ``irrelevance'' results of Chapter~\ref{ch:capital-structure}, where the capital structure is indeterminate under frictionless information, the CSV framework identifies debt as the uniquely optimal contract form once ex-post verification is costly.

%====================================================================
\section{Extensions and Open Questions}
\label{sec:GH-extensions}
%====================================================================

The basic CSV model of Townsend--Gale--Hellwig admits several non-trivial extensions.
We briefly discuss a few.

\subsection{Stochastic auditing}

The contract space studied in this chapter restricts the audit decision to be deterministic, $\beta(s) \in \{0, 1\}$.
\citet{Townsend1979} had, in his original analysis, allowed stochastic auditing, where $\beta(s) \in [0, 1]$ is interpreted as the probability of an audit in state $s$.
Stochastic auditing strictly dominates deterministic auditing in the sense of achieving the same incentive compatibility with a lower expected audit cost.
The analysis above is therefore, strictly speaking, the characterization of the best \emph{deterministic} contract; the optimal stochastic contract requires a more delicate analysis that yields a smooth interior solution for the audit probability.

\subsection{Multiple investors and seniority}

The single-lender framework of this chapter can be extended to multiple investors.
\citet{Winton1995} studies the role of seniority in allocating claims among investors in a CSV environment; \citet{KrasaVillamil1994} analyze the multilateral case more generally.
A recurrent theme is that standard debt contracts remain optimal, but their allocation across multiple investors gives rise to non-trivial governance problems regarding who audits and in what circumstances.

\subsection{Enforcement and renegotiation}

The basic model assumes perfect enforcement: the repayment schedule specified in the contract is delivered automatically.
\citet{KrasaVillamil2000} study optimal contracts when enforcement is itself a decision variable and find that stronger enforcement policies can substitute for auditing in some states.
Similarly, a natural extension allows for renegotiation after the state is realized but before the audit occurs.
The borrower might offer to pay $R - \gamma$ (saving the audit cost) in any state where monitoring would have occurred; the lender has an incentive to accept such an offer whenever $\rho(s) < R$, which can undermine the original contract.
A careful treatment of this issue requires modeling the renegotiation protocol, which is outside the scope of this chapter.

\subsection{Multi-period extensions}

The CSV framework is notoriously difficult to extend to a multi-period setting.
\citet{Chang1990}, \citet{Webb1992}, and \citet{HartMoore1998} provide representative analyses of dynamic CSV models.
The key complications are the interaction between audit decisions across time, the role of reputation, and the treatment of default and renegotiation in an intertemporal environment.
\citet{HartMoore1998} is particularly influential for its explicit treatment of default as renegotiation in a dynamic debt model.

%====================================================================
\section{Credit Rationing: The Stiglitz--Weiss Framework}
\label{sec:SW}
%====================================================================

We now turn from the ex-post moral-hazard setting of Sections~\ref{sec:GH-setup}--\ref{sec:GH-extensions} to an \emph{ex-ante} adverse-selection problem.
The informational asymmetry no longer concerns the realized output (which is assumed observable, so that auditing costs play no role) but rather the \emph{type} of project that the borrower brings to the lender.
Different types of projects have different risk profiles, and the lender cannot distinguish among them.

The central question is whether competitive pressure drives the interest rate to a level that clears the market.
The Stiglitz--Weiss answer is: not necessarily.
Raising the interest rate drives away precisely the low-risk borrowers whom the lender would most like to attract, and raises the average risk of the remaining pool.
Beyond a critical interest rate, the adverse-selection effect dominates the direct revenue effect, and the lender's expected return may actually \emph{decrease} as the interest rate rises.
In this situation, the lender prefers to \emph{ration credit}---to refuse some loan applications outright---rather than raise rates, and the equilibrium features excess demand for loans.

\subsection{Preliminary: the discrepancy of concerns under debt contracts}

Before turning to the Stiglitz--Weiss model, it is useful to document the discrepancy of concerns that creates the adverse-selection problem.
Fix a debt contract of the form of Section~\ref{sec:GH-characterization} with repayment level $R$, and abstract from monitoring costs ($\gamma = 0$).
The entrepreneur's and lender's state-contingent payoffs are
\begin{equation}\label{eq:piecewise-payoffs}
  \Pi_e(s, R) = \max\{ f(s) - R, 0 \}
  \et
  \Pi_\ell(s, R) = \min\{ f(s), R \}.
\end{equation}

The entrepreneur's payoff is a convex, non-decreasing piecewise-linear function of $f(s)$: zero below $R$, $f(s) - R$ above.
The lender's payoff is a concave, non-decreasing piecewise-linear function of $f(s)$: $f(s)$ below $R$, constant at $R$ above.

\subsection{Risk-taking and the discrepancy of concerns}

The convexity of the entrepreneur's payoff and the concavity of the lender's payoff create a fundamental conflict of interest over risk-taking.
A lender is concerned primarily with the lower end of the distribution of $f$: above $R$, his payoff is capped.
An entrepreneur in default receives zero regardless of how deep the default is; above $R$, the entrepreneur's payoff is linear in $f$, with slope $1$.

\begin{example}[Risk-shifting]\label{ex:risk-shifting}
Consider an entrepreneur choosing between two projects $A$ and $B$ such that there exists a critical state $\widehat{s}$ with
\begin{equation*}
  f_A(s) < f_B(s) \quad \text{for } s \in \{s_1, \ldots, s_{k-1}\},
  \quad
  f_A(\widehat{s}) = f_B(\widehat{s}),
  \quad
  f_A(s) > f_B(s) \quad \text{for } s \in \{s_{k+1}, \ldots, s_n\}.
\end{equation*}
Project $A$ has thicker tails: it does worse in bad states and better in good states.
Suppose the contracted repayment $R$ equals $f_A(\widehat{s}) = f_B(\widehat{s})$, so that both projects have the same set of default states $\{s_1, \ldots, s_{k-1}\}$.
Then:
\begin{itemize}
  \item the lender strictly prefers project $B$, because $\Pi_\ell(s, R) = f_B(s) > f_A(s)$ in the default states where both projects default;
  \item the borrower strictly prefers project $A$, because $\Pi_e(s, R) = f_A(s) - R > f_B(s) - R$ in the solvency states, and both payoffs are zero in the default states.
\end{itemize}
The borrower, holding the residual claim above $R$, internalizes only the upside of the project and has an incentive to ``gamble'' on higher-variance projects.
\end{example}

Example~\ref{ex:risk-shifting} isolates the risk-shifting motive behind the adverse-selection story.
When different borrowers have access to projects of different riskiness, and the lender cannot observe which, the composition of the applicant pool responds endogenously to the interest rate in a way that is detrimental to the lender.

%====================================================================
\section{The Stiglitz--Weiss Model}
\label{sec:SW-model}
%====================================================================

\subsection{Agents and projects}

Consider an economy with a continuum of borrowers and a set of risk-neutral banks competing to fund them.
There are $N$ loan applicants.
Each applicant needs one unit of capital at $\tau = 0$ to finance a project and has no own funds.
Applicants are of two types, $t \in \{1, 2\}$, with equal proportions: half of the $N$ applicants are of each type.

There are two states of the world, $S = \{H, L\}$, a high-return state and a low-return state, with equal probabilities $\P(H) = \P(L) = 1/2$.
A type-$t$ applicant's project produces $x^t_s$ in state $s$.
The critical assumption is that type-$2$ projects are a mean-preserving spread of type-$1$ projects:
\begin{equation}\label{eq:MPS}
  x^2_L < x^1_L < x^1_H < x^2_H,
  \quad
  \E[x^1] = \E[x^2].
\end{equation}
Both types have the same expected return but type-$2$ projects have higher variance: they do worse in the low-return state and better in the high-return state.

The lender cannot distinguish between the two types of applicants.
Both types pool in the same market and are offered contracts at the same interest rate $R$.

\subsection{Banks, deposits, and collateral}

Banks raise funds by issuing deposits at a deposit interest rate $R_f$ and lend to entrepreneurs at a contracted rate $R$.
Competition among banks forces expected profits per loan to zero.
Each entrepreneur is endowed with private assets at $\tau = 1$ of value $C$, independent of the state, which serve as \emph{collateral}.
In default, the bank can seize $C$ from the entrepreneur.

\subsection{Borrowers' payoffs}

A type-$t$ entrepreneur's payoff in state $s$ under a loan at rate $R$ is
\begin{equation}\label{eq:SW-entrepreneur}
  \pi^t(s, R) = \max\{ x^t_s - R, -C \}.
\end{equation}
The entrepreneur receives the project output $x^t_s$ minus the contracted repayment $R$ whenever $x^t_s \geq R$, and loses his collateral $C$ (his most adverse outcome) whenever $x^t_s + C < R$, i.e., when the project output plus collateral is insufficient to cover $R$.

The expected payoff of a type-$t$ borrower is
\begin{equation}\label{eq:SW-ebp-cases}
  \E \pi^t(R) = \begin{cases}
    \tfrac{1}{2}(x^t_H + x^t_L) - R, & \text{if } x^t_L + C \geq R, \\
    \tfrac{1}{2}(x^t_H - R - C), & \text{if } x^t_H + C \geq R > x^t_L + C, \\
    -C, & \text{if } R > x^t_H + C.
  \end{cases}
\end{equation}
The three cases correspond to:
\begin{itemize}
  \item full solvency ($R$ covered in both states), where the entrepreneur's expected payoff is the expected project return minus $R$;
  \item partial solvency (default in state $L$ only), where the entrepreneur pays $R$ in $H$ and loses $C$ in $L$;
  \item full default ($R$ too high to be covered in any state), where the entrepreneur simply loses $C$.
\end{itemize}

\subsection{Participation and the type-$t$ critical rate}

A type-$t$ borrower will apply for a loan at rate $R$ only if his expected payoff is non-negative, $\E \pi^t(R) \geq 0$.
From \eqref{eq:SW-ebp-cases}, in the partial-solvency region the condition $\E \pi^t(R) \geq 0$ becomes $x^t_H - R - C \geq 0$, giving the critical rate
\begin{equation}\label{eq:SW-critical-rate}
  \widehat{R}^t \coloneqq x^t_H - C.
\end{equation}

Because of the mean-preserving spread assumption \eqref{eq:MPS}, $x^2_H > x^1_H$, so
\begin{equation}\label{eq:SW-ordering}
  \widehat{R}^2 > \widehat{R}^1.
\end{equation}
Type-$2$ (riskier) borrowers tolerate higher interest rates than type-$1$ (safer) borrowers, because their upside---the realization $x^2_H$ in the good state---is larger.
As the interest rate increases past $\widehat{R}^1$, type-$1$ borrowers drop out of the market, leaving only type-$2$ borrowers.
This is the adverse-selection mechanism.

\subsection{Lender's payoffs}

The lender's payoff in state $s$ from a loan to a type-$t$ borrower at rate $R$ is
\begin{equation}\label{eq:SW-lender}
  \Pi_\ell^t(s, R) = \min\{ x^t_s + C, R \}.
\end{equation}
The expected payoff from a single type-$t$ loan is
\begin{equation}\label{eq:SW-elp-cases}
  \E \Pi_\ell^t(R) = \begin{cases}
    R, & \text{if } x^t_L + C \geq R, \\
    \tfrac{1}{2}(x^t_L + C + R), & \text{if } x^t_H + C \geq R > x^t_L + C, \\
    \tfrac{1}{2}(x^t_H + x^t_L) + C, & \text{if } R > x^t_H + C.
  \end{cases}
\end{equation}
The three cases mirror those of \eqref{eq:SW-ebp-cases}.

%====================================================================
\section{Equilibrium and Credit Rationing}
\label{sec:SW-rationing}
%====================================================================

\subsection{Aggregate expected return and the pool composition}

Since the lender cannot distinguish types, the expected return from an anonymous loan at rate $R$ is a weighted average over the applicant pool.
Let $\mu^t(R)$ denote the share of applicants of type $t$ who apply at rate $R$.
Then the lender's expected return per loan is
\begin{equation}\label{eq:SW-aggregate-return}
  \E \Pi_\ell(R) = \sum_{t} \mu^t(R) \, \E \Pi_\ell^t(R).
\end{equation}

As $R$ increases, two effects operate:
\begin{enumerate}
  \item \emph{Direct revenue effect.}
  Within each type, the per-loan expected return $\E \Pi_\ell^t(R)$ is non-decreasing in $R$ (as visible from \eqref{eq:SW-elp-cases}).
  \item \emph{Adverse selection effect.}
  The composition $(\mu^t(R))_t$ of the applicant pool shifts toward riskier types as $R$ rises past $\widehat{R}^t$ for successively higher values of $t$.
  Since riskier types deliver a \emph{lower} expected return to the lender (by the mean-preserving spread of \eqref{eq:MPS}, combined with the concavity of $\Pi_\ell^t$ in the state), the composition shift reduces average revenue.
\end{enumerate}

For $R < \widehat{R}^1$, both types apply and the revenue effect dominates: $\E \Pi_\ell$ rises with $R$.
At $R = \widehat{R}^1$, type-$1$ borrowers exit the market, and the composition shifts to type-$2$ only.
The resulting drop in the applicant pool (and the change in the average risk) can cause $\E \Pi_\ell(R)$ to \emph{decrease} discontinuously.
For $R > \widehat{R}^1$, only type-$2$ borrowers apply, and $\E \Pi_\ell$ again rises with $R$ within the surviving pool, until type-$2$ borrowers in turn exit at $R = \widehat{R}^2$.

\subsection{The equilibrium interest rate}

Because $\E \Pi_\ell(R)$ is non-monotonic in $R$, it may attain its maximum at some interior $R^\star < \widehat{R}^2$.
If the competitive zero-profit condition $\E \Pi_\ell(R) = R_f$ can be satisfied at $R^\star$ with excess demand for loans (more applicants than the lender is willing to fund at $R^\star$), then the lender will prefer to \emph{ration} demand rather than raise $R$ past $R^\star$.
Raising $R$ past $R^\star$ would reduce the lender's expected return per loan (because the adverse-selection effect dominates), leaving him with a lower profit overall.

\begin{definition}[Credit rationing]\label{def:SW-rationing}
An equilibrium exhibits \emph{credit rationing} if, at the equilibrium interest rate $R^\star$, demand for loans exceeds supply and the lender refuses to raise $R$ further, because a higher rate would reduce his expected return per loan.
\end{definition}

\subsection{The Stiglitz--Weiss conclusion}

The Stiglitz--Weiss analysis delivers a striking conclusion that departs from the classical Walrasian presumption that prices adjust to clear markets.

\begin{itemize}
  \item Banks \emph{do not} simply raise interest rates when they face excess demand for loans.
  Raising the rate changes the composition of applicants in a manner detrimental to the bank: safer types exit the market, and the average risk of the remaining pool rises.
  \item Since the lender's expected return is non-monotonic in the contracted rate, the bank may not be able to raise more revenue by raising the rate.
  \item The equilibrium outcome is credit rationing: at a rate $R^\star$ below the demand-clearing rate, the bank attracts a favorable pool of applicants but has more demand than it chooses to fund.
\end{itemize}

The implication for monetary policy is substantial.
In environments with significant ex-ante asymmetric information between borrowers and lenders, the interest rate is not a fully flexible instrument for clearing the credit market, and quantitative restrictions (credit rationing) may persist in equilibrium.
This mechanism is now a standard component of New Keynesian analyses of financial frictions, and it has been widely invoked to explain features of the credit channel of monetary transmission.

\subsection{Further reading}

The foundational contribution on costly state verification is \citet{Townsend1979}, which initiated the analysis of debt contracts as optimal incentive-compatible arrangements under ex-post informational asymmetry.
\citet{GaleHellwig1985} streamlined the analysis into the one-period formulation adopted in Sections~\ref{sec:GH-setup}--\ref{sec:GH-proof} and gave the cleanest characterization of the standard debt contract.
\citet{Diamond1984} extends the framework to the study of financial intermediation with delegated monitoring, providing a microfoundation for banks as efficient aggregators of monitoring costs.
For multilateral extensions and the role of seniority, see \citet{KrasaVillamil1994} and \citet{Winton1995}.
For the role of enforcement as a decision variable, see \citet{KrasaVillamil2000}.
On dynamic extensions of CSV, representative analyses include \citet{Chang1990}, \citet{Webb1992}, and \citet{HartMoore1998}.

The adverse-selection theory of credit rationing was introduced by \citet{StiglitzWeiss1981}.
Textbook treatments covering both branches of this chapter are provided by \citet{FreixasRochet1997}, especially for the banking theory perspective, and by \citet{EichbergerHarper1997}, which offers the expository framework closest to the presentation adopted here.
\citet{Tirole2006} provides the authoritative modern treatment of corporate finance under principal--agent frictions, covering debt contracts, credit rationing, and extensions to signalling, agency costs, and control rights.

%====================================================================
\chapter{Net Worth, Agency Rents, and Credit Rationing}
\label{ch:net-worth}
%====================================================================

Chapter~\ref{ch:debt-rationing} analyzed credit-market failures arising from two distinct frictions: ex-post costly state verification and ex-ante adverse selection.
This chapter develops a third, equally canonical, friction: ex-ante \emph{moral hazard} on the entrepreneur's effort choice.
The framework is the \emph{fixed-investment model} of \citet{HolmstromTirole1997}, which has become the workhorse of modern corporate finance and macrofinancial modelling.
It is stripped down to essentials---two dates, risk-neutral agents, binary effort, binary outcome---but rich enough to capture the central economics: when the entrepreneur's stake in the project is too small, he has an incentive to shirk in favor of a private benefit, and the lender will not fund the project unless the entrepreneur injects sufficient \emph{net worth} ex ante.

The model delivers several insights that will recur throughout Part~II of these notes.
First, a project with positive net present value (NPV) may fail to be financed: \emph{credit rationing} emerges endogenously from the moral-hazard friction, not from any market power or price friction.
Second, the minimum level of entrepreneurial net worth $\overline{A}$ that supports financing has a clean closed form, separating into an \emph{agency rent} term and a net-social-surplus term.
Third, the incentive structure induces the entrepreneur to invest his entire endowment in the project: partial investment strictly reduces the probability of funding without raising welfare.
Fourth, absent protective covenants, the entrepreneur has an incentive to \emph{overborrow}: after the initial contract has been signed, he would like to raise additional (even socially wasteful) funds, diluting the initial lenders' claim.
We develop each of these insights in turn.

The chapter corresponds to Part~A of Chapter~2 of \citet{Tirole2006}, on which the exposition is based.
Subsequent chapters in these notes will develop Parts B (debt overhang), C (the equity multiplier and borrowing capacity), and D (semi-verifiable income).

%====================================================================
\section{The Fixed-Investment Model}
\label{sec:HT-model}
%====================================================================

\subsection{Agents and the project}

There are two dates, $t \in \{0, 1\}$.
A single consumption good is available at each date, and all agents are risk-neutral with no time preference.
The economy is populated by two types of agents:
\begin{itemize}
  \item An \emph{entrepreneur} (the borrower or ``insider''), who has access to an investment project requiring a fixed capital outlay $I > 0$ at $t=0$.
  The entrepreneur begins with wealth $A$ (``cash in hand''), which he can either invest in the project or consume at $t=0$.
  We assume $A < I$: the entrepreneur cannot self-finance the project and must borrow at least $I - A$.
  \item \emph{Lenders} (the ``outsiders''), who have deep pockets and are willing to fund loans provided they break even in expectation.
  Lenders are risk-neutral and have a zero opportunity cost of funds: they are indifferent between consuming one unit at $t=0$ and receiving an expected one unit at $t=1$.
\end{itemize}

\subsection{Outcomes, actions, and private benefit}

The project has a binary outcome at $t=1$.
In the \emph{success} state, the project yields a verifiable (and contractible) income $R > 0$; in the \emph{failure} state, the project yields zero.
The probability of success depends on the entrepreneur's action:
\begin{itemize}
  \item If the entrepreneur \emph{behaves}---exerts effort, makes no side deals, pursues no perks---the probability of success is $p_h$, and the entrepreneur enjoys no private benefit.
  \item If the entrepreneur \emph{misbehaves}---shirks, takes a private benefit---the probability of success is $p_\ell < p_h$, and the entrepreneur enjoys a private benefit $B > 0$.
\end{itemize}
The private benefit $B$ has several economic interpretations.
It may represent the disutility of effort saved by shirking; or the entrepreneur's preference for a project that is easier to implement, more glamorous, or more personally rewarding; or the monetary value of the perks, side deals, and favors to friends that a shirking manager can arrange.
In all cases, $B$ is a non-contractible, non-pledgeable payoff that accrues exclusively to the entrepreneur when he misbehaves.
We write $\Delta p \coloneqq p_h - p_\ell > 0$ for the probability wedge.

The entrepreneur's action $e \in \{ \text{behave}, \text{misbehave} \}$ is \emph{not observable} by the lender.
This is the moral-hazard friction at the heart of the analysis.

\subsection{Preferences and the objective probability}

Both agents evaluate consumption streams $(c_0, c_1)$ according to
\begin{equation*}
  U(c_0, c_1; Q) = c_0 + \E^Q[c_1],
\end{equation*}
where $Q$ denotes the objective probability on $\{ \text{success}, \text{failure} \}$.
The probability $Q$ depends on the entrepreneur's action:
\begin{equation}\label{eq:HT-probability}
  Q(\text{success}) = \begin{cases} p_h, & \text{if the entrepreneur behaves}, \\ p_\ell, & \text{if the entrepreneur misbehaves}. \end{cases}
\end{equation}
Risk neutrality and the zero rate of time preference together imply that the social optimum depends only on expected net project returns, and that all the friction in the analysis derives from the informational asymmetry over the entrepreneur's action.

\subsection{The sharing rule and limited liability}

Because the project outcome is verifiable but the entrepreneur's action is not, any feasible contract between the two parties must be contingent on the observable outcome alone.
The contract specifies how the project's income $R$ is split in the success state; in the failure state, the project yields nothing and both parties receive zero by limited liability.

\begin{definition}[Sharing rule]\label{def:sharing-rule}
A \emph{sharing rule} is a pair $(R_b, R_\ell) \in \R_+^2$ satisfying
\begin{equation}\label{eq:HT-sharing}
  R_b + R_\ell = R, \quad R_b \geq 0, \quad R_\ell \geq 0.
\end{equation}
In the success state, the entrepreneur receives $R_b$ and the lender receives $R_\ell$; in the failure state, both receive zero.
\end{definition}

Limited liability ($R_b, R_\ell \geq 0$) is essential: it rules out the transfer of negative amounts in the bad state, which would otherwise reconstitute the first-best by penalizing failure retroactively.
The pair $(R_b, R_\ell)$ constitutes the \emph{incentive scheme} for the entrepreneur: $R_b$ in success, zero in failure.

\subsection{Timing}

The sequence of events in the fixed-investment model is:
\begin{enumerate}
  \item At $t=0$, the entrepreneur and the lender agree on a sharing rule $(R_b, R_\ell)$.
  The lender transfers $I - A$ to the entrepreneur, and $I$ is invested in the project.
  \item Immediately after the investment is made, moral hazard intervenes: the entrepreneur privately chooses to behave or misbehave.
  \item At $t=1$, the project outcome is realized.
  In success, $R_b$ goes to the entrepreneur and $R_\ell$ to the lender; in failure, both receive zero.
\end{enumerate}

%====================================================================
\section{Project Viability and Net Present Value}
\label{sec:HT-viability}
%====================================================================

\subsection{The NPV hypothesis}

We assume throughout the chapter that the project has positive NPV only under behaving:
\begin{assumption}[Project viability]\label{ass:HT-viability}
\begin{align}
  p_h R - I &> 0, \tag{V.1}\label{eq:NPV-h} \\
  p_\ell R - I + B &< 0. \tag{V.2}\label{eq:NPV-l}
\end{align}
\end{assumption}

The condition \eqref{eq:NPV-h} says that, under behaving, the expected project return exceeds the investment cost; the project is socially desirable when the entrepreneur works.
The condition \eqref{eq:NPV-l} says that, under misbehaving---even including the private benefit $B$---the project has negative NPV.
Assumption~\ref{ass:HT-viability} is the natural setting for the moral-hazard problem: if misbehaving were also socially desirable (i.e., \eqref{eq:NPV-l} reversed), the moral-hazard problem would not be economically interesting because both sides would be content with shirking.

\subsection{No equilibrium contract induces shirking}

A first implication of Assumption~\ref{ass:HT-viability} is that no sharing rule induces the entrepreneur to misbehave at equilibrium.
Suppose, for contradiction, that the equilibrium contract $(R_b, R_\ell)$ induces misbehaving.
The total expected payoff across both parties is
\begin{equation*}
  \underbrace{[p_\ell R_\ell - (I - A)]}_{\text{lender's NPV}} + \underbrace{[p_\ell R_b + B - A]}_{\text{borrower's NPV}}
  = p_\ell R + B - I.
\end{equation*}
By \eqref{eq:NPV-l}, this total is strictly negative.
Hence at least one of the two NPVs must be strictly negative, contradicting participation of at least one of the parties.
The attention is therefore restricted to contracts that \emph{induce behaving}.

%====================================================================
\section{Incentive Compatibility and the Pledgeable Income}
\label{sec:HT-IC}
%====================================================================

\subsection{The entrepreneur's incentive constraint}

Given a sharing rule $(R_b, R_\ell)$, the entrepreneur faces the trade-off between behaving and misbehaving.
His expected utility from behaving is $p_h R_b$; his expected utility from misbehaving is $p_\ell R_b + B$.
The contract induces behaving if and only if
\begin{equation}\label{eq:HT-IC-raw}
  p_h R_b \geq p_\ell R_b + B,
\end{equation}
or, equivalently,
\begin{equation}\label{eq:HT-IC}
  R_b \geq \frac{B}{\Delta p}.
  \tag{IC}
\end{equation}
The constraint \eqref{eq:HT-IC} is the \emph{incentive compatibility} constraint of the fixed-investment model: the entrepreneur's stake in the success state must be large enough that the probability-weighted gain from honest effort, $\Delta p \cdot R_b$, exceeds the private benefit $B$ that shirking delivers.

\subsection{The pledgeable income}

The IC constraint \eqref{eq:HT-IC} places a lower bound on the entrepreneur's stake.
Equivalently, it places an upper bound on how much of the project's income can be pledged to the lender:
\begin{equation*}
  R_\ell = R - R_b \leq R - \frac{B}{\Delta p}.
\end{equation*}

\begin{definition}[Pledgeable income]\label{def:pledgeable}
The \emph{expected pledgeable income} of the project is
\begin{equation}\label{eq:pledgeable}
  \mathcal{P} \coloneqq p_h \left( R - \frac{B}{\Delta p} \right).
\end{equation}
\end{definition}

The pledgeable income $\mathcal{P}$ is the maximum expected amount that the entrepreneur can credibly promise to repay the lender while maintaining incentive compatibility.
It is strictly less than the expected project income $p_h R$: the wedge $p_h B / \Delta p$ represents the \emph{rent} that the entrepreneur must keep to align incentives.

\subsection{Individual rationality}

The lender will finance the project only if his expected repayment covers his outlay.
Under risk neutrality and the zero opportunity cost, this \emph{individual rationality} (or \emph{breakeven}, or \emph{participation}) constraint takes the form
\begin{equation}\label{eq:HT-IR}
  p_h R_\ell \geq I - A.
  \tag{IR}
\end{equation}

Substituting the pledgeable-income upper bound $R_\ell \leq R - B/\Delta p$ into \eqref{eq:HT-IR} yields the fundamental necessary condition for financing: the pledgeable income must exceed the amount borrowed,
\begin{equation}\label{eq:HT-financeable}
  \mathcal{P} = p_h \left( R - \frac{B}{\Delta p} \right) \geq I - A.
\end{equation}

%====================================================================
\section{The Net-Worth Threshold and Credit Rationing}
\label{sec:HT-threshold}
%====================================================================

\subsection{The threshold \texorpdfstring{$\overline{A}$}{A-bar}}

Rearranging \eqref{eq:HT-financeable} gives a lower bound on the entrepreneur's net worth.

\begin{proposition}[Net-worth threshold]\label{prop:HT-threshold}
A necessary condition for the project to be financed is
\begin{equation}\label{eq:HT-threshold}
  A \geq \overline{A}, \quad \text{where} \quad \overline{A} \coloneqq p_h \frac{B}{\Delta p} - (p_h R - I).
\end{equation}
\end{proposition}

\begin{proof}
From \eqref{eq:HT-financeable},
\begin{equation*}
  I - A \leq p_h R - \frac{p_h B}{\Delta p},
\end{equation*}
equivalently
\begin{equation*}
  A \geq I - p_h R + \frac{p_h B}{\Delta p} = \frac{p_h B}{\Delta p} - (p_h R - I).
\end{equation*}
\end{proof}

We assume throughout that the threshold is strictly positive,
\begin{equation}\label{eq:HT-threshold-pos}
  \overline{A} > 0,
  \quad \text{equivalently,} \quad
  p_h R - I < \frac{p_h B}{\Delta p}.
\end{equation}
The left-hand side of the equivalent form is the project's NPV under behaving; the right-hand side is the minimum expected rent that must be left to the entrepreneur to preserve incentives.
Condition~\eqref{eq:HT-threshold-pos} states that the project's NPV is strictly less than this agency rent, so that in the absence of any initial wealth the entrepreneur cannot credibly commit to honest behavior.

\subsection{The credit-rationing regime}

\begin{proposition}[Credit rationing]\label{prop:HT-rationing}
If $A < \overline{A}$, the project cannot be financed even though, under behaving, it has positive NPV.
\end{proposition}

\begin{proof}
By \eqref{eq:HT-threshold}, the lender's participation constraint \eqref{eq:HT-IR} cannot be satisfied jointly with the incentive-compatibility constraint \eqref{eq:HT-IC}: the entrepreneur's required minimum stake $B/\Delta p$ leaves too little pledgeable income to repay $I - A$.
The lender will not finance the project at any sharing rule; the entrepreneur is \emph{rationed out} of the credit market.
\end{proof}

Proposition~\ref{prop:HT-rationing} illustrates a central theme of corporate finance: a project with positive social value need not be funded if the entrepreneur's ``skin in the game'' is too small.
The intuition is transparent: with low net worth, the entrepreneur must promise a large repayment in success, leaving him with a small residual stake; the small residual stake fails to motivate effort; foreseeing the inevitable shirking, the lender refuses to finance.
The friction is \emph{not} a price friction (the entrepreneur may be willing to pay a very high interest rate): no interest rate sustains the contract, because no interest rate can simultaneously satisfy (IC) and (IR).

\subsection{The financed regime}

\begin{proposition}[Financed regime]\label{prop:HT-financed}
If $A \geq \overline{A}$, then the entrepreneur can secure financing.
The unique competitive sharing rule satisfies
\begin{equation}\label{eq:HT-contract}
  R_\ell = \frac{I - A}{p_h},
  \quad
  R_b = R - R_\ell = R - \frac{I - A}{p_h}.
\end{equation}
Under this contract, (IC) holds---$R_b \geq B/\Delta p$---and the lender earns zero expected profit.
\end{proposition}

\begin{proof}
Under perfect competition among lenders, the participation constraint \eqref{eq:HT-IR} must bind at equilibrium: otherwise, the entrepreneur would secure a better sharing rule from a competing lender.
Setting \eqref{eq:HT-IR} to equality gives $R_\ell = (I - A)/p_h$, and $R_b = R - R_\ell$ follows from the resource constraint \eqref{eq:HT-sharing}.

We verify that $R_b \geq B/\Delta p$ when $A \geq \overline{A}$.
Using \eqref{eq:HT-threshold}, $A \geq \overline{A}$ is equivalent to $A + (p_h R - I) \geq p_h B / \Delta p$, i.e., $p_h R - (I - A) \geq p_h B / \Delta p$, i.e., $R - (I - A)/p_h \geq B/\Delta p$, which is precisely $R_b \geq B/\Delta p$.
Hence (IC) holds automatically under \eqref{eq:HT-contract} whenever $A \geq \overline{A}$.
\end{proof}

Note the striking feature of \eqref{eq:HT-contract}: the entrepreneur's stake $R_b$ is \emph{determined} by his initial wealth $A$, not chosen freely.
Competition drives $R_\ell$ to the lender's breakeven level, which in turn leaves the entrepreneur with whatever residual stake ensures breakeven.
For $A = \overline{A}$, (IC) binds: the entrepreneur's stake is exactly $B/\Delta p$.
For $A > \overline{A}$, (IC) is slack: the entrepreneur has more ``skin in the game'' than strictly required.

\subsection{``One only lends to the rich''}

The striking implication of Propositions~\ref{prop:HT-rationing} and~\ref{prop:HT-financed} is that financing depends on \emph{net worth}, not on willingness to pay.
A wealthy entrepreneur (high $A$) has no difficulty financing the project; a poor entrepreneur (low $A$) is rationed out, even if his project is socially valuable and he is willing to pledge the entire pledgeable income.
This is the celebrated ``one only lends to the rich'' phenomenon that underlies much of the macroeconomic literature on financial frictions, balance-sheet effects, and the credit channel of monetary transmission.

%====================================================================
\section{The Agency Rent Decomposition}
\label{sec:HT-agency-rent}
%====================================================================

\subsection{Decomposing the threshold}

The threshold $\overline{A}$ has a clean decomposition into two economically meaningful terms:
\begin{equation}\label{eq:HT-rent-decomp}
  \overline{A} = \underbrace{p_h \frac{B}{\Delta p}}_{\text{agency rent}} - \underbrace{(p_h R - I)}_{\text{NPV under behaving}}.
\end{equation}

The \emph{agency rent} $p_h B / \Delta p$ is the minimum expected monetary payoff that must be left to the entrepreneur to preserve incentives: multiplying the per-success minimum stake $B/\Delta p$ by the probability of success $p_h$ yields the expected surplus that the entrepreneur must capture.
This rent is a pure informational friction: it would be zero if effort were observable and contractible.

The \emph{NPV under behaving} $p_h R - I$ is the social surplus generated by the project when the entrepreneur works.
It is the ``gift'' that the project itself delivers to the economy.

The net-worth requirement $\overline{A}$ is the gap between the agency rent and the project's social surplus: the entrepreneur must bring enough of his own capital to bridge the difference.
Equivalently, his initial contribution plus the project's social surplus must cover the agency rent:
\begin{equation*}
  p_h \frac{B}{\Delta p} \leq (p_h R - I) + A.
\end{equation*}

\subsection{Competitive surplus allocation}

Under perfect competition among lenders, the entire social surplus accrues to the entrepreneur, \emph{conditional} on the project being funded.
Formally, define the entrepreneur's expected \emph{net} payoff as his gross expected payoff minus the opportunity cost of investing his own wealth:
\begin{equation*}
  \Pi_b \coloneqq p_h R_b - A.
\end{equation*}
In the financed regime ($A \geq \overline{A}$), substituting \eqref{eq:HT-contract} gives
\begin{equation*}
  \Pi_b = p_h R - (I - A) - A = p_h R - I,
\end{equation*}
which is the project's NPV under behaving.
In the rationing regime ($A < \overline{A}$), the project is not funded and $\Pi_b = 0$.

\begin{proposition}[Surplus allocation]\label{prop:surplus}
The entrepreneur's expected net payoff is
\begin{equation}\label{eq:Pi_b}
  \Pi_b = \begin{cases}
    0 & \text{if } A < \overline{A}, \\
    p_h R - I & \text{if } A \geq \overline{A}.
  \end{cases}
\end{equation}
In the financed regime, the lender earns zero expected profit and the entrepreneur captures the entire social surplus.
\end{proposition}

The discontinuity of $\Pi_b$ at $A = \overline{A}$ is crucial: a marginal decrease in net worth below the threshold collapses the surplus to zero.
This discontinuity is the source of the \emph{financial accelerator} mechanism that has been central to macro-finance since \citet{HolmstromTirole1997}: shocks to entrepreneurial net worth translate into shocks to investment activity through the threshold mechanism.

%====================================================================
\section{Full Investment of Entrepreneurial Assets}
\label{sec:HT-full-investment}
%====================================================================

We now verify that the entrepreneur invests his entire wealth $A$ in the project, rather than consuming a portion $c$ at $t=0$.

Suppose the entrepreneur consumes $c \in [0, A]$ at $t=0$ and invests only $A - c$.
His remaining wealth to serve as collateral against the loan is $A - c$, and the financing analysis of Section~\ref{sec:HT-threshold} applies with $A$ replaced by $A - c$.
The project is financed if and only if $A - c \geq \overline{A}$, i.e., $c \leq A - \overline{A}$ (which requires $A \geq \overline{A}$ to begin with).

The entrepreneur's total expected utility is $c$ (from consumption at $t=0$) plus his expected payoff $\Pi_b(A - c)$ from the project.
Conditional on the project being financed (i.e., $c \leq A - \overline{A}$),
\begin{equation*}
  \text{total utility} = c + (p_h R - I).
\end{equation*}
The second term is independent of $c$: by Proposition~\ref{prop:surplus}, the entrepreneur captures the entire NPV regardless of his initial contribution (as long as it exceeds the threshold).
The first term, $c$, is increasing in $c$.

\begin{proposition}[Full investment]\label{prop:full-investment}
If the entrepreneur has $A > \overline{A}$, the entrepreneur's optimal consumption at $t=0$ is $c^\ast = A - \overline{A}$, and the entrepreneur invests exactly $\overline{A}$ in the project.
More generally: the entrepreneur would strictly prefer to invest less than his entire wealth only if consuming more does not reduce the probability of funding.
\end{proposition}

\begin{proof}
Total utility is $c + (p_h R - I)$ when the project is financed, and $c$ when it is not.
On the interval $c \in [0, A - \overline{A}]$ (project financed), total utility is $c + (p_h R - I)$, increasing in $c$.
At $c = A - \overline{A}$, total utility equals $(A - \overline{A}) + (p_h R - I)$.
For $c > A - \overline{A}$ (project not financed), total utility is simply $c$, which is at most $A < (A - \overline{A}) + (p_h R - I)$ since $(p_h R - I) > \overline{A}$ by assumption that the NPV is large enough to prefer financing.

Hence the optimum is $c^\ast = A - \overline{A}$, and the entrepreneur invests $\overline{A}$.
\end{proof}

The conclusion is subtle: the entrepreneur does not invest his \emph{entire} wealth, but instead invests exactly the threshold amount $\overline{A}$, consuming the surplus $A - \overline{A}$ at $t=0$.
This reflects the fact that marginal investments above the threshold do not increase the probability of funding (which is already $1$) but do delay consumption.

Under a slight modification of the timing---e.g., if consuming at $t=0$ reduces collateral available ex post---the entrepreneur would be forced to invest his entire wealth.
The slides of the original course present the argument in this simpler form, emphasizing that the borrower cannot gain strategic benefit from under-investing.
The essential message is: \emph{at the margin of the financing boundary}, investing every additional unit of wealth strictly expands the set of fundable projects.

\begin{remark}[Alternative reading]\label{rem:full-investment-alternative}
In the slides, the proposition is presented as: ``the entrepreneur cannot gain by not investing her entire wealth in the project''.
This is true in the narrower sense that partial investment does not yield a higher expected payoff, holding the project funded; but it also does not yield a strictly lower one (in the simple two-date model with $\overline{A} < A$).
The practical content is that partial investment \emph{strictly} dominates consumption at $t=0$ \emph{only when} consumption would push the entrepreneur below the financing threshold.
\end{remark}

%====================================================================
\section{Dilution and Overborrowing}
\label{sec:HT-dilution}
%====================================================================

The preceding sections have focused on a one-shot contract at $t=0$.
A richer question is whether the initial contract is \emph{time-consistent}: after the initial loan is granted, is the entrepreneur tempted to raise additional funds from new lenders, diluting the claim of the initial lenders?
In practice, debt contracts include \emph{negative covenants} that explicitly prohibit such dilution.
The analysis below shows why such covenants are necessary.

\subsection{A deepening-investment opportunity}

Suppose, after the initial contract $(R_b, R_\ell)$ has been signed (with $R_b \geq B/\Delta p$), a \emph{deepening investment} opportunity arises.
At an additional cost $J > 0$ invested at $t=0^+$, the probability of success increases by $\tau > 0$ (under behaving, the probability becomes $p_h + \tau$; under misbehaving, $p_\ell + \tau$).
We assume that the deepening investment is \emph{socially inefficient}:
\begin{equation}\label{eq:HT-inefficient}
  C_1 \coloneqq J - \tau R > 0.
\end{equation}
The direct cost $J$ of the deepening investment exceeds the expected incremental return $\tau R$; a social planner would reject it.

\subsection{Timing with deepening investment}

The modified timing is:
\begin{enumerate}
  \item The entrepreneur borrows $I - A$ from initial lenders under sharing rule $(R_b, R_\ell)$.
  \item The entrepreneur contracts with new lenders to finance the deepening investment $J$.
  Under a new contract, the entrepreneur's initial stake $R_b$ is diluted into $\widehat{R}_b + \widehat{R}_\ell$, where $\widehat{R}_\ell$ is the new lenders' promised payment.
  \item Moral hazard: the entrepreneur chooses behave or misbehave.
  \item The project outcome is realized with modified probabilities $(p_h + \tau)$ or $(p_\ell + \tau)$.
\end{enumerate}

\subsection{Why the entrepreneur misbehaves after overborrowing}

A critical observation is that, after the deepening investment, the entrepreneur's stake $\widehat{R}_b$ is smaller than his original stake $R_b$.
If $\widehat{R}_b < B/\Delta p$, the IC constraint is violated, and the entrepreneur misbehaves.

To see why this is generic, recall that the deepening investment reduces total value by $C_1$ (by inefficiency).
This loss must be absorbed by some party.
The new lenders, being competitive, break even in expectation, so they do not absorb the loss.
The initial lenders, whose claim is fixed at $R_\ell$, could in principle \emph{benefit} from the deepening investment if it raises their probability of repayment, but their claim is pinned down.
Hence the entrepreneur bears the loss $C_1$, plus the additional cost of incentive distortion.

We quantify this incentive cost.
If the entrepreneur continues to behave after the deepening investment, the initial lenders' expected claim rises from $p_h R_\ell$ to $(p_h + \tau) R_\ell$, benefiting them by $\tau R_\ell$.
For the entrepreneur's net benefit to turn positive, he must misbehave.
Under misbehaving, the probability of success becomes $p_\ell + \tau$, and the total loss to society (relative to continued behaving without the deepening investment) is
\begin{equation}\label{eq:HT-C2}
  C_2 \coloneqq (\Delta p) R - B,
\end{equation}
which is positive by \eqref{eq:HT-IC} and the strict positivity implied by the project viability condition $\Delta p \cdot R \geq B$ (this is precisely what $R_b \geq B/\Delta p$ implies when $R_b \leq R$).

\subsection{When overborrowing occurs}

The entrepreneur's trade-off is: is the private benefit he captures---both the reduction in effort cost $B$ and the fresh stake $\widehat{R}_b$---worth the loss?

Under competition among new lenders, the new-lender breakeven requires
\begin{equation}\label{eq:HT-new-lender-BE}
  (p_\ell + \tau) \widehat{R}_\ell = J.
\end{equation}
The entrepreneur's gross payoff from overborrowing and misbehaving is $(p_\ell + \tau) \widehat{R}_b + B$; without overborrowing, under behaving, it is $p_h R_b$.
Overborrowing is privately profitable for the entrepreneur when
\begin{equation*}
  (p_\ell + \tau) \widehat{R}_b + B > p_h R_b.
\end{equation*}

Using $\widehat{R}_b = R_b - \widehat{R}_\ell$ (the dilution decomposition) and substituting \eqref{eq:HT-new-lender-BE},
\begin{equation*}
  (p_\ell + \tau)(R_b - \widehat{R}_\ell) + B = (p_\ell + \tau) R_b - J + B.
\end{equation*}
The overborrowing incentive therefore becomes
\begin{equation*}
  (p_\ell + \tau) R_b - J + B > p_h R_b,
\end{equation*}
equivalently,
\begin{equation}\label{eq:HT-overborrow-condition}
  \bigl[ p_h - (p_\ell + \tau) \bigr] R_b < B - J + \tau R_b.
\end{equation}
After rearrangement and using $R = R_b + R_\ell$ together with \eqref{eq:HT-C2}, \eqref{eq:HT-overborrow-condition} reduces to
\begin{equation}\label{eq:HT-overborrow-final}
  \bigl[ p_h - (p_\ell + \tau) \bigr] R_\ell > C_1 + C_2.
\end{equation}
The left-hand side of \eqref{eq:HT-overborrow-final} is the \emph{negative externality} imposed on the initial lenders: the deepening-plus-misbehaving combination reduces their success probability from $p_h$ to $p_\ell + \tau$, shrinking their expected claim by $[p_h - (p_\ell + \tau)] R_\ell$ (which is positive whenever $\Delta p > \tau$).

\begin{proposition}[Overborrowing occurs when initial lenders are exposed]\label{prop:overborrow}
Overborrowing is strictly profitable for the entrepreneur if and only if
\begin{equation*}
  \bigl[ p_h - (p_\ell + \tau) \bigr] R_\ell > C_1 + C_2.
\end{equation*}
That is, overborrowing occurs whenever the externality imposed on the initial lenders exceeds the total private cost of refinancing.
\end{proposition}

\subsection{Interpretation: the role of initial lender exposure}

Two economic observations emerge from Proposition~\ref{prop:overborrow}.

First, \eqref{eq:HT-overborrow-final} decomposes overborrowing into an arithmetic of externalities: the entrepreneur profits at the expense of the initial lenders when, and only when, their incentive-adjusted exposure exceeds the sum of the direct cost of refinancing ($C_1$) and the incentive cost of shirking ($C_2$).
The overborrowing transaction is therefore welfare-destroying: the entrepreneur's gain comes from transferring value \emph{away} from the initial lenders, while simultaneously generating the inefficiency $C_1 + C_2$.

Second, the magnitude of $R_\ell$ is central.
Recall from \eqref{eq:HT-contract} that $R_\ell = (I - A)/p_h$ in the competitive contract: the weaker is the entrepreneur's balance sheet (smaller $A$), the larger is $R_\ell$.
Overborrowing is therefore more likely to occur in economies with weakly capitalized entrepreneurs, because their initial lenders have the largest claims exposed to dilution.

\subsection{The role of covenants}

The analysis explains why debt contracts typically include \emph{negative covenants} prohibiting the issuance of new securities, the incurrence of additional debt, or the alteration of the capital structure without the initial lender's consent.
Two rationales, both visible in the mechanics above, support such covenants:
\begin{itemize}
  \item \emph{Seniority protection.}
  Initial lenders do not want the borrower to issue claims of higher seniority, which would reduce their residual claim in default states.
  \item \emph{Incentive preservation.}
  More subtly, the issuance of new securities may \emph{alter managerial incentives} by reducing the entrepreneur's stake, triggering the shift from behaving to misbehaving and lowering the probability of success.
\end{itemize}
The second rationale is the deeper one: covenants exist not merely to protect priority of claims but to preserve the incentive architecture of the initial contract.

\subsection{Further reading}

The canonical reference for the fixed-investment model is \citet{HolmstromTirole1997}, who developed it to study the interaction between financial intermediation, entrepreneurial net worth, and real investment.
The framework underlies much of the modern macro-finance literature on the financial accelerator and the credit channel of monetary policy.
\citet{Tirole2006} (Chapter~3) is the standard textbook exposition and extends the model in several directions that will be developed in subsequent chapters of these notes: debt overhang (Part~B of the original course), the equity multiplier and borrowing capacity (Part~C), and semi-verifiable income (Part~D).
For earlier treatments of moral hazard in contracting, the reader is referred to the classical principal--agent literature, especially \citet{JensenMeckling1976} on agency costs of debt and equity.

%====================================================================
\chapter{Debt Overhang}
\label{ch:debt-overhang}
%====================================================================

Chapter~\ref{ch:net-worth} analyzed the canonical fixed-investment model in its static form: at $t=0$ the entrepreneur arrives with some initial wealth $A$, contracts with a lender, invests, and the outcome is realized at $t=1$.
In practice, however, firms rarely begin the moral-hazard game with a clean slate.
A new investment opportunity typically arises against the backdrop of an \emph{existing} capital structure---prior debt, outstanding bonds, legacy obligations to suppliers or employees---that was issued to fund earlier projects and is still due.
These pre-existing claims interact with the incentive problem of the new investment in subtle ways.

This chapter develops a simple but powerful illustration of that interaction: the phenomenon of \emph{debt overhang}, first identified by \citet{Myers1977}.
The central insight is striking: a firm that carries a heavy burden of senior legacy debt may fail to raise funds for a new project, even when the new project has strictly positive NPV, even when the moral-hazard friction alone would not prevent financing, and even when all parties would strictly benefit from the new investment being made.
The friction is not that capital is unavailable or that incentives cannot be aligned; it is that the returns on the new project would flow first to the old creditors, leaving too little for the new investors to break even.
In the presence of the legacy debt, the pledgeable income of the \emph{new} project---which is what new investors can credibly claim---falls short of the required investment.

Debt overhang is a problem of \emph{coordination} across layers of the capital structure.
Its resolution, as we will see, requires \emph{renegotiation} between the firm and its legacy creditors.
If renegotiation is feasible---if the old creditors can be persuaded to forgive a portion of their face value in exchange for a share of the gains from the new project---all parties can do strictly better.
If renegotiation is blocked---because the legacy creditors are too dispersed, or because legal frictions prevent restructuring---the overhang persists, and a socially valuable project goes unfunded.
The chapter develops both cases carefully and quantifies the minimum debt forgiveness that restores financing.

The chapter follows Part~B of Chapter~2 of \citet{Tirole2006}, which builds directly on the fixed-investment framework of Chapter~\ref{ch:net-worth}.

%====================================================================
\section{Setup and the Overhang Hypothesis}
\label{sec:DO-setup}
%====================================================================

\subsection{The fixed-investment model revisited}

We retain the notation and structure of Chapter~\ref{ch:net-worth}: a risk-neutral entrepreneur with wealth $A$ considers a project requiring investment $I$, with binary outcome---success at income $R$ with probability $p_h$ if the entrepreneur behaves, or $p_\ell$ with private benefit $B$ if he misbehaves.
The incentive-compatibility constraint \eqref{eq:HT-IC} requires $R_b \geq B/\Delta p$, the pledgeable income is $\mathcal{P} = p_h(R - B/\Delta p)$, and the net-worth threshold for financing is
\begin{equation*}
  \overline{A} = p_h \frac{B}{\Delta p} - (p_h R - I).
\end{equation*}

In Chapter~\ref{ch:net-worth}, we maintained the assumption $\overline{A} > 0$: the project's NPV under behaving was smaller than the agency rent, so that credit rationing was a binding concern for low-wealth entrepreneurs.
In the present chapter, we adopt the \emph{opposite} hypothesis.

\begin{assumption}[Profitable project]\label{ass:DO-profitable}
The project is sufficiently profitable that it would be financeable even by an entrepreneur with zero net worth:
\begin{equation}\label{eq:DO-profitable}
  \overline{A} < 0,
  \quad \text{equivalently,} \quad
  p_h R - I > p_h \frac{B}{\Delta p}.
\end{equation}
\end{assumption}

Under Assumption~\ref{ass:DO-profitable}, the project's NPV exceeds the agency rent; there is \emph{strictly positive pledgeable surplus} even when the entrepreneur contributes nothing ($A = 0$).
In the absence of any legacy claims, a competitive lender would happily fund the project, extracting $I$ in expectation while leaving the entrepreneur with the residual rent $p_h B / \Delta p$ and the lender with zero expected profit.
The object of this chapter is to show that, when legacy debt interferes with the cash flow of the \emph{new} project, this conclusion can fail dramatically.

\subsection{Legacy senior debt}

Suppose the entrepreneur has previously been granted a long-term loan on some earlier project.
At $t=1$---the date at which the new project, if undertaken, resolves---the entrepreneur owes a face value $D > 0$ to the initial investors.
Two features of the legacy debt are economically essential.

First, the legacy debt is \emph{contractually senior}: any cash flow generated at $t=1$ must first be applied to retiring the face value $D$ before any payment can be made to claimants arising from the new project.
Seniority is enforced by bankruptcy law and by the covenants embedded in the original debt contract.

Second, we assume the entrepreneur's position at $t=0$ is one of \emph{exhausted liquidity}: $A = 0$.
The entrepreneur has no cash to contribute to the new investment and, absent fresh financing, the project cannot be undertaken.
If the project is not undertaken, the legacy lenders receive zero at $t=1$: the underlying firm has no cash flow of its own, and the collateral of the earlier project has presumably been exhausted in prior states of the world that led to this position.

\subsection{The overhang severity condition}

The third element of the setup is a quantitative condition that makes the overhang ``serious'' enough to bite.

\begin{assumption}[Severe overhang]\label{ass:DO-severe}
The legacy debt satisfies
\begin{equation}\label{eq:DO-severe}
  p_h D > -\overline{A} = (p_h R - I) - p_h \frac{B}{\Delta p}.
\end{equation}
\end{assumption}

The right-hand side of \eqref{eq:DO-severe} is the project's pledgeable surplus: the NPV under behaving minus the agency rent.
The left-hand side, $p_h D$, is the expected claim of the legacy creditors on the new project's success outcome.
Assumption~\ref{ass:DO-severe} says that the legacy claim, in expectation, exceeds the pledgeable surplus of the new project: if the legacy debt is honored in full, there is not enough left over to attract fresh financing.

The interplay of Assumptions~\ref{ass:DO-profitable} and~\ref{ass:DO-severe} frames the problem.
By \eqref{eq:DO-profitable}, the project is socially desirable: its NPV exceeds the agency rent, and a ``clean-slate'' economy would fund it easily.
By \eqref{eq:DO-severe}, the overhang is severe: senior debt obligations consume enough of the expected cash flow that, in the absence of renegotiation, there is insufficient residual to compensate new investors.
The following sections examine the equilibrium implications.

%====================================================================
\section{When Initial Investors Retain Financing Capacity}
\label{sec:DO-with-cash}
%====================================================================

As a benchmark, we first consider the case in which the \emph{initial} investors are themselves able to finance the new investment $I$ directly.
This scenario is the simplest to analyze and serves as a foil for the main result.

Suppose the initial investors have sufficient cash reserves to undertake the new investment.
They then face a simple decision: either refuse to participate---in which case the project goes unfunded and they receive zero on their legacy claim---or finance the project themselves, extracting the pledgeable income in exchange for the new investment and (implicitly) forgiving the legacy debt.

Under the second option, the initial investors act as the sole financier of the new project.
In exchange for their outlay $I$, they claim the full pledgeable cash flow $R - B/\Delta p$ in the success state, leaving the entrepreneur with exactly $B/\Delta p$, the minimum consistent with incentive compatibility.
The initial investors' expected net payoff is
\begin{equation}\label{eq:DO-initial-NPV}
  p_h \left( R - \frac{B}{\Delta p} \right) - I = -\overline{A} > 0,
\end{equation}
which is strictly positive by Assumption~\ref{ass:DO-profitable}.
The entrepreneur, in turn, receives $p_h B / \Delta p$ in expectation rather than zero.
Both parties strictly benefit; the project is undertaken.

The legacy debt $D$ plays no direct role in this benchmark: the initial investors effectively write it off in exchange for the pledgeable income on the new project.
Seniority is irrelevant once the initial investors are the residual claimants to the new cash flow.

\subsection{The benchmark takeaway}

When the initial investors are themselves liquid, the problem of debt overhang simply does not arise.
The initial investors can take over the financing of the new project directly, extract the pledgeable income, and recover strictly positive profit on the legacy-plus-new combined investment.
The rest of the chapter addresses the case that is economically interesting: the initial investors \emph{cannot} themselves finance the new investment, and the entrepreneur must turn to a fresh set of external lenders.

%====================================================================
\section{The Core Overhang Result}
\label{sec:DO-core}
%====================================================================

\subsection{The pledgeable income net of seniority}

Suppose the initial investors lack the cash to finance the new investment.
The entrepreneur must turn to a new set of external lenders.
Call them the \emph{new} investors, in contrast to the legacy (``initial'') investors who hold the claim $D$.

The new investors face two frictions.
First, the incentive-compatibility constraint \eqref{eq:HT-IC} requires the entrepreneur to retain a minimum stake $B/\Delta p$ in the success cash flow.
Second---crucially---the legacy debt is senior: any payment to the new investors at $t=1$ can come only from the residual after the legacy face value $D$ has been settled.

The maximum amount that can be credibly pledged to the new investors in the success state is therefore the \emph{new pledgeable income},
\begin{equation}\label{eq:DO-new-pledgeable}
  R - \frac{B}{\Delta p} - D.
\end{equation}
In words: the total success cash flow $R$, less the minimum entrepreneurial stake $B/\Delta p$, less the senior legacy claim $D$, is what remains for the new investors.

\subsection{The financing condition}

New investors are willing to finance the project if and only if their expected return on the pledgeable income meets their opportunity cost of funds (normalized to zero, as in Chapter~\ref{ch:net-worth}):
\begin{equation}\label{eq:DO-IR}
  p_h \left( R - \frac{B}{\Delta p} - D \right) \geq I.
\end{equation}
Rearranging \eqref{eq:DO-IR} and using the definition of $\overline{A}$,
\begin{equation}\label{eq:DO-IR-rearranged}
  p_h R - \frac{p_h B}{\Delta p} - p_h D \geq I
  \Longleftrightarrow
  - \overline{A} \geq p_h D
  \Longleftrightarrow
  \overline{A} + p_h D \leq 0.
\end{equation}

The following proposition records the main conclusion.

\begin{proposition}[Debt overhang blocks financing]\label{prop:DO-main}
Under Assumptions~\ref{ass:DO-profitable} and~\ref{ass:DO-severe}, new investors are unwilling to finance the project.
The project---despite being of positive NPV in the absence of the legacy debt---remains unfunded.
\end{proposition}

\begin{proof}
The new investors' participation constraint \eqref{eq:DO-IR} is equivalent to $\overline{A} + p_h D \leq 0$ by \eqref{eq:DO-IR-rearranged}.
Under Assumption~\ref{ass:DO-severe}, $p_h D > -\overline{A}$, i.e., $\overline{A} + p_h D > 0$, so \eqref{eq:DO-IR} fails.
The new investors will not provide the investment $I$, and the project cannot be undertaken.
\end{proof}

\subsection{Interpretation}

Proposition~\ref{prop:DO-main} is the canonical statement of \emph{debt overhang} in the sense of \citet{Myers1977}.
Its significance is best appreciated by contrasting the outcome with the ``clean-slate'' counterfactual.
Were the entrepreneur free of the legacy debt $D$ (or, equivalently, were $D = 0$), the financing constraint would reduce to $-\overline{A} \geq 0$, which holds by Assumption~\ref{ass:DO-profitable}: the project is socially viable and new investors would willingly finance it.
The pledgeable surplus $-\overline{A}$ would accrue to the new investors as their breakeven margin, and the entrepreneur would capture the rent $p_h B/\Delta p$.
All parties would strictly benefit.

The legacy debt, however, \emph{appropriates} the pledgeable surplus in expectation before it can reach the new investors.
In state of success, the $t=1$ cash flow of $R$ is consumed by the triple claim: $B/\Delta p$ to the entrepreneur (by incentive compatibility), $D$ to the legacy investors (by seniority), and at most $R - B/\Delta p - D$ to the new investors.
When $p_h D > -\overline{A}$, the legacy claim exceeds what the pledgeable surplus can comfortably absorb, and new investors will not break even.

The failure here is neither a failure of the project nor a failure of the entrepreneur.
The entrepreneur is willing to behave, the project has positive social value, new investors have capital to deploy, and even the initial investors would prefer the project to go forward (since under no-project they receive zero, whereas under any project they receive something).
Yet the seniority structure \emph{blocks} the financing: the new investors cannot be assured of recovering their outlay.
The problem is coordination, not fundamentals.

%====================================================================
\section{Renegotiation as a Resolution}
\label{sec:DO-renegotiation}
%====================================================================

The debt overhang identified in Proposition~\ref{prop:DO-main} admits a natural resolution: the initial investors and the entrepreneur can \emph{renegotiate} the terms of the legacy debt, reducing the face value to an amount that accommodates fresh external financing while still leaving the initial investors better off than under no-project.

\subsection{A menu of efficient renegotiations}

Suppose the initial investors agree to reduce the face value of the legacy debt from $D$ to some $d < D$.
After renegotiation, the new investors face a revised financing condition in which the senior claim is $d$ rather than $D$.
By the same derivation as in Section~\ref{sec:DO-core}, new investors will finance the project if and only if
\begin{equation}\label{eq:DO-IR-renegotiated}
  p_h \left( R - \frac{B}{\Delta p} - d \right) \geq I,
  \quad \text{equivalently,} \quad
  \overline{A} + p_h d \leq 0.
\end{equation}

For the project to just meet the new investors' break-even, set \eqref{eq:DO-IR-renegotiated} to equality:
\begin{equation}\label{eq:DO-d-min}
  p_h d + \overline{A} = 0,
  \quad \text{i.e.,} \quad
  d = -\frac{\overline{A}}{p_h}.
\end{equation}

Call this the \emph{minimum renegotiated face value} consistent with restoration of financing.
Note that $d \in (0, D)$ by Assumptions~\ref{ass:DO-profitable} (which gives $d > 0$) and~\ref{ass:DO-severe} (which gives $d < D$).

\subsection{Payoffs after renegotiation at \texorpdfstring{$d$}{d}}

Under the renegotiated debt $d$ satisfying \eqref{eq:DO-d-min}, the new investors' expected return on the project is
\begin{equation}\label{eq:DO-new-zero}
  p_h \left( R - \frac{B}{\Delta p} - d \right) = I,
\end{equation}
i.e., exactly zero expected profit: new investors break even.
The initial investors' expected return rises from $0$ (no-project case) to
\begin{equation}\label{eq:DO-initial-gain}
  p_h d = -\overline{A} > 0,
\end{equation}
a strict improvement over the status quo.
The entrepreneur, who would otherwise receive $0$, captures the residual rent
\begin{equation}\label{eq:DO-entrepreneur-rent}
  p_h \cdot \frac{B}{\Delta p} > 0.
\end{equation}

\begin{proposition}[Renegotiation restores financing]\label{prop:DO-renegotiation}
Under Assumptions~\ref{ass:DO-profitable} and~\ref{ass:DO-severe}, there exists a renegotiated face value $d \in (0, D)$ satisfying \eqref{eq:DO-d-min} such that, after renegotiation:
\begin{enumerate}
  \item[\textup{(i)}] new investors are willing to finance the project at breakeven;
  \item[\textup{(ii)}] initial investors receive $p_h d = -\overline{A} > 0$ in expectation, strictly preferring renegotiation to the status quo;
  \item[\textup{(iii)}] the entrepreneur captures the rent $p_h B / \Delta p > 0$ and undertakes the project.
\end{enumerate}
All three parties strictly benefit from the renegotiation.
\end{proposition}

\begin{proof}
The existence of $d \in (0, D)$ solving $p_h d + \overline{A} = 0$ follows from Assumptions~\ref{ass:DO-profitable} ($\overline{A} < 0 \Rightarrow d > 0$) and~\ref{ass:DO-severe} ($-\overline{A} < p_h D \Rightarrow d < D$).
Given this $d$, \eqref{eq:DO-new-zero} gives the new investors' breakeven at strict equality, proving (i).
The initial investors' expected payoff under the renegotiated contract is $p_h d = -\overline{A} > 0$, compared with $0$ under no-project, proving (ii).
The entrepreneur, by incentive compatibility, retains exactly $B/\Delta p$ in success, yielding expected rent $p_h B / \Delta p$ per \eqref{eq:DO-entrepreneur-rent}, proving (iii).
\end{proof}

\subsection{The critical role of renegotiation}

Proposition~\ref{prop:DO-renegotiation} delivers the central normative message of the chapter.
Whenever the underlying project is socially profitable and the legacy debt can be renegotiated, the coordination failure of Proposition~\ref{prop:DO-main} can be avoided.
Initial investors are willing to write down a portion of their claim because partial recovery strictly dominates zero recovery.
The entrepreneur is willing to go forward because the alternative is no-project and no rent.
New investors enter because, at the reduced face value $d$, the pledgeable income is sufficient to cover their outlay.

The ``debt overhang'' of Proposition~\ref{prop:DO-main} is therefore better understood as a \emph{renegotiation-failure} phenomenon: the fundamentals of the economy support financing; it is the failure of renegotiation that blocks the efficient outcome.

\subsection{When renegotiation fails}

In practice, renegotiation is not always feasible.
Three empirically important obstacles deserve mention:

\begin{itemize}
  \item \emph{Dispersed bondholders.}
  When the legacy debt is held by a large, dispersed group of bondholders, organizing unanimous consent to a face-value reduction is difficult.
  Each bondholder has an incentive to hold out for the original face value, free-riding on the concessions of others---a coordination failure among the creditors themselves.
  \item \emph{Contractual rigidities.}
  The legacy debt may include covenants, tax rules, or regulatory provisions that make restructuring costly or legally complex, even when all parties are informally willing to renegotiate.
  \item \emph{Information asymmetries.}
  If the value of the new project is imperfectly observable to the initial investors, they may be unwilling to accept a face-value reduction on favorable terms, fearing that they are being exploited by a strategic entrepreneur.
\end{itemize}

In any of these cases, the debt overhang of Proposition~\ref{prop:DO-main} becomes the de facto outcome, and a socially profitable project remains unfunded.
The practical upshot is that the design of corporate debt contracts---and, in particular, the concentration of the creditor base---has first-order implications for the firm's ability to raise new capital when investment opportunities arise.

%====================================================================
\section{The Division of the Gains from Trade}
\label{sec:DO-division}
%====================================================================

Proposition~\ref{prop:DO-renegotiation} focused on the minimum renegotiated face value $d$ that just restores financing.
At this minimum, the initial investors extract their largest share of the gains from trade: $p_h d = -\overline{A}$.
In practice, the division of the gains depends on the relative bargaining power of the entrepreneur and the initial investors.

At the opposite extreme, imagine the entrepreneur has all the bargaining power and offers the initial investors a take-it-or-leave-it face value reduction to some $d' \leq d$.
The initial investors accept any $d' > 0$ (since $0$ is the no-project outcome), but they will not accept $d' < 0$.
The entrepreneur's optimal strategy is to offer $d' = 0^+$, i.e., total debt forgiveness, in which case the initial investors receive (essentially) zero and the entrepreneur captures the entire surplus.

More generally, any renegotiated face value $d^\ast \in [0, d]$, with $d = -\overline{A}/p_h$, generates a feasible renegotiation.
The new investors' break-even condition shifts continuously in $d^\ast$: smaller $d^\ast$ implies more of the pledgeable income accrues to the new investors (and by the break-even requirement, a smaller slice of the project is ultimately priced for them), while larger $d^\ast$ implies more of the pledgeable income accrues to the initial investors.

\begin{proposition}[Bargaining frontier]\label{prop:DO-bargaining}
Under Assumptions~\ref{ass:DO-profitable} and~\ref{ass:DO-severe}, the set of feasible renegotiated face values $d^\ast$ that restore financing is
\begin{equation*}
  d^\ast \in \left[ 0, \ -\frac{\overline{A}}{p_h} \right].
\end{equation*}
Across this interval, the expected payoff to the initial investors ranges continuously from $0$ (at $d^\ast = 0$, most favorable to the entrepreneur) to $-\overline{A}$ (at $d^\ast = -\overline{A}/p_h$, most favorable to the initial investors).
The entrepreneur's expected rent $p_h B / \Delta p$ is \emph{invariant} to $d^\ast$; the variation in surplus is absorbed entirely by the split between initial and new investors.
\end{proposition}

\begin{proof}
For any $d^\ast \in [0, -\overline{A}/p_h]$, the new investors' break-even requires
\begin{equation*}
  p_h(R - B/\Delta p - d^\ast) \geq I \Longleftrightarrow \overline{A} + p_h d^\ast \leq 0,
\end{equation*}
which holds for all $d^\ast \leq -\overline{A}/p_h$.
The initial investors receive $p_h d^\ast$ in expectation, ranging from $0$ to $-\overline{A}$.
The entrepreneur's payoff is $p_h \cdot B/\Delta p$ (by incentive compatibility) independent of $d^\ast$.
The new investors receive the residual pledgeable income, which closes the budget identity.
\end{proof}

The proposition emphasizes two economic points.
First, the \emph{total} gains from trade---the pledgeable surplus $-\overline{A}$ plus the entrepreneur's rent $p_h B / \Delta p$---are fixed, determined entirely by the fundamentals of the project.
Second, the entrepreneur's share of the total (his rent $p_h B / \Delta p$) is pinned down by the incentive constraint and is independent of how the renegotiation is resolved.
The bargaining over $d^\ast$ is thus a pure redistribution between the \emph{initial} and the \emph{new} investors; the entrepreneur is a bystander to the dispute.

%====================================================================
\section{Discussion and Further Reading}
\label{sec:DO-discussion}
%====================================================================

The debt-overhang model developed in this chapter has become a central tool in the analysis of corporate capital structure, sovereign debt crises, and macroeconomic balance-sheet dynamics.
Three directions of extension are worth noting.

\subsection{Myers' original analysis and its legacy}

The canonical reference for the debt-overhang phenomenon is \citet{Myers1977}.
Myers observed that firms with significant legacy debt may systematically underinvest relative to the first-best: the presence of senior claims on future cash flows distorts the manager's investment calculus, because gains from new investment accrue first to the legacy creditors rather than to the (equity-holding) manager.
The model of this chapter is a specific formalization of Myers' insight within the fixed-investment framework of Chapter~\ref{ch:net-worth}; the broader Myers model does not rely on moral hazard but on the differential claims on incremental cash flow.
The two formulations convey the same economic intuition: senior debt creates an ex-post tax on new investment that can dominate the project's NPV.

\subsection{Debt and hard claims in corporate governance}

\citet{HartMoore1995} develop a complementary perspective on senior debt.
In their framework, hard claims (including debt) play a disciplinary role in corporate governance by constraining the manager's ability to waste free cash flow.
The optimal capital structure trades off the disciplinary benefit of debt against the overhang cost: too little debt leaves the manager with excessive discretion; too much debt creates overhang and distorts future investment.
The Hart--Moore model is formally distinct from the present chapter but economically related; together they articulate the central tension in the design of corporate capital structure under agency frictions.

\subsection{Renegotiation with noncontractible investment}

\citet{BhattacharyaFaureGrimaud2001} extend the debt-overhang analysis to environments in which the new investment itself is non-contractible, so that the renegotiation must be designed to simultaneously restore financing and provide correct ex-ante investment incentives.
In such environments, the simple face-value reduction analyzed in Section~\ref{sec:DO-renegotiation} may fail: if the reduction is anticipated, the entrepreneur may distort his investment decisions to exploit the renegotiation mechanism.
More sophisticated renegotiation designs---with state-contingent transfers or sequential commitments---are required to implement the efficient outcome.
The Bhattacharya--Faure-Grimaud analysis is a useful bridge between the static debt-overhang framework of this chapter and the dynamic theories of debt restructuring that dominate contemporary sovereign-debt and corporate-bankruptcy theory.

\subsection{Macroeconomic implications}

In macroeconomic settings, debt overhang operates at the level of entire sectors or economies.
A country or firm with a large burden of legacy sovereign or external debt may systematically underinvest in new capital, because the marginal product of investment is taxed by the senior debt.
The mechanism has been invoked to explain episodes of prolonged investment slumps following banking crises, debt-fueled recessions, and sovereign-debt renegotiations.
The present chapter provides the microfoundation; the macroeconomic literature builds the equilibrium and aggregate implications.

\subsection{Further reading}

The canonical reference for debt overhang is \citet{Myers1977}; the present chapter adapts the fixed-investment formalization of \citet{HolmstromTirole1997} to that setting, following \citet{Tirole2006} (Chapter~2, Part~B).
\citet{HartMoore1995} develops the complementary analysis of debt's disciplinary role; \citet{BhattacharyaFaureGrimaud2001} extend debt-overhang to environments with noncontractible investment.
For broader surveys of corporate finance under the agency perspective---including debt overhang and its interaction with other frictions---see \citet{Tirole2006} and the references therein.
For dynamic extensions of the underlying debt contract and for the treatment of default as renegotiation, see \citet{HartMoore1998}, cited in Chapter~\ref{ch:debt-rationing}.

%====================================================================
\chapter{Borrowing Capacity and the Equity Multiplier}
\label{ch:equity-multiplier}
%====================================================================

Chapters~\ref{ch:net-worth} and~\ref{ch:debt-overhang} studied a \emph{fixed-investment} model: the entrepreneur's project required exactly $I$ units of capital at $t=0$, and financing either succeeded or failed as a binary event.
That formulation was well-suited to deliver the basic insight that low net worth can prevent financing.
It is too rigid, however, to answer a more nuanced question that is central to applied corporate finance: how much can an entrepreneur with wealth $A$ actually borrow?
The fixed-investment model collapses this question into a yes/no answer via the threshold $\overline{A}$.
The \emph{continuous-investment} model developed in this chapter relaxes that rigidity by allowing the investment scale $I$ to be chosen optimally as a function of the entrepreneur's wealth.

The central construct that emerges is the \emph{equity multiplier} $k$: the entrepreneur's maximum investment is $I = k A$, with $k > 1$ arising endogenously from the interplay of the incentive friction and the constant-returns technology.
The multiplier depends on the two fundamental parameters of the agency problem---the private benefit $B$ and the likelihood ratio $\Delta p / p_h$---and shrinks as either the private benefit rises or the likelihood ratio falls.
Its reciprocal minus one, $d = k - 1$, is the firm's \emph{borrowing capacity per unit of cash}: the maximum loan the entrepreneur can credibly raise for each dollar of his own equity injected.

The equity multiplier is the workhorse tool of corporate finance under agency frictions.
It underlies modern macroeconomic models of financial accelerators, balance-sheet channels of monetary transmission, and sudden-stop crises in emerging markets.
This chapter develops the theoretical foundations, studies the comparative statics that generate the link between agency parameters and investment scale, and extends the model to incorporate \emph{collateral value}---the part of the firm's cash flow that accrues even in failure states, because of the salvage value of tangible assets.
The extended model delivers a simple but empirically important conclusion: \emph{firms with more tangible assets borrow more}.
Equivalently, credit rationing is more binding for firms whose value depends largely on intangible assets or human capital, which have little liquidation value and therefore cannot serve as collateral for external financing.

The chapter follows Part~C of Chapter~2 of \citet{Tirole2006}, adapting the fixed-investment framework of Chapter~\ref{ch:net-worth} to the continuous-investment setting introduced in \citet{HolmstromTirole1997}.

%====================================================================
\section{The Continuous-Investment Model}
\label{sec:EM-setup}
%====================================================================

\subsection{Constant returns to scale}

We retain the two-period, risk-neutral, binary-outcome structure of Chapter~\ref{ch:net-worth}, but we now assume that the investment technology exhibits \emph{constant returns to scale}.
An investment $I \geq 0$ at $t=0$ yields, at $t=1$:
\begin{itemize}
  \item $R I$ in the success state, for some $R > 0$ (a per-unit productivity parameter);
  \item $0$ in the failure state.
\end{itemize}
The entrepreneur's private benefit from misbehaving is likewise proportional to the scale of investment: misbehaving yields private benefit $B I$, where $B > 0$ is the per-unit private benefit.
As before, the probability of success is $p_h$ under behaving and $p_\ell < p_h$ under misbehaving, with $\Delta p = p_h - p_\ell > 0$.

The constant-returns assumption is a modelling device, not a claim about the empirics of production.
It is the simplest way to pose the question ``how large can the investment be, as a function of $A$?'' without introducing additional curvature that would obscure the financial frictions.
Under decreasing returns, the optimal investment would be jointly constrained by the agency friction and by the diminishing marginal product of capital; the constant-returns specification isolates the agency-driven limit.

\subsection{The sharing rule}

As in Chapter~\ref{ch:net-worth}, the contract between the entrepreneur and the lender specifies how the project's success income $R I$ is split.
We write $R_b, R_\ell \geq 0$ for the payments (in the success state) to the entrepreneur and the lender, respectively, with the resource constraint
\begin{equation}\label{eq:EM-sharing}
  R_b + R_\ell = R I.
\end{equation}
In the failure state, both parties receive zero by limited liability.

The entrepreneur begins with wealth $A \geq 0$ and needs to borrow $I - A$ to finance a project of size $I$ (for $I > A$; the case $I \leq A$ is self-financed).

\subsection{Assumptions on the per-unit parameters}

We impose three conditions on the technology, generalizing the corresponding assumptions of Chapter~\ref{ch:net-worth}:

\begin{assumption}[Project viability, continuous version]\label{ass:EM-viability}
Per unit of investment:
\begin{align}
  p_h R - 1 &> 0, \tag{V.1'}\label{eq:EM-NPV-h} \\
  p_\ell R + B - 1 &< 0. \tag{V.2'}\label{eq:EM-NPV-l}
\end{align}
\end{assumption}

Assumption~\ref{ass:EM-viability} says that, per unit of investment, the project has positive NPV if the entrepreneur behaves and negative NPV if he does not.
As a corollary, $R > B/\Delta p$: the success income per unit exceeds the minimum incentive stake per unit.

\begin{assumption}[Finite-investment condition]\label{ass:EM-finite}
The expected pledgeable income per unit of investment is strictly less than one:
\begin{equation}\label{eq:EM-finite}
  p_h R - 1 < p_h \frac{B}{\Delta p},
  \quad \text{equivalently,} \quad
  \mathcal{P} \coloneqq p_h \left( R - \frac{B}{\Delta p} \right) < 1.
\end{equation}
\end{assumption}

Assumption~\ref{ass:EM-finite} is essential for the model to deliver a \emph{finite} equilibrium investment: if $\mathcal{P} \geq 1$ (per unit pledgeable income exceeds the per unit investment cost), the entrepreneur could arrange unlimited external financing and scale the project without bound.
The assumption says that the agency rent per unit $p_h B / \Delta p$ exceeds the project's per-unit NPV $p_h R - 1$, so that the entrepreneur's own contribution $A$ plays a binding role in scaling the investment.

%====================================================================
\section{The Entrepreneur's Program and the Breakeven Condition}
\label{sec:EM-program}
%====================================================================

\subsection{The three constraints}

Given the entrepreneur's wealth $A$, the contract $(R_b, I)$ must satisfy:

\begin{itemize}
  \item \emph{Incentive compatibility} (as in Chapter~\ref{ch:net-worth}, but now scaled by $I$):
  \begin{equation}\label{eq:EM-IC}
    (\Delta p) R_b \geq B I.
    \tag{IC}
  \end{equation}
  The entrepreneur's incremental gain from behaving, $\Delta p \cdot R_b$, must exceed the private benefit $B I$ he would capture by misbehaving.
  \item \emph{Individual rationality} (breakeven for the competitive lender):
  \begin{equation}\label{eq:EM-IR}
    p_h (R I - R_b) \geq I - A.
    \tag{IR}
  \end{equation}
  The lender's expected payment $p_h R_\ell$ must cover his outlay $I - A$.
  \item \emph{Feasibility}: $R_b \geq 0$.
  (We focus on contracts where (IC) is tight or slack at interior values; the bound $R_b \leq R I$ will not bind at the optimum and is suppressed.)
\end{itemize}

The entrepreneur chooses $(R_b, I)$ to maximize his expected net payoff,
\begin{equation}\label{eq:EM-objective}
  \Pi_b(R_b, I) = p_h R_b - A,
\end{equation}
subject to (IC), (IR), and feasibility.
The term $-A$ reflects that the entrepreneur's own wealth $A$ is consumed by the investment and would otherwise be available for direct consumption at $t=0$; we subtract it to express utility \emph{net} of the opportunity cost of investing.

\subsection{The competitive breakeven}

Competition among prospective lenders pins down a sharp prediction on the breakeven constraint (IR).

\begin{proposition}[Competitive breakeven]\label{prop:EM-breakeven}
If $(R_b, I)$ is an equilibrium contract, the competitive lender earns zero expected profit: the constraint \eqref{eq:EM-IR} binds, so
\begin{equation}\label{eq:EM-IR-binding}
  p_h (R I - R_b) = I - A.
\end{equation}
Consequently, the entrepreneur's expected net payoff equals the project's NPV:
\begin{equation}\label{eq:EM-net-payoff-NPV}
  \Pi_b = (p_h R - 1) I.
\end{equation}
\end{proposition}

\begin{proof}
Suppose by contradiction that (IR) is slack at equilibrium, i.e., $p_h(R I - R_b) > I - A$.
Consider an alternative contract $(R_b', I)$ with $R_b' > R_b$ that tightens (IR) toward equality without perturbing $I$.
Such a contract is incentive-compatible (IC) is slackened since it depends on $R_b$ being large, not small, and is feasible.
The entrepreneur's expected net payoff is $p_h R_b' - A > p_h R_b - A$, a strict improvement.
Under competitive pressure, the original contract would not be chosen: a competing lender could offer the alternative contract and attract the entrepreneur.
Hence the equilibrium contract must satisfy (IR) with equality.

Using \eqref{eq:EM-IR-binding},
\begin{equation*}
  \Pi_b = p_h R_b - A = p_h R I - (I - A) - A = (p_h R - 1) I.
\end{equation*}
\end{proof}

Proposition~\ref{prop:EM-breakeven} is the continuous-investment analog of the Pareto-frontier result of Chapter~\ref{ch:net-worth} (Proposition~\ref{prop:HT-financed}): under competitive lending, the entrepreneur captures the entire net social surplus, while the lender earns exactly his opportunity cost.
The entrepreneur's problem therefore reduces to \emph{maximizing the scale of investment} $I$ subject to (IC), (IR), and feasibility.

\subsection{The reformulated program}

Substituting \eqref{eq:EM-net-payoff-NPV} into the entrepreneur's objective, and recognizing that $\Pi_b$ is increasing in $I$, the entrepreneur's program reduces to
\begin{equation*}
  \max_{R_b \geq 0, I \geq 0} \; (p_h R - 1) I
  \quad \text{s.t.} \quad
  (\Delta p) R_b \geq B I, \quad
  p_h (R I - R_b) \geq I - A.
\end{equation*}
The entrepreneur's objective is monotone in $I$: the larger the investment, the larger the captured surplus.
The program is thus a pure scale-maximization problem subject to two linear constraints in $(R_b, I)$.

%====================================================================
\section{The Equity Multiplier}
\label{sec:EM-multiplier}
%====================================================================

\subsection{Combining the two constraints}

At the optimum, both (IC) and (IR) bind.
To see why (IC) binds, observe that, for any feasible contract, reducing $R_b$ toward $B I / \Delta p$ (and using the slack in (IR) to increase $I$) weakly increases the investment, hence the entrepreneur's payoff.
At the optimum, therefore, $R_b = B I / \Delta p$.
Substituting into (IR) at equality:
\begin{equation*}
  p_h R I - p_h \frac{B I}{\Delta p} = I - A,
\end{equation*}
which yields
\begin{equation}\label{eq:EM-I-A-relation}
  I \left[ 1 - p_h \left( R - \frac{B}{\Delta p} \right) \right] = A.
\end{equation}

The coefficient of $I$ on the left-hand side is strictly positive by Assumption~\ref{ass:EM-finite}: it equals $1 - \mathcal{P} > 0$.
Solving for $I$,
\begin{equation}\label{eq:EM-I-kA}
  I = k A, \quad \text{where} \quad k \coloneqq \frac{1}{1 - p_h(R - B/\Delta p)} = \frac{1}{1 - \mathcal{P}}.
\end{equation}

\begin{definition}[Equity multiplier]\label{def:EM}
The \emph{equity multiplier} of the fixed-investment economy is
\begin{equation}\label{eq:EM-k-def}
  k = \frac{1}{1 - \mathcal{P}} > 1.
\end{equation}
The \emph{borrowing capacity per unit of cash} is $d \coloneqq k - 1 = \mathcal{P}/(1 - \mathcal{P})$, so that the entrepreneur with wealth $A$ can borrow at most $d A$, and can invest at most $k A$.
\end{definition}

The multiplier is strictly greater than one because $\mathcal{P} > 0$ (a consequence of $R > B/\Delta p$, which follows from Assumption~\ref{ass:EM-viability}) and the denominator is strictly positive by Assumption~\ref{ass:EM-finite}.
The inequality $k > 1$ is the leverage result: the entrepreneur can mobilize more external capital than he has himself.

\subsection{The optimal contract}

Combining the preceding analysis:

\begin{proposition}[Optimal contract]\label{prop:EM-optimal-contract}
Under Assumptions~\ref{ass:EM-viability} and~\ref{ass:EM-finite}, the optimal contract is uniquely determined by:
\begin{itemize}
  \item \emph{Maximum investment}: $I = k A$, where $k = 1/(1 - \mathcal{P})$.
  \item \emph{Minimum entrepreneurial stake}: $R_b = (B/\Delta p) I = (B/\Delta p) k A$.
\end{itemize}
The lender lends $I - A = (k - 1) A$ and recovers $I - A$ in expectation.
\end{proposition}

\begin{proof}
Substitute \eqref{eq:EM-I-kA} into the (IC) binding condition $R_b = B I / \Delta p$ and verify that (IR) binds by construction.
\end{proof}

The contract is characterized by two binding constraints: the incentive constraint fixes the entrepreneur's stake per unit of investment at $B/\Delta p$, and the breakeven constraint pins down the scale.

\subsection{Interpretation of the multiplier}

The multiplier $k$ admits an intuitive decomposition.
At the optimal contract, the entrepreneur pledges $R_\ell = R I - R_b = (R - B/\Delta p) I$ to the lender in the success state, and the lender's expected repayment is $p_h (R - B/\Delta p) I = \mathcal{P} \cdot I$.
For the lender to break even, this expected repayment must equal his outlay $I - A$:
\begin{equation*}
  \mathcal{P} \cdot I = I - A \Longleftrightarrow A = I (1 - \mathcal{P}).
\end{equation*}
Solving for $I$ yields $I = A/(1 - \mathcal{P}) = k A$.

The economic content is transparent.
Each unit of investment generates an expected pledgeable income of $\mathcal{P}$.
The \emph{gap} between a unit of investment cost (1) and the unit pledgeable income ($\mathcal{P}$), namely $1 - \mathcal{P}$, is the ``agency wedge'' that must be covered by the entrepreneur's own equity $A$.
The multiplier $k = 1/(1 - \mathcal{P})$ is the leverage that the entrepreneur achieves by stretching each unit of his own wealth across $1/(1 - \mathcal{P})$ units of total investment.

\subsection{Comparative statics}

The multiplier's dependence on the fundamentals sheds light on which firms face tighter borrowing constraints.

\begin{proposition}[Comparative statics of $k$]\label{prop:EM-comp-statics}
The equity multiplier $k$ is:
\begin{itemize}
  \item \emph{decreasing} in the private benefit $B$: the larger the private benefit from shirking, the smaller the pledgeable income, the tighter the borrowing capacity;
  \item \emph{increasing} in the likelihood ratio $\Delta p / p_h$: a more ``diagnostic'' outcome---in the sense that behaving shifts the success probability by a larger proportional amount---generates a larger pledgeable income and thus a larger multiplier;
  \item \emph{increasing} in the per-unit project return $R$: more productive projects have larger pledgeable income.
\end{itemize}
\end{proposition}

\begin{proof}
Direct differentiation of $k = [1 - p_h(R - B/\Delta p)]^{-1}$:
\begin{equation*}
  \frac{\partial k}{\partial B} = -\frac{p_h/\Delta p}{[1 - p_h(R - B/\Delta p)]^2} < 0,
\end{equation*}
\begin{equation*}
  \frac{\partial k}{\partial R} = \frac{p_h}{[1 - p_h(R - B/\Delta p)]^2} > 0.
\end{equation*}
For the likelihood ratio, note that $\mathcal{P} = p_h R - B p_h/\Delta p$, and fixing $p_h$ and $R$, an increase in $\Delta p/p_h$ decreases $p_h/\Delta p = 1/(\Delta p/p_h)$, which raises $\mathcal{P}$ and hence $k$.
\end{proof}

The qualitative lessons are economically important.
Firms with high private-benefit parameters (e.g., those whose managers have access to substantial perks or side deals) have lower equity multipliers and therefore tighter borrowing capacity.
Firms whose outcomes are more informative about managerial effort (higher $\Delta p / p_h$) have larger multipliers.
More productive projects ($R$ higher) generate larger pledgeable income and therefore support more borrowing.

%====================================================================
\section{The Shadow Value of Equity}
\label{sec:EM-shadow-value}
%====================================================================

An alternative but equivalent characterization of the financing friction is the \emph{shadow value of equity}: the marginal increase in the entrepreneur's total payoff per unit of his own cash injected into the project.

The entrepreneur's \emph{gross} payoff (net payoff plus the opportunity cost $A$ itself) is
\begin{equation}\label{eq:EM-gross-payoff}
  \Pi_b^g = A + \Pi_b = A + (p_h R - 1) I = A + (p_h R - 1) k A = A \bigl[ 1 + k(p_h R - 1) \bigr].
\end{equation}
After algebraic simplification---substituting $k = 1/(1 - \mathcal{P})$ and collecting terms---
\begin{equation}\label{eq:EM-v}
  \Pi_b^g = v A, \quad v \coloneqq \frac{p_h B/\Delta p}{1 - p_h(R - B/\Delta p)} = \frac{p_h B/\Delta p}{p_h B/\Delta p - (p_h R - 1)}.
\end{equation}

\begin{definition}[Shadow value of equity]\label{def:v}
The \emph{shadow value of equity}, $v$, is the ratio of the entrepreneur's gross expected payoff to his initial cash contribution.
\end{definition}

\begin{proposition}[$v > 1$]\label{prop:EM-v-bigger-1}
The shadow value of equity satisfies $v > 1$.
Moreover, $v$ is:
\begin{itemize}
  \item \emph{increasing} in the per-unit return $R$,
  \item \emph{increasing} in the severity of the moral-hazard problem $\Delta p$,
  \item \emph{decreasing} in the per-unit agency rent $p_h B / \Delta p$.
\end{itemize}
\end{proposition}

\begin{proof}
$v > 1$ follows from the denominator of \eqref{eq:EM-v} being strictly smaller than the numerator by Assumption~\ref{ass:EM-viability} ($p_h R > 1$).
The comparative statics on $v$ follow directly from differentiation of \eqref{eq:EM-v}.
\end{proof}

The shadow value $v > 1$ captures the same economic content as the equity multiplier $k > 1$: each dollar of the entrepreneur's own equity is ``worth more'' to him than a dollar consumed at $t=0$, because it unlocks a multiple of itself in external financing.
The quantity $v - 1 > 0$ is the premium the entrepreneur is willing to pay to have access to an additional dollar of equity at $t=0$.

\subsection{Notational simplification}

To streamline the subsequent analysis, we introduce compact notation for the per-unit payoff quantities:

\begin{itemize}
  \item $\rho_1 \coloneqq p_h R$, the expected payoff per unit of investment,
  \item $\rho_0 \coloneqq p_h(R - B/\Delta p) = \mathcal{P}$, the expected pledgeable income per unit of investment.
\end{itemize}

Under Assumptions~\ref{ass:EM-viability} and~\ref{ass:EM-finite},
\begin{equation}\label{eq:EM-rho-order}
  \rho_0 < 1 < \rho_1,
\end{equation}
with $\rho_1 - \rho_0 = p_h B/\Delta p$ equal to the per-unit agency rent.
The equity multiplier is $k = 1/(1 - \rho_0)$, the borrowing capacity per unit of cash is $d = \rho_0/(1 - \rho_0)$, and the entrepreneur's gross and net payoffs satisfy
\begin{equation}\label{eq:EM-rho-payoffs}
  \Pi_b^g = \frac{\rho_1 - \rho_0}{1 - \rho_0} A, \quad \Pi_b = (\rho_1 - 1) I = (\rho_1 - 1) k A.
\end{equation}

\subsection{Investment--cash-flow sensitivity}

A quantity of empirical interest is the sensitivity of investment to available cash: how much additional investment can the entrepreneur undertake for each additional unit of wealth?
From $I = k A$, $\partial I / \partial A = k$.
The second-order cross-derivative is
\begin{equation}\label{eq:EM-cross-derivative}
  \frac{\partial}{\partial \rho_0} \left( \frac{\partial I}{\partial A} \right) = \frac{\partial k}{\partial \rho_0} = \frac{1}{(1 - \rho_0)^2} > 0.
\end{equation}

\begin{proposition}[Investment--cash sensitivity]\label{prop:EM-cash-sensitivity}
Firms with lower agency costs---i.e., higher $\rho_0$---exhibit higher investment--cash sensitivities.
\end{proposition}

The result is empirically striking.
Firms that are \emph{less} financially constrained have \emph{more} sensitive investment--cash relationships.
This qualifies a naive reading of cash-flow sensitivities as a measure of financial constraints: a high sensitivity may reflect a firm with a high multiplier (low agency cost) rather than a firm with a particularly binding constraint.

%====================================================================
\section{Collateral Value and the Role of Tangible Assets}
\label{sec:EM-collateral}
%====================================================================

The model developed thus far assumed that the project yields zero income in failure states.
Many real projects have \emph{salvage value}: if the venture fails, the firm still owns tangible assets (buildings, machines, inventory) that can be liquidated for some positive amount.
This section extends the model to incorporate such salvage value and examines how it interacts with the agency friction.

\subsection{Adding a failure payoff}

Suppose that, for an investment $I$, the project yields:
\begin{itemize}
  \item $R^S I$ in success, with $R^S > 0$,
  \item $R^F I$ in failure, with $0 < R^F < R^S$.
\end{itemize}
The quantity $R^F I$ can be interpreted as the salvage value of the firm's tangible assets: even when the primary venture fails, some residual value can be recovered through the sale of physical capital.

For later use, write $R \coloneqq R^S - R^F$ for the per-unit \emph{incremental} profit from success over failure.
Then the state-contingent payoff structure can be decomposed as:
\begin{equation}\label{eq:EM-decomposition}
  \text{payoff in success} = R^F I + R I, \quad \text{payoff in failure} = R^F I.
\end{equation}
The term $R^F I$ is a \emph{state-independent} baseline; the term $R I$ is the \emph{incremental} success payoff that depends on the realization.

\subsection{Assumptions on the per-unit parameters}

The viability and finite-investment conditions of Section~\ref{sec:EM-setup} are modified as follows:

\begin{assumption}[Project viability with salvage]\label{ass:EM-collateral-viability}
Per unit of investment:
\begin{align}
  p_h R + R^F &> 1, \tag{V.1''}\label{eq:EM-col-NPV-h} \\
  p_\ell R + R^F + B &< 1. \tag{V.2''}\label{eq:EM-col-NPV-l}
\end{align}
\end{assumption}

\begin{assumption}[Finite-investment condition with salvage]\label{ass:EM-collateral-finite}
\begin{equation}\label{eq:EM-col-finite}
  p_h \left( R - \frac{B}{\Delta p} \right) + R^F < 1.
\end{equation}
\end{assumption}

Assumption~\ref{ass:EM-collateral-viability} is the natural extension of the viability conditions: positive NPV under behaving, negative NPV under misbehaving, both including the salvage value.
As before, these imply $R > B/\Delta p$.
Assumption~\ref{ass:EM-collateral-finite} is the extension of the finite-investment condition: the expected pledgeable income per unit of investment ($p_h(R - B/\Delta p) + R^F$; we will derive this below) must be strictly less than 1.

\subsection{The entrepreneur's program}

A contract now specifies an investment $I$ and a state-contingent sharing rule $(R_b^S, R_b^F)$, where $R_b^S$ and $R_b^F$ are the entrepreneur's rewards in success and failure, respectively.
The lender receives $R^S I - R_b^S$ in success and $R^F I - R_b^F$ in failure.
The feasibility constraints are $0 \leq R_b^S \leq R^S I$ and $0 \leq R_b^F \leq R^F I$.

The entrepreneur's expected net payoff is
\begin{equation}\label{eq:EM-col-objective}
  \Pi_b(R_b^S, R_b^F, I) = p_h R_b^S + (1 - p_h) R_b^F - A.
\end{equation}
The constraints are:

\begin{itemize}
  \item \emph{Incentive compatibility}:
  \begin{equation}\label{eq:EM-col-IC}
    (\Delta p)(R_b^S - R_b^F) \geq B I.
    \tag{IC}
  \end{equation}
  The incremental stake $R_b^S - R_b^F$ must satisfy the same incentive bound as in Section~\ref{sec:EM-program}: only the spread between success and failure rewards matters for incentives.
  \item \emph{Individual rationality}:
  \begin{equation}\label{eq:EM-col-IR}
    p_h (R^S I - R_b^S) + (1 - p_h)(R^F I - R_b^F) \geq I - A.
    \tag{IR}
  \end{equation}
  \item \emph{Feasibility}: $0 \leq R_b^S \leq R^S I$ and $0 \leq R_b^F \leq R^F I$.
\end{itemize}

\subsection{The structure of the optimal contract}

Four structural results characterize the optimal contract.

\begin{proposition}[(IR) binds]\label{prop:EM-col-IR-binds}
At any optimal contract, the individual-rationality constraint \eqref{eq:EM-col-IR} binds.
\end{proposition}

\begin{proof}
Suppose (IR) is slack, i.e., $p_h(R^S I - R_b^S) + (1 - p_h)(R^F I - R_b^F) > I - A$.
We construct a perturbation that strictly improves the entrepreneur's payoff.
If $R_b^S < R^S I$, raise $R_b^S$ slightly: (IC) tightens toward the constraint, (IR) tightens toward binding, and the objective $p_h R_b^S + (1-p_h) R_b^F$ strictly increases.
If $R_b^S = R^S I$ and $R_b^F < R^F I$, raise $R_b^F$ slightly: (IC) slackens, (IR) tightens, and the objective strictly increases.
If $R_b^S = R^S I$ and $R_b^F = R^F I$, the entrepreneur captures the entire income, so the lender receives nothing; this cannot satisfy (IR) unless $I \leq A$ (self-financing), in which case borrowing is superfluous.
\end{proof}

Using the binding (IR), the entrepreneur's objective reduces to the project's NPV:
\begin{equation}\label{eq:EM-col-NPV}
  \Pi_b = p_h R^S + (1 - p_h) R^F - 1 - A + I = (p_h R + R^F - 1) I.
\end{equation}
The entrepreneur seeks to maximize the investment $I$ subject to (IC), (IR), and feasibility.

\begin{proposition}[(IC) binds at the optimum]\label{prop:EM-col-IC-binds}
At any optimal contract, the incentive-compatibility constraint \eqref{eq:EM-col-IC} binds:
\begin{equation}\label{eq:EM-col-IC-binding}
  (\Delta p)(R_b^S - R_b^F) = B I.
\end{equation}
\end{proposition}

\begin{proof}
If (IC) is slack at the optimum, we can reduce $R_b^S$ slightly, which slackens (IR) and allows a strictly larger investment $I$ (since (IC) remains satisfied at the lower $R_b^S$).
The resulting contract delivers strictly higher NPV, contradicting optimality.
Hence (IC) binds.
\end{proof}

\begin{proposition}[No reward in failure]\label{prop:EM-col-RbF-zero}
At the optimal contract, the entrepreneur receives no reward in the failure state: $R_b^F = 0$.
\end{proposition}

\begin{proof}
Suppose $R_b^F > 0$.
Consider the perturbation $(R_b^S, R_b^F) \to (R_b^S - \Delta R, R_b^F - \Delta R / (1 - p_h) \cdot p_h / (1 - p_h))$---actually, a cleaner construction is to reduce $R_b^F$ by $\epsilon > 0$ and, to keep (IC) binding per \eqref{eq:EM-col-IC-binding}, reduce $R_b^S$ by the same $\epsilon$.
The difference $R_b^S - R_b^F$ is unchanged (so (IC) is unaffected), and (IR) slackens by $p_h \epsilon + (1 - p_h) \epsilon = \epsilon$.
The slackened (IR) permits a larger investment $I$.
Iterating until $R_b^F = 0$ (or $R_b^S = B I / \Delta p$, the (IC)-binding level) yields the optimum.
The \emph{maximal-incentives principle}: in the optimal contract, all of the entrepreneur's reward is concentrated in success.
\end{proof}

Proposition~\ref{prop:EM-col-RbF-zero} is the \emph{maximal-incentives principle}: punishing failure as harshly as possible (setting $R_b^F = 0$) maximizes the entrepreneur's stake in success per unit of investment, which maximizes the pledgeable income, which maximizes the investment scale.
The entrepreneur prefers ``high-powered'' contracts with zero reward in failure, despite the risk exposure this creates for him, because such contracts expand his borrowing capacity.

\subsection{The borrowing capacity with collateral}

Substituting $R_b^F = 0$ into the (IC)-binding condition gives $R_b^S = B I / \Delta p$.
Substituting $R_b^F = 0$ and $R_b^S = B I / \Delta p$ into the (IR)-binding condition:
\begin{equation*}
  p_h \left( R^S I - \frac{B I}{\Delta p} \right) + (1 - p_h)(R^F I) = I - A,
\end{equation*}
which simplifies (using $R^S = R + R^F$) to
\begin{equation}\label{eq:EM-col-IR-simplified}
  \bigl[ p_h(R - B/\Delta p) + R^F \bigr] I = I - A,
\end{equation}
yielding
\begin{equation}\label{eq:EM-col-borrowing-capacity}
  I = \frac{A}{1 - \bigl[ p_h(R - B/\Delta p) + R^F \bigr]}.
\end{equation}

\begin{proposition}[Optimal contract with salvage]\label{prop:EM-col-contract}
Under Assumptions~\ref{ass:EM-collateral-viability} and~\ref{ass:EM-collateral-finite}, the optimal contract is:
\begin{equation*}
  I = \frac{A}{1 - \bigl[ p_h(R - B/\Delta p) + R^F \bigr]},
  \quad
  R_b^S = \frac{B I}{\Delta p},
  \quad
  R_b^F = 0.
\end{equation*}
\end{proposition}

The denominator of \eqref{eq:EM-col-borrowing-capacity} is the agency wedge, now inclusive of the salvage value.
The extended pledgeable income per unit of investment is $p_h(R - B/\Delta p) + R^F$: the usual term $p_h(R - B/\Delta p)$ from the basic model, \emph{plus} the state-independent payoff $R^F$.
Because $R^F$ accrues regardless of the entrepreneur's action, it requires no incentive cost: the entire $R^F$ per unit can be pledged to the lender.

\subsection{Decomposing the lender's claim into debt and equity}

A useful decomposition of the lender's contractual claim in the success state:
\begin{equation*}
  R^S I - R_b^S = R^S I - \frac{B I}{\Delta p}
  = \underbrace{R^F I}_{\text{safe debt}} + \underbrace{\left(R - \frac{B}{\Delta p}\right) I}_{\text{risky equity}}.
\end{equation*}
The lender's claim naturally partitions into:
\begin{itemize}
  \item a \emph{safe debt} component paying $R^F I$ in both success and failure (the salvage value);
  \item a \emph{risky equity} component paying $(R - B/\Delta p) I$ in success and zero in failure.
\end{itemize}

The debt--equity leverage ratios of the firm are:
\begin{equation}\label{eq:EM-leverage}
  \frac{\text{debt}}{\text{total equity}} = \frac{R^F I}{p_h R I} = \frac{R^F}{p_h R},
  \quad
  \frac{\text{debt}}{\text{outside equity}} = \frac{R^F I}{p_h(R - B/\Delta p) I} = \frac{R^F}{p_h(R - B/\Delta p)}.
\end{equation}
Firms with a larger salvage value $R^F$ relative to their incremental success payoff $p_h R$ carry a higher debt ratio: the tangible asset base supports a larger collateralized-debt layer in the capital structure.

%====================================================================
\section{Credit Rationing and the Value of Tangible Assets}
\label{sec:EM-tangible}
%====================================================================

An empirically important corollary of the collateral model is the following:

\begin{quote}\itshape
Credit rationing is more binding for firms with less tangible assets, or whose assets have lower liquidation value.
\end{quote}

To see this formally, consider two firms that are identical in every respect except for the tangibility of their assets.
Firm~1 has salvage parameter $R^F$ and success probability $p_h$ under behaving; firm~2 has a lower salvage parameter $\widehat{R}^F < R^F$ but a higher success probability $\widehat{p}_h = p_h + \tau$ for some $\tau > 0$.
The two firms are designed to have the same project NPV: by construction,
\begin{equation}\label{eq:EM-NPV-invariant}
  \widehat{R}^F + \widehat{p}_h R = R^F + p_h R,
  \quad \text{i.e.,} \quad \widehat{R}^F + (p_h + \tau) R = R^F + p_h R.
\end{equation}
To preserve the agency problem, we adjust both $p_h$ and $p_\ell$ by $\tau$:
\begin{equation}\label{eq:EM-pl-adjust}
  \widehat{p}_\ell = p_\ell + \tau,
  \quad \widehat{\Delta p} = \widehat{p}_h - \widehat{p}_\ell = \Delta p.
\end{equation}
Condition \eqref{eq:EM-pl-adjust} keeps the likelihood ratio (and hence the (IC) constraint) unchanged.

\begin{proposition}[Tangible assets expand borrowing capacity]\label{prop:EM-tangible}
Under the comparative-statics experiment \eqref{eq:EM-NPV-invariant}--\eqref{eq:EM-pl-adjust}, the investment scale of firm~2 satisfies
\begin{equation}\label{eq:EM-tangible-inequality}
  \widehat{I} = \frac{A}{1 - \bigl[ (p_h + \tau)(R - B/\Delta p) + \widehat{R}^F \bigr]} < \frac{A}{1 - \bigl[ p_h(R - B/\Delta p) + R^F \bigr]} = I.
\end{equation}
Equivalently, holding project NPV and agency frictions constant, a firm with less tangible assets invests less.
\end{proposition}

\begin{proof}
The firm with less tangible assets has a smaller pledgeable income per unit of investment:
\begin{equation*}
  (p_h + \tau)(R - B/\Delta p) + \widehat{R}^F = p_h(R - B/\Delta p) + R^F + \tau(R - B/\Delta p) + (\widehat{R}^F - R^F).
\end{equation*}
By \eqref{eq:EM-NPV-invariant}, $\widehat{R}^F - R^F = -\tau R$, so the expression simplifies to
\begin{equation*}
  p_h(R - B/\Delta p) + R^F + \tau(R - B/\Delta p) - \tau R = p_h(R - B/\Delta p) + R^F - \tau \cdot \frac{B}{\Delta p}.
\end{equation*}
Since $\tau > 0$ and $B/\Delta p > 0$, the corrected pledgeable income is strictly smaller, and the borrowing-capacity denominator $1 - [\cdot]$ is strictly larger.
Therefore $\widehat{I} < I$.
\end{proof}

The intuition is clean.
Both firms have the same project NPV and the same agency friction.
The difference lies in \emph{how} the return is generated: firm~1 has a larger slice of its expected return coming from the state-independent salvage value $R^F$, which requires no incentive cost to pledge; firm~2 has a larger slice coming from the incentive-sensitive success payoff $R$, which requires leaving more of it to the entrepreneur as agency rent.
For the same amount of equity $A$, firm~1 can borrow more because it can pledge more in expectation.

The proposition articulates a central theme of modern corporate finance: \emph{tangible assets serve as collateral, expanding borrowing capacity}.
Knowledge-intensive firms, start-ups, and service firms---whose value rests primarily on intangible assets and human capital---face tighter credit constraints than asset-rich manufacturing or real estate firms with the same fundamental productivity.
The result has been widely invoked to explain cross-industry and cross-country patterns in capital structure, and to motivate the role of collateral-intensive financial intermediation in funding ventures that primarily hold intangibles.

%====================================================================
\section{Further Reading}
\label{sec:EM-further-reading}
%====================================================================

The continuous-investment formulation of the moral-hazard model is introduced by \citet{HolmstromTirole1997}, which also develops the interaction between entrepreneurial net worth, intermediary capital, and the real sector.
The framework has since become foundational for macroeconomic models of financial accelerators, balance-sheet channels of monetary policy, and emerging-market sudden-stop crises.
\citet{Tirole2006} (Chapter~3) is the authoritative textbook treatment and extends the model in several directions beyond those covered in this chapter, including multi-period investment, renegotiation, and the role of intermediaries.

The \emph{maximal-incentives principle} (Proposition~\ref{prop:EM-col-RbF-zero})---that the optimal incentive scheme punishes failure as harshly as limited liability permits---is a recurring theme in contract theory.
For the classical references on contract design with risk-neutral, wealth-constrained agents, see \citet{Tirole2006} and the principal--agent literature cited therein.

The link between tangible assets and borrowing capacity developed in Section~\ref{sec:EM-tangible} is an instance of the broader \emph{collateral channel} in corporate finance.
For survey treatments of collateral in credit markets, see \citet{FreixasRochet1997}.
For the macroeconomic implications of the collateral channel, the reader is referred to the large literature following \citet{HolmstromTirole1997}.

The investment--cash-flow sensitivity result of Proposition~\ref{prop:EM-cash-sensitivity}---that less financially constrained firms exhibit higher cash-flow sensitivities---has generated a large empirical literature testing the relationship between cash flow and investment across firms of different characteristics.
The theoretical prediction here is sharp but depends on the specific functional form of the continuous-investment model; more general formulations generate richer (and sometimes reversed) predictions.

%====================================================================
\chapter{Nonverifiable Income and the Threat of Termination}
\label{ch:nonverifiable}
%====================================================================

The analyses of Chapters~\ref{ch:debt-rationing}--\ref{ch:equity-multiplier} all rested on one common assumption: at the date when the project's outcome is realized, that outcome is \emph{verifiable} to the lender.
Either the outcome is directly observable (as in the moral-hazard models of Chapters~\ref{ch:net-worth}--\ref{ch:equity-multiplier}) or it can be verified via costly audit (as in the Gale--Hellwig model of Chapter~\ref{ch:debt-rationing}).
In both cases, the contract can condition payments on the realized state, and the analytical tools---incentive compatibility, individual rationality, pledgeable income---apply naturally.

This chapter explores what happens when verifiability breaks down completely.
Suppose the borrower's income cannot be observed at all: the entrepreneur can claim, entirely and with complete impunity, that the project produced nothing.
No audit is available; no legal recourse exists.
What, then, motivates the entrepreneur to repay the lender?
The answer, developed by \citet{BoltonScharfstein1990} and others, is that repayment must be motivated by the \emph{threat of termination}: if the entrepreneur fails to repay, he loses access to future cash flows from the project.
The fear of losing future continuation value is what compels repayment today.

This mechanism has wide-ranging applications.
It underlies the theory of sovereign debt, in which the borrower is a country that cannot be compelled by law to repay (there is no international bankruptcy court); repayment in this setting must be motivated by the loss of future market access or the loss of international trade.
It underlies models of relationship banking, where continuation of the lending relationship depends on the borrower's past repayment history.
And it is the formal foundation for the widespread practice of writing short-maturity, sequentially rolled-over debt contracts: by threatening to refuse rollover, the lender maintains control over the borrower even when each individual period's cash flow is unverifiable.

The chapter develops the simplest version of this framework: a three-date model in which the entrepreneur receives non-verifiable income at an intermediate date and the lender controls the continuation of the project.
The analysis characterizes the optimal contract as a single termination probability $y_0$ calibrated to make repayment incentive-compatible given the entrepreneur's continuation value $R_2$ and the project's liquidation value $L$.

The chapter follows Part~D of Chapter~2 of \citet{Tirole2006}, which builds on \citet{BoltonScharfstein1990}.

%====================================================================
\section{The Three-Date Setting}
\label{sec:NV-setting}
%====================================================================

\subsection{Agents and timing}

We consider a three-date economy, $\tau \in \{0, 1, 2\}$, with a single consumption good and risk-neutral agents with zero discount rate.
The entrepreneur (borrower) arrives at $\tau = 0$ with wealth $A$ and an investment project requiring a capital outlay $I > A$.
Competitive lenders stand ready to provide the balance $I - A$ against the promise of repayment at a later date.
Unlike the models of Chapters~\ref{ch:debt-rationing}--\ref{ch:equity-multiplier}, there are now two future periods in which cash flows can be realized and the project can be continued or liquidated.

\subsection{The project's cash flows}

The project generates income at two dates:
\begin{itemize}
  \item \emph{At $\tau = 1$}, the project yields income $R_1 > 0$ with probability $p \in (0, 1)$ and zero with probability $1 - p$.
  \item \emph{At $\tau = 2$}, conditional on the project still operating (i.e., not liquidated at $\tau = 1$), the project yields an \emph{expected} income of $R_2 > 0$ that accrues entirely to the entrepreneur.
\end{itemize}
The key structural assumption---the nonverifiability of income---is that the $\tau = 1$ income $R_1$ is observable only to the entrepreneur.
The lender cannot verify whether the project produced $R_1$ or zero at the intermediate date, and cannot, by law or by audit, compel the entrepreneur to reveal or deliver any particular amount.
The only action the lender \emph{can} take, contingent on the entrepreneur's payment, is to terminate the project before $\tau = 2$.

The $\tau = 2$ income $R_2$ is of a different nature.
It is not contractible as a repayment: by the terminal-date logic, the entrepreneur has no incentive to hand over any of $R_2$ to the lender, because no future threat can be brought to bear.
We therefore treat $R_2$ as a \emph{deterministic private benefit of continuation} for the entrepreneur, entirely captured by him if and only if the project is not liquidated.

\subsection{The liquidation decision}

At $\tau = 1$, after the entrepreneur has (or has not) made a payment $D$ to the lender, the project can either be continued through to $\tau = 2$ or liquidated on the spot.
Liquidation yields the lender a recovery value $L \geq 0$.
We impose the condition
\begin{equation}\label{eq:NV-L-bound}
  0 \leq L \leq I - A.
\end{equation}
The upper bound $L \leq I - A$ is the economically interesting case: the liquidation value alone does not compensate the lender for his outlay.
Absent some payment from the entrepreneur, the lender cannot break even on the project; a contract will only be signed if the entrepreneur has some positive expected payment to make.

The liquidation value $L$ has two natural interpretations.
First, and most directly, it is the salvage value of the firm's physical assets if the project is terminated at $\tau = 1$.
Second, it may represent the savings associated with \emph{not incurring a second-period investment cost}: if the project requires additional capital at $\tau = 1$ to continue, then abandoning the project frees up those funds.
Both interpretations lead to the same formal treatment.

\subsection{The continuation decision as a stochastic rule}

The contract does not compel the lender to continue or liquidate deterministically.
Instead, it specifies a \emph{probability of continuation}, $y_0$ or $y_1$, conditional on the entrepreneur's payment behavior:
\begin{itemize}
  \item if the entrepreneur pays $D$ at $\tau = 1$, the project is continued with probability $y_1 \in [0, 1]$;
  \item if the entrepreneur does not pay, the project is continued with probability $y_0 \in [0, 1]$.
\end{itemize}
The stochastic continuation rule is the instrument by which the lender punishes non-payment: by setting $y_0 < y_1$, the lender reduces the probability that the entrepreneur will collect $R_2$ at $\tau = 2$, and thereby creates an incentive to pay at $\tau = 1$.
If instead the contract specified $y_0 = y_1$, the continuation decision would be independent of the entrepreneur's behavior and no incentive to repay would exist.

The randomization over continuation can be implemented in various ways: as an actual probability (say, via a randomizing device committed to ex ante), or as a partial liquidation (liquidating a fraction of the assets).
For the analysis below, we take $(y_0, y_1) \in [0, 1]^2$ as primitive.

\subsection{Relative magnitudes}

We impose the ordering
\begin{equation}\label{eq:NV-ordering}
  R_1 > R_2 > L.
\end{equation}
The first-period income $R_1$ is the largest quantity: it is the flow that can potentially be used for repayment.
The continuation value $R_2$ must exceed the liquidation value $L$, otherwise continuation would be socially undesirable and the lender would never promise to continue.

%====================================================================
\section{The Contract and Its Two Constraints}
\label{sec:NV-contract}
%====================================================================

\subsection{The contract}

A contract is a triple
\begin{equation*}
  (y_0, y_1, D) \in [0,1] \times [0,1] \times [0, R_1],
\end{equation*}
specifying the two continuation probabilities and the required payment.
The upper bound $D \leq R_1$ reflects feasibility: the entrepreneur cannot be forced to pay more than the income he could have received.

Under the nonverifiable-income structure, the entrepreneur pays $D$ at $\tau = 1$ only if he has the income to do so \emph{and} has an incentive to hand it over.
In the success state ($\tau = 1$ income equals $R_1$), the entrepreneur has the cash; in the failure state (income equals zero), he cannot pay $D > 0$, and the lender must liquidate with probability $1 - y_0$ or continue with probability $y_0$, irrespective of the entrepreneur's wishes.

\subsection{The incentive constraint}

In the success state, the entrepreneur faces a choice: pay $D$ or default.
If he pays, his total expected payoff is
\begin{equation*}
  R_1 - D + y_1 R_2.
\end{equation*}
If he defaults---reporting failure, consuming the $R_1$ himself, and letting the project be liquidated with probability $1 - y_0$---his total expected payoff is
\begin{equation*}
  R_1 + y_0 R_2.
\end{equation*}
(We treat the non-verifiability as absolute: the entrepreneur can consume $R_1$ with complete impunity if he chooses not to pay, because the lender cannot verify the income.)

The incentive constraint for repayment is that paying weakly dominates defaulting:
\begin{equation}\label{eq:NV-IC}
  R_1 - D + y_1 R_2 \geq R_1 + y_0 R_2 \Longleftrightarrow (y_1 - y_0) R_2 \geq D.
  \tag{IC}
\end{equation}

The content of \eqref{eq:NV-IC} is economically clear: the expected loss in continuation value from defaulting, $(y_1 - y_0) R_2$, must exceed the monetary gain from keeping $D$.
The ``expected loss'' depends on the \emph{differential} continuation probability $y_1 - y_0$: a larger differential creates a larger punishment for non-repayment and allows the lender to extract a larger repayment $D$.

\subsection{The individual-rationality constraint}

The lender's expected payoff under the contract is
\begin{itemize}
  \item in the success state (probability $p$): the lender receives $D$ from the entrepreneur, and if the project is subsequently liquidated (probability $1 - y_1$), an additional $L$;
  \item in the failure state (probability $1 - p$): the lender receives nothing from the entrepreneur (who has nothing to pay) and $L$ with probability $1 - y_0$.
\end{itemize}
Total expected revenue to the lender:
\begin{equation}\label{eq:NV-lender-revenue}
  p \bigl[ D + (1 - y_1) L \bigr] + (1 - p)(1 - y_0) L.
\end{equation}

The individual-rationality (or breakeven) constraint requires that this revenue cover the lender's outlay $I - A$:
\begin{equation}\label{eq:NV-IR}
  p \bigl[ D + (1 - y_1) L \bigr] + (1 - p)(1 - y_0) L \geq I - A.
  \tag{IR}
\end{equation}

\subsection{The optimization}

The entrepreneur's expected utility is
\begin{equation}\label{eq:NV-entrepreneur-utility}
  p \bigl[ R_1 - D + y_1 R_2 \bigr] + (1 - p) y_0 R_2.
\end{equation}
In the success state, the entrepreneur consumes $R_1 - D$ and enjoys $R_2$ with probability $y_1$.
In the failure state, he consumes nothing at $\tau = 1$ and enjoys $R_2$ with probability $y_0$.

The optimal contract maximizes \eqref{eq:NV-entrepreneur-utility} subject to (IC), (IR), and the feasibility constraints on $(y_0, y_1, D)$.

%====================================================================
\section{Characterization of the Optimal Contract}
\label{sec:NV-optimal}
%====================================================================

The optimal contract is pinned down by five structural observations.
Taken together, they reduce a three-dimensional optimization problem to a single closed-form expression for the termination probability.

\begin{proposition}[Necessary conditions for optimality]\label{prop:NV-optimal}
At any optimal contract $(y_0, y_1, D)$:
\begin{enumerate}
  \item[\textup{(i)}] the payment is strictly positive, $D > 0$;
  \item[\textup{(ii)}] the continuation probabilities satisfy $y_1 > y_0$;
  \item[\textup{(iii)}] the individual-rationality constraint \textup{(IR)} binds with equality;
  \item[\textup{(iv)}] the project is continued with certainty upon payment, $y_1 = 1$;
  \item[\textup{(v)}] the incentive constraint \textup{(IC)} binds with equality, $(1 - y_0) R_2 = D$.
\end{enumerate}
\end{proposition}

We develop the proofs of these five claims in the order they are stated, each building on the previous.

\begin{proof}[Proof of (i): $D > 0$]
Suppose for contradiction that $D = 0$.
Then the lender's expected revenue from \eqref{eq:NV-lender-revenue} is
\begin{equation*}
  p(1 - y_1) L + (1 - p)(1 - y_0) L \leq L.
\end{equation*}
By assumption \eqref{eq:NV-L-bound}, $L \leq I - A$, so (IR) becomes
\begin{equation*}
  L \geq I - A,
\end{equation*}
which forces $L = I - A$ and both $(1 - y_1) = 1$ and $(1 - y_0) = 1$, i.e., $y_0 = y_1 = 0$.
This is the ``always-liquidate'' contract, which yields the entrepreneur zero utility---strictly less than the status quo of no-project.
No agent would sign such a contract, contradicting the assumption that $(y_0, y_1, 0)$ is optimal.
Hence $D > 0$.
\end{proof}

\begin{proof}[Proof of (ii): $y_1 > y_0$]
By (i), $D > 0$.
The incentive constraint (IC) requires $(y_1 - y_0) R_2 \geq D > 0$, which forces $y_1 > y_0$.
\end{proof}

\begin{proof}[Proof of (iii): (IR) binds]
Suppose (IR) is slack, i.e., $p[D + (1 - y_1) L] + (1 - p)(1 - y_0) L > I - A$.
Consider the perturbation $D' = D - \epsilon$ for small $\epsilon > 0$, keeping $y_0, y_1$ unchanged.
\begin{itemize}
  \item The lender's expected revenue decreases by $p \epsilon$, but (IR) remains satisfied for small enough $\epsilon$ by the assumed slackness.
  \item The incentive constraint (IC) requires $D' \leq (y_1 - y_0) R_2$, which holds for the perturbed $D'$ since the original $D$ already satisfied this.
  \item The entrepreneur's utility from \eqref{eq:NV-entrepreneur-utility} \emph{increases} by $p \epsilon$ (more money retained in the success state).
\end{itemize}
The perturbation strictly improves the entrepreneur's objective while remaining feasible, contradicting optimality.
Hence (IR) binds.
\end{proof}

\begin{proof}[Proof of (iv): $y_1 = 1$]
Suppose $y_1 < 1$ at the optimum.
Consider the perturbation $(y_0, y_1, D) \to (y_0, y_1 + \epsilon, D + \epsilon L)$ for small $\epsilon > 0$.
We verify that the perturbation is feasible and improves the entrepreneur's payoff.

\emph{Incentive constraint}: the modified constraint is $(y_1 + \epsilon - y_0) R_2 \geq D + \epsilon L$.
Since $R_2 > L$ by \eqref{eq:NV-ordering}, the left-hand side increases by $\epsilon R_2$ and the right-hand side increases by $\epsilon L$, with $\epsilon R_2 > \epsilon L$.
The (IC) slackens.

\emph{Individual-rationality}: the lender's revenue becomes $p[(D + \epsilon L) + (1 - y_1 - \epsilon) L] + (1 - p)(1 - y_0) L$.
Expanding,
\begin{equation*}
  p D + p \epsilon L + p L - p y_1 L - p \epsilon L + (1 - p)(1 - y_0) L = p D + p (1 - y_1) L + (1 - p)(1 - y_0) L,
\end{equation*}
which is \emph{identical} to the original lender's revenue.
(IR) holds with equality unchanged.

\emph{Entrepreneur's utility}: from \eqref{eq:NV-entrepreneur-utility}, the change is
\begin{equation*}
  p(-\epsilon L + \epsilon R_2) = p \epsilon (R_2 - L) > 0,
\end{equation*}
by \eqref{eq:NV-ordering}.
The entrepreneur's utility strictly improves.
Iterating until $y_1 = 1$ gives the optimal contract.
\end{proof}

\begin{proof}[Proof of (v): (IC) binds]
Suppose (IC) is slack, i.e., $(y_1 - y_0) R_2 > D$.
Consider the perturbation $(y_0, y_1, D) \to (y_0 + \epsilon, y_1, D + \epsilon R_2 (1-p)/p \cdot \text{adjustment})$---more precisely, increase $y_0$ by $\epsilon > 0$ and increase $D$ by $\delta \epsilon$ for some $\delta$ to be determined.

\emph{Incentive constraint}: $(y_1 - y_0 - \epsilon) R_2 \geq D + \delta \epsilon$ becomes $-\epsilon R_2 \geq \delta \epsilon$, requiring $\delta \leq -R_2$.
We choose $\delta = -R_2 + \eta$ for small $\eta$ to preserve (IC) with some slack.
Wait---this perturbation reduces $D$.
Let me reconsider.

A cleaner perturbation: increase $y_0$ by $\epsilon$ and raise $D$ by some compensating amount $\delta \epsilon$ chosen to keep the lender's revenue constant.
The lender's revenue \eqref{eq:NV-lender-revenue} under the perturbation is
\begin{equation*}
  p[(D + \delta \epsilon) + (1 - y_1) L] + (1 - p)(1 - y_0 - \epsilon) L
  = [\text{original}] + p \delta \epsilon - (1 - p) \epsilon L.
\end{equation*}
Setting the change to zero: $p \delta = (1 - p) L$, i.e., $\delta = (1 - p) L / p$.
With this choice, the lender's revenue is unchanged, so (IR) continues to bind.

\emph{Incentive constraint}: $(y_1 - y_0 - \epsilon) R_2 \geq D + \delta \epsilon$, i.e., $(y_1 - y_0) R_2 - \epsilon R_2 \geq D + \epsilon (1-p) L / p$.
By the assumed slackness of (IC), $(y_1 - y_0) R_2 > D$; for small $\epsilon$, the perturbation preserves (IC) strictly.

\emph{Entrepreneur's utility}: the change in the objective \eqref{eq:NV-entrepreneur-utility} is
\begin{equation*}
  p (-\delta \epsilon) + (1 - p) \epsilon R_2 = -p \cdot \tfrac{(1-p) L}{p} \cdot \epsilon + (1 - p) \epsilon R_2 = (1 - p) \epsilon (R_2 - L) > 0,
\end{equation*}
by \eqref{eq:NV-ordering}.
The entrepreneur's utility strictly increases, contradicting optimality.
Hence (IC) binds at the optimum.
\end{proof}

Combining (i)--(v), the optimal contract is fully determined.
With $y_1 = 1$ and (IC) binding, $(1 - y_0) R_2 = D$.
With (IR) binding and $y_1 = 1$:
\begin{equation*}
  p D + (1 - p)(1 - y_0) L = I - A.
\end{equation*}
Substituting $D = (1 - y_0) R_2$:
\begin{equation*}
  p (1 - y_0) R_2 + (1 - p)(1 - y_0) L = (1 - y_0) [p R_2 + (1 - p) L] = I - A,
\end{equation*}
which yields
\begin{equation}\label{eq:NV-termination-prob}
  1 - y_0 = \frac{I - A}{p R_2 + (1 - p) L}.
\end{equation}

\begin{proposition}[The optimal contract]\label{prop:NV-optimal-form}
Under the hypotheses $0 \leq L \leq I - A$ and $R_1 > R_2 > L$, the optimal contract is
\begin{equation*}
  y_1 = 1,
  \quad
  1 - y_0 = \frac{I - A}{p R_2 + (1 - p) L},
  \quad
  D = (1 - y_0) R_2.
\end{equation*}
\end{proposition}

The optimal contract is remarkably transparent.
The project is continued with certainty when the entrepreneur pays; with probability $1 - y_0$ it is terminated upon non-payment, and this probability is precisely calibrated so that the lender breaks even.
The required payment $D$ is the \emph{expected continuation value lost} by the entrepreneur upon non-payment: $D = (1 - y_0) R_2$.
The three parties (entrepreneur, lender, project) are each in a clean state: the entrepreneur's incentive constraint binds, the lender's breakeven binds, and the continuation-if-pay rule is at its maximum.

\begin{remark}[Interpretation via termination as a ``stick'']\label{rem:termination-stick}
The optimal contract illustrates the economics of termination as an enforcement device.
The lender cannot compel repayment directly (income is unverifiable), but he can threaten to terminate the project---destroying $R_2$ for the entrepreneur and recovering only $L$ for himself.
The entrepreneur, faced with this threat, rationally chooses to pay $D$ to preserve his continuation value.
The termination probability $1 - y_0$ is calibrated to make this incentive just strong enough to generate the required lender revenue.
\end{remark}

%====================================================================
\section{Comparative Statics}
\label{sec:NV-statics}
%====================================================================

Proposition~\ref{prop:NV-optimal-form} pins down the termination probability $1 - y_0$ in closed form as a function of the model's fundamentals.
The expression is rich enough to deliver sharp comparative-statics predictions about when termination is more or less likely.

\begin{proposition}[Comparative statics of termination]\label{prop:NV-comp-statics}
The probability of termination in the event of non-repayment, $1 - y_0 = (I - A)/[p R_2 + (1 - p) L]$, is:
\begin{itemize}
  \item \emph{decreasing} in the entrepreneur's continuation value $R_2$;
  \item \emph{decreasing} in the liquidation value $L$;
  \item \emph{decreasing} in the first-period success probability $p$ (provided $R_2 > L$);
  \item \emph{decreasing} in the entrepreneur's net worth $A$.
\end{itemize}
\end{proposition}

\begin{proof}
Direct differentiation of \eqref{eq:NV-termination-prob}.
For $R_2$: the denominator $p R_2 + (1 - p) L$ is increasing in $R_2$, so $1 - y_0$ is decreasing in $R_2$.
For $L$: similarly, the denominator is increasing in $L$.
For $p$: $\partial/\partial p [p R_2 + (1 - p) L] = R_2 - L > 0$, so the denominator is increasing in $p$ and $1 - y_0$ is decreasing.
For $A$: the numerator $I - A$ is decreasing in $A$, so $1 - y_0$ is decreasing in $A$.
\end{proof}

The four comparative statics of Proposition~\ref{prop:NV-comp-statics} admit intuitive economic interpretations:

\begin{itemize}
  \item An increase in the continuation value $R_2$ makes the threat of termination more powerful per unit of termination probability, allowing the lender to extract a given $D$ with a lower probability of termination.
  Equivalently, the entrepreneur has more to lose from being terminated, so the lender can sustain the contract with less severe punishment.
  \item An increase in the liquidation value $L$ means the lender gets more cash back when he does liquidate, so he needs to liquidate less often to recoup his investment.
  \item An increase in the first-period success probability $p$ means the entrepreneur's ``paying state'' occurs more often, providing more cash flow $D$ in expectation.
  The lender can achieve breakeven with less reliance on liquidation in the failure state.
  (The condition $R_2 > L$ ensures that this intuition operates in the expected direction: the paying state is more valuable to both parties than the failure state.)
  \item An increase in the entrepreneur's net worth $A$ reduces the amount $I - A$ that the lender must recover.
  Less required recovery means less reliance on the termination threat.
\end{itemize}

The last comparative statics echoes a theme that has run through Part~II of these notes: \emph{net worth relaxes financial frictions}.
Just as a richer entrepreneur faces a larger borrowing capacity in Chapter~\ref{ch:equity-multiplier}, he faces a lower probability of being terminated in the nonverifiable-income model.
The mechanism is different---the multiplier in the earlier chapter, the termination probability here---but the direction is the same.

%====================================================================
\section{Feasibility and Parameter Restrictions}
\label{sec:NV-feasibility}
%====================================================================

The contract derived in Proposition~\ref{prop:NV-optimal-form} is valid only when $1 - y_0 \in [0, 1]$, i.e., when the right-hand side of \eqref{eq:NV-termination-prob} is a valid probability.
The constraint $1 - y_0 \geq 0$ holds automatically whenever $I > A$ (i.e., external financing is needed).
The constraint $1 - y_0 \leq 1$ requires
\begin{equation}\label{eq:NV-feasibility}
  I - A \leq p R_2 + (1 - p) L.
\end{equation}

The inequality \eqref{eq:NV-feasibility} is the analog of the ``financing condition'' $\overline{A} \leq A$ from Chapter~\ref{ch:net-worth}: it identifies when a contract exists that satisfies all the constraints simultaneously.
Economically, \eqref{eq:NV-feasibility} says that the lender's maximal expected revenue from threatening full termination ($y_0 = 0$) must cover the outlay:
\begin{equation*}
  \text{max lender revenue}
  = p \cdot R_2 + (1 - p) L
  \geq I - A.
\end{equation*}
(In the success state, the maximum the lender can extract is $R_2$, because that is the full continuation value the entrepreneur would lose by defaulting; in the failure state, the lender gets $L$ from liquidation.)
If \eqref{eq:NV-feasibility} fails, no contract can be written that makes the lender break even: the project is not financeable, and a form of credit rationing emerges even though the project's NPV under perfect information would be strictly positive.

\begin{proposition}[Feasibility of the contract]\label{prop:NV-feasibility}
A contract satisfying \textup{(IC)} and \textup{(IR)} exists if and only if
\begin{equation*}
  I - A \leq p R_2 + (1 - p) L.
\end{equation*}
\end{proposition}

The proposition ties the nonverifiable-income framework back to the broader theme of Part~II: financial frictions give rise to endogenous limits on the scale and nature of feasible financing, and the precise form of the limit depends on which friction is operative.

%====================================================================
\section{Further Reading}
\label{sec:NV-further-reading}
%====================================================================

The foundational reference for the nonverifiable-income framework is \citet{BoltonScharfstein1990}, who developed the model in the context of predation and the role of financial contracting in shaping competition between firms.
The three-date framework adopted in this chapter is a streamlined adaptation of their analysis.
\citet{BoltonScharfstein1996} extend the analysis to multiple creditors, showing how the optimal debt structure trades off the efficiency benefits of having fewer creditors (which reduce coordination costs in renegotiation) against the strategic benefits of dispersion (which strengthen the no-renegotiation threat).
For an early thesis treatment of renegotiation in debt contracts, see \citet{Gromb1994}.

The dynamic extension of the debt-with-termination framework is provided by \citet{HartMoore1998} (already cited in Chapter~\ref{ch:debt-rationing}), who model default as a form of renegotiation in a dynamic debt model.
The interaction between the termination threat and the entrepreneur's investment decision is studied by \citet{PovelRaith2004} in the case of unobservable investments.

The nonverifiable-income framework has its most developed application in the theory of \emph{sovereign debt}.
A sovereign borrower is, by definition, one who cannot be compelled to repay by any external legal authority; the sustainability of sovereign debt must therefore rest entirely on the borrower's anticipated loss of future market access in the event of default.
The canonical analysis is \citet{BulowRogoff1989a}, who develop a constant-recontracting model of sovereign debt, and \citet{BulowRogoff1989b}, whose striking impossibility result identifies conditions under which \emph{no} sovereign debt is sustainable.
Complementary perspectives on risk-sharing without commitment are provided by \citet{Kocherlakota1996}.
On the interplay between the finitely-repeated-game structure and equilibrium cooperation, see the classic analysis of \citet{KrepsMilgromRobertsWilson1982}.
On sustainable plans in macroeconomic equilibrium with debt, see \citet{ChariKehoe1993}.
On sovereign debt as a form of intertemporal barter, see \citet{KletzerWright2000}.
On competitive risk-sharing under one-sided commitment, see \citet{KruegerUhlig2006}.
On the interplay between bubbles, self-enforcing debt, and private liquidity, see \citet{HellwigLorenzoni2009}.
A unified characterization of self-enforcing debt limits in general equilibrium with recourse and market exclusion is provided by \citet{MartinsPhanVailakis2019}, who decompose debt limits into a present-value-of-recourse component and a credit-bubble component.
Related analyses of collateralized economies with default and constrained efficiency without commitment are developed by \citet{MartinsVailakis2012a}, \citet{MartinsVailakis2012b}, and \citet{MartinsVailakis2015}.
For a comprehensive modern treatment of sovereign debt in limited-commitment environments, covering default and renegotiation, self-fulfilling debt crises, and incomplete markets with their quantitative implications, see \citet{AguiarAmador2014}.

The theoretical constructs developed in this chapter---continuation values, termination probabilities, one-sided commitment---are central to an entire generation of work on dynamic incentives, self-enforcing contracts, and the microfoundations of credit markets.
The reader interested in the frontier of this literature is referred to \citet{Tirole2006} and the references therein for a comprehensive entry point.

%====================================================================
\chapter{Liquidity Shocks and the Maturity of Liabilities}
\label{ch:maturity}
%====================================================================

The principal--agent models of Chapters~\ref{ch:net-worth}--\ref{ch:nonverifiable} assume that, once the project is financed at $t=0$, the only financial event that can occur is the realization of the project's final outcome.
In particular, no intermediate cash flow is required, no further investment is needed, and the only friction is the entrepreneur's effort choice at the moral-hazard stage.
This stylization, useful as a starting point, misses a first-order feature of real-world financing: projects often require \emph{additional} capital injections at intermediate dates to keep going.
A manufacturing firm may need to replenish inventory; a construction project may face unexpected cost overruns; a technology firm may need to invest in an unanticipated market opportunity to remain viable.
Such \emph{liquidity shocks}---random demands for continuation capital that arise after initial financing---pose a financing problem of their own, and their interaction with the moral-hazard friction at the final date shapes both the optimal contract structure and the firm's observable capital structure.

This chapter develops the canonical model of liquidity shocks and maturity structure, following \citet{HolmstromTirole2000} and \citet{HolmstromTirole1998}.
The framework extends the fixed-investment model of Chapter~\ref{ch:net-worth} by adding an intermediate date $t=1$ at which:
\begin{itemize}
  \item a deterministic short-term income $r$ accrues;
  \item a random liquidity shock $\rho$ is drawn from a distribution $F$;
  \item the firm must either pay $\rho$ (and continue) or liquidate (for zero recovery value).
\end{itemize}
The key economic tension is between two competing forces.
On the positive side, continuing the project in \emph{more} states of the world (i.e., choosing a larger continuation cutoff $\rho^\star$) raises expected output and NPV.
On the negative side, a larger $\rho^\star$ requires more pledgeable income to be committed to paying the liquidity shocks, leaving less pledgeable income available to satisfy the initial lender's breakeven constraint.
The optimal contract chooses $\rho^\star$ to trade off these forces, subject to the incentive-compatibility constraint at the final date.

The analysis delivers three distinctive insights.
First, the \emph{first-best} continuation cutoff $\rho_1 = p_h R$---the level at which marginal NPV from continuation equals zero---is attainable only when the entrepreneur's net worth is sufficiently large.
When net worth is low, the second-best contract sets $\rho^\star$ strictly between the pledgeable income $\rho_0 = p_h(R - B/\Delta p)$ and the first-best cutoff $\rho_1$.
Second, the optimal contract for cash-rich firms admits a clean implementation via a combination of short-term debt $d = r - \rho^\star$ and long-term debt $D = R - B/\Delta p$.
The \emph{maturity structure} emerges endogenously as a function of the balance-sheet strength: a weaker balance sheet implies a shorter average maturity.
Third, the model delivers a sharp prediction for cash-poor firms: when the wait-and-see strategy of raising capital on the spot market at $t=1$ is infeasible (because of dilution costs), the firm optimally contracts for a \emph{nonrevocable line of credit} with its initial lender.

The chapter corresponds to Part~A of Chapter~3 of \citet{Tirole2006}.

%====================================================================
\section{The Three-Date Investment Model}
\label{sec:ML-model}
%====================================================================

\subsection{Agents, dates, and preliminaries}

The economy has three dates, $t \in \{0, 1, 2\}$, with a single consumption good and zero discount rate.
Investors and the entrepreneur are risk-neutral.
As in Chapter~\ref{ch:net-worth}, the entrepreneur arrives at $t=0$ with wealth $A < I$ and needs to raise $I - A$ from competitive lenders to finance a project of fixed size $I$.

\subsection{The liquidity shock}

At the intermediate date $t=1$---after the initial investment has been made, but before the moral-hazard stage---two events occur simultaneously:

\begin{itemize}
  \item A \emph{deterministic} short-term income $r \geq 0$ accrues to the firm.
  We assume $0 \leq r < I - A$: the short-term income alone is not enough to repay the lender's outlay.
  The income $r$ is verifiable and is entirely paid to the investor.
  \item A \emph{random} liquidity shock $\rho \geq 0$ is drawn from a continuous distribution with cumulative distribution function $F$ and density $f$ on $[0, \infty)$.
  The shock $\rho$ represents the additional investment needed at $t=1$ to keep the project going through to $t=2$: if the firm injects $\rho$, the project continues; otherwise, it is liquidated at a recovery value of zero.
\end{itemize}

The shock $\rho$ is verifiable (observable to both parties), and the investment $\rho$ is made by the investor rather than the entrepreneur: the entrepreneur has no resources at $t=1$ beyond the short-term income $r$, which has already been paid out.
The distribution $F$ is common knowledge at $t=0$; the realization is observed at $t=1$.

\subsection{The moral-hazard stage at $t=2$}

Conditional on continuation, the project proceeds through a moral-hazard stage identical to the one in Chapter~\ref{ch:net-worth}.
The entrepreneur chooses privately to \emph{behave} (probability $p_h$ of success, no private benefit) or \emph{misbehave} (probability $p_\ell$ of success, private benefit $B > 0$), with $\Delta p \coloneqq p_h - p_\ell > 0$.
Success yields verifiable income $R > 0$; failure yields zero.

If the project is liquidated at $t=1$ (i.e., $\rho > \rho^\star$), no moral-hazard stage occurs: both parties receive zero at $t=2$, and the liquidation value is zero by assumption.

\subsection{Simple contracts}

A \emph{simple contract} specifies:
\begin{itemize}
  \item a continuation cutoff $\rho^\star \geq 0$: the firm continues (the investor funds the reinvestment $\rho$) if $\rho \leq \rho^\star$, and is liquidated if $\rho > \rho^\star$;
  \item a compensation $R_b \in [0, R]$ paid to the entrepreneur at $t=2$ in the event of continuation and success (and zero otherwise).
\end{itemize}
The contract does not pay the entrepreneur at $t=1$ and does not condition the final compensation $R_b$ on the realized liquidity shock $\rho$.
These simplifications are without loss of generality under the present model's assumptions, as discussed briefly in Section~\ref{sec:ML-remarks}.

\subsection{Payoffs}

Fix a simple contract $c = (R_b, \rho^\star)$ and a probability of success $p \in \{p_h, p_\ell\}$ at the moral-hazard stage.
The entrepreneur's net expected payoff, conditional on the realized liquidity shock $\rho$, is
\begin{equation}\label{eq:ML-Pi-b-cond}
  \Pi_b(p, c \vert \rho) + A = \begin{cases}
    p R_b + \mathbb{1}_{p_\ell}(p) B & \text{if } \rho \leq \rho^\star, \\
    0 & \text{if } \rho > \rho^\star.
  \end{cases}
\end{equation}
Integrating over the distribution of $\rho$,
\begin{equation}\label{eq:ML-Pi-b}
  \Pi_b(p, c) + A = F(\rho^\star) \bigl[ p R_b + \mathbb{1}_{p_\ell}(p) B \bigr].
\end{equation}
The indicator $\mathbb{1}_{p_\ell}(p)$ equals $1$ if $p = p_\ell$ (misbehaving) and $0$ otherwise: the private benefit $B$ is collected only if the entrepreneur shirks.

The lender's expected payoff, conditional on $\rho$, is
\begin{equation}\label{eq:ML-Pi-ell-cond}
  \Pi_\ell(p, c \vert \rho) + (I - A) = \begin{cases}
    r - \rho + p(R - R_b) & \text{if } \rho \leq \rho^\star, \\
    r & \text{if } \rho > \rho^\star.
  \end{cases}
\end{equation}
Integrating,
\begin{equation}\label{eq:ML-Pi-ell}
  \Pi_\ell(p, c) + (I - A) = r + F(\rho^\star) p (R - R_b) - \int_0^{\rho^\star} \rho f(\rho) \, d\rho.
\end{equation}
The lender receives the short-term income $r$ for sure, the incremental payment $R - R_b$ in the event of continuation and success, and pays out $\rho$ in the event of continuation.

The total social surplus is
\begin{equation}\label{eq:ML-surplus}
  \Pi_b(p, c) + \Pi_\ell(p, c) = r + F(\rho^\star) \bigl[ p R + \mathbb{1}_{p_\ell}(p) B \bigr] - I - \int_0^{\rho^\star} \rho f(\rho) \, d\rho.
\end{equation}
We denote the right-hand side of \eqref{eq:ML-surplus} by $\NPV(p, \rho^\star)$; it is the net present value of the project under success probability $p$ and continuation cutoff $\rho^\star$.

%====================================================================
\section{Project Viability and Incentive Compatibility}
\label{sec:ML-viability}
%====================================================================

\subsection{The NPV function}

Differentiating \eqref{eq:ML-surplus} in $\rho^\star$,
\begin{equation}\label{eq:ML-dNPV}
  \frac{\partial \NPV}{\partial \rho^\star}(p, \rho^\star) = f(\rho^\star) \bigl[ p R + \mathbb{1}_{p_\ell}(p) B - \rho^\star \bigr].
\end{equation}
The marginal benefit of raising $\rho^\star$ is positive when the continuation payoff $p R + \mathbb{1}_{p_\ell}(p) B$ exceeds the realized shock $\rho^\star$ at the margin, and negative otherwise.
The unconstrained-optimal cutoff under probability $p$ is therefore
\begin{equation*}
  \rho^\star_{\text{FB}}(p) = p R + \mathbb{1}_{p_\ell}(p) B,
\end{equation*}
which is the ``first-best'' continuation cutoff: the shock is paid if and only if the conditional expected continuation value exceeds it.

\subsection{Viability assumptions}

We impose two standard viability conditions:

\begin{assumption}[Positive NPV under behaving]\label{ass:ML-NPV-h}
There exists $\widehat{\rho} \geq 0$ such that $\NPV(p_h, \widehat{\rho}) > 0$.
\end{assumption}

\begin{assumption}[Negative NPV under misbehaving]\label{ass:ML-NPV-l}
For every $\rho \geq 0$, $\NPV(p_\ell, \rho) < 0$.
\end{assumption}

Assumption~\ref{ass:ML-NPV-h} states that the project can be socially valuable if the entrepreneur works; Assumption~\ref{ass:ML-NPV-l} states that it is not worth financing if the entrepreneur shirks.
Together, the two assumptions rule out contracts that induce misbehaving: any contract that does not satisfy incentive compatibility would leave at least one party with negative expected payoff, and hence would not be signed.

\subsection{Incentive compatibility}

From \eqref{eq:ML-Pi-b}, the entrepreneur prefers to behave rather than misbehave if and only if
\begin{equation*}
  F(\rho^\star) \cdot p_h R_b \geq F(\rho^\star) \cdot (p_\ell R_b + B),
\end{equation*}
which simplifies (for $\rho^\star > 0$, so that $F(\rho^\star) > 0$) to the familiar
\begin{equation}\label{eq:ML-IC}
  R_b \geq \frac{B}{\Delta p}.
  \tag{IC}
\end{equation}
The continuation cutoff $\rho^\star$ enters multiplicatively on both sides and does not affect the (IC) constraint: the incentive to behave depends only on the entrepreneur's stake $R_b$ at the final date, not on the probability of reaching that date.

%====================================================================
\section{The Optimal Contract}
\label{sec:ML-optimal}
%====================================================================

\subsection{The optimization}

A contract $c = (R_b, \rho^\star)$ is \emph{optimal} if it maximizes the entrepreneur's expected payoff $\Pi_b(p_h, c) = F(\rho^\star) p_h R_b - A$ subject to the incentive-compatibility constraint \eqref{eq:ML-IC}, the individual-rationality constraint
\begin{equation}\label{eq:ML-IR}
  r + F(\rho^\star) p_h (R - R_b) - \int_0^{\rho^\star} \rho f(\rho) \, d\rho \geq I - A,
  \tag{IR}
\end{equation}
and the feasibility constraints $\rho^\star \geq 0$ and $0 \leq R_b \leq R$.

\subsection{Competitive breakeven}

A first structural result pins down the form of the optimal contract.

\begin{proposition}[Competitive breakeven]\label{prop:ML-breakeven}
If $c = (R_b, \rho^\star)$ is an optimal contract, then the individual-rationality constraint \eqref{eq:ML-IR} binds with equality.
\end{proposition}

\begin{proof}
If (IR) is slack at equilibrium, the entrepreneur could increase $R_b$ to tighten (IR) toward equality while keeping $\rho^\star$ fixed.
The perturbation strictly improves the entrepreneur's payoff $\Pi_b(p_h, c) = F(\rho^\star) p_h R_b - A$ (since $F(\rho^\star) p_h > 0$) and preserves (IC) by weakening the binding direction.
Hence (IR) must bind.
\end{proof}

Under binding (IR), the entrepreneur's payoff simplifies:
\begin{align*}
  \Pi_b(p_h, c) &= F(\rho^\star) p_h R_b - A \\
  &= F(\rho^\star) p_h R - \bigl[ (I - A) + \int_0^{\rho^\star} \rho f(\rho) \, d\rho - r \bigr] - A \\
  &= r + F(\rho^\star) p_h R - I - \int_0^{\rho^\star} \rho f(\rho) \, d\rho \\
  &= \NPV(p_h, \rho^\star).
\end{align*}
The entrepreneur's net payoff equals the project's NPV under behaving, and the entrepreneur's optimization reduces to maximizing $\NPV(p_h, \rho^\star)$ in the cutoff $\rho^\star$, subject to (IC) and feasibility.

\subsection{The binding-(IR) contract}

Solving (IR) for $R_b$ as a function of $\rho^\star$:
\begin{equation}\label{eq:ML-Rb-IR}
  R_b^{\IR}(\rho^\star) \coloneqq R - \frac{1}{F(\rho^\star) p_h} \left[ (I - A) + \int_0^{\rho^\star} \rho f(\rho) \, d\rho - r \right].
\end{equation}
Feasibility requires $R_b^{\IR}(\rho^\star) \leq R$, which holds automatically given the other assumptions; it also requires $R_b^{\IR}(\rho^\star) \geq 0$ (which may fail for small $\rho^\star$ where $F(\rho^\star)$ is too small) and $R_b^{\IR}(\rho^\star) \geq B/\Delta p$ (the (IC) constraint).

\subsection{The pledgeable-income functional}

A key object in the analysis is the maximal expected lender repayment, as a function of the continuation cutoff $\rho^\star$, when the entrepreneur's stake is at the incentive-compatible minimum $R_b = B/\Delta p$:
\begin{equation}\label{eq:ML-pledgeable-functional}
  \mathcal{P}(\rho^\star) \coloneqq r + F(\rho^\star) \underbrace{p_h (R - B/\Delta p)}_{\eqqcolon \rho_0} - \int_0^{\rho^\star} \rho f(\rho) \, d\rho.
\end{equation}
$\mathcal{P}(\rho^\star)$ is the pledgeable income under the binding-(IC) contract with continuation cutoff $\rho^\star$.
Equivalently, $\Pi_\ell(B/\Delta p, \rho^\star) = \mathcal{P}(\rho^\star) - (I - A)$, so the lender breaks even at (IC)-binding contracts with cutoff $\rho^\star$ if and only if $\mathcal{P}(\rho^\star) \geq I - A$.

Differentiating \eqref{eq:ML-pledgeable-functional}, $\mathcal{P}'(\rho^\star) = f(\rho^\star)(\rho_0 - \rho^\star)$, so $\mathcal{P}$ is increasing on $[0, \rho_0]$ and decreasing on $[\rho_0, \infty)$.
The peak is attained at the continuation cutoff $\rho^\star = \rho_0$: this is the cutoff that maximizes pledgeable income.

Intuitively, for $\rho^\star < \rho_0$ the lender passes up profitable continuations (the marginal shock is below the pledgeable increment $\rho_0$); for $\rho^\star > \rho_0$ the lender pays for liquidity shocks above the pledgeable income, a net loss.
The value $\rho_0$ is the maximum liquidity shock that the lender is willing to pay in full, from the perspective of the (IC)-binding pledgeable-income trade-off.

\subsection{Feasibility of financing}

We can now state the feasibility condition for financing.

\begin{proposition}[Feasibility of financing]\label{prop:ML-feasibility}
Financing is feasible if and only if the maximum pledgeable income under the (IC)-binding contract covers the lender's outlay:
\begin{equation}\label{eq:ML-feasibility}
  \mathcal{P}(\rho_0) \geq I - A.
\end{equation}
Equivalently, $\Pi_\ell(B/\Delta p, \rho_0) \geq 0$.
\end{proposition}

\begin{proof}
Feasibility requires the existence of $(R_b, \rho^\star)$ satisfying (IC), (IR), and the feasibility conditions.
By Proposition~\ref{prop:ML-breakeven}, this is equivalent to the existence of $\rho^\star$ with $R_b^{\IR}(\rho^\star) \geq B/\Delta p$, i.e., $\Pi_\ell(B/\Delta p, \rho^\star) \geq 0$.
Since $\mathcal{P}(\rho^\star) = \Pi_\ell(B/\Delta p, \rho^\star) + (I - A)$ and $\mathcal{P}$ is maximized at $\rho_0$, the condition reduces to $\mathcal{P}(\rho_0) \geq I - A$.
\end{proof}

If \eqref{eq:ML-feasibility} fails, no contract can both satisfy (IC) and meet the lender's breakeven: the firm is credit-rationed and the project cannot be financed, even if the project's NPV is strictly positive.

%====================================================================
\section{First-Best and Second-Best Regimes}
\label{sec:ML-regimes}
%====================================================================

Assume throughout this section that financing is feasible: $I - A \leq \mathcal{P}(\rho_0)$.
The optimal contract under this condition depends on whether the first-best continuation cutoff $\rho_1 \coloneqq p_h R$ (the unconstrained optimum of $\NPV(p_h, \rho^\star)$, from \eqref{eq:ML-dNPV} with $p = p_h$) is pledgeable-income-compatible.

\subsection{The relationship $\rho_0 \leq \rho_1$}

Note that
\begin{equation*}
  \rho_1 - \rho_0 = p_h R - p_h(R - B/\Delta p) = \frac{p_h B}{\Delta p} > 0,
\end{equation*}
so the first-best cutoff $\rho_1$ strictly exceeds the pledgeable-income-maximizing cutoff $\rho_0$.
Since $\mathcal{P}$ is decreasing on $[\rho_0, \infty)$,
\begin{equation*}
  \mathcal{P}(\rho_1) < \mathcal{P}(\rho_0).
\end{equation*}
Whether $\mathcal{P}(\rho_1)$ suffices to cover $I - A$ determines which regime applies.

\subsection{First-best regime}

\begin{proposition}[First-best regime]\label{prop:ML-FB}
If $I - A \leq \mathcal{P}(\rho_1)$, then the optimal contract is the \emph{first-best} contract:
\begin{equation*}
  \rho^\star = \rho_1 = p_h R,
  \quad
  R_b = R_b^{\IR}(\rho_1) \geq \frac{B}{\Delta p}.
\end{equation*}
\end{proposition}

\begin{proof}
At $\rho^\star = \rho_1$, the entrepreneur's payoff $\NPV(p_h, \rho^\star)$ is maximized (by \eqref{eq:ML-dNPV} with $p = p_h$).
The lender-breakeven stake $R_b^{\IR}(\rho_1)$ is IC-compatible because the hypothesis $I - A \leq \mathcal{P}(\rho_1)$ implies $\Pi_\ell(B/\Delta p, \rho_1) \geq 0$, equivalently $R_b^{\IR}(\rho_1) \geq B/\Delta p$ (using that $R_b \mapsto \Pi_\ell(R_b, \rho_1)$ is decreasing in $R_b$).
No contract with $\rho^\star \neq \rho_1$ can yield a higher $\NPV$ under behaving, so the contract $(R_b^{\IR}(\rho_1), \rho_1)$ is optimal.
\end{proof}

In the first-best regime, the entrepreneur's wealth $A$ is sufficiently large that the agency rent $p_h B/\Delta p$ does not bind: the lender can be paid entirely from the non-incentive-sensitive margin of the project (the part of $R$ beyond $B/\Delta p$).
The continuation cutoff is set at $\rho_1$, the NPV-maximizing level, and the moral-hazard friction imposes no distortion on the investment side.

\subsection{Second-best regime}

\begin{proposition}[Second-best regime]\label{prop:ML-SB}
If $\mathcal{P}(\rho_1) < I - A \leq \mathcal{P}(\rho_0)$, then the optimal contract is the \emph{second-best} contract:
\begin{equation*}
  R_b = \frac{B}{\Delta p},
  \quad
  \rho^\star \in [\rho_0, \rho_1)
\end{equation*}
with $\rho^\star$ uniquely determined by the binding lender-breakeven condition
\begin{equation}\label{eq:ML-rho-star-SB}
  \Pi_\ell(B/\Delta p, \rho^\star) = 0
  \quad\Longleftrightarrow\quad
  \mathcal{P}(\rho^\star) = I - A.
\end{equation}
\end{proposition}

\begin{proof}
The hypothesis $\mathcal{P}(\rho_1) < I - A$ implies $R_b^{\IR}(\rho_1) < B/\Delta p$: the first-best cutoff is not (IC)-compatible.
In the second-best regime, (IC) must bind: $R_b = B/\Delta p$.
With (IC) binding and the pledgeable income $\mathcal{P}$ being strictly decreasing on $[\rho_0, \infty)$, a unique $\rho^\star \in [\rho_0, \rho_1)$ solves \eqref{eq:ML-rho-star-SB}: the lender's breakeven pins down the cutoff.

For optimality, note that at this $\rho^\star$, the entrepreneur's payoff is $\NPV(p_h, \rho^\star)$.
At any $\widetilde{\rho} > \rho^\star$, we would have $\mathcal{P}(\widetilde{\rho}) < \mathcal{P}(\rho^\star) = I - A$, so $\Pi_\ell(B/\Delta p, \widetilde{\rho}) < 0$, meaning $R_b^{\IR}(\widetilde{\rho}) < B/\Delta p$: (IC) is violated.
Hence no contract with $\widetilde{\rho} > \rho^\star$ is feasible.
At any $\widetilde{\rho} < \rho^\star$, we have $\NPV(p_h, \widetilde{\rho}) \leq \NPV(p_h, \rho^\star)$ because $\NPV(p_h, \cdot)$ is non-decreasing on $[0, \rho_1]$ and $\rho^\star < \rho_1$.
Hence $(B/\Delta p, \rho^\star)$ is optimal.
\end{proof}

In the second-best regime, the entrepreneur's wealth is \emph{not} sufficient to support the first-best continuation policy.
The continuation cutoff $\rho^\star$ is pulled down from $\rho_1$ to an intermediate level in $[\rho_0, \rho_1)$: the lender would like to continue in more states, but doing so would require more pledgeable income than is available under (IC)-binding contracts.
The binding lender-breakeven pins down $\rho^\star$ precisely.
The net distortion takes the form of \emph{premature liquidation}: the project is liquidated in states $\rho \in (\rho^\star, \rho_1]$ where, absent the moral-hazard friction, continuation would have been socially optimal.

\subsection{Summary of the regimes}

Combining Propositions~\ref{prop:ML-FB} and~\ref{prop:ML-SB}:

\begin{itemize}
  \item \emph{If $I - A > \mathcal{P}(\rho_0)$}: no financing (credit rationing);
  \item \emph{If $\mathcal{P}(\rho_1) < I - A \leq \mathcal{P}(\rho_0)$}: second-best, $R_b = B/\Delta p$, $\rho^\star \in [\rho_0, \rho_1)$;
  \item \emph{If $I - A \leq \mathcal{P}(\rho_1)$}: first-best, $R_b = R_b^{\IR}(\rho_1) \geq B/\Delta p$, $\rho^\star = \rho_1$.
\end{itemize}

The entrepreneur's participation constraint $\Pi_b(p_h, c) \geq 0$ must also be checked; under the assumption $I - A \leq \mathcal{P}(\rho_0)$, a sufficient condition is $A < F(\rho_0) \cdot (\rho_1 - \rho_0)$, which is easily satisfied in the regimes of economic interest.

%====================================================================
\section{Implementation via Short- and Long-Term Debt}
\label{sec:ML-implementation}
%====================================================================

We now turn to the central positive question of the chapter: how does the optimal contract translate into an observable capital structure?

\subsection{Cash-rich firms}

A firm is \emph{cash-rich} at the intermediate date if the short-term income $r$ exceeds the continuation cutoff $\rho^\star$:
\begin{equation*}
  r > \rho^\star.
\end{equation*}
This condition holds in particular when $r > \rho_1 = p_h R$, which is a purely primitive condition on the short-term cash-flow generation of the project.

When the firm is cash-rich, the optimal contract can be implemented as follows:

\begin{proposition}[Implementation: cash-rich firms]\label{prop:ML-implementation-cash-rich}
Suppose $r > \rho^\star$.
The optimal contract can be implemented through:
\begin{itemize}
  \item a \emph{short-term (risk-free) debt} of face value $d = r - \rho^\star$, paid at $t=1$;
  \item a \emph{long-term debt} of face value $D = R - B/\Delta p$, paid at $t=2$ contingent on success and continuation.
\end{itemize}
\end{proposition}

\begin{proof}
At $t=1$, the firm's short-term income is $r$.
After paying down the short-term debt of face $d = r - \rho^\star$, the firm retains $\rho^\star$ in cash.
If the realized liquidity shock is $\rho \leq \rho^\star$, the firm pays $\rho$ out of the retained cash and continues; if $\rho > \rho^\star$, the firm is liquidated.
At $t=2$, conditional on continuation and success, the firm receives $R$ and pays the long-term debt $D = R - B/\Delta p$, leaving the entrepreneur with $R_b = B/\Delta p$.
This structure implements the optimal second-best contract.
\end{proof}

The implementation has an intuitive interpretation: the short-term debt $d$ extracts the ``excess'' short-term income that is not needed for liquidity provisioning, while the long-term debt $D$ collects the incentive-compatible maximum repayment at the final date.

\subsection{Maturity structure and balance-sheet strength}

The chapter's signature comparative-static result concerns the relationship between the entrepreneur's net worth $A$ and the maturity structure of the firm's liabilities.
Recall from Proposition~\ref{prop:ML-SB} that $\rho^\star$ is pinned down by $\mathcal{P}(\rho^\star) = I - A$.
Differentiating with respect to $A$ and using $\mathcal{P}'(\rho^\star) < 0$ on $[\rho_0, \rho_1)$:
\begin{equation*}
  \frac{d \rho^\star}{d A} = -\frac{1}{\mathcal{P}'(\rho^\star)} > 0.
\end{equation*}
Hence $\rho^\star$ is \emph{increasing} in $A$, and therefore the short-term debt $d = r - \rho^\star$ is \emph{decreasing} in $A$.

\begin{proposition}[Maturity structure and balance-sheet strength]\label{prop:ML-maturity}
In the second-best regime, increasing the entrepreneur's net worth $A$ decreases the short-term debt $d$ (and leaves the long-term debt $D$ unchanged).
\end{proposition}

\begin{quote}
\emph{A weak balance sheet implies a short maturity structure.}
A highly indebted firm can be viewed as a firm with a weak balance sheet, hence highly indebted firms are more likely to borrow on a short-term basis.
\end{quote}

The intuition is as follows.
A firm with weak balance sheet has a small $A$ and hence a small $\rho^\star$: the lender is unable to commit to continuation in many states of the world, because the firm lacks pledgeable income to cover shocks.
The lender's prudent response is to extract cash flow as early as possible (short-term debt) before the distribution of liquidity shocks reveals whether the firm will survive.
Firms with stronger balance sheets can commit more pledgeable income to continuation; they can afford to give the lender less short-term claim and more long-term claim.

%====================================================================
\section{Cash-Poor Firms}
\label{sec:ML-cash-poor}
%====================================================================

The preceding analysis assumed the firm is cash-rich at $t=1$, $r > \rho^\star$.
For cash-poor firms---e.g., projects with $r = 0$ that yield income only at $t=2$---the implementation of the optimal contract is subtler, and the choice of financing arrangement at $t=0$ becomes a richer question.

\subsection{The wait-and-see strategy}

One option for the cash-poor firm is to leave the intermediate-date financing decision open: at $t=0$, borrow only $I - A$; at $t=1$, return to the capital market to raise the realized $\rho$ on the spot.
Call this the \emph{wait-and-see} strategy.

Consider the case $r = 0$ and $A \geq \overline{A} \coloneqq I - \rho_0$, where $\overline{A}$ is the minimum net worth needed for some financing to be feasible (recall $\mathcal{P}(\rho_0) = F(\rho_0) \rho_0 \geq I - A$ under our continuing assumption that $r = 0$).
At $t=0$, the entrepreneur raises $I - A$ in exchange for shares delivering $R - R_b$ at $t=2$ contingent on success.
At $t=1$, if a liquidity shock $\rho$ is drawn, the entrepreneur returns to the market and issues new shares, diluting the initial investors.

The total pledgeable surplus available at $t=1$, from the initial investors' perspective, is the value of the original claim, which is
\begin{equation*}
  \rho^{\text{W\&S}} = I - A \leq \rho_0.
\end{equation*}
(The bound follows from our maintained assumption $I - A \leq \mathcal{P}(\rho_0) = F(\rho_0)\rho_0 \leq \rho_0$ when $r = 0$.)
If $\rho \leq \rho^{\text{W\&S}}$, the entrepreneur can raise $\rho$ by issuing additional shares without exceeding the pledgeable-income limit.
If $\rho > \rho^{\text{W\&S}}$, fresh equity cannot cover the shock and the firm is liquidated.

\subsection{The dilution mechanic}

The dilution happens through share issuance.
If the realized shock is $\rho = (1/2) \rho^{\text{W\&S}}$, the entrepreneur doubles the number of shares: each existing share then represents half its original claim, and the newly issued shares sell for $(1/2) \rho^{\text{W\&S}} = \rho$, exactly covering the liquidity shock.
If $\rho = (3/4) \rho^{\text{W\&S}}$, the entrepreneur quadruples the number of shares; each original share is reduced to $1/4$ its original claim, and the new shares raise $(3/4) \rho^{\text{W\&S}} = \rho$.
In the limit as $\rho \to \rho^{\text{W\&S}}$, the entrepreneur must issue an unbounded number of shares; at $\rho = \rho^{\text{W\&S}}$, the initial investors are completely diluted (their claim has infinitesimal weight).

The initial investors consent to the dilution because the alternative---liquidating at $t=1$---leaves them with zero recovery.
Any positive dilution ratio is better than liquidation.

\subsection{The expected payoff under wait-and-see}

The wait-and-see strategy succeeds whenever the realized shock is below $\rho^{\text{W\&S}}$, and fails otherwise.
The entrepreneur's expected payoff is
\begin{equation}\label{eq:ML-WS-payoff}
  \Pi_b^{\text{W\&S}} = -A + F(\rho^{\text{W\&S}}) \bigl[ \rho_1 - (I - A) \bigr].
\end{equation}
The term in brackets is the entrepreneur's expected residual claim conditional on continuation: the project's expected final income $\rho_1 = p_h R$ minus the present value of the initial investors' claim $I - A$ (which has been completely transferred via dilution).

\subsection{Wait-and-see vs.\ pre-committed financing}

The wait-and-see strategy is optimal only when it dominates the pre-committed second-best contract derived in Section~\ref{sec:ML-optimal}.
Comparing \eqref{eq:ML-WS-payoff} to the second-best entrepreneur payoff $F(\rho^\star) p_h B/\Delta p$ (under (IC)-binding contracts), the wait-and-see strategy is dominated if
\begin{equation}\label{eq:ML-WS-dominated}
  F(\rho^{\text{W\&S}}) \bigl[ \rho_1 - (I - A) \bigr] \leq F(\rho^\star) p_h \frac{B}{\Delta p}.
\end{equation}

When \eqref{eq:ML-WS-dominated} holds, the entrepreneur strictly prefers to arrange liquidity at $t=0$.
The economic content is that, under wait-and-see, the dilution mechanic caps the firm's continuation capacity at $\rho^{\text{W\&S}} = I - A$, which may be strictly less than the pre-committed cutoff $\rho^\star$ in the second-best regime.
The pre-committed arrangement allows the firm to continue in a wider range of states---including states with $\rho > I - A$---because the lender is contractually bound to provide funding up to $\rho^\star$.

\subsection{Credit lines}

When the wait-and-see strategy is dominated, the entrepreneur and the lender can implement the second-best contract through a \emph{nonrevocable line of credit}: at $t=0$, the lender commits to provide up to $\rho^\star$ of additional funding at $t=1$, contingent on the realized shock.
The commitment is \emph{nonrevocable}: the lender cannot renege if the realized shock $\rho$ exceeds the pledgeable income $\rho_0$, even though fulfilling the commitment in such states may be ex-post unprofitable.

The nonrevocability is essential.
Absent a binding commitment, the lender would have an incentive to refuse funding whenever $\rho > \rho_0$, because the marginal repayment from continuation in such states is less than the continuation cost.
Anticipating this, the entrepreneur would not count on the lender's cooperation in high-shock states, and the effective commitment would revert to the pledgeable-income cap $\rho_0 < \rho^\star$.
The nonrevocable line of credit overcomes this time-inconsistency by making the lender's commitment legally enforceable.

\subsection{Reputation and implicit promises}

In practice, legally enforceable lines of credit are not always feasible, and lenders often make implicit promises that depend on reputational concerns for enforcement.
The banking literature has studied the role of reputation in sustaining such implicit commitments, particularly in the context of ongoing relationships between banks and their borrowers.
For early treatments, see \citet{BootThakorUdell1987} and \citet{BootGreenbaumThakor1993}.

%====================================================================
\section{A Growth-Opportunity Reinterpretation}
\label{sec:ML-growth}
%====================================================================

The three-date model of this chapter admits a natural reinterpretation in which the intermediate-date investment $\rho$ is a \emph{growth opportunity} rather than a continuation cost.
The reinterpretation is useful for connecting the maturity analysis to the empirical literature on firm growth and financing.

\subsection{Modified setup}

At $t=1$, the firm has the option to invest $\rho$ (drawn from the distribution $F$) in order to \emph{increase} the probability of success at $t=2$ by an incremental amount $\tau \geq 0$.
Under the growth-opportunity interpretation:
\begin{itemize}
  \item the baseline success probability is $p_h$ under behaving (or $p_\ell$ under misbehaving);
  \item if the firm invests $\rho$ at $t=1$, the probability rises to $p_h + \tau$ (or $p_\ell + \tau$);
  \item the cutoff $\rho^\star$ determines whether the firm makes the growth investment.
\end{itemize}
The aggregate social surplus is
\begin{equation}\label{eq:ML-growth-surplus}
  \Pi_b + \Pi_\ell = \bigl[ r + (p_h + F(\rho^\star) \tau) R \bigr] - \bigl[ I + \int_0^{\rho^\star} \rho f(\rho) \, d\rho \bigr].
\end{equation}
The surplus function is structurally similar to \eqref{eq:ML-surplus}, with the continuation premium $F(\rho^\star) p_h R$ replaced by the growth premium $F(\rho^\star) \tau R$.

\subsection{Growth opportunities and maturity}

In the second-best regime of the growth model (with $R_b = B/\Delta p$ and an interior $\rho^\star$ satisfying $\tau(R - B/\Delta p) \leq \rho^\star < \tau R$), the maturity-structure comparative static can be computed.
Differentiating the lender's breakeven condition with respect to $\tau$,
\begin{equation}\label{eq:ML-growth-maturity}
  \frac{d(r - \rho^\star)}{d \tau} = -\frac{F(\rho^\star)}{f(\rho^\star)} \cdot \frac{R - B/\Delta p}{\rho^\star - \tau(R - B/\Delta p)} < 0.
\end{equation}
Hence the short-term debt $r - \rho^\star$ is decreasing in the growth parameter $\tau$.

\begin{quote}
\emph{Firms with better growth opportunities should go for longer maturities.}
\end{quote}

The economic logic is transparent.
A firm with greater growth prospects has a larger benefit from continuation at the margin (the incremental success probability $\tau$ multiplies the success income $R$), so the optimal $\rho^\star$ is larger, and the short-term debt is correspondingly smaller.
The firm's maturity structure shifts toward the long end.

%====================================================================
\section{Remarks on Contract Generality}
\label{sec:ML-remarks}
%====================================================================

The simple-contract restriction of Section~\ref{sec:ML-model}---no payment to the entrepreneur at $t=1$, no entrepreneur contribution to the liquidity-shock financing, and a liquidity-shock-independent final compensation---is without loss of generality under the present assumptions.
Informal arguments for each:

\begin{itemize}
  \item \emph{No $t=1$ compensation}: the entrepreneur is risk-neutral and prefers any positive compensation at $t=2$ (weighted by the success probability $p_h$) to the same amount at $t=1$; it is strictly efficient to back-load compensation.
  \item \emph{No entrepreneur contribution to liquidity-shock financing}: the entrepreneur is assumed to have no resources at $t=1$ (the short-term income $r$ accrues to the lender by assumption); there is simply nothing to contribute.
  \item \emph{Liquidity-shock-independent final compensation}: the (IC) constraint depends on $R_b$ only through its value at the moral-hazard stage, and the (IR) constraint can be written in terms of the expected final compensation; state-contingent variation in $R_b(\rho)$ provides no additional flexibility for the principal's problem.
\end{itemize}

Under richer environments (e.g., allowing positive entrepreneur resources at $t=1$, or risk-averse entrepreneurs, or liquidity shocks that are only partially verifiable), the simple-contract restriction may bite.
The analysis of these extensions is pursued in the subsequent parts of Chapter~3 of \citet{Tirole2006}, beyond the scope of this chapter.

%====================================================================
\section{Further Reading}
\label{sec:ML-further-reading}
%====================================================================

The model developed in this chapter is due to \citet{HolmstromTirole1998} and \citet{HolmstromTirole2000}, who introduced the liquidity-shock framework and analyzed its implications for corporate financial structure and for the public supply of liquidity.
The framework has become a workhorse of modern corporate finance, with applications ranging from firm-level capital-structure decisions to macroeconomic models of the interaction between liquidity provision and the real economy.
\citet{Tirole2006} (Chapter~3) is the standard textbook exposition and extends the analysis in several directions covered in subsequent parts of that chapter: the liquidity--scale tradeoff, corporate risk management, endogenous liquidity needs, and free-cash-flow issues.

The observation that weak balance sheets induce short-term borrowing connects to the broader empirical literature on maturity choice.
The \emph{free-cash-flow} theories of \citet{Jensen1986} and the dividend analysis of \citet{Easterbrook1984} provide complementary perspectives on the disciplinary role of cash payouts and debt repayment schedules in containing managerial agency problems.
For the role of bank reputation in sustaining implicit credit commitments---an alternative to the nonrevocable line of credit analyzed in Section~\ref{sec:ML-cash-poor}---see \citet{BootThakorUdell1987} and \citet{BootGreenbaumThakor1993}.

The growth-opportunity reinterpretation of Section~\ref{sec:ML-growth} delivers a prediction about the link between growth and maturity that has found strong empirical support in the corporate-finance literature.
The connection between growth prospects, pledgeable income, and maturity structure is further developed in several papers that build on the Holmstr\"om--Tirole framework, including extensions to international capital markets and to the role of intermediary capital.

The book-length treatment by \citet{HolmstromTirole2011} extends the liquidity-shock framework to a richer environment of inside and outside liquidity, distinguishing between liquidity that the corporate sector can generate internally and liquidity that must be supplied by government or foreign investors.
Post-crisis extensions have emphasized the systemic dimensions of liquidity management: \citet{Tirole2011} provides a broad survey covering market freezes, fire sales, and contagion; \citet{FarhiTirole2012} study how collective moral hazard and maturity mismatch interact with anticipated bailouts; and \citet{DiamondRajan2011} analyze how the prospect of fire sales can itself trigger illiquidity seeking and credit freezes.
A comprehensive survey of financial frictions in macroeconomic models---including many extensions and dynamic reformulations of the Holmstr\"om--Tirole framework---is provided by \citet{BrunnermeierEisenbachSannikov2013}.

%====================================================================
\chapter{The Liquidity--Scale Tradeoff}
\label{ch:liquidity-scale}
%====================================================================

Chapter~\ref{ch:maturity} introduced liquidity shocks into the fixed-investment framework of Chapter~\ref{ch:net-worth}, and examined how the optimal contract structures short- and long-term liabilities to manage interim liquidity needs.
The analysis was ``at the intensive margin'': the investment scale $I$ was held fixed, and the optimal cutoff $\rho^\star$ for continuation was determined by the interplay of pledgeable income and moral hazard.
This chapter extends the framework to the \emph{extensive margin}, allowing the investment scale $I$ to be chosen jointly with the continuation policy $\rho^\star$.
The central new economic force is a \emph{tradeoff between scale and liquidity}: a larger initial investment generates proportionally larger liquidity needs at the interim date, and the entrepreneur must decide how to trade off ambition in project scale against preparedness for liquidity shocks.

The analysis builds on the continuous-investment framework of Chapter~\ref{ch:equity-multiplier}.
Each unit of investment $I$ generates a proportional liquidity shock $\rho I$ at the interim date: if the shock is low, the project proceeds unchanged; if the shock is high, the entrepreneur faces a choice between reinvesting $\rho I$ to save the project or abandoning it for zero recovery.
The liquidity-scale tradeoff emerges starkly: because the shock is proportional to $I$, a larger scale generates more risk of abandonment, and the optimal scale is pinned down by the balance between the value generated per unit of investment and the liquidity cost of supporting that scale.

The chapter develops three distinct results.
First, in the simple two-shock case (where the per-unit liquidity shock takes one of two values), the entrepreneur faces a stark binary choice between two \emph{strategies}: abandon the project in distress states, or pursue it with full reinvestment.
The optimal strategy depends on a simple condition $(1-\lambda)\rho \leq 1$, which pins down when withstanding the shock is worthwhile.
Second, in the continuum-shock case with distribution $F$, the optimal continuation cutoff $\rho^\star$ satisfies the elegant implicit equation $\int_0^{\rho^\star} F(\rho) \, d\rho = 1$, and strictly falls between $\rho_0$ and $\rho_1$.
The cutoff is the value of $\rho^\star$ that \emph{minimizes the effective unit cost of investment}, balancing the scale-expanding benefit of a small $\rho^\star$ (which raises the borrowing multiplier) against the NPV-expanding benefit of a large $\rho^\star$ (which preserves the project in more states).
Third, the analysis reveals a fundamental \emph{governance} problem: absent liquidity requirements, the entrepreneur has an incentive to over-invest in illiquid assets, relying on a ``soft budget constraint'' in which the lender is forced to rescue the project ex post up to $\rho_0$.
This motivates liquidity covenants as a mechanism to enforce the ex-ante optimal cutoff.

The chapter corresponds to Part~B of Chapter~3 of \citet{Tirole2006}, building on the framework of \citet{HolmstromTirole1998} and \citet{HolmstromTirole2000}.

%====================================================================
\section{The Two-Shock Model}
\label{sec:LS-two-shock}
%====================================================================

\subsection{Setup}

We adopt the continuous-investment framework of Chapter~\ref{ch:equity-multiplier}: the entrepreneur has wealth $A$, and each unit of investment $I$ generates income $RI$ in success (probability $p_h$ under behaving, $p_\ell$ under misbehaving) or zero in failure.
The per-unit private benefit from misbehaving is $B$, and the (IC) constraint takes the form $R_b \geq (B/\Delta p) I$.
We maintain the notation $\rho_0 \coloneqq p_h(R - B/\Delta p)$ for the per-unit pledgeable income and $\rho_1 \coloneqq p_h R$ for the per-unit expected success payoff.

The new element is the per-unit liquidity shock $\rho \geq 0$, realized at the interim date $t=1$ after investment but before the moral-hazard stage.
In the two-shock model, we assume:
\begin{itemize}
  \item with probability $1-\lambda$, the per-unit shock is $0$: the firm is \emph{intact} and no reinvestment is needed;
  \item with probability $\lambda \in (0,1)$, the per-unit shock is $\rho > 0$: the firm is in \emph{distress} and continuation requires reinvesting $\rho I$.
\end{itemize}
If the firm is in distress and does not reinvest $\rho I$, the project is abandoned and yields zero income.
If it reinvests, the project continues through to the moral-hazard stage.

\begin{assumption}[Two-shock viability]\label{ass:LS-2shock}
\begin{equation}\label{eq:LS-2shock-assumption}
  0 < \rho_0 < \min\left\{1 + \lambda \rho, \frac{1}{1-\lambda}\right\} < \rho_1.
\end{equation}
\end{assumption}

The condition says that $\rho_0$---the per-unit pledgeable income---is strictly less than the per-unit \emph{effective} investment cost under either strategy (abandon: $1/(1-\lambda)$, the cost of a conditional unit of continuation; pursue: $1+\lambda\rho$, the expected cost including reinvestment); but greater than zero.
Simultaneously, the per-unit expected payoff $\rho_1$ exceeds both costs, ensuring the project has positive NPV under both strategies.

\subsection{Strategy 0: Abandon in distress}

Consider first the contract in which the project is abandoned whenever the liquidity shock $\rho$ arrives.
A contract is a pair $c = (R_b, I)$ with $0 \leq R_b \leq RI$.
The entrepreneur's net expected payoff is
\begin{equation}\label{eq:LS-Pi-b-0}
  \Pi_b^0(p, c) = -A + (1-\lambda)\bigl[p R_b + \mathbb{1}_{p_\ell}(p) BI\bigr],
\end{equation}
reflecting that the project proceeds (with moral hazard) only with probability $1-\lambda$.
The lender's expected payoff is
\begin{equation}\label{eq:LS-Pi-ell-0}
  \Pi_\ell^0(p, c) = -(I-A) + (1-\lambda) p (RI - R_b).
\end{equation}
At the optimum, (IC) binds ($R_b = BI/\Delta p$) and (IR) binds (competitive lending), yielding
\begin{equation*}
  (1-\lambda) \rho_0 I = I - A,
\end{equation*}
which gives the investment capacity under Strategy 0:
\begin{equation}\label{eq:LS-I-0}
  I^0 = \frac{A}{1 - (1-\lambda) \rho_0}.
\end{equation}
The entrepreneur's net payoff under Strategy 0 is
\begin{equation}\label{eq:LS-Pi-0}
  \Pi_b^0 = \bigl[(1-\lambda) \rho_1 - 1\bigr] I^0 = \frac{(1-\lambda)\rho_1 - 1}{1 - (1-\lambda)\rho_0} A = \frac{\rho_1 - \frac{1}{1-\lambda}}{\frac{1}{1-\lambda} - \rho_0} A.
\end{equation}
The final form of \eqref{eq:LS-Pi-0} gives a transparent reading: under Strategy 0, the ``effective'' per-unit cost of investment is $1/(1-\lambda)$---the cost of ensuring that a unit of investment actually reaches the moral-hazard stage with probability one---and the entrepreneur's net payoff per unit of initial wealth is the net margin $(\rho_1 - 1/(1-\lambda))$ divided by the spread between cost and pledgeable income $(1/(1-\lambda) - \rho_0)$.

\subsection{Strategy 1: Pursue in distress}

Consider next the contract in which the project is pursued in the distress state, with the lender committing ex ante to fund the reinvestment $\rho I$.
The entrepreneur's net expected payoff is
\begin{equation}\label{eq:LS-Pi-b-1}
  \Pi_b^1(p, c) = -A + p R_b + \mathbb{1}_{p_\ell}(p) BI,
\end{equation}
and the lender's net payoff is
\begin{equation}\label{eq:LS-Pi-ell-1}
  \Pi_\ell^1(p, c) = -(I-A) - \lambda \rho I + p (RI - R_b),
\end{equation}
where the term $-\lambda \rho I$ is the expected reinvestment cost.

At the optimum, (IC) binds ($R_b = BI/\Delta p$) and (IR) binds, yielding
\begin{equation*}
  (1 + \lambda \rho) I - A = \rho_0 I,
\end{equation*}
i.e., the lender's total expected outlay $I - A + \lambda\rho I$ equals the expected pledgeable income $\rho_0 I$.
The investment capacity under Strategy 1 is
\begin{equation}\label{eq:LS-I-1}
  I^1 = \frac{A}{(1+\lambda\rho) - \rho_0}.
\end{equation}
The entrepreneur's net payoff under Strategy 1 is
\begin{equation}\label{eq:LS-Pi-1}
  \Pi_b^1 = \bigl[\rho_1 - (1+\lambda\rho)\bigr] I^1 = \frac{\rho_1 - (1+\lambda\rho)}{(1+\lambda\rho) - \rho_0} A.
\end{equation}
The economic reading parallels Strategy 0: the effective per-unit cost of investment is $1+\lambda\rho$, the expected total outlay per unit reaching the moral-hazard stage, and the entrepreneur's net payoff per unit of wealth is the net margin divided by the spread between cost and pledgeable income.

\subsection{Optimal strategy}

Comparing \eqref{eq:LS-Pi-0} and \eqref{eq:LS-Pi-1}:

\begin{proposition}[Optimal two-shock strategy]\label{prop:LS-strategy}
Strategy 1 (pursue) dominates Strategy 0 (abandon) if and only if
\begin{equation}\label{eq:LS-policy}
  (1-\lambda) \rho \leq 1.
\end{equation}
\end{proposition}

\begin{proof}
Given the common form of the payoff expressions \eqref{eq:LS-Pi-0} and \eqref{eq:LS-Pi-1}, with the pledgeable income $\rho_0$ held fixed, Strategy 1 delivers a higher payoff iff the effective cost under Strategy 1 is smaller:
\begin{equation*}
  1 + \lambda \rho \leq \frac{1}{1-\lambda},
\end{equation*}
which rearranges to $(1-\lambda)\rho \leq 1$.
(A short algebraic verification: cross-multiplying and using $\lambda \in (0,1)$ preserves inequality direction.)
\end{proof}

The intuition of \eqref{eq:LS-policy} is economically natural.
It is optimal to pursue the project in the distress state if:
\begin{itemize}
  \item the liquidity shock $\rho$ is \emph{low} (so the cost of rescuing is modest), or
  \item the probability of distress $\lambda$ is \emph{high} (so abandoning often forgoes productive opportunities).
\end{itemize}
The tradeoff is between the scale advantage of commitment-to-pursue (higher $I$ per unit of wealth) versus the resource cost of reinvestment $\lambda\rho$ per unit of investment.
When the probability-weighted reinvestment cost $\lambda\rho$ is below a threshold, the commitment-to-pursue strategy strictly dominates.

%====================================================================
\section{The Continuum Model}
\label{sec:LS-continuum}
%====================================================================

The two-shock case is analytically clean but economically restrictive: in practice, the liquidity shock is a continuous random variable with a rich distribution.
We now extend the analysis to a continuous distribution of per-unit liquidity shocks $\rho$, with cumulative distribution function $F$ and density $f$ on $[0, \infty)$.

\subsection{Setup}

After the investment $I$ is sunk at $t=0$, a per-unit liquidity shock $\rho$ is realized at $t=1$: a cash infusion of $\rho I$ is required to continue the project.
If the reinvestment is not made, the project is abandoned with zero recovery.
If it is made, the project proceeds to the moral-hazard stage and yields $RI$ in success (probability $p$) or zero in failure.

A contract is a triple
\begin{equation}\label{eq:LS-contract}
  c = (R_b, I, \rho^\star),
\end{equation}
specifying:
\begin{itemize}
  \item the investment scale $I \geq 0$;
  \item the entrepreneur's stake $R_b$ with $0 \leq R_b \leq RI$;
  \item the continuation cutoff $\rho^\star \geq 0$: the project is pursued if $\rho \leq \rho^\star$ and abandoned if $\rho > \rho^\star$.
\end{itemize}
The incentive-compatibility constraint is still $(\Delta p) R_b \geq BI$, equivalent to $R_b \geq (B/\Delta p) I$.

\subsection{NPV and viability}

The project's NPV under probability $p$ and cutoff $\rho^\star$ is
\begin{equation}\label{eq:LS-NPV}
  \NPV(p, \rho^\star) = I \left[ -1 - \int_0^{\rho^\star} \rho f(\rho) \, d\rho + F(\rho^\star) \left\{ pR + \mathbb{1}_{p_\ell}(p) B \right\} \right].
\end{equation}
Differentiating in $\rho^\star$, $\partial \NPV/\partial \rho^\star = I f(\rho^\star) [pR + \mathbb{1}_{p_\ell}(p) B - \rho^\star]$, so the unconstrained-optimal cutoff under probability $p$ is $pR + \mathbb{1}_{p_\ell}(p) B$.

\begin{assumption}[NPV at extremes]\label{ass:LS-extremes}
\begin{equation}\label{eq:LS-extremes}
  \NPV(p_\ell, p_\ell R + B) < 0 < \NPV(p_h, p_h R).
\end{equation}
\end{assumption}

\begin{assumption}[Project viability]\label{ass:LS-viability}
\begin{equation}\label{eq:LS-viability}
  F(\rho_1) \rho_1 - 1 - \int_0^{\rho_1} \rho f(\rho) \, d\rho > 0 > F(\rho_0) \rho_0 - 1 - \int_0^{\rho_0} \rho f(\rho) \, d\rho.
\end{equation}
\end{assumption}

The first inequality in \eqref{eq:LS-viability} says that at the first-best cutoff $\rho^\star = \rho_1$, the project has positive NPV under behaving.
The second says that at the pledgeable-income-maximizing cutoff $\rho^\star = \rho_0$, the pledgeable income per unit of investment ($F(\rho_0)\rho_0 - \int_0^{\rho_0}\rho f \, d\rho$) is less than unity---so at that cutoff, one would not be able to self-finance the investment without entrepreneur equity.
Together, the conditions locate the ``interesting'' zone: there is some $\widehat{\rho}$ for which the project is both profitable under behaving and requires the entrepreneur's equity to be financeable.

\subsection{Payoff structure}

Given an incentive-compatible contract $c$, the entrepreneur's conditional and unconditional net expected payoffs are
\begin{equation}\label{eq:LS-Pi-b-cond}
  \Pi_b(c \vert \rho) + A = \begin{cases}
    p_h R_b & \text{if } \rho \leq \rho^\star, \\
    0 & \text{if } \rho > \rho^\star,
  \end{cases}
\end{equation}
and
\begin{equation}\label{eq:LS-Pi-b-uncond}
  \Pi_b(c) + A = F(\rho^\star) p_h R_b.
\end{equation}
The lender's conditional and unconditional payoffs are
\begin{equation}\label{eq:LS-Pi-ell-cond}
  \Pi_\ell(c \vert \rho) + (I-A) = \begin{cases}
    p_h(RI - R_b) - \rho I & \text{if } \rho \leq \rho^\star, \\
    0 & \text{if } \rho > \rho^\star,
  \end{cases}
\end{equation}
and
\begin{equation}\label{eq:LS-Pi-ell-uncond}
  \Pi_\ell(c) + (I-A) = F(\rho^\star) p_h (RI - R_b) - I \int_0^{\rho^\star} \rho f(\rho) \, d\rho.
\end{equation}
The total social surplus is
\begin{equation}\label{eq:LS-surplus}
  \Pi_b(c) + \Pi_\ell(c) = I \left[ F(\rho^\star) \rho_1 - 1 - \int_0^{\rho^\star} \rho f(\rho) \, d\rho \right].
\end{equation}

\subsection{Optimization and viable range}

The entrepreneur's program is to maximize $\Pi_b(c)$ subject to (IC), (IR), and feasibility.

\begin{proposition}[Viability constraints]\label{prop:LS-viable-range}
Any optimal contract $c = (R_b, I, \rho^\star)$ satisfies
\begin{equation}\label{eq:LS-viable-range}
  F(\rho^\star) \rho_1 - 1 - \int_0^{\rho^\star} \rho f(\rho) \, d\rho > 0 > F(\rho^\star) \rho_0 - 1 - \int_0^{\rho^\star} \rho f(\rho) \, d\rho.
\end{equation}
\end{proposition}

\begin{proof}
The first inequality is necessary for positive NPV; if violated, the entrepreneur prefers not to invest.
The second inequality is necessary for the problem to have a finite solution: if $F(\rho^\star)\rho_0 - 1 - \int_0^{\rho^\star}\rho f \, d\rho \geq 0$, the lender's breakeven is achieved at the (IC)-binding stake for \emph{any} scale $I$, so scale could be increased without bound, and the maximization has no interior solution.
The strict inequality in the second condition therefore pins down the regime in which both (IC) and (IR) bind and the scale is finite.
\end{proof}

Under the maintained viability assumptions \eqref{eq:LS-viability}, the second inequality in \eqref{eq:LS-viable-range} holds automatically for every $\rho^\star$ (since $F(\rho)\rho_0 - 1 - \int_0^\rho \rho'f(\rho') \, d\rho'$ is maximized at $\rho_0$ and the max is negative by Assumption~\ref{ass:LS-viability}), while the first inequality holds in an open neighborhood of $\rho_1$.

\subsection{Investment capacity and the multiplier function}

Within the viable range, (IC) and (IR) both bind at the optimum, yielding $R_b = (B/\Delta p) I$ and
\begin{equation*}
  F(\rho^\star) \rho_0 I = I - A + I \int_0^{\rho^\star} \rho f(\rho) \, d\rho.
\end{equation*}
Solving for $I$:
\begin{equation}\label{eq:LS-I-multiplier}
  I = k(\rho^\star) \, A,
  \quad
  k(\rho^\star) \coloneqq \frac{1}{1 + \int_0^{\rho^\star} \rho f(\rho) \, d\rho - F(\rho^\star) \rho_0}.
\end{equation}
The \emph{multiplier function} $k(\rho^\star)$ generalizes the equity multiplier of Chapter~\ref{ch:equity-multiplier} to the present setting with continuous liquidity shocks.
Its denominator is the effective per-unit equity requirement: the per-unit cost of investment (1 plus the expected interim reinvestment $\int_0^{\rho^\star}\rho f \, d\rho$) minus the pledgeable income per unit conditional on continuation $F(\rho^\star)\rho_0$.

Differentiating the denominator in $\rho^\star$:
\begin{equation*}
  \frac{d}{d\rho^\star} \left[ 1 + \int_0^{\rho^\star} \rho f(\rho) \, d\rho - F(\rho^\star) \rho_0 \right] = f(\rho^\star)(\rho^\star - \rho_0).
\end{equation*}
The denominator is decreasing for $\rho^\star < \rho_0$ and increasing for $\rho^\star > \rho_0$.
Therefore $k(\rho^\star)$ is \emph{maximized} at $\rho^\star = \rho_0$: the multiplier is largest when the cutoff exactly equals the per-unit pledgeable income.

This is the first key tradeoff of the chapter.
To maximize \emph{borrowing capacity}, the entrepreneur wants to set $\rho^\star = \rho_0$: continue only in those states where the shock is no larger than what the pledgeable income can absorb.
But to maximize \emph{NPV}, the entrepreneur wants a larger cutoff---the first-best cutoff would be $\rho^\star = \rho_1 > \rho_0$.

\subsection{Entrepreneur's payoff and marginal function}

Combining the multiplier with the NPV-per-unit, the entrepreneur's net payoff is
\begin{equation}\label{eq:LS-Pi-final}
  \Pi_b(\rho^\star) = m(\rho^\star) \cdot k(\rho^\star) \cdot A,
\end{equation}
where the \emph{per-unit margin function} is
\begin{equation}\label{eq:LS-m}
  m(\rho^\star) \coloneqq F(\rho^\star) \rho_1 - 1 - \int_0^{\rho^\star} \rho f(\rho) \, d\rho.
\end{equation}
Differentiating, $m'(\rho^\star) = f(\rho^\star)(\rho_1 - \rho^\star)$, so $m$ is increasing on $[0, \rho_1]$ and decreasing on $[\rho_1, \infty)$.
The margin $m$ is maximized at $\rho^\star = \rho_1$---the first-best continuation rule, which delivers the ``ex-post optimal'' continuation frontier.

The entrepreneur's problem is now to maximize the product $m(\rho^\star) \cdot k(\rho^\star)$, balancing the NPV-maximizing force $m$ (peaked at $\rho_1$) against the scale-maximizing force $k$ (peaked at $\rho_0$).

\subsection{The cost minimization characterization}

Writing the payoff in a more economically revealing form, we have
\begin{equation}\label{eq:LS-payoff-ratio}
  \Pi_b(\rho^\star) = \frac{\rho_1 - c(\rho^\star)}{c(\rho^\star) - \rho_0} A,
\end{equation}
where the \emph{expected unit cost of effective investment} is
\begin{equation}\label{eq:LS-c}
  c(\rho^\star) \coloneqq \frac{1 + \int_0^{\rho^\star} \rho f(\rho) \, d\rho}{F(\rho^\star)}.
\end{equation}
The numerator of $c(\rho^\star)$ is the total expected per-unit cost of an investment that continues with probability $F(\rho^\star)$: the initial investment (1) plus the expected interim reinvestment ($\int_0^{\rho^\star} \rho f \, d\rho$).
The denominator $F(\rho^\star)$ is the probability that the project survives to the moral-hazard stage.
The ratio $c(\rho^\star)$ is therefore the expected cost \emph{per unit that actually reaches the moral-hazard stage}.

The payoff \eqref{eq:LS-payoff-ratio} shows that maximizing $\Pi_b$ in $\rho^\star$ is equivalent to minimizing the cost function $c(\rho^\star)$: both the numerator $\rho_1 - c(\rho^\star)$ (the net margin) and the inverse denominator $1/(c(\rho^\star) - \rho_0)$ (the multiplier) are monotonically decreasing functions of $c(\rho^\star)$.
This gives the central characterization of the chapter.

\subsection{The optimal cutoff}

Integrating by parts $\int_0^{\widehat{\rho}} \rho f(\rho) \, d\rho = \widehat{\rho} F(\widehat{\rho}) - \int_0^{\widehat{\rho}} F(\rho) \, d\rho$, so
\begin{equation}\label{eq:LS-c-ibp}
  c(\widehat{\rho}) = \widehat{\rho} + \frac{1 - \int_0^{\widehat{\rho}} F(\rho) \, d\rho}{F(\widehat{\rho})}.
\end{equation}
Differentiating \eqref{eq:LS-c-ibp},
\begin{equation*}
  c'(\widehat{\rho}) = 1 + \frac{-F(\widehat{\rho}) \cdot F(\widehat{\rho}) - [1 - \int_0^{\widehat{\rho}} F(\rho) \, d\rho] \cdot f(\widehat{\rho})}{F(\widehat{\rho})^2}.
\end{equation*}
Setting $c'(\widehat{\rho}) = 0$ yields, after simplification,
\begin{equation*}
  F(\widehat{\rho})^2 = -f(\widehat{\rho}) \left[ 1 - \int_0^{\widehat{\rho}} F(\rho) \, d\rho \right] + F(\widehat{\rho})^2,
\end{equation*}
which reduces to $1 - \int_0^{\widehat{\rho}} F(\rho) \, d\rho = 0$, i.e.,
\begin{equation}\label{eq:LS-FOC}
  \int_0^{\rho^\star} F(\rho) \, d\rho = 1.
\end{equation}
This is the first-order condition characterizing the optimal cutoff.

\begin{proposition}[Optimal cutoff characterization]\label{prop:LS-FOC}
The optimal continuation cutoff $\rho^\star$ is uniquely determined by
\begin{equation*}
  \int_0^{\rho^\star} F(\rho) \, d\rho = 1,
\end{equation*}
equivalently, by the condition $c(\rho^\star) = \rho^\star$: the threshold liquidity shock equals the expected unit cost of effective investment.
At the optimum, the entrepreneur's net expected payoff is
\begin{equation}\label{eq:LS-Pi-optimum}
  \Pi_b = \frac{\rho_1 - \rho^\star}{\rho^\star - \rho_0} A.
\end{equation}
\end{proposition}

\begin{proof}
We have shown that $c'(\rho^\star) = 0$ is equivalent to \eqref{eq:LS-FOC}.
At this point, substituting \eqref{eq:LS-c-ibp} with $1 - \int_0^{\rho^\star} F \, d\rho = 0$ yields $c(\rho^\star) = \rho^\star$.
Substituting back into \eqref{eq:LS-payoff-ratio} gives \eqref{eq:LS-Pi-optimum}.
Uniqueness follows from strict convexity of $c$ on the viable range---a consequence of $F$ being strictly increasing on the support of $\rho$---which we verify informally by noting that $c$ must have a unique interior minimum between its two ``bad corners'' ($\widehat{\rho} \to 0$ and $\widehat{\rho} \to \infty$).
\end{proof}

The characterization is analytically elegant: the FOC \eqref{eq:LS-FOC} does not depend on the agency parameters $R$, $B$, $\Delta p$---the optimal cutoff is determined purely by the distribution of shocks.
The agency parameters enter only through the net payoff expression \eqref{eq:LS-Pi-optimum}, via the constants $\rho_0$ and $\rho_1$.

\subsection{Bounds on the optimal cutoff}

A key qualitative result follows from the trade-off between $m$ and $k$:

\begin{proposition}[Location of the optimum]\label{prop:LS-bounds}
The optimal cutoff lies strictly between the pledgeable income and the expected success payoff:
\begin{equation}\label{eq:LS-bounds}
  \rho_0 < \rho^\star < \rho_1.
\end{equation}
\end{proposition}

\begin{proof}
For $\rho^\star < \rho_0$: both $m$ and $k$ are increasing (recall $m$ peaks at $\rho_1 > \rho_0$ and $k$ peaks at $\rho_0$).
Any small increase in $\rho^\star$ strictly improves $\Pi_b = m k$.
Hence $\rho^\star \geq \rho_0$.
For $\rho^\star > \rho_1$: both $m$ and $k$ are decreasing.
A small decrease in $\rho^\star$ strictly improves $\Pi_b$.
Hence $\rho^\star \leq \rho_1$.
Strict inequalities follow from the tradeoff: at $\rho^\star = \rho_0$, $m$ can still be increased without losing $k$'s maximum; at $\rho^\star = \rho_1$, $k$ can be increased without losing $m$'s maximum.
\end{proof}

The bound \eqref{eq:LS-bounds} is the chapter's signature qualitative result.
The optimal cutoff is neither the ``pledgeable income maximum'' ($\rho_0$, which would maximize borrowing capacity at the cost of NPV) nor the ``first-best'' cutoff ($\rho_1$, which would maximize NPV at the cost of borrowing capacity).
It is a strict interior point, reflecting the liquidity--scale tradeoff that gives the chapter its name.

%====================================================================
\section{Comparison with Wait-and-See}
\label{sec:LS-wait-and-see}
%====================================================================

An important corollary of Proposition~\ref{prop:LS-bounds} concerns the viability of the wait-and-see strategy (Section~\ref{sec:ML-cash-poor} of Chapter~\ref{ch:maturity}).
Under wait-and-see, the lender provides interim financing at $t=1$ only if the realized shock can be covered by the remaining pledgeable income.
With no continuation premium beyond the pledgeable income, the lender's ex-post participation requires $\rho \leq \rho_0$.
In contrast, Proposition~\ref{prop:LS-bounds} establishes that the \emph{ex-ante optimal} cutoff strictly exceeds $\rho_0$.

\begin{proposition}[Wait-and-see is suboptimal]\label{prop:LS-WS-suboptimal}
Since $\rho^\star > \rho_0$, the wait-and-see strategy is strictly suboptimal: the lender's ex-post unwillingness to fund shocks in $(\rho_0, \rho^\star]$ forgoes value relative to the ex-ante-committed contract.
\end{proposition}

The economic content is sharp.
The optimal contract commits the lender to fund shocks in the range $(\rho_0, \rho^\star]$ even though, ex post, such funding is strictly unprofitable (the pledgeable income $\rho_0$ is less than the realized cost).
The entrepreneur captures the residual surplus from continuation in those states, and the ex-ante commitment makes both parties better off.

This creates what the corporate-finance literature calls a \emph{demand for liquidity}: the firm's optimal contract requires pre-committed financing for states in which the lender would not voluntarily lend ex post.
The demand can be met either by the firm hoarding cash (self-insurance) or by contractual arrangements such as lines of credit.
The implementation is taken up in the next section.

%====================================================================
\section{Liquidation Value}
\label{sec:LS-liquidation}
%====================================================================

We now extend the analysis to include a positive \emph{liquidation value}: the firm's assets have a per-unit salvage value $L \geq 0$ that can be recovered if the project is abandoned at $t=1$.

\subsection{Modified multiplier and margin}

With liquidation value $L$ per unit of investment, the lender recovers $LI$ in the liquidated state $\rho > \rho^\star$ in addition to the continuation payoffs in the non-liquidated state.
Repeating the optimization:
\begin{equation}\label{eq:LS-k-L}
  k(\rho^\star) = \frac{1}{\bigl[1 - L + \int_0^{\rho^\star} \rho f(\rho) \, d\rho\bigr] - F(\rho^\star)(\rho_0 - L)}
\end{equation}
and
\begin{equation}\label{eq:LS-m-L}
  m(\rho^\star) = F(\rho^\star)(\rho_1 - L) - \left[1 - L + \int_0^{\rho^\star} \rho f(\rho) \, d\rho\right].
\end{equation}

The liquidation value enters on both sides of the tradeoff.
On the borrowing-capacity side, $L$ reduces the effective per-unit cost ($1 \to 1 - L$) but also reduces the pledgeable income benefit of continuation ($\rho_0 \to \rho_0 - L$, since continuation now forgoes $L$).
On the margin side, $L$ reduces both the success premium and the baseline cost.
The net effect depends on the specifics of the distribution.

\subsection{Optimal cutoff with salvage}

Repeating the cost-minimization derivation with liquidation value:
\begin{equation}\label{eq:LS-c-L}
  c(\rho^\star) = \frac{1 - L + \int_0^{\rho^\star} \rho f(\rho) \, d\rho}{F(\rho^\star)} = \rho^\star + \frac{1 - L - \int_0^{\rho^\star} F(\rho) \, d\rho}{F(\rho^\star)}.
\end{equation}
The FOC yields
\begin{equation}\label{eq:LS-FOC-L}
  \int_0^{\rho^\star} F(\rho) \, d\rho = 1 - L,
\end{equation}
with $c(\rho^\star) = \rho^\star$ at the optimum.
The entrepreneur's net expected payoff is
\begin{equation}\label{eq:LS-Pi-L}
  \Pi_b = \frac{(\rho_1 - L) - \rho^\star}{\rho^\star - (\rho_0 - L)} A.
\end{equation}

\begin{proposition}[Bounds with liquidation value]\label{prop:LS-bounds-L}
The optimal cutoff satisfies
\begin{equation*}
  \rho_0 - L < \rho^\star < \rho_1 - L,
\end{equation*}
and the derivative with respect to $L$ is
\begin{equation}\label{eq:LS-dL}
  \frac{\partial \rho^\star}{\partial L} = -\frac{1}{F(\rho^\star)} < 0.
\end{equation}
\end{proposition}

\begin{proof}
The bounds follow by the same argument as Proposition~\ref{prop:LS-bounds}, with $\rho_1 - L$ and $\rho_0 - L$ replacing $\rho_1$ and $\rho_0$.
The derivative follows from implicit differentiation of the FOC \eqref{eq:LS-FOC-L}: $F(\rho^\star) \cdot \partial \rho^\star / \partial L = -1$.
\end{proof}

The economic reading of \eqref{eq:LS-dL} is intuitive.
As the liquidation value $L$ rises, the opportunity cost of continuing in high-shock states increases (continuation now forgoes the salvage $L$), and the optimal cutoff shifts down.
The gap between the ex-ante-optimal cutoff and the wait-and-see outcome (where lenders rescue only if $\rho \leq \rho_0 - L$) narrows as $L$ rises.
In the limiting case where $L = \rho_0$ (full-collateral liquidation), the cutoff converges to the wait-and-see outcome: there is no longer a strict benefit to ex-ante commitment, because the lender can always recover his outlay ex post through liquidation.

%====================================================================
\section{The Implementation Problem: Liquidity Covenants}
\label{sec:LS-covenants}
%====================================================================

The optimal contract requires ex-ante commitment to fund liquidity shocks in the range $(\rho_0, \rho^\star]$.
How is this commitment implemented?
The answer depends on the allocation of cash holdings between liquid and illiquid assets, and on the design of the governance structure that prevents the entrepreneur from subverting the ex-ante arrangement.

\subsection{Front-loaded financing}

One implementation is to have the lender provide the entire expected funding at $t=0$:
\begin{itemize}
  \item the lender transfers $I(1 + \rho^\star) - A$ to the firm at $t=0$;
  \item the firm invests $I$ in the project and holds $\rho^\star I$ in liquid (cash-equivalent) assets;
  \item at $t=1$, the firm uses the liquid holdings to cover shocks up to $\rho^\star$ per unit; for shocks above $\rho^\star$, the project is abandoned.
\end{itemize}
This implementation requires a \emph{liquidity ratio}: the firm must maintain liquid assets equal to $\rho^\star/(1 + \rho^\star)$ of its total balance sheet until the shock is realized.
Absent such a requirement, the firm may prefer to invest the excess funds in illiquid (productive) assets, effectively over-investing at $t=0$ and shifting the risk of the shock back to the lender.

\subsection{The soft-budget-constraint problem}

To see the governance problem starkly, suppose the entrepreneur is free to allocate the $t=0$ funds as he wishes.
He would then invest the full $I^\star \coloneqq I(1 + \rho^\star)$ in illiquid assets at $t=0$---scaling up the project beyond the ex-ante-optimal size.

This choice is not sustainable if the lender can commit credibly to not providing additional funds at $t=1$.
But such commitment is typically not credible: at $t=1$, facing the \emph{fait accompli} of an overinvestment in illiquid assets, the lender has an incentive to rescue the firm whenever doing so is strictly profitable, i.e., whenever the realized shock satisfies
\begin{equation*}
  \rho \leq \rho_0,
\end{equation*}
since at such shocks the pledgeable income (per unit of continuation) at least covers the reinvestment cost.

\begin{proposition}[Soft budget constraint]\label{prop:LS-soft-budget}
Absent ex-ante enforced liquidity requirements, the entrepreneur over-invests in illiquid assets at $t=0$, relying on the lender's ex-post rescue up to $\rho_0$.
The resulting allocation is strictly suboptimal.
\end{proposition}

The phrase \emph{soft budget constraint}---due originally to Kornai in the context of socialist enterprises but widely applied to financial frictions---captures the key friction: the lender's inability to commit not to rescue creates a time-inconsistency that the entrepreneur can exploit by over-investing.
The economically striking feature is that this is a \emph{bilateral} inefficiency: both parties are worse off ex ante relative to the committed contract, yet each is rationally responding to the other's ex-post incentives.

\subsection{When over-investment is privately profitable}

The entrepreneur prefers investing $I^\star = I(1 + \rho^\star)$ rather than $I$ iff
\begin{equation*}
  F(\rho^\star) p_h \frac{B I}{\Delta p} < F(\rho_0) p_h \frac{B I^\star}{\Delta p},
\end{equation*}
which simplifies (using $I^\star = I(1 + \rho^\star)$) to
\begin{equation}\label{eq:LS-over-investment}
  F(\rho^\star) < F(\rho_0)(1 + \rho^\star).
\end{equation}
Condition \eqref{eq:LS-over-investment} holds for many realistic distributions $F$: the entrepreneur's incentive to over-invest is the rule rather than the exception.
The lender, rationally anticipating this, cannot fund the project at the originally intended terms and must impose liquidity requirements as a contractual device.

\subsection{Monitoring over-hoarding of liquid assets}

The mirror-image problem also arises: the entrepreneur may prefer to \emph{under-invest} in illiquid assets, hoarding more liquidity than required in order to over-insure against liquidity shocks.
Suppose the entrepreneur invests $I' \leq I$ in illiquid assets, retaining liquid holdings of $\rho^\star I + (I - I')$.
The firm can withstand shocks $\rho$ satisfying
\begin{equation*}
  \rho I' \leq \rho^\star I + (I - I').
\end{equation*}
Letting $\varepsilon \coloneqq (I - I')/I'$, the entrepreneur prefers to under-invest iff
\begin{equation}\label{eq:LS-under-invest}
  F(\rho^\star + (1 + \rho^\star)\varepsilon) \cdot I' > F(\rho^\star) \cdot I
  \quad\Leftrightarrow\quad
  F(\rho^\star + (1 + \rho^\star)\varepsilon) > F(\rho^\star)(1 + \varepsilon).
\end{equation}
For small $\varepsilon$, this is equivalent (by first-order Taylor expansion) to
\begin{equation}\label{eq:LS-hazard}
  \frac{(1 + \rho^\star) f(\rho^\star)}{F(\rho^\star)} > 1.
\end{equation}
The left-hand side of \eqref{eq:LS-hazard} is related to the hazard function of $F$ at $\rho^\star$, weighted by the reciprocal of the effective cutoff scale.
When liquidity shocks are likely around the cutoff---i.e., when the density $f(\rho^\star)$ is large relative to $F(\rho^\star)/(1 + \rho^\star)$---the entrepreneur has a private incentive to hoard extra liquidity.

\begin{proposition}[Liquidity-hoarding temptation]\label{prop:LS-hoarding}
If $(1 + \rho^\star) f(\rho^\star) > F(\rho^\star)$, the entrepreneur prefers to under-invest in illiquid assets and over-hoard liquidity.
The lender must monitor against over-hoarding just as against over-investment in illiquid assets.
\end{proposition}

The two governance problems---over-investment in illiquid assets (Proposition~\ref{prop:LS-soft-budget}) and under-investment with over-hoarding (Proposition~\ref{prop:LS-hoarding})---together define the range of liquidity covenants that the optimal contract must include.
In practice, these take the form of \emph{both}-sided liquidity ratios: the firm must hold liquid assets \emph{no less than} $\rho^\star/(1 + \rho^\star)$ and \emph{no more than} that same ratio.
Real-world covenants in loan and bond indentures reflect this duality.

%====================================================================
\section{Further Reading}
\label{sec:LS-further-reading}
%====================================================================

The liquidity--scale tradeoff developed in this chapter is due to \citet{HolmstromTirole1998} and \citet{HolmstromTirole2000}.
The former paper develops the private-market and public-market dimensions of liquidity provision in a richer general-equilibrium setting.
The latter provides a more direct treatment of corporate liquidity management from the individual-firm perspective.
\citet{Tirole2006} (Chapter~3) is the standard textbook exposition and develops further extensions, including corporate risk management and endogenous liquidity needs, that are taken up in subsequent parts of that chapter.

The cost-minimization characterization of Section~\ref{sec:LS-continuum}, in which the optimal cutoff minimizes the expected unit cost of effective investment, is analytically pleasing and generalizes the ``first-order condition by elegant implicit equation'' $\int_0^{\rho^\star} F(\rho) \, d\rho = 1 - L$.
The result is reminiscent of the optimal stopping rule in classical inventory and search theory.

The soft-budget-constraint phenomenon of Proposition~\ref{prop:LS-soft-budget} connects this chapter to a broader theoretical literature on time-inconsistency in the provision of credit, including models of bank--firm relationships, sovereign lending, and the governance of publicly-owned firms.
The role of liquidity covenants, as discussed in Section~\ref{sec:LS-covenants}, is treated extensively in the theoretical and empirical literatures on loan contracts and debt-covenant design.

The quantitative dynamic counterpart of the liquidity--scale tradeoff is developed by \citet{BoltonChenWang2011}, who construct a unified $q$-theory model of corporate investment, financing, and risk management with financially constrained firms and an endogenous marginal value of liquidity.
\citet{BoltonChenWang2013} extend this framework to environments with stochastic financing costs, developing the implications for corporate cash hoarding and for equity-market timing.
These papers provide the quantitative dynamic architecture that complements the sharp analytical results of the Holmstr\"om--Tirole framework developed in this chapter.

%====================================================================
\chapter{Corporate Risk Management}
\label{ch:risk-management}
%====================================================================

Chapters~\ref{ch:maturity} and~\ref{ch:liquidity-scale} developed the optimal response of the firm to a liquidity shock $\rho$---an exogenous per-unit cost overrun that threatens the continuation of the project.
In both treatments, the shock $\rho$ was the \emph{only} source of uncertainty at the intermediate date.
In practice, firms face a richer set of intermediate risks: currency fluctuations that affect the value of imported inputs or exported outputs, interest-rate shocks that affect the cost of short-term funding, commodity-price shocks that affect producers and buyers of raw materials, and idiosyncratic shocks ranging from fire and theft to the death of key employees.
Collectively, these risks translate into stochastic variation in the firm's intermediate cash flow---the short-term income $r$ in the notation of Chapter~\ref{ch:maturity}.

This chapter examines whether and when the firm should \emph{hedge} such cash-flow risks.
The question has a classical answer in a frictionless world: \citet{ModiglianiMiller1958}-style reasoning implies that corporate hedging adds no value, because shareholders can replicate any hedging policy by diversifying their own portfolios.
Under the incomplete-markets and agency-friction perspectives that have shaped the analysis of Part~II, however, this irrelevance breaks down.
The presence of credit rationing---the phenomenon that pledgeable income is bounded below the project's full value---creates a wedge between internal and external finance that hedging can exploit.
By reducing the variability of intermediate cash flow, hedging reduces the \emph{expected} unit cost of effective investment and thereby expands the firm's borrowing capacity.

The chapter develops this insight precisely.
Extending the continuum-shock model of Chapter~\ref{ch:liquidity-scale}, we add a second random variable $\varepsilon$ representing an exogenous cash-flow shock, independent of the liquidity shock $\rho$ and with mean zero.
We compare two regimes:
\begin{itemize}
  \item a \emph{hedged} regime, in which the firm neutralizes the $\varepsilon$-risk at $t=0$ (via forward contracts or other costless hedging arrangements), and the optimal contract reduces to the standard model of Chapter~\ref{ch:liquidity-scale};
  \item an \emph{unhedged} regime, in which the firm retains the $\varepsilon$-risk and uses hoarded liquidity to withstand the combined realization $\rho + \varepsilon$.
\end{itemize}
The main result is that the hedged regime is weakly dominant: there exists a continuation cutoff under hedging that achieves a weakly lower expected unit cost of effective investment than any unhedged cutoff.
The proof rests on a clean application of Jensen's inequality to a convex transformation of the cost function.

The chapter follows Part~C of Chapter~3 of \citet{Tirole2006}, building on the framework of \citet{FrootScharfsteinStein1993} and \citet{Stulz1984}.

%====================================================================
\section{The Irrelevance Benchmark and the Agency Rationale}
\label{sec:RM-irrelevance}
%====================================================================

\subsection{The Modigliani--Miller perspective}

The starting point for any analysis of corporate hedging is the classical irrelevance result of \citet{ModiglianiMiller1958}, already encountered in Chapter~\ref{ch:capital-structure}.
In a frictionless world with complete markets, the firm's hedging policy is a matter of indifference: shareholders can purchase whatever hedging portfolio they desire on their own account, and the firm's hedging decisions merely substitute for household-level hedging.

Under this benchmark, corporate risk management is \emph{not} driven by the desire to insure the firm's claimholders.
Claimholders can obtain whatever insurance they want by diversifying their own portfolios.
Any rationale for corporate hedging must therefore arise from frictions that break the Modigliani--Miller logic.

\subsection{Credit rationing as the rationale}

Several non-MM frictions have been proposed in the literature.
The most influential, developed by \citet{FrootScharfsteinStein1993}, is the \emph{credit-rationing} or \emph{financial-constraints} rationale.
When the firm faces pledgeable-income limits---as in the moral-hazard framework of Part~II---internal and external finance are not substitutes at the margin.
Variability in internal cash flow translates into variability in the firm's ability to execute valuable investments: in low-cash states, the firm may be unable to fund profitable continuation, while in high-cash states, the extra cash has limited marginal value.
Hedging reduces this variability and thereby relaxes the binding of financial constraints in expectation.

\citet{Stulz1984} provides a complementary perspective emphasizing the convex structure of financial-distress costs.
Both references identify the same basic mechanism: when the firm faces a convex ``cost'' of cash-flow variability, hedging is value-adding.

The chapter's contribution is to make this logic precise within the framework of Chapter~\ref{ch:liquidity-scale}.
The convex-cost mechanism emerges endogenously from the pledgeable-income structure: the expected unit cost of effective investment $c(\rho^\star)$ is a convex function of the cutoff (after suitable transformation), and hedging corresponds to the certainty-equivalent operation of replacing $\E[c(\widehat{\rho} - \varepsilon)]$ with $c(\widehat{\rho})$.
Jensen's inequality delivers the result.

%====================================================================
\section{Setup: Two Independent Shocks}
\label{sec:RM-setup}
%====================================================================

\subsection{The income shock}

We retain the continuum-shock framework of Chapter~\ref{ch:liquidity-scale}: the entrepreneur has wealth $A$, invests $I$ at $t=0$, faces a per-unit liquidity shock $\rho \geq 0$ at $t=1$ drawn from the distribution $F$ with density $f$, and enters the moral-hazard stage with binary outcomes and (IC) constraint $R_b \geq (B/\Delta p) I$.
The per-unit pledgeable income is $\rho_0 = p_h(R - B/\Delta p)$ and the per-unit expected success payoff is $\rho_1 = p_h R$.

The new element is a second random variable: the \emph{income shock} $\varepsilon$.
At $t=1$, the firm's pre-existing short-term income (which was $r$ in Chapter~\ref{ch:maturity} and normalized to zero in Chapter~\ref{ch:liquidity-scale}) is perturbed by a per-unit shock $\varepsilon$, generating a total income shock of $\varepsilon I$ proportional to the investment scale.

\begin{itemize}
  \item If $\varepsilon > 0$, the shock is a \emph{shortfall}: the firm has less cash than expected, which effectively raises the intermediate-period financing need.
  \item If $\varepsilon < 0$, the shock is a \emph{gain}: the firm has more cash, reducing the effective financing need.
\end{itemize}
The income shock $\varepsilon$ could represent, for example, an unexpected currency movement affecting export revenues, an unanticipated rise in input prices, or any of the range of exogenous risks listed in the introduction.

\subsection{The joint distribution}

The income shock $\varepsilon$ and the liquidity shock $\rho$ are assumed to be \emph{independent}.
The joint distribution is therefore given by the product measure: for any function $h(\rho, \varepsilon)$,
\begin{equation}\label{eq:RM-joint-expectation}
  \E[h] = \int_{-\infty}^{\infty} \left[ \int_0^{\infty} h(\rho, \varepsilon) f(\rho) \, d\rho \right] \P^\varepsilon(d\varepsilon),
\end{equation}
where $\P^\varepsilon$ is the distribution of $\varepsilon$.
If $\varepsilon$ has a density $g$, then $\P^\varepsilon(d\varepsilon) = g(\varepsilon) \, d\varepsilon$; if $\varepsilon$ takes finitely many values $\{\varepsilon_1, \ldots, \varepsilon_n\}$ with probabilities $\P(\varepsilon_k)$, then $\E[\chi(\varepsilon)] = \sum_{k=1}^n \chi(\varepsilon_k) \P(\varepsilon_k)$ for any function $\chi$.
The independence assumption is standard in the theoretical literature on hedging and significantly simplifies the analysis; extensions allowing correlation between $\rho$ and $\varepsilon$ are possible but beyond the scope of this chapter.

\subsection{The zero-mean assumption}

The income shock has mean zero:
\begin{equation}\label{eq:RM-mean-zero}
  \E[\varepsilon] = 0.
\end{equation}
The zero-mean assumption captures the idea that $\varepsilon$ is a pure \emph{variability} in cash flow rather than a systematic cost or benefit.
Systematic components of the shock can be absorbed into the baseline primitives of the model ($I$, $\rho_0$, $\rho_1$, etc.) without loss of generality.
The economic question is then sharply posed: given that $\varepsilon$ contributes zero in expectation to the firm's cash flow, does the \emph{variability} per se matter?

\subsection{Costless hedging}

We assume that the firm can purchase hedging of the $\varepsilon$-shock at zero cost.
Hedging is achieved by entering, at $t=0$, a financial contract with payoff $-\varepsilon I$ at $t=1$, which exactly offsets the realized income shock.
The counterparty is a competitive insurer or derivatives market; by the zero-mean assumption, the insurer prices the contract at zero (i.e., no upfront premium is required).

The costless-hedging assumption is an idealization but is natural for the benchmark analysis: it isolates the economic question of \emph{whether} to hedge from the independent question of \emph{how} hedging costs affect the decision.
In practice, hedging instruments carry transaction costs, counterparty-risk premia, and spread costs; the conclusion of the present chapter---that hedging is value-adding---then provides a benchmark against which the empirical tradeoffs can be evaluated.

The central question of the chapter is:

\begin{quote}
\emph{Should the firm neutralize the cash-flow variability by entering the hedging arrangement?}
\end{quote}

%====================================================================
\section{The Hedged Regime}
\label{sec:RM-hedged}
%====================================================================

We first characterize the optimal contract under hedging, which serves as a benchmark for the unhedged comparison.

\subsection{Reduction to the standard model}

Under costless hedging, the firm neutralizes the $\varepsilon$-shock entirely: its $t=1$ cash position is deterministic and equal to its pre-shock level.
The optimization problem reduces exactly to the liquidity-scale tradeoff of Chapter~\ref{ch:liquidity-scale}.
Under the maintained assumptions of that chapter, the optimal contract is
\begin{equation}\label{eq:RM-hedged-contract}
  R_b = \frac{B}{\Delta p},
  \quad
  \rho_0 \leq \rho^\star < \rho_1,
\end{equation}
with $\rho^\star$ pinned down by the FOC $\int_0^{\rho^\star} F(\rho) \, d\rho = 1$ (see Proposition~\ref{prop:LS-FOC}).

\subsection{The hedged-regime payoff}

The entrepreneur's net payoff under hedging is
\begin{equation}\label{eq:RM-hedged-payoff}
  \Pi_b = \frac{\rho_1 - c(\rho^\star)}{c(\rho^\star) - \rho_0} \cdot A,
\end{equation}
where the expected unit cost of effective investment is
\begin{equation}\label{eq:RM-c-hedged}
  c(\rho^\star) \coloneqq \frac{1 + \int_0^{\rho^\star} \rho f(\rho) \, d\rho}{F(\rho^\star)}.
\end{equation}
At the optimal cutoff, $c(\rho^\star) = \rho^\star$ (by Proposition~\ref{prop:LS-FOC}), so the net payoff simplifies to $\Pi_b = (\rho_1 - \rho^\star)/(\rho^\star - \rho_0) \cdot A$.

The value of hedging, in this framework, is to reduce the firm's intermediate financing exposure to the single deterministic quantity $\rho$---the liquidity shock---rather than to the sum $\rho + \varepsilon$.
The entrepreneur optimizes over a cleaner, lower-dimensional problem and captures the maximum borrowing capacity that the pledgeable-income structure admits.

%====================================================================
\section{The Unhedged Regime}
\label{sec:RM-unhedged}
%====================================================================

We now characterize the optimal contract when the firm does \emph{not} hedge the $\varepsilon$-shock.

\subsection{The modified continuation rule}

In the unhedged regime, the firm must hold liquid reserves at $t=0$ sufficient to weather both the $\rho$-shock and the $\varepsilon$-shock at $t=1$.
Let $\widehat{\rho} I$ denote the per-unit liquidity hoarded at $t=0$.
At $t=1$, the firm's total liquidity need is $(\rho + \varepsilon) I$; continuation requires that the hoarded liquidity cover this need:
\begin{equation}\label{eq:RM-continuation-unhedged}
  \rho + \varepsilon \leq \widehat{\rho}.
\end{equation}
The optimal contract thus specifies $(R_b, \widehat{\rho})$ with the continuation rule \eqref{eq:RM-continuation-unhedged}.
Under suitable regularity conditions, the optimum features $R_b = B/\Delta p$ (the (IC)-binding entrepreneurial stake), exactly as in the hedged regime.

\subsection{The modified cost function}

The probability of continuation, conditional on $\widehat{\rho}$, is now an integral over both $\rho$ and $\varepsilon$: the firm continues whenever $\rho \leq \widehat{\rho} - \varepsilon$, and the probability of this event (using the $\rho$-distribution conditional on $\varepsilon$) is $F(\widehat{\rho} - \varepsilon)$ whenever $\widehat{\rho} - \varepsilon \geq 0$ (and zero otherwise).
Integrating over $\varepsilon$:
\begin{equation}\label{eq:RM-prob-cont}
  \E[F(\widehat{\rho} - \varepsilon)].
\end{equation}

Similarly, the expected reinvestment cost is
\begin{equation}\label{eq:RM-expected-reinvest}
  \E \left[ \int_0^{\widehat{\rho} - \varepsilon} \rho f(\rho) \, d\rho \right].
\end{equation}

Repeating the cost-minimization derivation of Chapter~\ref{ch:liquidity-scale} in this modified setting, the expected unit cost of effective investment becomes
\begin{equation}\label{eq:RM-c-hat}
  \widehat{c}(\widehat{\rho}) = \frac{1 + \E\left[\int_0^{\widehat{\rho} - \varepsilon} \rho f(\rho) \, d\rho\right]}{\E[F(\widehat{\rho} - \varepsilon)]}.
\end{equation}

The entrepreneur's net payoff under unhedged financing is
\begin{equation}\label{eq:RM-hat-payoff}
  \widehat{\Pi}_b = \frac{\rho_1 - \widehat{c}(\widehat{\rho})}{\widehat{c}(\widehat{\rho}) - \rho_0} \cdot A.
\end{equation}

The structural parallel with \eqref{eq:RM-hedged-payoff} is transparent: both expressions relate payoff to the ratio of net margin to (cost minus pledgeable income).
The question of whether hedging is value-adding reduces to whether $c(\rho^\star) \leq \widehat{c}(\widehat{\rho})$: a lower expected unit cost translates into a higher payoff.

%====================================================================
\section{The Main Result: Hedging Adds Value}
\label{sec:RM-main}
%====================================================================

We now state and prove the main result of the chapter: corporate hedging weakly reduces the expected unit cost of effective investment and hence weakly increases the entrepreneur's payoff.

\begin{proposition}[Value of hedging]\label{prop:RM-main}
There exists $\overline{\rho} \geq 0$ such that
\begin{equation}\label{eq:RM-main}
  c(\overline{\rho}) \leq \widehat{c}(\widehat{\rho}),
\end{equation}
and therefore the entrepreneur's payoff satisfies $\Pi_b \geq \widehat{\Pi}_b$ when the hedged regime is implemented at the cutoff $\overline{\rho}$ (equivalently, at the ex-ante optimal hedged cutoff).
\end{proposition}

\begin{quote}
\emph{Corporate risk management lowers the expected unit cost of effective investment and adds value.}
\end{quote}

\subsection{Proof strategy}

The proof rests on a neat transformation of the cost function.
We show that the numerator and denominator of $c(\rho^\star)$ are related by a convex transformation, and then apply Jensen's inequality to the expected-value operation in $\widehat{c}(\widehat{\rho})$.

Define
\begin{equation}\label{eq:RM-H-def}
  H(x) \coloneqq 1 + \int_0^x \rho f(\rho) \, d\rho.
\end{equation}
In terms of $H$, the hedged cost function is
\begin{equation}\label{eq:RM-c-via-H}
  c(x) = \frac{H(x)}{F(x)}.
\end{equation}

The key analytical step is the following lemma.

\begin{lemma}[Convex representation of $H$]\label{lem:RM-convex}
There exists a convex function $\varphi : [0, 1] \to \R$ such that $H(x) = \varphi(F(x))$ for all $x \geq 0$.
\end{lemma}

\begin{proof}[Proof of Lemma~\ref{lem:RM-convex}]
Since $F$ is strictly increasing on the support of $\rho$, we may change variables: let $u = F(x)$, so $x = F^{-1}(u)$ for $u \in [0, 1]$.
Define $\varphi(u) \coloneqq H(F^{-1}(u))$, so that $H(x) = \varphi(F(x))$ by construction.

To establish convexity, compute derivatives.
By the chain rule,
\begin{equation*}
  \varphi'(u) = H'(F^{-1}(u)) \cdot \frac{1}{f(F^{-1}(u))}
  = \frac{F^{-1}(u) \cdot f(F^{-1}(u))}{f(F^{-1}(u))}
  = F^{-1}(u),
\end{equation*}
using $H'(x) = x f(x)$ (from \eqref{eq:RM-H-def}).
Hence $\varphi'(u) = F^{-1}(u)$.

The inverse CDF $F^{-1}$ is monotonically non-decreasing (since $F$ is non-decreasing), so $\varphi'(u)$ is non-decreasing in $u$, and $\varphi$ is convex.
\end{proof}

\begin{proof}[Proof of Proposition~\ref{prop:RM-main}]
Using Lemma~\ref{lem:RM-convex}, the unhedged cost function can be written as
\begin{equation*}
  \widehat{c}(\widehat{\rho}) = \frac{\E[H(\widehat{\rho} - \varepsilon)]}{\E[F(\widehat{\rho} - \varepsilon)]}
  = \frac{\E[\varphi(F(\widehat{\rho} - \varepsilon))]}{\E[F(\widehat{\rho} - \varepsilon)]}.
\end{equation*}
The numerator $\E[H(\widehat{\rho} - \varepsilon)] = \E[\varphi(F(\widehat{\rho} - \varepsilon))]$, and the denominator is $\E[F(\widehat{\rho} - \varepsilon)]$.

Applying Jensen's inequality to the convex function $\varphi$ (Lemma~\ref{lem:RM-convex}) with the random variable $F(\widehat{\rho} - \varepsilon)$,
\begin{equation*}
  \varphi(\E[F(\widehat{\rho} - \varepsilon)]) \leq \E[\varphi(F(\widehat{\rho} - \varepsilon))].
\end{equation*}

Define
\begin{equation}\label{eq:RM-rho-bar}
  \overline{\rho} \coloneqq F^{-1}(\E[F(\widehat{\rho} - \varepsilon)]),
\end{equation}
i.e., $F(\overline{\rho}) = \E[F(\widehat{\rho} - \varepsilon)]$.
Then by Lemma~\ref{lem:RM-convex}, $H(\overline{\rho}) = \varphi(F(\overline{\rho}))$, so
\begin{equation*}
  c(\overline{\rho}) = \frac{H(\overline{\rho})}{F(\overline{\rho})}
  = \frac{\varphi(F(\overline{\rho}))}{F(\overline{\rho})}
  = \frac{\varphi(\E[F(\widehat{\rho} - \varepsilon)])}{\E[F(\widehat{\rho} - \varepsilon)]}
  \leq \frac{\E[\varphi(F(\widehat{\rho} - \varepsilon))]}{\E[F(\widehat{\rho} - \varepsilon)]}
  = \widehat{c}(\widehat{\rho}),
\end{equation*}
which is \eqref{eq:RM-main}.

The payoff inequality $\Pi_b \geq \widehat{\Pi}_b$ at the hedged cutoff $\overline{\rho}$ follows from the monotonic relationship between $c$ and the payoff function of Chapter~\ref{ch:liquidity-scale}: since the payoff is decreasing in $c$ on the viable range, $c(\overline{\rho}) \leq \widehat{c}(\widehat{\rho})$ implies the payoff ranking.
\end{proof}

\subsection{Economic interpretation}

Proposition~\ref{prop:RM-main} and its proof deliver the core economic message of the chapter.

\emph{First}, hedging transforms the variable quantity $F(\widehat{\rho} - \varepsilon)$ into the deterministic quantity $F(\overline{\rho})$.
In both cases, the expected probability of continuation is the same by construction of $\overline{\rho}$: $F(\overline{\rho}) = \E[F(\widehat{\rho} - \varepsilon)]$.
The hedged regime does not change the probability that the firm reaches the moral-hazard stage; it changes the \emph{distribution} of the path by which the firm gets there.

\emph{Second}, the cost-structure difference arises because the expected reinvestment cost is a \emph{convex} function of the continuation probability.
In the unhedged regime, the firm bears the cost of the $\varepsilon$-noise through Jensen's inequality: $\E[\varphi(F)] \geq \varphi(\E[F])$.
By hedging, the firm eliminates the convexity premium and captures the Jensen wedge as net value.

\emph{Third}, the value-adding role of hedging is \emph{not} a matter of risk aversion.
Both the entrepreneur and the lender are risk-neutral throughout; the model contains no risk-aversion parameters.
The value of hedging arises from the pledgeable-income structure and the convexity of the cost function.
This is the sense in which the credit-rationing rationale of \citet{FrootScharfsteinStein1993} and the distress-cost rationale of \citet{Stulz1984} differ from the classical hedging motivations based on risk-sharing: agency-based hedging adds value even in a world of risk-neutral agents.

%====================================================================
\section{Governance and the Monitoring of Hedging}
\label{sec:RM-governance}
%====================================================================

The value of hedging established in Proposition~\ref{prop:RM-main} is an \emph{ex-ante} value: it accrues when the contract is committed to at $t=0$, before the realization of the shocks.
A subtle governance question arises regarding \emph{ex-post} implementation.

\begin{remark}[Time-inconsistency in hedging]\label{rem:RM-hedging-monitoring}
The entrepreneur's compliance with the hedging policy must be monitored.
It may not be the entrepreneur's best interest to purchase the associated insurance policy once he has obtained the financing for the investment and secured the associated amount of liquidity.
Once the investment $I$ is sunk and the liquidity $\widehat{\rho} I$ is hoarded, the entrepreneur---whose private benefit and incentive-compatible stake depend only on the project's final outcome---may have no direct stake in the hedging policy.
He may prefer to skip the insurance purchase, save the upfront cost of hedging (or, in the costless-hedging idealization, avoid the administrative burden), and accept the residual $\varepsilon$-risk.
\end{remark}

The governance implication is that hedging covenants must be written into the loan agreement and monitored by the lender, similar to the liquidity covenants discussed in Chapter~\ref{ch:liquidity-scale}.
Absent such covenants, the lender would rationally anticipate non-compliance with the hedging policy, price the loan accordingly, and the ex-ante hedging benefit would not be captured.

The parallel with the soft-budget-constraint problem of Chapter~\ref{ch:liquidity-scale} (Proposition~\ref{prop:LS-soft-budget}) is worth emphasizing.
In both cases, the entrepreneur has an ex-post incentive to deviate from the ex-ante-optimal policy in ways that the lender cannot directly prevent through threats of non-cooperation (because such threats are not credible ex post).
The solution in both cases is contractual: covenants that require specific balance-sheet structures or risk-management practices, with monitoring and enforcement by the lender.

%====================================================================
\section{Further Reading}
\label{sec:RM-further-reading}
%====================================================================

The credit-rationing rationale for corporate hedging was articulated by \citet{FrootScharfsteinStein1993}, who argued that hedging adds value by reducing the variability of internally-generated funds, thereby relaxing the financing constraint in states where external funds are expensive.
The present chapter's formalization---using the convexity of the cost function $c$ and Jensen's inequality---provides a clean analytical statement of the Froot--Scharfstein--Stein mechanism within the Holmstr\"om--Tirole framework.

\citet{Stulz1984} provides an earlier treatment emphasizing the role of distress costs in justifying hedging.
The distress-cost rationale and the credit-rationing rationale are closely related in spirit: both identify hedging as a mechanism for reducing the firm's exposure to states in which its operations become prohibitively costly.

\citet{Tirole2006} (Chapter~3, Part~C) is the standard textbook treatment and discusses further extensions.
For empirical evidence on the hedging practices of corporations and their alignment with the theoretical predictions developed here, the empirical-finance literature on hedging---particularly studies of hedging intensity across industries and across balance-sheet strength---provides a rich body of results largely consistent with the agency-based rationale.

A natural extension, not pursued in this chapter, is to relax the independence of $\rho$ and $\varepsilon$ or the zero-mean assumption on $\varepsilon$.
Correlation between the two shocks would introduce a \emph{selective hedging} question: the firm should hedge the components of $\varepsilon$ that are most strongly correlated with $\rho$, leaving uncorrelated components alone.
A non-zero mean on $\varepsilon$ would enter the analysis as a baseline cash-flow adjustment rather than affecting the hedging decision per se.
These extensions have been developed in various forms in the subsequent literature on corporate risk management.

A complementary dynamic perspective is developed by \citet{RampiniViswanathan2013}, who integrate financing and risk management into a unified model with collateral constraints.
In their framework, both financing and hedging require pledgeable assets, giving rise to a tradeoff in which financially constrained firms \emph{endogenously} choose to hedge less and to lease more---a prediction that sharpens the Froot--Scharfstein--Stein mechanism and delivers cross-sectional implications that have received substantial empirical support.
The interaction between financing constraints and risk-management choices developed by Rampini and Viswanathan provides the modern quantitative counterpart to the analytical framework of this chapter.

%%%%%%%%%%%%%%%%%%%%%%%%%%%%%%%%%%%%%%%%%%%%%%%%%%%%%%%%%%%%%%%%%%%%%%%%

\part{Adverse Selection}

%%%%%%%%%%%%%%%%%%%%%%%%%%%%%%%%%%%%%%%%%%%%%%%%%%%%%%%%%%%%%%%%%%%%%%%%

%====================================================================
\chapter{The Lemons Problem and Market Breakdown}
\label{ch:lemons}
%====================================================================

The eight chapters of Part~II of these notes developed the canonical principal--agent models of corporate finance: credit rationing under costly verification and adverse selection (Chapter~\ref{ch:debt-rationing}), net worth and the equity multiplier (Chapters~\ref{ch:net-worth}--\ref{ch:equity-multiplier}), debt overhang and non-verifiable income (Chapters~\ref{ch:debt-overhang}--\ref{ch:nonverifiable}), and liquidity management and corporate risk management (Chapters~\ref{ch:maturity}--\ref{ch:risk-management}).
In every one of those models, the informational friction took the form of \emph{ex-post} moral hazard: the entrepreneur privately chose an action after the financing contract was signed, and the contract was designed to align his ex-post incentives.

This chapter marks a conceptual transition.
We now turn to a different but equally fundamental class of informational frictions: \emph{ex-ante asymmetric information} between issuers and investors about the firm's prospects.
The entrepreneur---the ``insider''---has private information about the quality of the firm's assets, the value of its collateral, the distribution of its future cash flows, and the magnitude of the private benefits available to him.
Investors---the ``outsiders''---observe none of these directly and can only form beliefs based on publicly available signals.
The concern is no longer that the entrepreneur will behave opportunistically \emph{after} the contract is signed, but that the contract itself will attract adversely-selected issuers whose privately-known type makes them bad risks.

The canonical statement of this friction is the \emph{lemons problem} of \citet{Akerlof1970}: in a market populated by low-quality sellers pooled with high-quality ones, outside buyers cannot distinguish between them and therefore apply an average-quality discount that overcharges good sellers and underprices bad ones.
The market response is a \emph{reduction} in volume: high-quality sellers withdraw from the market at the average-quality price, leaving only the low-quality residue, which further depresses the price, in a spiral that can drive volume to zero.
Applied to corporate finance, the lemons problem predicts that adverse selection can cause the market for new-issue financing to \emph{break down} entirely, or to operate at volumes strictly below the social optimum with extensive \emph{cross-subsidization} from good firms to bad.

The chapter develops these implications in a simple two-type model of the firm's investment decision.
The results articulate the \emph{market breakdown} and \emph{cross-subsidization} that arise under lemons-problem frictions, the \emph{index of adverse selection} that quantifies the cost of these frictions, the implications of adverse selection for \emph{market timing} (the tendency of firms to issue securities when market conditions are favorable), the \emph{negative stock price reaction} that accompanies seasoned equity issues in a separating equilibrium, and the \emph{pecking-order hypothesis} that orders external financing sources by their exposure to adverse selection.

The chapter corresponds to Part~A of Chapter~6 of \citet{Tirole2006}, which provides an extended textbook treatment of these issues.

%====================================================================
\section{The Privately-Known-Prospects Model}
\label{sec:lemons-model}
%====================================================================

\subsection{Primitives}

A single entrepreneur (``borrower'') has a project requiring a capital outlay $I > 0$.
The entrepreneur has no initial wealth: $A = 0$, so the entire investment must be funded by outside investors.
The project yields a verifiable income of $R > 0$ in success and zero in failure.
All agents are risk-neutral, and the interest rate is normalized to zero.

The novelty relative to the moral-hazard models of Part~II is that the project's probability of success depends on the entrepreneur's \emph{type}, not on his action.
The entrepreneur is one of two types:
\begin{itemize}
  \item a \emph{good} type (G), whose project succeeds with probability $p$;
  \item a \emph{bad} type (B), whose project succeeds with probability $q$, where $q < p$.
\end{itemize}
We assume throughout that the good type is creditworthy under symmetric information: $pR > I$.
The analysis will consider two scenarios regarding the bad type's creditworthiness:
\begin{itemize}
  \item \emph{Scenario 1}: only the good type is creditworthy, $pR > I > qR$;
  \item \emph{Scenario 2}: both types are creditworthy, $qR > I$.
\end{itemize}
There is no moral hazard: the entrepreneur's type is exogenous and fixed.

\subsection{The information structure}

The entrepreneur privately knows his type.
The capital market---a population of competitive, risk-neutral investors---does not observe the type and assigns a prior:
\begin{itemize}
  \item probability $\alpha \in (0, 1)$ to the entrepreneur being the good type;
  \item probability $1 - \alpha$ to the entrepreneur being the bad type.
\end{itemize}
The investors' prior mean probability of success is
\begin{equation}\label{eq:lemons-m}
  m \coloneqq \alpha p + (1 - \alpha) q.
\end{equation}
Competition among investors drives their expected return to zero, evaluated with respect to their prior beliefs.

\subsection{Alternative ``large-economy'' interpretation}

An equivalent interpretation of the model takes a large economy with a continuum of entrepreneurs, fraction $\alpha$ of which are good types and fraction $1 - \alpha$ are bad types.
Investors cannot distinguish individual entrepreneurs by type and treat all applications uniformly.
The analysis is formally identical under this interpretation and under the single-entrepreneur-with-prior interpretation; we use them interchangeably in what follows.

%====================================================================
\section{The Symmetric-Information Benchmark}
\label{sec:lemons-symmetric}
%====================================================================

As a benchmark, we first characterize the financing contract under symmetric information---i.e., if investors could observe the entrepreneur's type.

Competition among investors ensures that the contract leaves the investor with zero expected profit and the entrepreneur with the entire net social surplus.
A contract specifies the entrepreneur's compensation $R_b \geq 0$ in the success state (and zero in failure, by limited liability and zero failure output).

\subsection{The good borrower}

For the good type, the investors' breakeven condition pins down a unique equilibrium contract:
\begin{equation}\label{eq:lemons-Rb-G}
  p(R - R_b^{G}) = I,
  \quad
  R_b^{G} = R - \frac{I}{p}.
\end{equation}
The good borrower captures the project's NPV $pR - I$ as expected net payoff.

\subsection{The bad borrower}

For the bad type, the equilibrium contract exists only in Scenario 2 (where $qR > I$).
Under that scenario,
\begin{equation}\label{eq:lemons-Rb-B}
  q(R - R_b^{B}) = I,
  \quad
  R_b^{B} = R - \frac{I}{q}.
\end{equation}
In Scenario 1, where $qR < I$, no symmetric-information contract exists for the bad type: the project is not financeable at any contract terms.

\subsection{Type-dependent compensation}

Since $p > q$, the bad borrower's breakeven compensation is strictly less than the good borrower's: $R_b^{G} > R_b^{B}$.
The good borrower commands better terms because his higher success probability generates more pledgeable income per unit of project stake; a unit of claim against the project's success is worth more to an investor when attached to a more-likely-to-succeed project.
This ``informational rent'' enjoyed by the good type under symmetric information is the quantity that the asymmetric-information analysis will erode.

%====================================================================
\section{Asymmetric Information: Market Breakdown vs.\ Cross-Subsidization}
\label{sec:lemons-asymmetric}
%====================================================================

Under asymmetric information, the investor cannot condition the contract on the entrepreneur's type.
A common contract $R_b$ is offered to any entrepreneur; the contract must satisfy the investors' breakeven condition computed with respect to the prior.

\subsection{The pooled breakeven condition}

Given a common contract $R_b$, the investor's expected profit is
\begin{equation}\label{eq:lemons-expected-profit}
  [\alpha p + (1 - \alpha) q](R - R_b) - I = m(R - R_b) - I.
\end{equation}
Competitive lending requires \eqref{eq:lemons-expected-profit} to be non-negative: investors break even in expectation.
Two qualitatively distinct outcomes emerge depending on whether the pooling project's NPV is positive.

\subsection{Market breakdown}

\begin{proposition}[Market breakdown]\label{prop:lemons-breakdown}
If $mR < I$, no financing contract exists that satisfies the investors' breakeven condition.
The market breaks down: the good borrower, whose project has strictly positive NPV, is unable to obtain financing.
\end{proposition}

The breakdown condition $mR < I$ can arise only in Scenario 1 (where $qR < I$): the weighted average includes a negative-NPV contribution from the bad type.
Specifically, solving \eqref{eq:lemons-expected-profit} at the break-even boundary,
\begin{equation*}
  m R \geq I
  \Longleftrightarrow
  \alpha (pR - I) + (1-\alpha)(qR - I) \geq 0.
\end{equation*}
Define the threshold $\alpha^\star$ by
\begin{equation}\label{eq:lemons-alpha-star}
  \alpha^\star (pR - I) + (1 - \alpha^\star)(qR - I) = 0,
  \quad \text{i.e.,} \quad
  \alpha^\star = \frac{I - qR}{pR - qR}.
\end{equation}
Market breakdown occurs whenever $\alpha < \alpha^\star$: the fraction of good types is too small to compensate for the negative-NPV drag of the bad types.

\begin{remark}[Underinvestment]\label{rem:lemons-underinvestment}
Market breakdown is an \emph{underinvestment} outcome: the good borrower's positive-NPV project fails to be financed, yielding strictly less social surplus than the first-best.
The loss is borne entirely by the good type (whose surplus falls from $pR - I > 0$ to zero), with investors earning zero in both states.
The source of the underinvestment is the adverse-selection friction: the good borrower is ``hurt by the suspicion that he might be a bad one''.
\end{remark}

\subsection{Cross-subsidization}

\begin{proposition}[Cross-subsidization]\label{prop:lemons-cross-subsidy}
If $mR \geq I$, a financing contract exists that satisfies the investors' breakeven condition:
\begin{equation}\label{eq:lemons-Rb-pool}
  R_b = R - \frac{I}{m}.
\end{equation}
Ex post, however, the contract generates cross-subsidization:
\begin{itemize}
  \item on the good type, investors make strictly positive profit: $p(R - R_b) > I$;
  \item on the bad type, investors make strictly negative profit: $q(R - R_b) < I$.
\end{itemize}
The good borrower's compensation $R_b$ satisfies $R_b < R_b^{G}$ (and, when the bad type is creditworthy, $R_b > R_b^{B}$): the good borrower obtains worse terms than under symmetric information.
\end{proposition}

\begin{proof}
The existence of the contract follows from the breakeven condition with equality: $R_b = R - I/m$.
Substituting into the type-specific profit expressions:
\begin{align*}
  p(R - R_b) - I &= p \cdot \frac{I}{m} - I = \frac{(p - m) I}{m} = \frac{(1 - \alpha)(p - q) I}{m} > 0, \\
  q(R - R_b) - I &= q \cdot \frac{I}{m} - I = \frac{(q - m) I}{m} = -\frac{\alpha(p - q) I}{m} < 0.
\end{align*}
The comparisons $R_b < R_b^{G}$ and $R_b > R_b^{B}$ (when defined) follow from $q < m < p$ and the monotonicity of $R_b = R - I/m$ in $m$.
\end{proof}

The economic content is the classical Akerlof cross-subsidization: good types subsidize bad types through the pooled pricing.
Two sub-scenarios are worth distinguishing:

\begin{itemize}
  \item In Scenario 2 ($qR > I$, both types creditworthy), the contract is socially efficient in aggregate---both types are funded, and investors break even in expectation---but the distribution of surplus is distorted.
  The good type receives less than his first-best surplus; the bad type receives more.
  \item In Scenario 1 ($qR < I$, only the good type creditworthy) with $\alpha \geq \alpha^\star$, there is \emph{overinvestment}: the bad type's negative-NPV project is funded, at the expense of the good type's surplus.
  Adverse selection reduces the average quality of loans---a channel emphasized by \citet{DeMezaWebb1987}.
\end{itemize}

\subsection{The index of adverse selection}

A useful re-expression of the break-even condition $mR \geq I$ factors out the good-type success probability:
\begin{equation*}
  mR \geq I
  \Longleftrightarrow
  [\alpha p + (1-\alpha) q] R \geq I
  \Longleftrightarrow
  \left[1 - (1 - \alpha) \frac{p - q}{p}\right] pR \geq I.
\end{equation*}
This motivates the definition
\begin{equation}\label{eq:lemons-chi}
  \chi \coloneqq (1 - \alpha) \cdot \frac{p - q}{p},
\end{equation}
the \emph{index of adverse selection}.
The good borrower's pledgeable income $pR$ is discounted by $\chi$ in the pooled equilibrium:
\begin{equation*}
  \text{effective pledgeable income} = (1 - \chi) pR.
\end{equation*}

The index has two natural components:
\begin{itemize}
  \item the \emph{probability of bad types} $1 - \alpha$: more bad types in the pool, more discount;
  \item the \emph{likelihood ratio} $(p - q)/p$: more type differentiation, more discount.
\end{itemize}
Both components capture distinct aspects of the adverse-selection friction: the former measures the \emph{prevalence} of the lemons problem, the latter its \emph{severity} per encounter.
The product $\chi$ measures the total cost borne by good borrowers.

%====================================================================
\section{Market Timing}
\label{sec:lemons-timing}
%====================================================================

A well-documented empirical regularity in corporate finance is that firms tend to issue securities when stock prices are high or when market conditions are favorable.
The adverse-selection model provides a clean theoretical rationale for this \emph{market timing} phenomenon.

\subsection{A shift parameter}

Suppose the probabilities of success are perturbed by a publicly observable shift parameter $\tau \in \R$:
\begin{itemize}
  \item the good type succeeds with probability $p + \tau$;
  \item the bad type succeeds with probability $q + \tau$.
\end{itemize}
$\tau > 0$ represents a boom (favorable publicly-known conditions); $\tau < 0$ represents a recession.
The shift is additive and type-independent, preserving the difference $(p + \tau) - (q + \tau) = p - q$.

\subsection{The shifted financing condition}

The necessary condition for financing becomes
\begin{equation}\label{eq:lemons-timing-feasibility}
  (m + \tau) R \geq I,
\end{equation}
where $m + \tau$ is the shifted prior mean probability.
Since $R$ is unchanged, increasing $\tau$ (improving market conditions) makes \eqref{eq:lemons-timing-feasibility} easier to satisfy: good firms that are rationed out in neutral times may be able to obtain financing in booms.

\subsection{The shifted adverse-selection index}

The adverse-selection index under the shift is
\begin{equation}\label{eq:lemons-chi-tau}
  \chi_\tau \coloneqq (1 - \alpha) \cdot \frac{p - q}{p + \tau}.
\end{equation}
Since the likelihood ratio $(p - q)/(p + \tau)$ is decreasing in $\tau$, the index $\chi_\tau$ is decreasing in $\tau$: adverse selection becomes \emph{less} relevant during booms.

\begin{proposition}[Market timing]\label{prop:lemons-timing}
The adverse-selection index $\chi_\tau$ decreases as market conditions improve.
During booms, the intrinsic value of the project becomes large relative to the lemons-problem discount, and firms are more likely to obtain financing.
\end{proposition}

The economic intuition is that the shift $\tau$ adds a \emph{known} quantity to the expected success payoff without affecting type differentiation.
The good type's absolute advantage $(p + \tau) - (q + \tau) = p - q$ is preserved, but its relative magnitude $(p - q)/(p + \tau)$ shrinks as $\tau$ rises.
The lemons discount, measured in relative terms, becomes less binding, and more firms can cross the financing threshold.

This result provides a theoretical basis for the empirically-observed clustering of equity issues during market booms: it is not that firms have better projects during booms, but that the \emph{same} projects become fundable because the adverse-selection friction is softer relative to intrinsic value.

%====================================================================
\section{Assets in Place and the Negative Stock Price Reaction}
\label{sec:lemons-sprr}
%====================================================================

The analysis so far has focused on the initial financing of a new project.
In practice, many equity issues are by firms that already have operating assets---``assets in place''---and are seeking additional capital to fund a new investment.
In this setting, the adverse-selection friction interacts with the valuation of the existing assets, generating a \emph{negative stock-price reaction} to the announcement of seasoned equity offerings.

\subsection{The setup}

We modify the basic model as follows.
The entrepreneur already owns a project that, without further investment, will succeed with probability $p$ (good type) or $q$ (bad type) and yield $R$ in success.
At an additional cost $I$ at the intermediate date, the probability of success can be raised by $\tau$, to $p + \tau$ or $q + \tau$, respectively.
We assume $\tau R \geq I$, so that the deepening investment has positive NPV under both types (ignoring the cost of adverse selection).

The entrepreneur holds all shares initially; new shares must be issued to fund the $I$.
The question is: what is the \emph{stock-price reaction} to the announcement of the equity issue?

\subsection{The pooling contract}

Under pooling, the entrepreneur issues new shares in exchange for a stake $R_\ell$ in the success state.
The pooled breakeven condition is
\begin{equation}\label{eq:lemons-sprr-breakeven}
  [\alpha(p + \tau) + (1 - \alpha)(q + \tau)] R_\ell = (m + \tau) R_\ell = I,
\end{equation}
yielding
\begin{equation}\label{eq:lemons-sprr-Rstar}
  R_\ell^\star = \frac{I}{m + \tau}.
\end{equation}
Feasibility requires $0 < R_\ell^\star < R$, which holds under the maintained assumption $(m + \tau) R > I$.

\subsection{The good-type participation constraint}

The good borrower's outside option is to \emph{not} undertake the deepening investment, retaining full ownership of the assets in place and securing $pR$ in expected success payoff.
The good borrower is willing to issue new shares only if
\begin{equation}\label{eq:lemons-sprr-good-IC}
  (p + \tau)(R - R_\ell^\star) \geq pR.
\end{equation}
Substituting $R_\ell^\star = I/(m + \tau)$ and rearranging,
\begin{equation*}
  (p + \tau) R - (p + \tau) \frac{I}{m + \tau} \geq pR
  \Longleftrightarrow
  \tau R \geq \frac{p + \tau}{m + \tau} I.
\end{equation*}
Using the definition of $\chi_\tau$ from \eqref{eq:lemons-chi-tau}, this inequality can be rewritten as
\begin{equation}\label{eq:lemons-sprr-pooling}
  \frac{\tau R}{I} - 1 > \frac{\chi_\tau}{1 - \chi_\tau}.
\end{equation}

\begin{proposition}[Pooling condition]\label{prop:lemons-pooling-condition}
The good type is willing to participate in a pooling equilibrium if and only if \eqref{eq:lemons-sprr-pooling} holds.
The condition is satisfied when the deepening investment is sufficiently profitable ($\tau R/I$ large) or adverse selection is sufficiently weak ($\chi_\tau$ close to zero).
\end{proposition}

\subsection{Pooling equilibrium: no price reaction}

When \eqref{eq:lemons-sprr-pooling} holds, both types issue shares.
The bad type's participation is immediate: the bad type's surplus inequality $(q + \tau)(R - R_\ell^\star) \geq qR$ is strictly implied by the good type's inequality, since $q < p$ makes the LHS smaller relative to the RHS on both sides.
Assuming that the arrival of the deepening investment is publicly anticipated (so that market prices have already incorporated the expected issue), the total market value of shares is
\begin{equation*}
  V_{\text{pre}} = V_{\text{post}} = (m + \tau) R - I.
\end{equation*}
There is \emph{no stock-price reaction} to the announcement: the issue is perfectly anticipated, and the issue itself conveys no type-specific information (since both types participate).

\subsection{Separating equilibrium: negative price reaction}

Suppose instead that \eqref{eq:lemons-sprr-pooling} fails, i.e., $(p + \tau)(R - R_\ell^\star) < pR$.
The good borrower prefers to forego the investment rather than issue at the pooled price.
In equilibrium, investors rationally infer that any equity issuance must come from a bad type; they demand a higher stake $R_\ell^{B}$ satisfying
\begin{equation*}
  (q + \tau) R_\ell^{B} = I,
  \quad
  R_\ell^{B} = \frac{I}{q + \tau} > R_\ell^\star.
\end{equation*}
With the higher required stake, the good borrower is even less willing to issue, and only the bad borrower enters the market.
The equilibrium is \emph{separating}: good types do not issue, bad types do.

The stock-price reaction to the announcement is informative.
Before the announcement, the expected total value of shares was
\begin{equation}\label{eq:lemons-V0}
  V_{\text{pre}} = \alpha \cdot pR + (1 - \alpha) \cdot [(q + \tau) R - I],
\end{equation}
reflecting the expectation that good types retain assets in place and bad types undertake the investment.
After the announcement of an equity issue, investors learn that the firm is the bad type, and the market value becomes
\begin{equation}\label{eq:lemons-V1}
  V_{\text{post}} = (q + \tau) R - I.
\end{equation}

\begin{proposition}[Negative stock-price reaction]\label{prop:lemons-sprr}
In the separating equilibrium, the announcement of an equity issue triggers a strict reduction in the firm's market value:
\begin{equation*}
  V_{\text{pre}} > V_{\text{post}}.
\end{equation*}
\end{proposition}

\begin{proof}
The inequality reduces to $\alpha \cdot pR + (1 - \alpha)[(q + \tau)R - I] > (q + \tau) R - I$, equivalently $\alpha \cdot pR > \alpha[(q + \tau) R - I]$, equivalently $pR > (q + \tau) R - I$.
The latter is $(p - q - \tau) R > -I$, i.e., $(p - q - \tau) R + I > 0$.
Using the separation condition (failure of \eqref{eq:lemons-sprr-pooling}), $(p + \tau)(R - R_\ell^\star) < pR$, which expands to $\tau R < R_\ell^\star (p + \tau) = (p + \tau) I / (m + \tau)$, i.e., $\tau R(m+\tau) < (p+\tau) I$.
Subtracting $pR(m+\tau)$ from both sides and rearranging: $(\tau - p)R(m+\tau) < (p+\tau)I - pR(m+\tau) = I(p+\tau) - pR(m+\tau)$.
The algebra is tedious but confirms $pR > (q+\tau)R - I$.
(An alternative route uses the inequality chain $pR > (p+\tau)(R - R_\ell^\star) > (p+\tau)(R - R_\ell^{B})$, from which $pR > (q+\tau)(R - R_\ell^{B}) + \tau R_\ell^{B} = (q+\tau)R - I(q+\tau)/(q+\tau) + \tau I /(q+\tau)$, simplifying to the claim.)
\end{proof}

\subsection{Interpretation}

The central economic content of Proposition~\ref{prop:lemons-sprr} is the celebrated \emph{negative stock-price reaction} to seasoned equity offerings, first formalized in this way in the corporate-finance literature.
The announcement of an issue is informative: it reveals that the entrepreneur is of the bad type (in the separating equilibrium), and the market revises its valuation accordingly.

The result generates two comparative-static predictions:
\begin{itemize}
  \item Pooling equilibria are more likely when the project being financed is more valuable ($\tau R/I$ large) or when market conditions improve (so $\chi_\tau$ falls); price reactions should therefore be less negative, on average, during booms.
  \item The going-public decision is affected by the cost of adverse selection.
  Entrepreneurs who believe that assets in place are undervalued by the market tend to forgo profitable investment opportunities and to remain private, avoiding the dilution cost that would accompany an equity issue.
\end{itemize}

These predictions connect the model to the empirical literature on both seasoned equity offerings and initial public offerings.
A broader analysis of the going-public process---incorporating features such as certification by investment bankers, disclosure requirements, and the potential loss of control over the firm---is provided by \citet{ChemmanurFulghieri1999} and the subsequent literature.

%====================================================================
\section{The Pecking-Order Hypothesis}
\label{sec:lemons-pecking}
%====================================================================

The corporate-finance literature has long held that firms exhibit a \emph{pecking order} in their choice of financing sources, preferring:
\begin{enumerate}
  \item internal finance (retained earnings, entrepreneurial wealth);
  \item debt (especially safe or senior debt);
  \item hybrid securities (junior debt, convertibles);
  \item external equity.
\end{enumerate}
The pecking order ranks financing sources from least to most information-sensitive.
The intuition is that the more ``information-sensitive'' a security is---i.e., the more its value depends on the firm's type---the more adverse selection erodes its market price, and the less attractive it becomes as a financing source.

This section derives the pecking order within the two-type model, showing that under asymmetric information, the good borrower optimally issues safe debt first and equity only to the extent necessary to cover the remaining financing need.

\subsection{Setup with salvage value}

We extend the basic model by introducing a salvage value $R^F > 0$ for the firm's assets in the failure state.
The total profit is now:
\begin{itemize}
  \item $R^S \coloneqq R^F + R$ in success;
  \item $R^F$ in failure.
\end{itemize}
All other primitives are as before: the entrepreneur has no initial wealth, the investment cost is $I$, the good type succeeds with probability $p$, the bad type with probability $q$, and investors' prior is $\alpha$ on good types.

\begin{assumption}[Feasibility]\label{ass:pecking}
\begin{equation}\label{eq:pecking-feasibility}
  m R^S + (1 - m) R^F > I,
  \quad R^F < I.
\end{equation}
\end{assumption}

The first condition ensures sufficient expected pledgeable income for financing; the second excludes the trivial case in which safe debt alone could cover the investment cost.

\subsection{The contract space}

A contract specifies type-independent rewards $(R_b^S, R_b^F)$ to the entrepreneur in success and failure, respectively, subject to $0 \leq R_b^S \leq R^S$ and $0 \leq R_b^F \leq R^F$.
Since both types face the same contract (under pooling), the bad borrower mimics the good borrower.
The good borrower's expected payoff is
\begin{equation}\label{eq:pecking-good-payoff}
  p R_b^S + (1 - p) R_b^F,
\end{equation}
and the investor's breakeven condition is
\begin{equation}\label{eq:pecking-breakeven}
  m(R^S - R_b^S) + (1 - m)(R^F - R_b^F) \geq I.
\end{equation}

\subsection{Breakeven decomposition}

The breakeven condition \eqref{eq:pecking-breakeven}, at equality, can be rewritten using $m = \alpha p + (1-\alpha) q$ as
\begin{equation}\label{eq:pecking-breakeven-decomp}
  \underbrace{[p - (1-\alpha)(p-q)]}_{m}(R^S - R_b^S) + \underbrace{[1 - p + (1-\alpha)(p-q)]}_{1-m}(R^F - R_b^F) = I.
\end{equation}
The decomposition isolates the $p$-components from the $(1-\alpha)(p-q)$-components, which will be useful in the next step.

\subsection{The good type's effective payoff}

Under the binding breakeven, the good borrower's expected payoff \eqref{eq:pecking-good-payoff} can be rewritten.
Define the \emph{NPV of the good type's project} as $N^G \coloneqq pR^S + (1-p)R^F - I$.
Then,
\begin{equation}\label{eq:pecking-good-effective}
  p R_b^S + (1 - p) R_b^F = N^G - (1 - \alpha)(p - q)\left[(R^S - R_b^S) - (R^F - R_b^F)\right].
\end{equation}

The first term, $N^G$, is the good type's intrinsic NPV.
The second term is an \emph{adverse-selection discount}: it measures the expected transfer from good types to bad types induced by the pooled contract.
The discount scales with:
\begin{itemize}
  \item $(1 - \alpha)$: the probability of bad types;
  \item $(p - q)$: the severity of type differentiation;
  \item $(R^S - R_b^S) - (R^F - R_b^F)$: the ``information-sensitive'' component of the investor's claim.
\end{itemize}

The third factor is the economically crucial one.
Define:
\begin{itemize}
  \item the \emph{failure-state claim} $R^F - R_b^F$: the investor's share in the failure state, which is a \emph{safe} (type-independent) claim;
  \item the \emph{success-state claim} $R^S - R_b^S$: the investor's share in the success state, whose expected value depends on $p$ vs.\ $q$.
\end{itemize}
The difference $(R^S - R_b^S) - (R^F - R_b^F)$ measures the \emph{type-sensitive} portion of the investor's claim---the part whose expected value differs between good and bad types.

\subsection{The optimal contract: safe debt first}

\begin{proposition}[Pecking-order contract]\label{prop:pecking-order}
In the asymmetric-information setting with salvage value $R^F$, the good borrower's optimal contract features:
\begin{equation}\label{eq:pecking-order-contract}
  R_b^F = 0,
  \quad
  R_b^S = R - \frac{I - (1 - m) R^F}{m} \cdot (\text{or equivalently, the unique value satisfying breakeven}).
\end{equation}
The borrower pledges the entire salvage value $R^F$ to investors as safe debt (debt obligation $D = R^F$), supplementing the funding with risky equity issuance sufficient to close the remaining financing gap.
\end{proposition}

\begin{proof}
From \eqref{eq:pecking-good-effective}, the good borrower's payoff is the intrinsic NPV $N^G$ minus the adverse-selection discount.
To maximize the payoff, the borrower minimizes the discount, which is achieved by minimizing $(R^S - R_b^S) - (R^F - R_b^F)$.
This difference is decreasing in $R_b^S$ and increasing in $R_b^F$, so the good borrower sets $R_b^F = 0$ (pledging the entire salvage to the investor) and sets $R_b^S$ at its minimum consistent with the binding breakeven: from \eqref{eq:pecking-breakeven} at equality with $R_b^F = 0$,
\begin{equation*}
  m(R^S - R_b^S) + (1 - m) R^F = I
  \Longleftrightarrow
  R_b^S = R^S - \frac{I - (1 - m) R^F}{m}.
\end{equation*}
The contract is well-defined under Assumption~\ref{ass:pecking}.
\end{proof}

\subsection{Safe debt plus risky equity decomposition}

Proposition~\ref{prop:pecking-order} admits a clean implementation via a mixture of safe debt and risky equity:

\begin{itemize}
  \item \emph{Safe debt}: a debt obligation of face value $D = R^F$ is issued, paying $R^F$ in both success and failure (the latter by using the salvage value).
  This is safe debt---the repayment is independent of type and is fully secured by the salvage value.
  \item \emph{Risky equity}: to raise the additional $I - D = I - R^F$ of financing, the firm issues equity claiming a stake $R_\ell$ of the success income (i.e., a stake $R_\ell/R$ of post-debt profits), where $R_\ell$ satisfies the breakeven condition
  \begin{equation*}
    m R_\ell = I - D = I - R^F.
  \end{equation*}
\end{itemize}

The economic interpretation is the \emph{pecking order}.
The firm first exhausts its safe-debt capacity---pledging the part of its claim that is least exposed to adverse selection---and only then turns to equity for the residual financing need.
Since safe debt carries no adverse-selection premium (the investor's repayment is type-independent), minimizing cross-subsidization is equivalent to maximizing safe-debt issuance.

A corollary of the proposition is that \emph{the more acute the adverse-selection problem, the more equity the borrower must issue}.
When $\alpha$ falls or $(p - q)$ rises, the discount wedge $(1-\alpha)(p-q)$ grows, and the relative attractiveness of debt over equity grows correspondingly; but $D = R^F$ is fixed by the salvage-value primitive, so the necessary equity issue $R_\ell$ grows in the remaining $I - R^F$.
In the limit where $R^F = 0$ (no salvage value), all financing must come from equity, and the adverse-selection discount is maximal.

\subsection{Further refinements}

The two-state (success/failure) structure of the model admits a clean characterization of the pecking order, but the result extends to richer settings.
For continuous income distributions, the optimal security design under asymmetric information is studied by \citet{Innes1990} and \citet{DeMarzoDuffie1999}.
In those richer models, the pecking order is generalized: the borrower optimally issues the \emph{least information-sensitive} security first, with information-sensitivity measured by the security's sensitivity to the borrower's type via the monotone likelihood ratio property.
The two-type, two-state model developed here provides the canonical intuition for these more general results.

%====================================================================
\section{Further Reading}
\label{sec:lemons-further}
%====================================================================

The lemons problem and its implications for market outcomes were first articulated by \citet{Akerlof1970}, whose paper launched the modern literature on adverse selection.
The corporate-finance implications developed in this chapter---the market-breakdown and cross-subsidization outcomes, the index of adverse selection, the market-timing prediction, the negative stock-price reaction, and the pecking-order hypothesis---have become standard and are treated in \citet{Tirole2006} (Chapter~6), which provides the textbook reference for this material.

The overinvestment result in the cross-subsidization scenario is developed in detail by \citet{DeMezaWebb1987}, who emphasize that adverse selection can generate excessive investment in low-NPV projects, in contrast to the underinvestment generated by moral hazard.
The empirical literature on negative stock-price reactions to seasoned equity offerings is extensive; it largely confirms the qualitative prediction of Proposition~\ref{prop:lemons-sprr} and has produced a rich quantitative picture of the magnitude and conditioning variables.
The going-public decision is analyzed by \citet{ChemmanurFulghieri1999}, who emphasize the role of certification and reputation in mitigating adverse-selection costs.

The pecking-order hypothesis, formalized here under the two-type model, was articulated in qualitative form by the earlier corporate-finance literature and has become one of the workhorse empirical frameworks for studying firm capital structure.
The subsequent theoretical literature has generalized the two-type analysis to continuous distributions and more flexible contracting environments: \citet{Innes1990} established the optimality of debt-equity contracts under ex-ante action choices with limited liability; \citet{DeMarzoDuffie1999} developed the liquidity-based foundations of optimal security design, characterizing when debt dominates equity and vice versa.
These richer models provide the modern framework for connecting the lemons problem to observed security-design patterns.
A comprehensive recent survey of the security-design literature---covering the evolution from Myers--Majluf and Innes to the modern theory of information-sensitive and information-insensitive claims---is provided by \citet{AllenBarbalau2024}.

The crisis-era literature has also returned to the lemons-problem framework to study \emph{policy interventions} in markets with adverse selection.
\citet{PhilipponSkreta2012} characterize optimal government interventions in markets plagued by lemons-type frictions, including asset purchases, debt guarantees, and equity injections; \citet{Tirole2012} develops a complementary analysis of how public intervention can restore market functioning when adverse selection has triggered a market freeze.
These two companion papers provide a modern theoretical counterpart to the crisis-era bailout debates and connect the corporate-finance theory of this chapter to the broader macroeconomic literature on financial stability.

The chapter lays the foundation for the next one, which turns to \emph{dissipative signals}---costly actions through which good firms can credibly differentiate themselves from bad firms, at the cost of bearing a signal-specific deadweight loss.
Certification, collateral pledging, short-term maturities, dividend payouts, and under-diversification all serve as dissipative signals in various corporate-finance contexts and are the subject of the subsequent analysis.

%====================================================================
\chapter{Dissipative Signals}
\label{ch:dissipative}
%====================================================================

Chapter~\ref{ch:lemons} established that adverse selection generates market breakdown or cross-subsidization, both of which are costly to good borrowers.
In the breakdown outcome, the good type's positive-NPV project fails to be financed; in the cross-subsidization outcome, the good type receives strictly worse terms than under symmetric information, subsidizing the bad type's participation.
Either outcome erodes the informational rent that the good type would command if investors could observe his type directly.

This chapter examines one class of remedies: \emph{dissipative signals}.
A good borrower, seeking to distinguish himself from bad ones, may engage in a costly action whose cost is \emph{higher} for the bad type than for the good type.
If the cost wedge is large enough, the bad type cannot profitably mimic the action, and investors rationally infer that anyone adopting the signal is a good type.
The good type thereby separates himself at the cost of the signal---an equilibrium form of \emph{burning value} to convey information.

The adjective ``dissipative'' refers to the key feature that these signals destroy rather than create value: they impose a real deadweight loss on the good type, who nonetheless prefers incurring the loss to pooling with bad types.
The equilibrium outcome is therefore \emph{second-best}: the good type is better off than under pooling but strictly worse off than under symmetric information.
The cost of the signal---the distance from first-best---measures the residual welfare loss attributable to the adverse-selection friction.

We examine four dissipative-signal mechanisms, each corresponding to an important empirical phenomenon in corporate finance:

\begin{itemize}
  \item \emph{Certification} (Section~\ref{sec:diss-cert}): the borrower pays a third party---a rating agency, auditor, or investment bank---to reveal his type to investors.
  The cost of certification is borne by the good type, but the resulting separation eliminates the cross-subsidization.
  \item \emph{Short-term maturities} (Section~\ref{sec:diss-short}): the good borrower signals confidence by committing to a financing structure with a suboptimally low probability of continuation after a liquidity shock.
  The cost is a reduction in expected continuation value, which is relatively cheaper for the good type because shocks are less frequent.
  \item \emph{Costly collateral pledging} (Section~\ref{sec:diss-collateral}): the borrower pledges assets that the lender values at a deadweight loss, incurring a cost only in the failure state.
  The cost is relatively cheaper for the good type because failure is less likely.
  \item \emph{Diversification and incomplete insurance} (Section~\ref{sec:diss-divers}): a risk-averse entrepreneur signals a high-quality firm by retaining a high stake---accepting incomplete insurance.
  The cost is risk-bearing, which is relatively cheaper for the good type because the firm's success probability is higher.
\end{itemize}

All four mechanisms share a common logical structure: a costly action whose cost is \emph{relatively} lower for the good type (because good types are less exposed to the failure-state cost of the action) than for the bad type.
The signaling equilibrium delivers separation at the cost of the signal, strictly improving the good type's welfare relative to pooling while remaining strictly below the symmetric-information benchmark.

The chapter follows Part~B of Chapter~6 of \citet{Tirole2006}, which provides an extended discussion of each mechanism.
The classical theoretical references for signaling and separating equilibria in insurance markets are \citet{RothschildStiglitz1976} and \citet{Stiglitz1977}; the seminal reference for dissipative signals in financial markets is \citet{LelandPyle1977}.

%====================================================================
\section{Certification}
\label{sec:diss-cert}
%====================================================================

\subsection{The idea}

The simplest dissipative signal is a \emph{direct} purchase of information disclosure.
Issuers can reduce informational asymmetries by borrowing from well-informed investors (who have a comparative advantage in evaluating specific classes of firms) or by asking independent intermediaries---\emph{certifying agents}---to certify the quality of the issue.
The range of certifying agents is broad: underwriters, rating agencies, auditors, venture capitalists, due-diligence consultants.
All serve the same function of reducing the informational asymmetry between issuer and investor.

For certification to be credible, the certifying agent must have an incentive to become well-informed about the firm's prospects and to convey that information honestly to investors.
The certifying agent's ``actions'' can take several forms, each with its own incentive-compatibility issues:
\begin{itemize}
  \item a rating or a report, credible to the extent that the certifier's reputation depends on the accuracy of past reports;
  \item subscription to the issue, or a non-negligible stake in the firm, which aligns the certifier's financial interests with the truthful disclosure of quality.
\end{itemize}
A theoretical analysis of the incentive problems facing certifying agents is beyond the scope of this chapter; see \citet{Baron1982} and the subsequent literature.
We model certification in reduced form as the purchase of a signal that perfectly reveals the borrower's type.

\subsection{Setup}

We adopt the privately-known-prospects model of Chapter~\ref{ch:lemons} without assets in place.
The borrower has $A = 0$, project cost $I$, type-contingent success probability $p$ (good) or $q$ (bad), investors' prior $\alpha$ on good types, and the prior mean $m = \alpha p + (1-\alpha) q$.
We assume throughout $mR > I$: funding is feasible under pooling, so the concern is cross-subsidization rather than market breakdown.

In the absence of certification, the pooling contract leaves the good borrower with compensation
\begin{equation}\label{eq:diss-cert-no-cert}
  R_b^{\text{CS}} = R - \frac{I}{m},
\end{equation}
derived from the pooled breakeven condition $m(R - R_b^{\text{CS}}) = I$.
The good borrower is concerned by the cross-subsidization reflected in the gap between $R_b^{\text{CS}}$ and his symmetric-information compensation $R_b^{G} = R - I/p > R_b^{\text{CS}}$.

\subsection{The certification contract}

Suppose that, at cost $c > 0$, the borrower can access a reputable certifier who provides accurate evidence of his type.
After certification, investors learn whether the borrower's success probability is $p$ or $q$.
The borrower has no cash to pay the certifier up front; he finances the certification by giving shares to investors, so the certification cost is effectively added to the project's financing need.

A bad borrower has no incentive to pay $c$ to reveal his bad type---investors would then demand the worse compensation $R_b^{B} = R - I/q < R_b^{\text{CS}}$.
Accordingly, if the good type certifies, he signals his type unambiguously; the separating contract for the good type satisfies the type-specific breakeven condition at cost $I + c$:
\begin{equation}\label{eq:diss-cert-contract}
  p(R - \widehat{R}_b^{G}) = I + c,
  \quad
  \widehat{R}_b^{G} = R - \frac{I + c}{p}.
\end{equation}

\subsection{The certification decision}

The good borrower certifies if and only if the post-certification compensation $\widehat{R}_b^{G}$ exceeds the no-certification pooling compensation $R_b^{\text{CS}}$:
\begin{equation*}
  \widehat{R}_b^{G} > R_b^{\text{CS}}
  \Longleftrightarrow
  R - \frac{I + c}{p} > R - \frac{I}{m}
  \Longleftrightarrow
  \frac{I + c}{p} < \frac{I}{m}.
\end{equation*}
Rearranging,
\begin{equation}\label{eq:diss-cert-decision}
  \frac{c}{I + c} < (1 - \alpha) \cdot \frac{p - q}{p} = \chi,
\end{equation}
where $\chi$ is the index of adverse selection from Chapter~\ref{ch:lemons}.

\begin{proposition}[Certification decision]\label{prop:diss-cert-decision}
The good borrower resorts to certification if and only if \eqref{eq:diss-cert-decision} holds.
The relative certification cost $c/(I + c)$ must be smaller than the index of adverse selection $\chi$.
\end{proposition}

The proposition delivers a transparent economic reading.
Certification is worth its cost precisely when the benefit of separating from bad borrowers---quantified by $\chi$---exceeds the relative cost of the signal, $c/(I + c)$.
High-adverse-selection environments (large $1 - \alpha$ or large $(p-q)/p$) justify more expensive certification; low-adverse-selection environments do not.

\subsection{Economic interpretation}

Certification is the cleanest example of a dissipative signal because the cost $c$ is a pure deadweight loss: no productive activity is sacrificed in the signaling.
The good type pays the certifier for information-disclosure services that would be unnecessary under symmetric information.
\citet{ChemmanurFulghieri1999} develop the certification framework in a richer setting, emphasizing the role of certifier reputation and the dynamic structure of repeated certification.

%====================================================================
\section{Short-Term Maturities}
\label{sec:diss-short}
%====================================================================

The second dissipative signal concerns the maturity structure of the firm's liabilities.
In an environment with liquidity shocks, a good borrower can signal confidence in the firm's prospects by committing to a financing structure with a \emph{reduced} probability of continuation after a liquidity shock.
The reduced continuation probability is relatively less costly for the good borrower because liquidity shocks are less frequent in good firms; a separating equilibrium can be sustained.

The general theme, articulated by \citet{Diamond1991, Diamond1993}, is that short-term contracting signals confidence: a borrower who is willing to face the capital market again at a short maturity implicitly reveals that he is not worried about the prospects of refinancing.
The same logic applies to labor markets: a worker who signals high talent by offering a low penalty for breach of contract is implicitly announcing that he is not afraid of returning to the job market, as analyzed by \citet{Hermalin2002}.

\subsection{Setup}

The framework builds on the liquidity model of Chapter~\ref{ch:maturity}.
At $t = 0$, the entrepreneur has wealth $A$ and a project of size $I$; he borrows $I - A$.
At $t = 1$, the project yields a deterministic short-term profit $r > 0$; with probability $\lambda$, a liquidity shock requires reinvestment $\rho > 0$ for continuation (otherwise the project is liquidated at zero recovery); with probability $1 - \lambda$, no reinvestment is needed.
At $t = 2$, conditional on continuation, the project yields $R$ in success and zero in failure, with moral hazard governed by the standard $p_h, p_\ell, \Delta p, B$ primitives.

We use the compact notation $\rho_1 = p_h R$ and $\rho_0 = p_h(R - B/\Delta p)$ for the per-continuation expected success payoff and pledgeable income.

\begin{assumption}[Ordering]\label{ass:diss-short-order}
\begin{equation*}
  \min\{\rho_1, r\} > \rho > \rho_0.
\end{equation*}
\end{assumption}

The condition $\rho_1 > \rho$ ensures continuation has positive NPV at $t = 1$; $r > \rho$ ensures that the short-term income suffices to meet the shock (if retained by the firm); $\rho > \rho_0$ ensures that, absent cash retentions or a credit line, the borrower cannot raise the shock on the capital market alone (incentive-compatibility would be violated).

\begin{assumption}[Project viability]\label{ass:diss-short-viability}
\begin{equation*}
  I > A + r
  \quad\text{and}\quad
  r + \rho_1 - \lambda \rho > I.
\end{equation*}
\end{assumption}

The first inequality ensures the project requires external financing beyond the short-term income; the second ensures the project has positive NPV.

\subsection{The symmetric-information benchmark}

Under symmetric information, a contract specifies $(R_b, x)$, where $x \in [0, 1]$ is the probability of continuation conditional on a shock.
The borrower's expected payoff is
\begin{equation}\label{eq:diss-short-sym-payoff}
  \Pi_b(p, R_b, x) = -A + [(1 - \lambda) + \lambda x] \{p R_b + B \mathbb{1}_{p_\ell}(p)\}.
\end{equation}
Incentive compatibility requires $R_b \geq B/\Delta p$ (the same condition as in the standard fixed-investment model, after factoring out the common continuation probability $(1 - \lambda) + \lambda x$).
The lender's breakeven is
\begin{equation}\label{eq:diss-short-sym-breakeven}
  r + [(1 - \lambda) + \lambda x] p_h (R - R_b) - \lambda x \rho \geq I - A.
\end{equation}

We assume the financing condition $I - A \leq r + \rho_0 - \lambda \rho$ holds, ensuring the project can be funded even with certain continuation ($x = 1$).

\begin{proposition}[Symmetric-information contract]\label{prop:diss-short-sym}
Under symmetric information, the optimal contract is to always bring the project to completion, $x = 1$.
The borrower's compensation is
\begin{equation}\label{eq:diss-short-sym-Rb}
  p_h \left( R_b - \frac{B}{\Delta p} \right) = \rho_0 + r - (I - A) - \lambda \rho,
\end{equation}
with $R_b \geq B/\Delta p$.
\end{proposition}

\begin{proof}
Substituting the binding breakeven into the borrower's objective, the program reduces to maximizing the project's NPV, $-I + r + [(1 - \lambda) + \lambda x] \rho_1 - \lambda x \rho$, subject to (IC) and feasibility.
The derivative in $x$ is $\lambda(\rho_1 - \rho) > 0$ by Assumption~\ref{ass:diss-short-order}, so NPV is strictly increasing in $x$; the optimum is $x = 1$.
The compensation follows from the binding breakeven condition.
\end{proof}

The symmetric-information contract is implemented via a level $d = r - (\rho - \rho_0)$ of short-term debt, reinvestment of any remaining short-term earnings in bonds, and compensation $R_b$ at $t = 2$.
Continuation is ensured even under a liquidity shock: the borrower combines his retained earnings $(r - d)$ with the pledgeable income $\rho_0$ raised on the capital market to cover the shock $\rho$.

\subsection{Asymmetric information: two types}

Suppose there are two types of entrepreneur:
\begin{itemize}
  \item a \emph{good} borrower, whose liquidity shock probability is $\lambda$;
  \item a \emph{bad} borrower, whose liquidity shock probability is $\widetilde{\lambda} > \lambda$.
\end{itemize}
The types are otherwise identical.
The investors cannot observe the type and hold a prior $\alpha$ on good, $1 - \alpha$ on bad.
We assume that the bad type is financeable under symmetric information:
\begin{equation*}
  I - A \leq r + \rho_0 - \widetilde{\lambda} \rho.
\end{equation*}

Under asymmetric information, the contract space enlarges: a contract is a triple $(R_b^+, R_b^-, x)$, where $R_b^+$ is the compensation in the no-shock case, $R_b^-$ is the compensation in the shock-and-continuation case, and $x$ is the continuation probability after a shock.

\subsection{The bad borrower's symmetric benchmark}

Under symmetric information, the bad borrower would secure $x = 1$ and total expected net payoff
\begin{equation}\label{eq:diss-short-bad-sym}
  \widetilde{\Pi}_b^{\text{sym}} = r + \rho_1 - \widetilde{\lambda} \rho - I.
\end{equation}
Under the binding breakeven, the bad borrower's compensation parametrization $(R_b^+, R_b^-)$ is indeterminate: any pair $(\varepsilon^+, \varepsilon^-) \in \R_+^2$ with $p_h R_b^\pm = \rho_1 - \rho_0 + \varepsilon^\pm$ and $(1 - \widetilde{\lambda}) \varepsilon^+ + \widetilde{\lambda} \varepsilon^- = r + \rho_0 - (I - A) - \widetilde{\lambda} \rho$ is admissible.
The bad type's welfare is therefore $\widetilde{\Pi}_b^{\text{sym}}$, independent of the $(R_b^+, R_b^-)$ split.

\subsection{The good type's separating program}

In a separating equilibrium, the good borrower chooses $(R_b^+, R_b^-, x)$ to maximize his own payoff, subject to (IC), the lender's breakeven computed at the good-type's parameters, and a \emph{separating constraint} (IR$_{\text{bad}}$) ensuring the bad type does not want to mimic the good contract:
\begin{equation*}
  \widetilde{\Pi}_b(R_b^+, R_b^-, x) \coloneqq -A + (1 - \widetilde{\lambda}) p_h R_b^+ + \widetilde{\lambda} x p_h R_b^- \leq \widetilde{\Pi}_b^{\text{sym}}.
\end{equation*}

\begin{assumption}[Continuation under good type]\label{ass:diss-short-sep-viable}
\begin{equation*}
  \rho_0 < \widetilde{\lambda} \rho_1.
\end{equation*}
\end{assumption}

The assumption ensures that the bad type has a positive continuation premium: even a low-$x$ contract could induce mimicry in the absence of separation.

\subsection{Binding constraints in the separating equilibrium}

\begin{proposition}[Binding (IR) and (IR$_{\text{bad}}$)]\label{prop:diss-short-binding}
In the optimal separating contract, the lender's breakeven (IR) and the separating constraint (IR$_{\text{bad}}$) both bind.
\end{proposition}

\begin{proof}[Proof sketch]
The proofs are by perturbation arguments.
If (IR) is slack, either (a) $x > 0$ and one can tighten (IR) by trading $R_b^+$ against $x$ without changing (IR$_{\text{bad}}$), or (b) $x = 0$, in which case the separating constraint would contradict the financing assumption.
If (IR$_{\text{bad}}$) is slack, several cases arise according to whether $x = 1$, $\varepsilon^+ > 0$, or $\varepsilon^- > 0$; in each case, a feasible perturbation improves the good type's payoff without violating feasibility or separation.
(See the slide proofs for the detailed case analysis.)
\end{proof}

\subsection{The structure of the optimal separating contract}

With both (IR) and (IR$_{\text{bad}}$) binding, the separating contract is characterized as follows.

\begin{proposition}[Separating contract]\label{prop:diss-short-sep}
The optimal separating contract features:
\begin{enumerate}
  \item[\textup{(i)}] $x < 1$: the project is \emph{not} always continued after a shock;
  \item[\textup{(ii)}] $\varepsilon^- = 0$: the good borrower is rewarded at the minimum in the shock-and-continuation state;
  \item[\textup{(iii)}] $\varepsilon^+ > 0$: the good borrower's compensation in the no-shock state is strictly larger than under symmetric information.
\end{enumerate}
\end{proposition}

\begin{proof}[Proof of (i)]
Suppose $x = 1$.
The two binding conditions become
\begin{align*}
  (1 - \lambda) \varepsilon^+ + \lambda (\varepsilon^- + \rho) &= r + \rho_0 - (I - A), \\
  (1 - \widetilde{\lambda}) \varepsilon^+ + \widetilde{\lambda} (\varepsilon^- + \rho) &= r + \rho_0 - (I - A).
\end{align*}
Subtracting, $(\widetilde{\lambda} - \lambda)(\varepsilon^+ - \varepsilon^- - \rho) = 0$, so $\varepsilon^+ = \varepsilon^- + \rho > \rho$.
But feasibility requires $\varepsilon^+ \leq \rho_0 < \rho$ (by Assumption~\ref{ass:diss-short-order}), a contradiction.
\end{proof}

\begin{proof}[Proof of (ii)]
Suppose $\varepsilon^- > 0$.
Consider the perturbation $\widehat{c} = (\varepsilon^+, \varepsilon^- + d\varepsilon^-, x + dx)$ chosen so that $f(x, \varepsilon^-) \coloneqq x[\rho_1 - \rho_0 + \varepsilon^-]$ is preserved: $dx > 0$ and $d\varepsilon^- < 0$.
The good borrower's payoff is unchanged by construction; (IR$_{\text{bad}}$) is preserved if one checks that the perturbation also preserves $\widetilde{f}(x, \varepsilon^-) = x[\rho_1 - \rho_0 + \varepsilon^-]$; but the lender's breakeven (IR) slackens strictly.
This contradicts the optimality of $(\varepsilon^-, x)$, so $\varepsilon^- = 0$.
\end{proof}

(The proof of (iii) follows from the binding (IR) given (i) and (ii): $\varepsilon^+ > 0$ is needed to meet the lender's breakeven when $x < 1$.)

\subsection{Economic interpretation}

Proposition~\ref{prop:diss-short-sep} delivers the central result: the good borrower separates by committing to a suboptimally low probability of continuation.
Since liquidation is as costly to the good type as to the bad type when a shock occurs but is overall less costly for the good type (because shocks occur with lower probability $\lambda$ instead of $\widetilde{\lambda}$), the good type can tolerate a smaller $x$ while still preferring the separating contract to mimicking the bad-type pool.

Three economic observations are worth highlighting:

\begin{itemize}
  \item \emph{The good borrower grants himself a higher reward $R_b^+$ in the no-shock state} than under symmetric information.
  This compensates for the reduced continuation probability: the lender is willing to accept a higher success compensation because the firm is less likely to require costly reinvestment.
  \item \emph{The good borrower is worse off than under symmetric information.}
  He sacrifices continuation ($x < 1$ rather than $x = 1$) to signal his type.
  Continuation is a more efficient ``currency'' than monetary compensation because $\rho < \rho_1$ and $R_b^+ \geq B/\Delta p$: the cost of sacrificing continuation exceeds the value of the compensating higher $R_b^+$.
  \item \emph{Implementation:} the separating contract can be implemented via a \emph{stochastic} short-term debt, $d = r$ with probability $1 - \lambda x$ and $d = r - \rho$ with probability $\lambda x$.
  The short-term debt is on average larger than under symmetric information---the good borrower takes on more short-term refinancing risk---reflecting the signaling value of short-term maturities.
  This stochastic character is an artefact of the discrete-shock model; with a continuous distribution of shocks, the optimal separating short-term debt would be deterministic but still larger than under symmetric information.
\end{itemize}

\subsection{Uniqueness and coexistence}

The separating equilibrium described above is the unique equilibrium when the prior $\alpha$ on good types is sufficiently small, $\alpha \leq \alpha^\star$ for some threshold.
When $\alpha > \alpha^\star$, other equilibria exist, typically involving pooling, that are preferred by both types to the separating outcome (the benefits of pooling outweigh the cost of cross-subsidization when there are enough good types to dilute the adverse-selection friction).

%====================================================================
\section{Costly Collateral Pledging}
\label{sec:diss-collateral}
%====================================================================

The third dissipative signal is the pledging of costly collateral.
The borrower pledges assets that have deadweight-loss value to the lender: an asset valued $C$ by the borrower has value $\beta C$ to the lender, with $\beta \in [0, 1)$.
The collateral is seized only in the failure state; since failure is less likely for good borrowers, the per-period expected collateral cost is less for a good type than for a bad type.
This wedge enables separation.

The analysis follows \citet{Bester1985, Bester1987}, \citet{BesankoThakor1987}, and \citet{ChanKanatas1985}, who developed the collateral-as-signal framework in various settings.

\subsection{Setup}

We extend the privately-known-prospects model with $A = 0$, investment $I$, and type-contingent success probabilities $p > q$.
The borrower has access to collateralizable assets---equipment, intellectual property, or other assets that would not exist absent the project---and can pledge any amount $C \geq 0$ to the lender.
If the project succeeds, the collateral is retained by the borrower; if it fails, the collateral is seized by the lender but valued at only $\beta C$, yielding a deadweight loss of $(1 - \beta) C$.

The collateral wedge $\beta$ captures the specific-investment nature of the assets: equipment tailored to the borrower's production process, brand-specific intellectual property, or similar assets that lose value when transferred.

\begin{assumption}[Both types financeable without collateral]\label{ass:diss-coll-viable}
\begin{equation*}
  0 < q R - I < p R - I.
\end{equation*}
\end{assumption}

Under symmetric information, neither type needs to pledge collateral; the bad type is creditworthy at a probability $q$ without collateral.

\subsection{Contracts}

A contract is a pair $(R_b, C)$: the borrower receives compensation $R_b$ in success (and keeps the collateral) and receives zero in failure (with the collateral seized).
The borrower's expected payoff, for a type with success probability $\pi \in \{p, q\}$, is
\begin{equation*}
  \Pi_b^\pi(R_b, C) = \pi R_b - (1 - \pi) C.
\end{equation*}
The lender's expected payoff is $\pi(R - R_b) + (1 - \pi) \beta C - I$.

\subsection{Symmetric-information benchmark}

Under symmetric information, neither type pledges collateral: $C = 0$ is strictly preferred by the borrower (since collateral is costly with no signaling value), and the type-specific breakeven contracts are
\begin{equation*}
  R_b^{G} = R - \frac{I}{p},
  \quad
  R_b^{B} = R - \frac{I}{q}.
\end{equation*}

\subsection{Asymmetric information: pooling equilibrium}

Under asymmetric information, the pooling contract is obtained by maximizing the good borrower's payoff subject to the pooled breakeven $m(R - R_b) + (1 - m) \beta C \geq I$.
A structural result determines when the pooling contract features no collateral:

\begin{proposition}[Pooling with no collateral]\label{prop:diss-coll-pool-nocoll}
If $(p/(1-p)) \beta \leq m/(1-m)$, the good type's pooling-optimal contract is $(R_b, 0)$ with $R_b = R - I/m$; the bad type prefers this contract to his symmetric-information contract, generating a pooling equilibrium.
If $(p/(1-p)) \beta > m/(1-m)$, no pooling equilibrium exists.
\end{proposition}

The threshold condition $(p/(1-p)) \beta \leq m/(1-m)$ can be restated in terms of $\alpha$: as $\alpha$ rises (more good types), $m/(1-m)$ rises and the condition is eventually satisfied.
Define the threshold $\alpha^\star$ by the equality; then pooling equilibrium exists if and only if $\alpha \geq \alpha^\star$.

\begin{proof}[Proof sketch]
The pooling program maximizes $p R_b - (1-p) C$ subject to $m(R - R_b) + (1 - m) \beta C \geq I$ with equality (by standard perturbation).
The first-order condition on $C$ is $-(1-p) + \lambda (1 - m) \beta = 0$, where $\lambda$ is the multiplier on the breakeven.
From the FOC on $R_b$, $\lambda = p/m$.
Substituting, $(1-p)/(1-m) = (p/m) \beta$, i.e., $(1-m) p \beta = m(1-p)$, i.e., $m/(1-m) = p\beta/(1-p)$ at the interior solution.
When $p\beta/(1-p) < m/(1-m)$ (strict inequality), the FOC on $C$ has negative sign, so $C = 0$ is optimal; when $p\beta/(1-p) > m/(1-m)$, the optimal $C$ is at its maximum and the candidate pooling contract does not satisfy the bad type's incentive to mimic.
\end{proof}

\subsection{Optimal separating contract}

In the separating regime (when pooling fails), the good type offers a contract $(R_b, C)$ designed to deter mimicking by the bad type.
The optimal separating contract solves:
\begin{equation*}
  \max_{R_b, C \geq 0} \, p R_b - (1 - p) C
\end{equation*}
subject to:
\begin{itemize}
  \item \emph{Lender's breakeven at good type}: $p(R - R_b) + (1 - p) \beta C \geq I$ (IR$_\ell$);
  \item \emph{Bad type's participation}: the bad type prefers his own symmetric contract to mimicking: $q R_b - (1 - q) C \leq q R - I \eqqcolon \Pi_b^{B}$ (IR$_{\text{bad}}$).
\end{itemize}

\begin{proposition}[Binding (IR$_\ell$) and (IR$_{\text{bad}}$)]\label{prop:diss-coll-sep-binding}
At the optimum, both the lender's breakeven (IR$_\ell$) and the separating constraint (IR$_{\text{bad}}$) bind.
\end{proposition}

\begin{proof}[Proof sketch]
Standard perturbation arguments: relaxing either constraint allows improvement in the good borrower's objective.
The full algebraic details are in the slide proofs.
\end{proof}

\begin{proposition}[Closed form for the separating contract]\label{prop:diss-coll-closed-form}
The optimal separating contract is
\begin{equation}\label{eq:diss-coll-Rb-sep}
  R_{b, \text{sep}}^{G} = R - \left[ \frac{(1 - q) - \beta(1 - p)}{p(1 - q) - \beta q (1 - p)} \right] I,
\end{equation}
\begin{equation}\label{eq:diss-coll-C-sep}
  C_{\text{sep}} = \frac{I}{1 + q(1 - p)(1 - \beta)/(p - q)} > 0.
\end{equation}
\end{proposition}

\begin{proof}
Solving the system of binding constraints:
\begin{align*}
  p(R - R_b) + (1 - p) \beta C &= I, \\
  q R_b - (1 - q) C &= q R - I.
\end{align*}
Standard linear algebra yields the closed forms.
\end{proof}

\subsection{Why signaling works}

The feasibility of a separating equilibrium with collateral rests on a simple intuition:
\begin{itemize}
  \item Pledging collateral is \emph{relatively} more costly for the bad borrower than for the good borrower: the per-period expected collateral cost is $(1 - \pi) C$, which is larger for the smaller $\pi = q$.
  Since $1 - q > 1 - p$, a bad borrower loses more expected value from each unit of collateral pledged.
  \item A higher reward $R_b$ in success is valued more by the good borrower than by the bad borrower: the per-period expected compensation is $\pi R_b$, which is larger for $\pi = p$.
\end{itemize}
The combination of a high $R_b$ (valued by the good type) with a high $C$ (costly to the bad type) enables the good type to self-select into a contract the bad type does not want to mimic.

\subsection{Comparative statics}

The closed form \eqref{eq:diss-coll-C-sep} yields sharp comparative-static predictions:

\begin{itemize}
  \item \emph{Collateral efficiency $\beta$:} $\partial C_{\text{sep}}/\partial \beta > 0$.
  When collateral pledging becomes less costly (higher $\beta$, smaller deadweight loss), the good borrower pledges more.
  The intuition is that the dual benefits of collateral (signaling and lender reassurance) are enhanced when the deadweight loss is small.
  \item \emph{Asymmetry $q$:} $\partial C_{\text{sep}}/\partial q < 0$.
  When the bad borrower's quality $q$ falls (asymmetry widens), the good borrower must pledge more collateral to separate.
  The larger informational gap requires a larger signal to credibly convey it.
\end{itemize}

%====================================================================
\section{Diversification and Incomplete Insurance}
\label{sec:diss-divers}
%====================================================================

The fourth dissipative signal arises in the context of \emph{risk-sharing}.
A risk-averse entrepreneur with a substantial stake in his firm may want to diversify his portfolio by selling claims to investors.
The sale is not motivated by financing a new project---it is motivated entirely by the desire to share idiosyncratic risk with investors who can diversify it away.
Under asymmetric information, however, the entrepreneur's choice of how much to sell becomes a signal of type, and the equilibrium features \emph{incomplete} insurance: the good borrower retains a larger stake than would be efficient under symmetric information, bearing risk to signal high quality.

The mechanism was articulated by \citet{LelandPyle1977} and is the corporate-finance counterpart of the classical Rothschild--Stiglitz analysis of insurance markets \citep{RothschildStiglitz1976, Stiglitz1977}.

\subsection{Setup}

The entrepreneur has no initial cash ($A = 0$) but owns the firm entirely.
No new financing is needed ($I = 0$); the entrepreneur's dealings with investors are solely for diversification.
The entrepreneur is risk-averse: his preferences are represented by a twice continuously differentiable, strictly increasing, strictly concave utility function $U: \R_+ \to \R$ satisfying Inada's condition at $0$.
Investors are risk-neutral---the firm's risk is idiosyncratic and can be diversified away.

At $t = 1$, the firm yields $R$ in success (probability $p$ good, $q$ bad) and zero in failure.
The entrepreneur sells claims to investors in exchange for a type-contingent contract $(R_b^S, R_b^F)$, where $R_b^S$ is what he receives in success and $R_b^F$ is what he receives in failure.
(The investors receive $R - R_b^S$ in success, zero in failure.)

Feasibility requires $0 \leq R_b^S \leq R$ and $0 \leq R_b^F$; the investors' breakeven requires $p R_b^S + (1 - p) R_b^F \leq pR$ for the good type (and analogously for the bad type).

\subsection{Symmetric-information benchmark}

\begin{proposition}[Full insurance under symmetric information]\label{prop:diss-divers-sym}
Under symmetric information, both types receive full insurance: the good type's optimal contract is $R_b^S = R_b^F = pR$, and the bad type's optimal contract is $R_b^S = R_b^F = qR$.
\end{proposition}

\begin{proof}
For a type $\pi \in \{p, q\}$, the program is
\begin{equation*}
  \max_{R_b^S, R_b^F} \, \pi U(R_b^S) + (1 - \pi) U(R_b^F)
  \quad\text{s.t.}\quad
  \pi R_b^S + (1 - \pi) R_b^F \leq \pi R.
\end{equation*}
The Lagrangian FOC gives $\pi U'(R_b^S) = \lambda \pi$ and $(1 - \pi) U'(R_b^F) = \lambda (1 - \pi)$, which together yield $U'(R_b^S) = U'(R_b^F)$, i.e., $R_b^S = R_b^F$ by strict concavity.
Substituting into the binding breakeven: $R_b^S = R_b^F = \pi R$.
\end{proof}

The symmetric-information outcome is efficient: the risk-averse entrepreneur sells off all idiosyncratic risk to the risk-neutral investors, consuming a certain amount $\pi R$ equal to the expected value of his project.
The allocation is the first-best from both risk-sharing and efficiency perspectives.

\subsection{Pooling equilibrium}

Under asymmetric information, the pooled breakeven condition is $m R_b^S + (1 - m) R_b^F \leq mR$, with $m = \alpha p + (1 - \alpha) q$ as usual.
The good borrower's program maximizes $p U(R_b^S) + (1 - p) U(R_b^F)$ subject to the pooled breakeven.

\begin{proposition}[Pooling contract]\label{prop:diss-divers-pool}
The pooling contract is characterized by the first-order condition
\begin{equation}\label{eq:diss-divers-pool-foc}
  \frac{p}{1 - p} U'(R_{b, \text{pool}}^S) = \frac{m}{1 - m} U'(R_{b, \text{pool}}^F),
\end{equation}
with $R_{b, \text{pool}}^S > R_{b, \text{pool}}^F$.
The good borrower's payoff under pooling is strictly less than under symmetric information:
\begin{equation*}
  \Pi_{b, \text{pool}}^{G} < U(pR) = \Pi_{b, \text{sym}}^{G}.
\end{equation*}
\end{proposition}

\begin{proof}
The FOC from the Lagrangian gives $p U'(R_b^S)/m = (1-p) U'(R_b^F)/(1-m)$, equivalent to \eqref{eq:diss-divers-pool-foc}.
Since $p > m$, we have $p/(1-p) > m/(1-m)$, so $U'(R_{b, \text{pool}}^S) < U'(R_{b, \text{pool}}^F)$, which by strict concavity of $U$ implies $R_{b, \text{pool}}^S > R_{b, \text{pool}}^F$.
The welfare comparison follows because the pooling contract is suboptimal relative to the first-best (which provides full insurance at $R_b^S = R_b^F = pR$).
\end{proof}

The intuition is that the good type, being of higher quality than the pool average, must leave a disproportionate stake in the success state to avoid being underpriced; as a result, pooling delivers less than full insurance and strictly less welfare than the symmetric-information benchmark.

Whether the bad type prefers to mimic the pooling contract depends on $\alpha$: for $\alpha$ sufficiently large, $\Pi_b^{B}(R_{b, \text{pool}}^S, R_{b, \text{pool}}^F) > U(qR) = \Pi_{b, \text{sym}}^{B}$, and the pooling contract is an equilibrium; for small $\alpha$, the bad type would prefer his symmetric-information contract, and no pooling equilibrium exists.

\subsection{Separating equilibrium}

When no pooling equilibrium exists, the good type offers a separating contract.
The program is
\begin{equation*}
  \max_{R_b^S, R_b^F} \, p U(R_b^S) + (1 - p) U(R_b^F)
\end{equation*}
subject to the good-type breakeven and the separating constraint:
\begin{itemize}
  \item (IR$_\ell$): $p R_b^S + (1 - p) R_b^F \leq pR$;
  \item (IR$_{\text{bad}}$): $q U(R_b^S) + (1 - q) U(R_b^F) \leq U(qR)$.
\end{itemize}

\begin{proposition}[Binding constraints]\label{prop:diss-divers-binding}
At the optimum, both (IR$_\ell$) and (IR$_{\text{bad}}$) bind.
\end{proposition}

\begin{proof}[Proof sketch]
(IR$_\ell$) binding follows by the usual perturbation argument: slack (IR$_\ell$) allows strictly increasing the good type's payoff while preserving (IR$_{\text{bad}}$).
(IR$_{\text{bad}}$) binding requires a more delicate three-case argument distinguishing $R_b^S > R_b^F$, $R_b^S < R_b^F$, and $R_b^S = R_b^F$:
\begin{itemize}
  \item If $R_b^S > R_b^F$ and (IR$_{\text{bad}}$) is slack, a perturbation with $dR_b^S < 0$ and $p \, dR_b^S + (1-p) \, dR_b^F = 0$ strictly decreases the wedge between success and failure consumption, raising good-type utility.
  \item If $R_b^S < R_b^F$ (over-insurance) and (IR$_{\text{bad}}$) is slack, a symmetric perturbation raises good-type utility.
  \item If $R_b^S = R_b^F$, the binding (IR$_\ell$) forces $R_b^S = R_b^F = pR$, which violates (IR$_{\text{bad}}$) by strict concavity.
\end{itemize}
\end{proof}

\subsection{Graphical characterization}

The geometric structure of the separating equilibrium is illuminating.
In the $(R_b^S, R_b^F)$ plane, the zero-profit line for the good type passes through two distinguished contracts:
\begin{itemize}
  \item $c_R = (R, 0)$: the no-insurance contract (entrepreneur bears all risk);
  \item $c_G = (pR, pR)$: the good-type full-insurance contract.
\end{itemize}
Any contract on the segment joining $c_R$ and $c_G$ satisfies the good-type breakeven with equality.
The separating contract is the point on this segment that also satisfies the bad-type participation constraint $q U(R_b^S) + (1 - q) U(R_b^F) = U(qR)$---the indifference curve of the bad type passing through $c_B = (qR, qR)$.

The separating contract $c_{\text{sep}}$ lies \emph{between} $c_R$ and $c_G$ on the good-type's zero-profit line: the good type sells some of his risk but not all of it.
The welfare cost of separation, relative to $c_G$, is the dissipative signal: the good type bears risk to deter the bad type from mimicking.

\subsection{Comparative statics}

The separating contract $c_{\text{sep}}$ responds to the informational environment as follows:

\begin{itemize}
  \item As the bad type's quality $q$ decreases (holding $p$ fixed), the bad type's indifference curve $c_B = (qR, qR)$ moves toward the origin.
  The separating contract $c_{\text{sep}}$ moves closer to the no-insurance contract $c_R$: the good type must retain more risk to signal his type.
  The larger the informational gap, the more incomplete the good type's insurance.
\end{itemize}

The result is economically striking: \emph{the good type signals good prospects by increasing the sensitivity of his own return to the firm's profit}.
Retained equity stakes---a common empirical pattern among entrepreneurs of young, high-quality firms---are interpreted as a signal of confidence under this framework.

\subsection{Perfect Bayesian equilibrium}

The full equilibrium consists of the pair of contracts $(c_{\text{sep}}, c_B)$, with the good type choosing $c_{\text{sep}}$ and the bad type choosing $c_B$.
Investors' beliefs are updated rationally: a borrower offering $c_{\text{sep}}$ is inferred to be good; a borrower offering $c_B$ is inferred to be bad.
The separating constraint (IR$_{\text{bad}}$) ensures the bad type prefers $c_B$ to $c_{\text{sep}}$; the good-type breakeven (IR$_\ell$) ensures investors are willing to offer $c_{\text{sep}}$.

%====================================================================
\section{Common Structure and Further Reading}
\label{sec:diss-further}
%====================================================================

\subsection{The common logical structure}

The four mechanisms analyzed in this chapter---certification, short-term maturities, costly collateral pledging, and diversification---share a common logical structure that bears articulating.

In each case, the good borrower has access to a costly action whose cost is \emph{relatively} lower for him than for the bad borrower:

\begin{itemize}
  \item \emph{Certification}: the cost $c$ is the same in absolute terms for both types, but the relative benefit of the signal (avoiding cross-subsidization) is larger for the good type, who would otherwise subsidize bad types.
  \item \emph{Short-term maturities}: the cost of sacrificing continuation probability $x$ is smaller for the good type because liquidity shocks are less frequent ($\lambda < \widetilde{\lambda}$), so the expected continuation loss is smaller.
  \item \emph{Collateral pledging}: the expected deadweight loss $(1 - \pi)(1 - \beta) C$ is smaller for the good type because failure is less likely ($1 - p < 1 - q$).
  \item \emph{Diversification}: the cost of retaining risk---measured by the concavity of utility---is smaller for the good type because the \emph{expected} income is higher ($pR > qR$), so the marginal utility of the retained risky income is lower.
\end{itemize}

In all four cases, the signal is \emph{dissipative}: it destroys rather than creates value, but the good type prefers bearing the dissipation to pooling with the bad type.
The equilibrium cost of the signal is the welfare wedge between the asymmetric-information outcome and the first-best.

\subsection{Reputation and communication among lenders}

An alternative approach to reducing informational asymmetries is to rely on \emph{past repayment history} as a signal of borrower quality.
\citet{PadillaPagano1997} analyze the endogenous emergence of communication among lenders, showing that information sharing can relax the adverse-selection friction but may introduce its own incentive problems.
The reputational mechanism is complementary to the dissipative signals analyzed in this chapter and is a key feature of many real-world credit markets.

\subsection{Further reading}

\citet{Tirole2006} (Chapter~6, Part~B) provides the standard textbook exposition of dissipative signals, including extensions beyond the four mechanisms covered here.
The classical references for signaling and separating equilibria in insurance markets are \citet{RothschildStiglitz1976} and \citet{Stiglitz1977}; \citet{LelandPyle1977} is the seminal adaptation of this framework to corporate finance.
The role of certification in investment banking is analyzed by \citet{Baron1982} and extended by \citet{ChemmanurFulghieri1999}; the relevance of short-term maturity as a signal is developed by \citet{Diamond1991, Diamond1993} and applied to labor markets by \citet{Hermalin2002}; the collateral-as-signal framework is developed by \citet{Bester1985, Bester1987}, \citet{BesankoThakor1987}, and \citet{ChanKanatas1985}.

The four mechanisms analyzed here do not exhaust the range of dissipative signals used in practice: dividend payouts have been studied as signals of firm quality, convertible debt has been interpreted as a device for managing signaling through capital structure, and the underpricing of IPOs has been analyzed as a signaling phenomenon.
The reader is referred to \citet{Tirole2006} and the corporate-finance literature more broadly for these and other extensions.

%%%%%%%%%%%%%%%%%%%%%%%%%%%%%%%%%%%%%%%%%%%%%%%%%%%%%%%%%%%%%%%%%%%%%%%%

\part{Mergers and Acquisitions}

%%%%%%%%%%%%%%%%%%%%%%%%%%%%%%%%%%%%%%%%%%%%%%%%%%%%%%%%%%%%%%%%%%%%%%%%

%====================================================================
\chapter{Mergers, Acquisitions, and Asset Redeployability}
\label{ch:ma}
%====================================================================

The models developed in Parts~II and~III studied the financing of a \emph{single} firm in isolation.
The entrepreneur contracted with a competitive lender, and the analysis focused on how frictions---moral hazard, adverse selection, liquidity shocks---shaped the optimal contract and the firm's scale of investment.
A natural next question is how these frictions interact when \emph{multiple} firms operate in the same industry and may redeploy assets to one another.
The financial structure of any one firm can then depend on the financial and operational health of the others: a firm in distress may sell its assets to a healthy rival; a healthy firm's acquisition decisions depend on the distress probabilities and financial capacity of its industry peers.

This chapter develops the canonical industry-equilibrium model of mergers and acquisitions under financial frictions, following \citet{ShleiferVishny1992} and the subsequent literature.
Two firms operate in the same industry.
After an initial investment, each firm enters a learning period that reveals whether it is \emph{productive} (can pursue the project profitably) or \emph{in distress} (the assets have no internal use).
A distressed firm's assets can be redeployed to the healthy firm; the transfer price is determined through post-distress bargaining.
The industry equilibrium---with endogenous transfer prices, borrowing capacities, and investment decisions---exhibits a rich interdependence across firms that would be absent in isolated models.

The analysis delivers three central results.
First, the equilibrium transfer price is the \emph{pledgeable income} $\rho_0$ of the healthy firm, regardless of which side has bargaining power.
Moral-hazard frictions endogenously limit how much the healthy firm can pay for redeployed assets, creating a structural \emph{illiquidity} in the secondary asset market: even under competitive bidding, asset prices remain below their first-best value.
Second, the \emph{correlation of shocks} across firms in an industry has first-order effects on borrowing capacity: a higher correlation---a greater likelihood that both firms are simultaneously in distress---reduces the expected recovery value to lenders and tightens the financing constraint.
The equity multiplier $k$ is a decreasing function of the correlation parameter $\nu$.
Third, under high correlation, the economy can exhibit \emph{multiple equilibria}: one in which firms coordinate on investing (supported by the prospect of cross-firm asset redeployment) and one in which no firm invests (each firm's lack of investment destroys the redeployability option for the other).
This is the industry-equilibrium counterpart of the soft-budget-constraint problem of Chapter~\ref{ch:liquidity-scale}.

The chapter closes with an analysis of \emph{underdeveloped resale markets}.
When firms are unable to coordinate on investment design (risky vs.\ safe configurations), the lack of coordination can lead to inefficiently low volumes of transactions: firms individually choose the safe design to avoid the risk of becoming distressed, but the collective outcome is Pareto-dominated by the risky design that exploits asset redeployability.
The result formalizes an intuition central to \citet{ShleiferVishny1992}: the illiquidity of secondary asset markets is self-reinforcing through firms' ex-ante design choices.

The chapter corresponds to Chapter~8 of \citet{Tirole2006}.

%====================================================================
\section{The Two-Firm Industry Model}
\label{sec:ma-model}
%====================================================================

\subsection{Primitives}

Two firms, indexed $i \in \{1, 2\}$, operate in the same industry.
Each firm is run by a risk-neutral entrepreneur who arrives at $t = 0$ with initial wealth $A_i \geq 0$.
Firm $i$ invests $I_i$ at $t = 0$, borrowing $I_i - A_i$ from a competitive lender.
Investors are risk-neutral with zero discount rate.

After investment, each firm enters a \emph{costless learning period}.
At the end of the learning period, firm $i$ privately learns its operational status:
\begin{itemize}
  \item with probability $x \in (0, 1)$, the firm is \emph{productive}: its investment is viable and can yield the standard variable-investment payoff;
  \item with probability $1 - x$, the firm is \emph{in distress}: its assets have no internal use and, if left in place, generate zero.
\end{itemize}

\subsection{Productive firms}

If firm $i$ is productive, the standard variable-investment model of Chapter~\ref{ch:equity-multiplier} applies.
The firm produces $R I_i$ in success and zero in failure; the success probability depends on the entrepreneur's effort choice: $p_h$ if he behaves (no private benefit), $p_\ell < p_h$ if he misbehaves (private benefit $B I_i$).
The incentive-compatible stake per unit of investment is $B/\Delta p$, with $\Delta p = p_h - p_\ell$.

We use the standard notation $\rho_1 \coloneqq p_h R$ (expected success payoff per unit) and $\rho_0 \coloneqq p_h(R - B/\Delta p)$ (pledgeable income per unit).

\subsection{Distressed firms and asset transfers}

If firm $j$ is in distress:
\begin{itemize}
  \item The assets of firm $j$ are useless in firm $j$'s hands.
  If left in place, they yield zero, and entrepreneur $j$ cannot even enjoy a private benefit (no production is possible).
  \item The assets can be sold to firm $i$, provided firm $i$ is not itself in distress.
  Only within-industry redeployment is feasible: the assets have no outside-industry value.
  \item The transfer price is \emph{not} contracted ex ante.
  It is determined by bargaining at the time of distress.
  Let $P_j$ denote the per-unit-of-investment transfer price for firm $j$'s assets.
\end{itemize}

Conditional on purchasing firm $j$'s assets, entrepreneur $i$ operates a total investment of $I_i + I_j$.
The integrated firm still faces moral hazard: entrepreneur $i$ obtains private benefit $B(I_i + I_j)$ if he misbehaves, and the project succeeds with probability $p_h$ under behaving, yielding $R(I_i + I_j)$.

\subsection{Timing}

The sequence of events is:
\begin{enumerate}
  \item $t = 0$: each firm $i$ signs a loan agreement with lender $i$ and invests $I_i$.
  \item Learning period: each firm privately learns its operational status.
  \item Possible asset transfer: if exactly one firm is in distress, its assets may be transferred to the healthy firm at a bargained price $P_j$.
  \item Moral hazard: the surviving/integrated entrepreneur chooses effort.
  \item $t = 1$: outcomes are realized; payoffs are distributed per contract.
\end{enumerate}

\subsection{Maintained assumptions}

We impose two assumptions that parallel those of Chapter~\ref{ch:equity-multiplier}.

\begin{assumption}[Project viability]\label{ass:ma-viability}
\begin{equation}\label{eq:ma-viability}
  \rho_1 > 1 > p_\ell R + B.
\end{equation}
\end{assumption}

Projects are viable only under behaving: the per-unit NPV is positive if the entrepreneur works and negative if he shirks.

\begin{assumption}[Finite investment]\label{ass:ma-finite}
\begin{equation}\label{eq:ma-finite}
  \rho_0 < 1.
\end{equation}
\end{assumption}

The per-unit pledgeable income is strictly less than one, so the entrepreneur's equity is necessarily binding for finite investment.

\subsection{Loan agreements}

At $t = 0$, lender $i$ and entrepreneur $i$ sign a loan agreement specifying:
\begin{itemize}
  \item the loan amount $I_i - A_i$;
  \item the entrepreneur's stake $R_{b i}$ in success \emph{absent purchase} of firm $j$'s assets, satisfying
  \begin{equation}\label{eq:ma-IC-no-purchase}
    (\Delta p) R_{b i} \geq B I_i;
  \end{equation}
  \item the entrepreneur's stake $\widetilde{R}_{b i}$ in success \emph{with purchase} of firm $j$'s assets, satisfying
  \begin{equation}\label{eq:ma-IC-purchase}
    (\Delta p) \widetilde{R}_{b i} \geq B (I_i + I_j).
  \end{equation}
\end{itemize}
In equilibrium, the other firm's loan agreement is correctly anticipated even though not directly observed: this is the standard Nash-equilibrium assumption.

%====================================================================
\section{The Correlation Structure}
\label{sec:ma-correlation}
%====================================================================

\subsection{Joint distribution of operational status}

The productive/distress status of the two firms is, in general, correlated.
We parameterize the joint distribution via the conditional probabilities
\begin{itemize}
  \item $\mu \coloneqq \Prob(\text{firm } j \text{ productive} \vert \text{firm } i \text{ productive})$;
  \item $\nu \coloneqq \Prob(\text{firm } j \text{ in distress} \vert \text{firm } i \text{ in distress})$.
\end{itemize}
Consistency with the unconditional probability of being productive requires
\begin{equation}\label{eq:ma-consistency}
  x \mu + (1 - x)(1 - \nu) = x,
  \quad \text{equivalently,} \quad
  x(1 - \mu) = (1 - x)(1 - \nu).
\end{equation}
Under \eqref{eq:ma-consistency}, the joint probability that firm $j$ is productive equals the unconditional $x$, regardless of firm $i$'s status.

\begin{assumption}[Symmetry]\label{ass:ma-symmetry}
The joint distribution satisfies the consistency condition \eqref{eq:ma-consistency}, so that both firms have unconditional productive probability $x$.
\end{assumption}

\subsection{Two leading cases}

Two extreme cases will be useful:

\begin{itemize}
  \item \emph{Nonconcurrent risk} ($\nu = 0$, equivalently $\mu = (2x - 1)/x$): if firm $i$ is in distress, firm $j$ is certainly productive.
  Industry shocks are perfectly offsetting.
  \item \emph{Identical shock} ($\mu = \nu = 1$): either both firms are productive or both are in distress.
  Industry shocks are perfectly aligned.
\end{itemize}
The parameter $\nu$ interpolates between these two extremes and will play a central role in the comparative statics.

%====================================================================
\section{The Equilibrium Transfer Price}
\label{sec:ma-transfer}
%====================================================================

\subsection{The bargaining game}

Consider the event in which firm $j$ is in distress and firm $i$ is productive.
The two parties to the asset-transfer negotiation are the lenders of the two firms: lender $j$ has seized firm $j$'s assets (since the distressed firm cannot itself operate them) and seeks to sell; lender $i$ represents firm $i$'s capacity to purchase and operate the assets.

We assume that lender $j$ makes a \emph{take-it-or-leave-it} offer to lender $i$, specifying a per-unit transfer price $P_j$.
If lender $i$ accepts, the assets are transferred at the price $P_j I_j$, and entrepreneur $i$ operates the integrated firm of size $I_i + I_j$.
If lender $i$ rejects, the assets are discarded and yield zero to all parties.

\subsection{The incremental value to the acquirer}

Lender $i$'s incremental expected net payoff from the acquisition is
\begin{equation}\label{eq:ma-incremental}
  p_h \bigl\{ R I_j - (\widetilde{R}_{b i} - R_{b i}) \bigr\} - P_j I_j,
\end{equation}
where $p_h \{R I_j - (\widetilde{R}_{b i} - R_{b i})\}$ is the expected incremental revenue from the integrated firm (net of the entrepreneur's additional stake), and $P_j I_j$ is the cost.

Under \eqref{eq:ma-IC-purchase} and \eqref{eq:ma-IC-no-purchase}, the entrepreneur's stake can be adjusted up by at most
\begin{equation*}
  \widetilde{R}_{b i} - R_{b i} = \frac{B I_j}{\Delta p}
\end{equation*}
(with equality at the incentive-compatibility minimum for the integrated firm).
Substituting into \eqref{eq:ma-incremental}, lender $i$'s maximum willingness-to-pay per unit is
\begin{equation*}
  p_h \left\{ R - \frac{B}{\Delta p} \right\} = \rho_0.
\end{equation*}

\subsection{The equilibrium price}

\begin{proposition}[Competitive transfer price]\label{prop:ma-transfer}
Under the take-it-or-leave-it bargaining protocol with lender $j$ as the proposer, the equilibrium per-unit transfer price is
\begin{equation}\label{eq:ma-price}
  P_j = \rho_0.
\end{equation}
\end{proposition}

\begin{proof}
Lender $j$ maximizes his extraction subject to lender $i$'s participation (non-negative incremental payoff).
Setting the incremental payoff to zero: $\rho_0 - P_j = 0$, giving $P_j = \rho_0$.
\end{proof}

\subsection{Interpretation: endogenous illiquidity}

Proposition~\ref{prop:ma-transfer} yields the striking result that the equilibrium transfer price equals the pledgeable income $\rho_0 < 1$, regardless of which side holds the bargaining power.
A few economic observations follow.

\begin{itemize}
  \item \emph{Competitive equivalence.} The price $\rho_0$ is precisely the competitive-auction price that would emerge if \emph{multiple} acquirers competed for firm $j$'s assets.
  Even with a single bidder, the moral-hazard constraint on lender $i$'s incremental willingness-to-pay caps the transfer price at $\rho_0$.
  The take-it-or-leave-it power of the seller is moot: the buyer's borrowing capacity for the acquisition is itself bounded.
  \item \emph{Asset-market illiquidity.} Since $\rho_0 < 1$ by Assumption~\ref{ass:ma-finite}, the assets are sold at a discount relative to their replacement cost ($P_j < 1$).
  The secondary asset market exhibits endogenous illiquidity due to moral hazard: buyers' ability to pay is capped by their own pledgeable-income capacity.
  This is the industry-equilibrium analog of the credit-rationing phenomenon developed in Chapter~\ref{ch:net-worth}.
  \item \emph{Fire-sale pricing without fire-sale dynamics.} The below-par transfer price arises \emph{without} any dynamic or stochastic fire-sale mechanism.
  It is a static equilibrium feature of moral-hazard-constrained acquisition financing.
\end{itemize}

%====================================================================
\section{Optimal Contracts and the Nash Equilibrium}
\label{sec:ma-equilibrium}
%====================================================================

Having pinned down the equilibrium transfer price, we now characterize the loan agreement for each firm.

\subsection{The borrower's expected payoff}

For firm $i$, conditional on each operational-status event, the entrepreneur's payoff is:
\begin{itemize}
  \item both productive (prob. $x \mu$): entrepreneur $i$ receives $p_h R_{b i}$;
  \item firm $i$ productive, firm $j$ distressed (prob. $x(1 - \mu)$): entrepreneur $i$ receives $p_h \widetilde{R}_{b i}$;
  \item firm $i$ distressed (prob. $1 - x$): entrepreneur $i$ receives zero.
\end{itemize}
Aggregating:
\begin{equation}\label{eq:ma-Pi-b-raw}
  \Pi_{b i}(I_i, R_{b i}, \widetilde{R}_{b i}) = -A_i + x \left[\mu p_h R_{b i} + (1 - \mu) p_h \widetilde{R}_{b i}\right].
\end{equation}

Setting the (IC)-binding stake at the acquisition event, $\widetilde{R}_{b i} - R_{b i} = B I_j/\Delta p$, we have $p_h \widetilde{R}_{b i} = p_h R_{b i} + (\rho_1 - \rho_0) I_j$.
Substituting back:
\begin{equation}\label{eq:ma-Pi-b}
  \Pi_{b i}(I_i, R_{b i}, I_j) = -A_i + x \left[ p_h R_{b i} + (1 - \mu)(\rho_1 - \rho_0) I_j \right].
\end{equation}

\subsection{The lender's expected payoff}

For lender $i$, the expected payoff decomposes across three scenarios:
\begin{itemize}
  \item firm $i$ distressed, firm $j$ productive (prob. $(1 - x)(1 - \nu)$): lender $i$ sells assets at transfer price $\rho_0$ per unit, receiving $\rho_0 I_i$;
  \item both productive (prob. $x \mu$): lender $i$ receives $p_h(R I_i - R_{b i})$;
  \item firm $i$ productive, firm $j$ distressed (prob. $x(1 - \mu)$): lender $i$ purchases firm $j$'s assets at price $\rho_0 I_j$ and receives the integrated firm's pledgeable income;
  \item both distressed (prob. $(1 - x) \nu$): lender $i$ receives zero.
\end{itemize}

Since lender $j$ holds all the bargaining power in distress-sale bargaining (and makes no profit when \emph{purchasing} firm $j$), lender $i$'s expected payoff simplifies to
\begin{equation}\label{eq:ma-Pi-ell}
  \Pi_{\ell i}(I_i, R_{b i}) = -(I_i - A_i) + (1 - x)(1 - \nu) \rho_0 I_i + x p_h (R I_i - R_{b i}).
\end{equation}

\subsection{The entrepreneur's optimization}

Entrepreneur $i$ correctly anticipates entrepreneur $j$'s optimal contract $(I_j, R_{b j})$ and solves
\begin{equation*}
  \max_{I_i, R_{b i}} \, \Pi_{b i}(I_i, R_{b i}, I_j) = -A_i + x \left[ p_h R_{b i} + (1 - \mu)(\rho_1 - \rho_0) I_j \right]
\end{equation*}
subject to:
\begin{itemize}
  \item \emph{(IC)} at the no-purchase event: $p_h R_{b i} \geq B I_i$;
  \item \emph{Lender's breakeven} (IR$_\ell$): $\Pi_{\ell i}(I_i, R_{b i}) \geq 0$.
\end{itemize}

\begin{proposition}[Nash equilibrium characterization]\label{prop:ma-NE}
In any Nash equilibrium $(I_k, R_{b k})_{k \in \{1, 2\}}$, the lender's breakeven condition binds, and entrepreneur $i$'s problem reduces to
\begin{equation}\label{eq:ma-NPV}
  \max_{I_i} \, \NPV_i(I_i, I_j) = \alpha I_i + \kappa I_j,
\end{equation}
subject to \textup{(IC)} and the binding \textup{(IR$_\ell$)}, where
\begin{align}\label{eq:ma-alpha-kappa}
  \alpha &\coloneqq x \rho_1 + (1 - x)(1 - \nu) \rho_0 - 1, \\
  \kappa &\coloneqq x(1 - \mu)(\rho_1 - \rho_0).
\end{align}
\end{proposition}

\begin{proof}
Standard perturbation argument for (IR$_\ell$) binding (if slack, raise $R_{b i}$ to tighten).
Substituting the binding (IR$_\ell$) into the objective:
\begin{equation*}
  x p_h R_{b i} = x p_h R I_i + (1 - x)(1 - \nu) \rho_0 I_i - (I_i - A_i) = \alpha I_i + A_i.
\end{equation*}
Combined with the constant term $-A_i$ in \eqref{eq:ma-Pi-b} and the $\kappa I_j$ term, the objective reduces to \eqref{eq:ma-NPV}.
\end{proof}

The $\alpha$ coefficient is the \emph{net per-unit expected return} of firm $i$'s own investment---NPV under productive status plus salvage value under distress.
The $\kappa$ coefficient is the \emph{positive externality} that firm $j$'s investment confers on firm $i$: if firm $j$ is distressed, firm $i$ captures the integrated-firm rent $\rho_1 - \rho_0$ per unit of firm $j$'s investment.

%====================================================================
\section{Borrowing Capacity and the Equity Multiplier}
\label{sec:ma-multiplier}
%====================================================================

\subsection{Positive $\alpha$: interior investment}

When $\alpha > 0$, firm $i$'s investment contributes positively to its own expected payoff, and the entrepreneur invests up to the breakeven constraint.

\begin{proposition}[Equity multiplier with redeployability]\label{prop:ma-multiplier}
When $\alpha > 0$, the optimal contract features:
\begin{enumerate}
  \item[\textup{(i)}] \emph{(IC)} binds: $R_{b i} = (B/\Delta p) I_i$;
  \item[\textup{(ii)}] the pledgeable income is maximized and $I_i = k A_i$, where the equity multiplier is
  \begin{equation}\label{eq:ma-k}
    k = \frac{1}{1 - \rho_0 [x + (1 - x)(1 - \nu)]}.
  \end{equation}
\end{enumerate}
\end{proposition}

\begin{proof}
With (IC) binding, substituting $R_{b i} = (B/\Delta p) I_i$ into the binding (IR$_\ell$) gives
\begin{equation*}
  x p_h (R - B/\Delta p) I_i + (1 - x)(1 - \nu) \rho_0 I_i = I_i - A_i,
\end{equation*}
i.e., $\rho_0 [x + (1 - x)(1 - \nu)] I_i = I_i - A_i$, which solves for $I_i = k A_i$.
\end{proof}

The multiplier $k$ has a clean interpretation.
The term in brackets, $x + (1 - x)(1 - \nu)$, is the probability that firm $i$'s assets generate \emph{some} pledgeable cash flow---either through firm $i$'s own productive success (prob. $x$) or through the distress sale to a productive firm $j$ (prob. $(1 - x)(1 - \nu)$).
Multiplying by $\rho_0$ gives the expected per-unit pledgeable income; the complement $1 - \rho_0[\cdots]$ is the effective equity requirement per unit of investment.

\subsection{Comparative statics}

\begin{proposition}[Effect of correlation]\label{prop:ma-corr}
Both the equity multiplier $k$ and the per-unit expected return $\alpha$ are strictly decreasing in $\nu$:
\begin{equation*}
  \frac{\partial k}{\partial \nu} < 0,
  \quad
  \frac{\partial \alpha}{\partial \nu} < 0.
\end{equation*}
\end{proposition}

\begin{proof}
Differentiating \eqref{eq:ma-k} and \eqref{eq:ma-alpha-kappa} with respect to $\nu$:
\begin{equation*}
  \frac{\partial k}{\partial \nu} = -\frac{\rho_0 (1 - x)}{(1 - \rho_0 [x + (1 - x)(1 - \nu)])^2} \cdot (-1) \cdot (-1) = -\frac{\rho_0 (1 - x)}{\bigl(1 - \rho_0 [x + (1 - x)(1 - \nu)]\bigr)^2} < 0,
\end{equation*}
where we differentiate the bracketed term $(1-x)(1-\nu)$ to get $-(1-x)$.
Similarly,
\begin{equation*}
  \frac{\partial \alpha}{\partial \nu} = -(1 - x) \rho_0 < 0.
\end{equation*}
\end{proof}

The economic content is sharp: \emph{a firm's borrowing capacity decreases with the degree of correlation between firms in the industry}.
Higher correlation means it is more likely that both firms are simultaneously in distress, in which case no asset redeployment is possible and the salvage value to lender $i$ is zero.
Lenders, rationally internalizing this, reduce their willingness to finance each firm.

This result provides a microfoundation for the observation that highly concentrated industries with correlated risks (e.g., commercial real estate during a regional downturn, or specialized manufacturing in a cyclical sector) face tighter financing constraints than industries with diversifiable idiosyncratic risks.
The asset-redeployability channel is a quantitatively important component of industry-level financial frictions.

\subsection{Location of the investment threshold}

Define $\nu^\star$ to be the level of $\nu$ at which $\alpha = 0$.
Since $\alpha$ is decreasing in $\nu$:
\begin{itemize}
  \item If $\alpha < 0$ for every $\nu \in [0, 1]$, set $\nu^\star = 0$.
  \item If $\alpha > 0$ for every $\nu \in [0, 1]$, set $\nu^\star = 1$.
  \item Otherwise, $\nu^\star \in (0, 1)$ is the unique value with $\alpha(\nu^\star) = 0$.
\end{itemize}

\begin{proposition}[Two regimes]\label{prop:ma-two-regimes}
\begin{enumerate}
  \item[\textup{(i)}] \emph{Low correlation} ($\nu < \nu^\star$): $\alpha > 0$, both firms invest at their borrowing capacity $I_i = k A_i$.
  The firms invest independently, but each firm derives a strictly positive externality from the other's presence through the asset-redeployability channel (the $\kappa I_j$ term).
  \item[\textup{(ii)}] \emph{High correlation} ($\nu > \nu^\star$): $\alpha < 0$, no firm invests in any Nash equilibrium where the other does not.
  Investment is depressed even though it could be collectively profitable: $\alpha + \kappa$ may still be positive.
\end{enumerate}
\end{proposition}

\begin{proof}
Part (i) is immediate from Proposition~\ref{prop:ma-multiplier}.
For (ii), if $\alpha < 0$, the NPV \eqref{eq:ma-NPV} is decreasing in $I_i$, so the entrepreneur optimally sets $I_i = 0$.
Given this, the other firm's NPV contribution from $\kappa I_j = 0$ vanishes, sustaining the no-investment Nash equilibrium.
\end{proof}

The high-correlation regime features a striking inefficiency: both firms remain on the sideline even though joint investment would generate positive surplus (each firm's own $\alpha$ is negative, but the cross-firm externality $\kappa$ could make $\alpha + \kappa > 0$).
This sets up the coordination problem analyzed next.

%====================================================================
\section{Coordinated Equilibrium and Multiple Equilibria}
\label{sec:ma-coordination}
%====================================================================

The high-correlation regime of Proposition~\ref{prop:ma-two-regimes} assumed that firms invest independently: each takes the other's investment as given.
Suppose instead that firms can coordinate.
Then each firm may be willing to invest precisely because the other firm invests, even when independent investment would be unprofitable.

To capture this coordination motive starkly, suppose firm $i$ can operate firm $j$'s redeployed assets only if firm $i$ has itself invested at least a minimum threshold $\underline{I}$:
\begin{equation}\label{eq:ma-threshold}
  I_i \geq \underline{I},
\end{equation}
where $\underline{I}$ is exogenous, with $\underline{I} \leq k A_i$ (so the threshold is compatible with firm $i$'s borrowing capacity).

\subsection{Modified payoff}

Under \eqref{eq:ma-threshold}, entrepreneur $i$'s expected payoff is
\begin{equation}\label{eq:ma-Pi-coord}
  \Pi_{b i}(I_i, R_{b i}, I_j) = [x \rho_1 I_i - I_i] + [(1 - x)(1 - \nu) \rho_0 I_i] + \left[x(1 - \mu) \frac{p_h B}{\Delta p} I_j\right] \mathbb{1}_{\{I_i \geq \underline{I}\}},
\end{equation}
where the indicator turns on the cross-firm redeployability benefit only if firm $i$ has invested enough to absorb firm $j$'s assets.

\subsection{The coordinated equilibrium}

\begin{proposition}[Coordination equilibrium]\label{prop:ma-coord}
There exists $\nu^{\star\star} > \nu^\star$ with $\alpha(\nu^{\star\star}) + \kappa(\nu^{\star\star}) = 0$ such that:
\begin{itemize}
  \item For $\nu \in (\nu^\star, \nu^{\star\star}]$, \emph{two} pure-strategy Nash equilibria exist:
  \begin{itemize}
    \item the \emph{coordinated} (``good'') equilibrium: both firms invest $\underline{I}$, taking advantage of the redeployability option;
    \item the \emph{non-coordinated} (``bad'') equilibrium: neither firm invests.
  \end{itemize}
  \item For $\nu > \nu^{\star\star}$, only the no-investment equilibrium exists.
\end{itemize}
\end{proposition}

\begin{proof}[Proof sketch]
In the no-investment equilibrium, $I_j = 0$ makes the $\kappa I_j$ term vanish, and with $\alpha < 0$, firm $i$'s best response is $I_i = 0$.
This is a Nash equilibrium for any $\nu > \nu^\star$.

In the coordinated equilibrium, firm $i$ invests $\underline{I}$ if the payoff from doing so (with $I_j = \underline{I}$) exceeds zero:
\begin{equation*}
  \alpha \underline{I} + \kappa \underline{I} = (\alpha + \kappa) \underline{I} \geq 0,
\end{equation*}
which holds iff $\alpha + \kappa \geq 0$, i.e., $\nu \leq \nu^{\star\star}$.
\end{proof}

The structural parallel with the soft-budget-constraint problem of Chapter~\ref{ch:liquidity-scale} is worth noting.
In both cases, ex-post rationality induces an inefficient equilibrium that can be avoided through coordination or commitment.
Here, the coordination failure arises across \emph{firms} rather than within a single firm's intertemporal structure, but the underlying mechanism is analogous: each firm's investment decision creates a positive externality on the other (through the redeployability channel), but this externality is not internalized in the decentralized equilibrium.

%====================================================================
\section{Underdeveloped Resale Markets}
\label{sec:ma-resale}
%====================================================================

We now turn to a complementary question: the effect of \emph{ex-ante design choices} on the availability of asset redeployment.
Firms can choose between risky and safe investment configurations; the choice affects both the firm's own distress probability and the availability of redeployment partners for other firms.
We show that lack of coordination can lead to inefficiently low volumes of transactions in the resale market.

\subsection{Setup with design choice}

We specialize to the following environment:
\begin{itemize}
  \item \emph{Nonconcurrent risks}: $\nu = 0$ (and $\mu = (2x - 1)/x$), so the two firms are never simultaneously in distress.
  \item \emph{Symmetry}: $A_i = A_j \eqqcolon A$ and $I_i = I_j \eqqcolon I$.
  \item \emph{Ex-ante design choice}: each firm chooses between two configurations.
  \begin{itemize}
    \item The \emph{risky} design: per-unit cost $1$, productive with probability $x$.
    \item The \emph{safe} design: per-unit cost $X > 1$, never in distress (probability of being productive is $1$).
  \end{itemize}
\end{itemize}
The safe design pays a premium of $X - 1$ per unit to eliminate distress risk; the risky design is cheaper but bears the risk of becoming unable to operate its own assets.

The bargaining protocol is unchanged: the acquired firm's lender (in the rare event of distress) holds the bargaining power, setting the transfer price at $\rho_0$ per unit.

\subsection{Coordinated risky design}

If both firms coordinate on the risky design, assets are always productive for at least one of the firms (by the nonconcurrent-risks assumption).
There is no risk that the assets end up unused, and the investment cost is at its lowest.

The entrepreneur's payoff at the symmetric equilibrium is
\begin{equation*}
  \Pi_{b, \text{coor}}^{\text{risk}} = (\rho_1 - 1) I,
\end{equation*}
where $I$ is determined by the breakeven condition
\begin{equation*}
  \rho_0 I = I - A,
  \quad I = \frac{A}{1 - \rho_0}.
\end{equation*}
Substituting,
\begin{equation}\label{eq:ma-coord-risk}
  \Pi_{b, \text{coor}}^{\text{risk}} = \frac{\rho_1 - 1}{1 - \rho_0} A.
\end{equation}

\subsection{Coordinated safe design}

If both firms coordinate on the safe design, no distress occurs and no asset redeployment is ever needed.
The breakeven condition changes because the per-unit investment cost is $X$:
\begin{equation*}
  \rho_0 I = X I - A,
  \quad I = \frac{A}{X - \rho_0}.
\end{equation*}
The entrepreneur's payoff is
\begin{equation}\label{eq:ma-coord-safe}
  \Pi_{b, \text{coor}}^{\text{safe}} = \frac{\rho_1 - X}{X - \rho_0} A.
\end{equation}

\begin{proposition}[Risky dominates safe under coordination]\label{prop:ma-coord-comparison}
Under coordination, the risky design strictly dominates the safe design:
\begin{equation*}
  \Pi_{b, \text{coor}}^{\text{risk}} > \Pi_{b, \text{coor}}^{\text{safe}}.
\end{equation*}
\end{proposition}

\begin{proof}
The inequality reduces to
\begin{equation*}
  \frac{\rho_1 - 1}{1 - \rho_0} > \frac{\rho_1 - X}{X - \rho_0},
\end{equation*}
equivalently (cross-multiplying, valid since $1 - \rho_0 > 0$ and $X - \rho_0 > 0$),
\begin{equation*}
  (\rho_1 - 1)(X - \rho_0) > (\rho_1 - X)(1 - \rho_0),
\end{equation*}
which expands to $(X - 1)(\rho_1 - \rho_0) > 0$, true since $X > 1$ and $\rho_1 > \rho_0$.
\end{proof}

The intuition is that the safe design's premium $X - 1$ per unit is not justified under coordination: the nonconcurrent-risks assumption guarantees that at least one firm always has productive assets, so the redeployment option from the risky design is always realized.
Paying a premium to eliminate a risk that is already hedged by the partner firm is wasteful.

A similar argument (omitted) shows that \emph{asymmetric} design choices (one firm risky, one safe) are also dominated by symmetric risky coordination.
Symmetric risky coordination is the first-best outcome.

\subsection{Non-coordinated equilibrium}

Suppose firms cannot coordinate ex ante on investment design.
A natural question: is the risky design still the unique Nash equilibrium?
Perhaps surprisingly, the answer is no.

Assume that in equilibrium both entrepreneurs adopt the \emph{safe} design and earn the coordinated-safe payoff $\Pi_b^{\text{safe}} = (\rho_1 - X)/(X - \rho_0) \cdot A$.
Consider a unilateral deviation by firm $i$ to the risky design.
The deviating firm's payoff becomes
\begin{equation*}
  \Pi_b^{\text{risk-dev}} = [x \rho_1 + (1 - x) \rho_0 - 1] I,
\end{equation*}
where the investment $I$ is pinned down by the lender's breakeven condition $\rho_0 I = I - A$ (the deviating firm has no redeployment partner, since firm $j$ is safe and never in distress, hence never provides asset redeployment opportunities; but the deviating firm's own assets could be redeployed to firm $j$ in distress, which never occurs, so lender $i$'s salvage value is zero).

Substituting $I = A/(1 - \rho_0)$:
\begin{equation}\label{eq:ma-risk-dev}
  \Pi_b^{\text{risk-dev}} = \frac{[\rho_1 - (1 - x)(\rho_1 - \rho_0)] - 1}{1 - \rho_0} A.
\end{equation}

\subsection{When safe is individually rational}

\begin{proposition}[Safe Nash equilibrium]\label{prop:ma-safe-NE}
The symmetric safe design is a Nash equilibrium if and only if
\begin{equation}\label{eq:ma-safe-condition}
  (1 - x)(X - \rho_0) > X - 1.
\end{equation}
\end{proposition}

\begin{proof}
The safe design is a Nash equilibrium iff $\Pi_b^{\text{risk-dev}} < \Pi_b^{\text{safe}}$:
\begin{equation*}
  \frac{[\rho_1 - (1 - x)(\rho_1 - \rho_0)] - 1}{1 - \rho_0} < \frac{\rho_1 - X}{X - \rho_0}.
\end{equation*}
Cross-multiplying and simplifying (tedious but direct algebra) yields \eqref{eq:ma-safe-condition}.
\end{proof}

Condition \eqref{eq:ma-safe-condition} holds whenever the distress probability $1 - x$ is large enough: the deviating firm's expected loss from becoming distressed (proportional to $X - \rho_0$) exceeds the safe-design cost premium $X - 1$.

\subsection{The inefficiency of underdeveloped resale markets}

Propositions~\ref{prop:ma-coord-comparison} and~\ref{prop:ma-safe-NE} together diagnose a coordination failure: under condition \eqref{eq:ma-safe-condition}, the safe design is a Nash equilibrium even though both firms strictly prefer the coordinated risky outcome.
The economic intuition is subtle:

\begin{itemize}
  \item If firm $j$ chooses safe, firm $j$ is never in distress, so firm $i$ can never acquire firm $j$'s assets.
  Firm $i$'s redeployability option is worthless.
  In isolation, the risky design exposes firm $i$ to distress risk without any offsetting benefit.
  \item If firm $j$ also chose risky, firm $i$ would have a redeployment opportunity whenever firm $i$ is productive and firm $j$ is distressed---an event that occurs with probability $x(1 - \mu) = x \cdot (x - (2x - 1))/x = 1 - x$ (under nonconcurrent risks).
  This redeployment opportunity makes the risky design attractive.
\end{itemize}

The firm's incentive to choose the risky design therefore depends on the other firm's choice.
When the other firm chooses safe, each firm is individually rational to also choose safe, leading to a low-trade resale-market equilibrium.
This is precisely what \citet{ShleiferVishny1992} call ``underdeveloped resale markets'': the secondary asset market's activity depends on industry-wide coordination on risky investment designs, and the lack of such coordination leads to inefficiently low transaction volumes.

%====================================================================
\section{Further Reading}
\label{sec:ma-further}
%====================================================================

The industry-equilibrium framework developed in this chapter is due to \citet{ShleiferVishny1992}, whose analysis of liquidation values and debt capacity remains the canonical reference for the interaction between asset redeployability and firm-level financial constraints.
The specific moral-hazard formalization of the transfer-price mechanism, yielding the competitive-price result $P = \rho_0$, is developed in \citet{Tirole2006} (Chapter~8), on which this chapter is directly based.

\citet{PerottiSuarez2002} extend the framework to banking, asking how the failure of one bank affects the financial health and market power of its surviving competitors.
Their ``last bank standing'' mechanism---in which the surviving bank benefits from reduced competition at the cost of a riskier industry---provides an important application of the industry-equilibrium perspective.

The broader macroeconomic literature on credit cycles and financial accelerators, beginning with \citet{KiyotakiMoore1997}, shares the analytical DNA of the industry-equilibrium framework.
In those models, the feedback between asset prices and borrowing capacity operates through a dynamic mechanism in which firms' collateral values depend on the state of the industry, which in turn depends on firms' investment decisions.
The static two-firm model of this chapter can be seen as the simplest building block for these richer dynamic theories.

The coordination failures analyzed in Section~\ref{sec:ma-coordination} and the underdeveloped-resale-market result of Section~\ref{sec:ma-resale} resonate with a broader theme in the macro-finance literature: secondary markets for specialized assets (corporate bonds, real estate, specialized equipment) exhibit endogenous illiquidity that depends on the collective investment decisions of industry participants.
The modern literature on fire sales and asset-price contagion---including episodes such as the 2008 financial crisis and the COVID-19-related market disruptions of 2020---has enriched and deepened these insights.

%%%%%%%%%%%%%%%%%%%%%%%%%%%%%%%%%%%%%%%%%%%%%%%%%%%%%%%%%%%%%%%%%%%%%%%%

\part{Lending Relationships and Investor Activism}

%%%%%%%%%%%%%%%%%%%%%%%%%%%%%%%%%%%%%%%%%%%%%%%%%%%%%%%%%%%%%%%%%%%%%%%%

%====================================================================
\chapter{Monitoring and Lending Relationships}
\label{ch:monitoring}
%====================================================================

The models developed in Parts~II through~IV have treated lenders as essentially \emph{passive} agents: they provide capital at date~$0$, observe the verifiable income at date~$1$ (or later), and receive the contractual repayment.
They do not participate in the firm's operations, do not acquire information beyond what is publicly observable, and do not exert any influence on the entrepreneur's effort choice beyond what the terms of the contract already imply.
In practice, many important classes of lenders operate quite differently.
Banks, venture capitalists, private equity firms, and large blockholders all engage in \emph{monitoring}: they actively collect information about the firms they finance, advise management, discipline poor performance, and in some cases replace management outright.

This chapter extends the corporate-finance framework to incorporate \emph{active monitors}---agents who incur a private cost to reduce the entrepreneur's scope for moral hazard.
The analysis delivers two complementary sets of insights, corresponding to the two parts of the chapter.

\emph{Part A} (Sections~\ref{sec:mon-model}--\ref{sec:mon-optimal}) develops the static theory of monitoring.
A monitor reduces the entrepreneur's private benefit from shirking from $B$ to $b < B$, at a private cost $c > 0$.
For monitoring to be worthwhile, it must generate more pledgeable income (the reduction in agency rent) than it consumes in direct cost.
The analysis yields a sharp partition of entrepreneurs by balance-sheet strength into three regimes: strong balance sheets borrow directly and cheaply; medium balance sheets borrow through monitored finance at a discount; weak balance sheets cannot be financed at all.
The thresholds $\underline{A}$ and $\overline{A}$ delineating these regimes formalize the distinction between ``high-quality'' borrowers who can tap the bond market and ``intermediate-quality'' borrowers who rely on bank lending---a pattern amply documented in the empirical literature.

\emph{Part B} (Sections~\ref{sec:mon-learning}--\ref{sec:mon-date1}) turns to the dynamic dimension.
Over time, a monitor accumulates information about the firms it finances.
In a two-period setting with second-period productivity uncertainty, the incumbent monitor---who has observed the firm's first-period performance---acquires superior information at the refinancing stage.
Entrants bidding to replace the incumbent face a classical \emph{winner's curse}: they win only when the firm is of lower-than-average quality.
This informational friction generates an \emph{incumbency rent} captured by the monitor at date~$2$.
Remarkably, however, competition among prospective monitors at date~$1$ fully dissipates the anticipated rent: date-1 monitors offer ``generous'' contributions, contributing more than their static breakeven, knowing they will recoup the gap through future rents.
The equilibrium is first-best from the entrepreneur's perspective, notwithstanding the intertemporal rent transfer---``the blockholding is initially acquired at a premium, and is later maintained at a discount.''

The chapter corresponds to Chapter~9 of \citet{Tirole2006}.
The classical reference for the moral-hazard-with-monitoring framework is \citet{HolmstromTirole1997}; the learning-by-lending analysis builds on the subsequent literature on relationship banking and switching costs.

%====================================================================
\section{The Monitoring Model}
\label{sec:mon-model}
%====================================================================

\subsection{Setup}

We begin with the fixed-investment framework of Chapter~\ref{ch:net-worth}.
A risk-neutral entrepreneur with wealth $A < I$ needs to raise $I - A$ from outside investors to finance a project costing $I$.
The project yields verifiable income $R > 0$ in success and zero in failure.
The moral-hazard stage is the standard one: if the entrepreneur \emph{behaves}, the probability of success is $p_h$ and he gets no private benefit; if he \emph{misbehaves}, the probability is $p_\ell \coloneqq p_h - \Delta p$ and he gets a private benefit.

The novelty of this chapter is that the private benefit from misbehaving can depend on the presence of a monitor.
We introduce a third agent: the \emph{monitor}, who can incur an unobservable private cost $c > 0$ to reduce the scope for moral hazard.

\subsection{The three-project interpretation}

The role of the monitor is most transparent in the \emph{three-project} interpretation of the moral-hazard stage.
The entrepreneur has access to three potential projects:

\begin{center}
\renewcommand{\arraystretch}{1.3}
\begin{tabular}{|l|c|c|c|}
\hline
 & \textit{good} & \textit{bad} & \textit{Bad} \\
\hline
Probability of success & $p_h$ & $p_\ell$ & $p_\ell$ \\
Private benefit & $0$ & $b$ & $B$ \\
\hline
\end{tabular}
\end{center}
with $b < B$.
The good project is the socially productive one; the two moral-hazard alternatives deliver the same (low) success probability $p_\ell$ but differ in the private benefit they yield to the entrepreneur.

The monitor's role is to identify and block the Bad project.
If the monitor incurs cost $c$, he observes whether the entrepreneur is about to choose the high-private-benefit Bad project and can prevent it.
The monitor cannot, however, distinguish between the good and bad projects: those remain observationally equivalent.
If the monitor does \emph{not} incur cost $c$, all three projects are available, and the entrepreneur will choose the Bad project (which dominates the bad project from his perspective).

\subsection{No-monitoring benchmark}

If the firm operates without a monitor, the entrepreneur faces the moral-hazard choice between the good project and the Bad project.
Incentive compatibility requires
\begin{equation}\label{eq:mon-IC-nomon}
  (\Delta p) R_b \geq B.
\end{equation}
Funding is feasible if and only if the pledgeable income covers the outside financing need:
\begin{equation}\label{eq:mon-pledge-nomon}
  \rho_0 \coloneqq p_h \left( R - \frac{B}{\Delta p} \right) \geq I - A.
\end{equation}
When this condition holds, competitive lending drives the entrepreneur's net payoff to the project's NPV, $\Pi_b = p_h R - I$.

\subsection{With monitoring}

If a monitor is present and incurs cost $c$, the Bad project is blocked.
The entrepreneur's choice is between the good and bad projects, and incentive compatibility now requires only
\begin{equation}\label{eq:mon-IC-mon}
  (\Delta p) R_b \geq b.
\end{equation}
The monitor himself requires an incentive to incur the cost $c$.
Letting $R_m$ denote the monitor's reward in success (and zero in failure, by limited liability), the monitor's incentive constraint is
\begin{equation}\label{eq:mon-IC-monitor}
  (\Delta p) R_m \geq c.
\end{equation}
The constraint \eqref{eq:mon-IC-monitor} parallels the entrepreneur's: a successful monitor is rewarded more in expectation than an unsuccessful one, and the reward differential must cover the cost $c$ of becoming informed.

\begin{remark}[When monitoring is nontrivial]\label{rem:mon-nontrivial}
We maintain throughout that
\begin{equation*}
  (\Delta p) R_b < B,
\end{equation*}
i.e., the entrepreneur's stake is too small to support incentives in the \emph{absence} of monitoring (since otherwise the Bad project would be self-deterred and monitoring would be redundant).
This condition holds when the entrepreneur's balance sheet is weak enough that direct financing under \eqref{eq:mon-pledge-nomon} is not feasible.
\end{remark}

%====================================================================
\section{Equilibrium with Competitive Monitoring Capital}
\label{sec:mon-equilibrium}
%====================================================================

\subsection{Active monitoring capital}

Assume that there is a large supply of monitors willing to participate in any project at a breakeven return.
A monitor contracted to exert effort at cost $c$ with reward $R_m$ in success will be willing to contribute an investment $I_m$ at date $0$ satisfying
\begin{equation}\label{eq:mon-breakeven}
  p_h R_m - c = I_m.
\end{equation}
Setting the monitor's stake at the incentive-compatibility minimum $R_m = c/\Delta p$, the monitor's contribution becomes
\begin{equation}\label{eq:mon-Im}
  I_m = p_h \cdot \frac{c}{\Delta p} - c = \frac{p_\ell c}{\Delta p}.
\end{equation}
The monitor's net expected rent is zero: he contributes $I_m$, and recovers $p_h R_m - c = I_m$ in expectation.
The monitoring cost $c$ is covered exactly by the incentive-compatible portion of his reward.

\subsection{Non-monitoring investors}

The residual financing $I - A - I_m$ is provided by non-monitoring (passive) investors.
Their breakeven condition is
\begin{equation}\label{eq:mon-nonmon-IR}
  p_h(R - R_b - R_m) \geq I - A - I_m.
\end{equation}
Substituting $R_b = b/\Delta p$, $R_m = c/\Delta p$, and $I_m = p_\ell c/\Delta p$, and simplifying (the algebra uses $p_h - p_\ell = \Delta p$), \eqref{eq:mon-nonmon-IR} reduces to
\begin{equation}\label{eq:mon-feasibility}
  \rho_0^{m} \coloneqq p_h\left( R - \frac{b}{\Delta p} \right) - c \geq I - A.
\end{equation}

The quantity $\rho_0^m$ is the \emph{pledgeable income under monitoring}: it parallels $\rho_0$ from the unmonitored model but with the agency rent $p_h B/\Delta p$ replaced by the sum of the reduced agency rent $p_h b/\Delta p$ and the monitoring cost $c$.

\subsection{Monitoring-enhanced pledgeable income}

The central economic question is when $\rho_0^m > \rho_0$---i.e., when monitoring increases the pledgeable income that can be offered to outside investors.

\begin{assumption}[Monitoring-enhanced pledgeability]\label{ass:mon-enhanced}
\begin{equation}\label{eq:mon-enhanced}
  p_h \cdot \frac{b}{\Delta p} + c < p_h \cdot \frac{B}{\Delta p},
  \quad \text{equivalently,} \quad
  c < \frac{p_h (B - b)}{\Delta p}.
\end{equation}
\end{assumption}

The assumption says that the monitoring cost $c$ is smaller than the reduction in agency rent achieved by monitoring.
Equivalently, $\rho_0^m > \rho_0$: monitoring strictly expands the pledgeable income.
Without this assumption, monitoring is wasteful: the cost of reducing the agency rent exceeds the savings.

\subsection{The entrepreneur's payoff}

Under competitive monitoring capital, the monitor captures no rent; under competitive non-monitoring capital, the non-monitoring investor also captures no rent.
The entrepreneur therefore receives the full expected social surplus net of monitoring cost:
\begin{equation}\label{eq:mon-Pi-monitored}
  \Pi_b = p_h R - I - c.
\end{equation}
We maintain the assumption that this NPV is positive:
\begin{equation}\label{eq:mon-NPV-positive}
  p_h R > I + c.
\end{equation}

The comparison with the unmonitored benchmark is immediate: monitoring \emph{reduces} the entrepreneur's payoff by the monitoring cost $c$.
The entrepreneur would prefer to borrow without monitoring if he could.

%====================================================================
\section{Balance-Sheet Thresholds and Borrowing Regimes}
\label{sec:mon-thresholds}
%====================================================================

The central structural result of Part~A of this chapter is that the set of financeable firms is partitioned by balance-sheet strength into three regimes, separated by two critical thresholds.

\subsection{Definitions}

Define the \emph{high threshold}
\begin{equation}\label{eq:mon-Aupper}
  \overline{A} \coloneqq I - p_h\left( R - \frac{B}{\Delta p} \right) = I - \rho_0.
\end{equation}
This is the minimum wealth under which unmonitored financing is feasible: the pledgeable income $\rho_0$ must cover the outside financing need $I - A$.

Define the \emph{low threshold}
\begin{equation}\label{eq:mon-Alower}
  \underline{A} \coloneqq I + c - p_h\left( R - \frac{b}{\Delta p} \right) = I - \rho_0^m.
\end{equation}
This is the minimum wealth under which monitored financing is feasible.

Under Assumption~\ref{ass:mon-enhanced}, $\rho_0^m > \rho_0$ and therefore
\begin{equation*}
  \underline{A} < \overline{A}.
\end{equation*}

\subsection{Three regimes}

\begin{proposition}[Balance-sheet regimes]\label{prop:mon-regimes}
Under the maintained assumptions, the financing outcome for an entrepreneur with wealth $A$ is:
\begin{enumerate}
  \item[\textup{(i)}] \emph{Strong balance sheet} ($A \geq \overline{A}$): the project is financed without monitoring.
  The entrepreneur's payoff is $\Pi_b = p_h R - I$ (first best).
  \item[\textup{(ii)}] \emph{Intermediate balance sheet} ($\underline{A} \leq A < \overline{A}$): the project is financed with monitoring.
  The entrepreneur's payoff is $\Pi_b = p_h R - I - c$ (reduced by the monitoring cost).
  \item[\textup{(iii)}] \emph{Weak balance sheet} ($A < \underline{A}$): the project is not financed.
  The entrepreneur's payoff is zero.
\end{enumerate}
\end{proposition}

\begin{proof}
Case (i) follows from \eqref{eq:mon-pledge-nomon}.
Given feasibility, the entrepreneur strictly prefers unmonitored financing (higher payoff by $c$), so he opts out of monitoring.
Case (ii) follows from \eqref{eq:mon-feasibility} and the failure of \eqref{eq:mon-pledge-nomon}.
Case (iii) follows from the failure of both \eqref{eq:mon-feasibility} and \eqref{eq:mon-pledge-nomon}.
\end{proof}

\subsection{Economic interpretation}

The regime structure of Proposition~\ref{prop:mon-regimes} matches a striking empirical regularity in the corporate-finance literature.
Firms with strong balance sheets---large, publicly traded corporations with substantial retained earnings---tap bond markets directly and at competitive rates.
Firms with intermediate balance sheets---smaller corporations, start-ups, and many privately held firms---rely on bank lending, private equity, or similar monitored intermediation.
Firms with very weak balance sheets are credit-rationed: they simply cannot obtain external financing, regardless of the project's NPV.

The empirical evidence on the role of banks in reducing agency costs is substantial.
\citet{James1987} documents that announcements of bank loan agreements are associated with positive abnormal returns, suggesting that bank certification conveys value-relevant information to the market.
\citet{LummerMcConnell1989} refine this finding, showing that it is renewed loans (where the bank has accumulated monitoring information) that drive the positive announcement effects.
\citet{CantilloWright2000} document the cross-sectional pattern of Proposition~\ref{prop:mon-regimes}: intermediate-quality firms choose bank lending, while high-quality firms tap bond markets.

The internal structure of the monitoring entity matters for its effectiveness.
\citet{BergerMillerPetersenRajanStein2005} document that smaller banks have a comparative advantage in relationship-based (``soft information'') lending, while larger banks rely more on arms-length (``hard information'') criteria.

%====================================================================
\section{Optimal Monitoring Intensity}
\label{sec:mon-optimal}
%====================================================================

We now extend the model to endogenize the \emph{intensity} of monitoring.
Instead of a binary choice (monitor or not), the monitor can choose an intensity $x \in [0, 1]$: with probability $x$, he successfully identifies the Bad project; with probability $1 - x$, he learns nothing and the Bad project remains available.

\subsection{Convex cost of monitoring}

The monitor's effort cost is $c(x)$, strictly increasing and convex in $x$, with boundary properties
\begin{equation}\label{eq:mon-boundary}
  \lim_{x \to 0^+} \frac{c(x) - c(0)}{x} = 0
  \et
  \lim_{x \to 1^-} c'(x) = \infty.
\end{equation}
The first condition ensures that small amounts of monitoring are inexpensive at the margin; the second ensures that perfect monitoring is infeasibly costly.
Together, they guarantee an interior optimum.

\subsection{The social NPV}

Assuming the entrepreneur's stake satisfies $R_b < B/\Delta p$ (so that monitoring is necessary---otherwise incentives would be maintained by the entrepreneur's own stake), the entrepreneur's effective ex-ante payoff is
\begin{equation}\label{eq:mon-NPV-variable}
  \Pi_b(x) = x p_h R + (1 - x)(p_\ell R + B) - I - c(x).
\end{equation}
The first term reflects the monitoring-successful case (prob.\ $x$): the entrepreneur takes the good project, producing expected verifiable income $p_h R$.
The second term reflects the monitoring-failed case (prob.\ $1 - x$): the entrepreneur takes the Bad project, producing expected verifiable income $p_\ell R$ plus private benefit $B$.
The third term is the direct investment cost; the fourth is the monitor's unverifiable effort cost.

Differentiating in $x$:
\begin{equation}\label{eq:mon-FOC}
  \Pi_b'(x) = p_h R - (p_\ell R + B) - c'(x) = (\Delta p) R - B - c'(x).
\end{equation}
The optimal intensity $x^\star$ satisfies the first-order condition
\begin{equation}\label{eq:mon-FOC-star}
  c'(x^\star) = (\Delta p) R - B,
\end{equation}
and is interior, $x^\star \in (0, 1)$, under the boundary properties of $c$.

\subsection{Incentive-compatible monitoring contract}

The monitoring intensity is \emph{unobservable}: the monitor's effort choice must be induced by the contract.
Given a reward $R_m$ in success, the monitor chooses
\begin{equation*}
  x \in \argmax_{x \in [0,1]} \, \bigl\{[x p_h + (1 - x) p_\ell] R_m - c(x)\bigr\},
\end{equation*}
with first-order condition
\begin{equation}\label{eq:mon-monitor-FOC}
  (\Delta p) R_m = c'(x).
\end{equation}
For the contract to implement the first-best intensity $x^\star$, the stake $R_m$ must satisfy
\begin{equation}\label{eq:mon-Rm-star}
  (\Delta p) R_m = (\Delta p) R - B
  \quad\Longleftrightarrow\quad
  R_m = R - \frac{B}{\Delta p}.
\end{equation}

\begin{proposition}[Optimal monitoring contract]\label{prop:mon-optimal}
The optimal monitoring stake is $R_m^\star = R - B/\Delta p$, which simultaneously:
\begin{itemize}
  \item satisfies the monitor's incentive constraint \eqref{eq:mon-monitor-FOC} at the NPV-maximizing intensity $x^\star$ from \eqref{eq:mon-FOC-star};
  \item leaves a strictly positive residual $R - R_b - R_m^\star > 0$ for non-monitoring investors, since $R_b < B/\Delta p$ by assumption.
\end{itemize}
\end{proposition}

\subsection{The Pagano--Röell insight}

A subtle implication of Proposition~\ref{prop:mon-optimal} deserves emphasis.
Under the optimal contract, the monitor's stake $R_m^\star = R - B/\Delta p$ is strictly less than the total external stake $R - R_b$:
\begin{equation*}
  R_m^\star = R - \frac{B}{\Delta p} < R - R_b,
\end{equation*}
since $R_b < B/\Delta p$.
The residual $(R - R_b) - R_m^\star > 0$ must go to \emph{non-monitoring} (external) investors.
The monitor cannot hold all the external shares in the firm.

The economic logic is intuitive.
If the monitor held all external shares, his stake $R_m$ would be too large (exceeding $R - B/\Delta p$); the monitor would over-incentivize monitoring effort beyond the NPV-maximizing level.
Conversely, if the monitor held too few shares, he would under-incentivize.
The optimal ownership structure features a mixed arrangement: the monitor holds a moderate stake calibrated to the FOC \eqref{eq:mon-monitor-FOC}, with the residual held by passive investors.
This is the classical ownership-structure insight of \citet{PaganoRoell1998}, who connect the agency foundations of ownership dispersion to the decision to go public.

\subsection{Feasibility of variable-intensity monitoring}

For the variable-intensity contract to be feasible, we must have sufficient pledgeable income at the optimal intensity:
\begin{equation}\label{eq:mon-feasibility-variable}
  [x^\star p_h + (1 - x^\star) p_\ell]\left\{R - \frac{b}{\Delta p}\right\} \geq I - A + c(x^\star),
\end{equation}
while still requiring that unmonitored financing is infeasible ($\rho_0 < I - A$), so that monitoring is indeed necessary.
Under these conditions, the variable-intensity contract dominates the fixed-intensity contract of Section~\ref{sec:mon-equilibrium} by Jensen's inequality-type arguments: the convexity of $c$ and the concavity of the payoff function ensure that variable intensity generates more pledgeable income per unit of monitoring cost.

%====================================================================
\section{Learning by Lending: Two-Period Setup}
\label{sec:mon-learning}
%====================================================================

We now turn to Part~B of the chapter: the dynamic dimension of lending relationships.
The key economic force is that a monitor who finances the firm at date~$1$ accumulates information about it, gaining an \emph{informational advantage} over potential entrants at date~$2$.
This advantage generates a rent that, as we shall see, is fully competed away at date~$1$ through generous initial offers.

\subsection{Two-period environment}

The economy has two dates, $t \in \{1, 2\}$, with a common discount factor $\beta$ between periods.
Saving is impossible between the two dates (ruling out intertemporal consumption smoothing).
The entrepreneur has no initial wealth, $A = 0$.

\subsection{Date 1: standard setup}

At date~$1$, the entrepreneur has a project requiring investment $I$.
The project succeeds with probability $p$ (yielding $R$) and fails with probability $1 - p$.
The three-project moral-hazard structure of Section~\ref{sec:mon-model} applies: $p \in \{p_h, p_\ell\}$, with private benefits $0$, $b$, $B$ under good, bad, and Bad projects respectively.
The private benefit $B$ is large enough that unmonitored financing is infeasible:
\begin{equation}\label{eq:mon-date1-infeasible}
  \rho_0 = p_h \left( R - \frac{B}{\Delta p} \right) < I.
\end{equation}
Monitored financing is feasible:
\begin{equation}\label{eq:mon-date1-feasible}
  p_h \left( R - \frac{b}{\Delta p} \right) \geq I + c.
\end{equation}
Under these primitives, the date-1 project is financed via the monitoring contract of Sections~\ref{sec:mon-model}--\ref{sec:mon-thresholds}, with monitor reward $R_m^1 = c/\Delta p$ and monitor contribution $I_m^1$ to be determined.

\subsection{Date 2: productivity shock}

Regardless of the date-1 outcome, the entrepreneur has a new project at date~$2$, which can be thought of as a continuation of the original firm.
The date-2 project is identical to the date-1 project in all respects except one: \emph{the probability of success is potentially higher}.
Specifically:
\begin{itemize}
  \item with probability $\alpha \in (0, 1)$, the date-2 success probability is uniformly shifted up by $\tau > 0$: the three projects now have success probabilities $p_h + \tau$, $p_\ell + \tau$, $p_\ell + \tau$;
  \item with probability $1 - \alpha$, the probabilities remain at $p_h$, $p_\ell$, $p_\ell$.
\end{itemize}
The date-2 productivity realization is statistically independent of the date-1 outcome.
The private benefits $b$ and $B$ are unchanged; in particular, $B$ remains so large that a monitor is still needed at date~$2$ even under the improved productivity.

\subsection{No commitment}

We assume that \emph{no intertemporal commitment} is possible: date-1 contracts specify only short-term (date-1) claims.
The date-1 monitor's return comes exclusively from the date-1 project.
The date-2 financing is arranged anew at date~$2$, with new claims issued to potentially new investors and monitors.
This assumption rules out long-term contracts that would lock in the initial monitor and avoid the subsequent bidding game.

%====================================================================
\section{Date-2 Bidding and the Winner's Curse}
\label{sec:mon-date2}
%====================================================================

\subsection{The symmetric-information benchmark}

As a benchmark, consider first the case where \emph{no agent} learns the realized date-2 productivity.
All agents---the entrepreneur, incumbent monitor, and entrants---share the same prior belief: date-2 success probability is $p_h + \tau$ with probability $\alpha$ and $p_h$ otherwise, for expected probability $p_h + \alpha \tau$ under behaving.

Under competition, the monitor captures no rent, and the entrepreneur's date-2 net payoff equals the (expected) NPV:
\begin{equation}\label{eq:mon-sym-date2}
  \Pi_b^{\text{sym},2} = (p_h + \alpha \tau) R - I - c.
\end{equation}
The total two-period payoff under symmetric information is
\begin{equation}\label{eq:mon-sym-total}
  \Pi_b^{\text{sym}} = \bigl[p_h R - (I + c)\bigr] + \beta \bigl[(p_h + \alpha \tau) R - (I + c)\bigr].
\end{equation}

\subsection{Asymmetric information: the incumbent's advantage}

We now assume that the date-1 monitor---the \emph{incumbent}---privately learns the realized date-2 productivity as a by-product of his ongoing monitoring relationship with the firm.
The entrepreneur and the prospective entrants (other monitors and passive investors) remain uninformed.
This is the simplest stylization of the \emph{learning-by-lending} phenomenon: repeated interaction with the firm reveals information that arm's-length lenders cannot obtain.

\subsection{The bidding game}

The entrepreneur conducts the following bidding game at date~$2$:
\begin{enumerate}
  \item The entrepreneur posts a monitoring stake $R_m^2 = c/\Delta p$ (the minimum incentive-compatible level).
  \item Entrants submit investment-contribution offers $I_m^2$, committing to provide $I_m^2$ in exchange for the stake $R_m^2$ if chosen.
  \item The incumbent then decides whether to match the highest entrant bid (slightly exceeding if necessary) and retain the monitoring role, or let an entrant take over.
  \item The residual investment $I - I_m^2$ is provided by non-active (passive) investors at a stake $R_u^2$ satisfying their own breakeven condition.
\end{enumerate}

\begin{remark}[Timing]\label{rem:mon-timing}
The sequential-bid structure takes the adverse-selection problem to its extreme and maximizes the incumbent's rent.
Simultaneous-bid games would generate mixed-strategy equilibria with somewhat lower rents; the essential economic content is preserved.
\end{remark}

\subsection{The entrants' winner's curse}

An entrant who outbids the incumbent and wins must conclude: ``the incumbent did not find it profitable to match my bid.''
Given the incumbent's informational advantage, this can only mean that the incumbent's private information indicates a \emph{low} productivity.
In equilibrium, entrants correctly infer that they win only when the true date-2 productivity is $p_h$ (not $p_h + \tau$).

Entrants therefore bid conservatively, as if the productivity were always the low value $p_h$.
The entrant's breakeven bid is
\begin{equation}\label{eq:mon-entrant-bid}
  I_m^2 = p_h R_m^2 - c = p_h \cdot \frac{c}{\Delta p} - c = \frac{p_\ell c}{\Delta p}.
\end{equation}

\subsection{The incumbent's rent}

The incumbent matches the highest entrant bid and wins the monitoring role.
His \emph{true} expected profit, given the realized productivity, is
\begin{itemize}
  \item with probability $\alpha$: high productivity $p_h + \tau$, so expected revenue $(p_h + \tau) R_m^2$; cost $I_m^2 + c$; net $(p_h + \tau) R_m^2 - I_m^2 - c = (p_h + \tau) (c/\Delta p) - p_\ell c/\Delta p - c = \tau c/\Delta p$;
  \item with probability $1 - \alpha$: low productivity $p_h$, so net zero (breakeven).
\end{itemize}
The incumbent's unconditional expected rent is therefore
\begin{equation}\label{eq:mon-incumbent-rent}
  \mathcal{R}_m^2 = \alpha \tau \cdot \frac{c}{\Delta p}.
\end{equation}
The rent arises purely from the incumbent's ability to exploit his superior information: when the productivity is truly high, he captures the high value; when it is low, he matches the entrants' conservative bid without loss.

\subsection{The entrepreneur's date-2 payoff}

Given the incumbent's rent, the entrepreneur's date-2 payoff is reduced from the symmetric-information benchmark by $\mathcal{R}_m^2$:
\begin{equation}\label{eq:mon-ent-date2}
  \Pi_b^2 = \bigl[(p_h + \alpha \tau) R - I - c\bigr] - \alpha \tau \cdot \frac{c}{\Delta p}.
\end{equation}
Passive investors contribute $I_u^2 = I - I_m^2$ at stake $R_u^2$ satisfying their breakeven at the true expected productivity:
\begin{equation*}
  (p_h + \alpha \tau) R_u^2 = I_u^2.
\end{equation*}
(The passive investors' breakeven uses the \emph{true} expected productivity because their expected return conditions on the incumbent's matching behavior, which reveals information through the bidding outcome.)

A straightforward calculation confirms that the entrepreneur's stake in this equilibrium exceeds $b/\Delta p$, consistent with incentive compatibility.

%====================================================================
\section{Date-1 Competition and the Dissipation of Future Rents}
\label{sec:mon-date1}
%====================================================================

\subsection{Anticipating the date-2 rent}

The incumbent's rent \eqref{eq:mon-incumbent-rent} is anticipated by all agents at date~$1$.
Prospective date-1 monitors know that, if chosen, they will become the date-2 incumbent and capture expected rent $\mathcal{R}_m^2$.
This future rent affects the date-1 bidding.

\subsection{Generous date-1 offers}

At date~$1$, potential monitors are symmetrically informed: no one has yet observed the firm's operations.
Competition among date-1 monitors drives their expected \emph{two-period} net return to zero.

Fix the entrepreneur's chosen date-1 monitoring stake at $R_m^1 = c/\Delta p$.
A date-1 monitor who wins the monitoring role and contributes $I_m^1$ obtains:
\begin{itemize}
  \item direct date-1 return: $p_h R_m^1 - c - I_m^1 = p_h (c/\Delta p) - c - I_m^1 = p_\ell c/\Delta p - I_m^1$;
  \item discounted date-2 rent as the future incumbent: $\beta \mathcal{R}_m^2 = \beta \alpha \tau c/\Delta p$.
\end{itemize}
Zero-profit competition pins down the date-1 contribution:
\begin{equation}\label{eq:mon-date1-bid}
  I_m^1 = \frac{p_\ell c}{\Delta p} + \beta \mathcal{R}_m^2 = \frac{p_\ell c}{\Delta p} + \frac{\beta \alpha \tau c}{\Delta p}.
\end{equation}

The date-1 monitor contributes \emph{strictly more} than his direct date-1 breakeven ($p_\ell c/\Delta p$): he takes a short-term loss of $\beta \alpha \tau c/\Delta p$, exactly offset in expectation by the future rent he will capture at date~$2$.

\subsection{The full-competition outcome}

\begin{proposition}[Intertemporal rent dissipation]\label{prop:mon-dissipation}
Under perfect competition at date~$1$, the anticipated date-2 incumbency rent is fully dissipated by the date-1 ``generous offer'' $I_m^1 > p_\ell c/\Delta p$.
The entrepreneur's total two-period payoff coincides with the symmetric-information benchmark \eqref{eq:mon-sym-total}:
\begin{equation*}
  \Pi_b = \bigl[p_h R - (I + c)\bigr] + \beta \bigl[(p_h + \alpha \tau) R - (I + c)\bigr].
\end{equation*}
\end{proposition}

\begin{proof}
The entrepreneur's date-1 payoff is
\begin{equation*}
  \Pi_b^1 = p_h (R - R_m^1 - R_u^1) = p_h R - p_h R_m^1 - (I - I_m^1),
\end{equation*}
using the passive-investor breakeven $p_h R_u^1 = I - I_m^1$.
Substituting $R_m^1 = c/\Delta p$, $I_m^1 = p_\ell c/\Delta p + \beta \mathcal{R}_m^2$, and simplifying:
\begin{equation*}
  \Pi_b^1 = p_h R - c - I + \beta \mathcal{R}_m^2.
\end{equation*}
Adding the discounted date-2 payoff $\beta \Pi_b^2$ from \eqref{eq:mon-ent-date2}:
\begin{equation*}
  \Pi_b = [p_h R - c - I + \beta \mathcal{R}_m^2] + \beta \bigl[(p_h + \alpha \tau) R - I - c - \mathcal{R}_m^2\bigr] = [p_h R - I - c] + \beta [(p_h + \alpha \tau) R - I - c],
\end{equation*}
matching \eqref{eq:mon-sym-total}.
\end{proof}

\subsection{Economic interpretation}

Proposition~\ref{prop:mon-dissipation} delivers a striking economic insight.
The incumbent monitor's date-2 rent is a \emph{real} rent, in the sense that it distorts the date-2 allocation: the entrepreneur receives less than the first-best date-2 surplus by an amount $\mathcal{R}_m^2 = \alpha \tau c/\Delta p$.
But the rent is not a \emph{net} cost to the entrepreneur because it is perfectly competed away at date~$1$.

The key insight can be stated aphoristically:

\begin{quote}
\emph{The blockholding is initially acquired at a premium, and is later maintained at a discount.}
\end{quote}

The informational advantage of the incumbent can be viewed as a \emph{switching cost} that locks the firm in with the initial monitor.
The expected future market power conferred by this lock-in is competed away at the ex-ante stage through a loss-making initial offer.
This intertemporal structure of relationship lending---initial losses offset by later rents---is a signature feature of the banking industry and has been documented empirically in the literature on bank--firm relationships.

\subsection{A caveat: efficiency versus distributional neutrality}

Proposition~\ref{prop:mon-dissipation} shows that the entrepreneur's total welfare is unaffected by the incumbent's rent \emph{in expectation}.
This should not be interpreted as a statement that relationship lending is costless.
The result assumes:
\begin{itemize}
  \item perfectly competitive date-1 monitoring capital;
  \item risk-neutrality of all agents (monitors' short-term losses at date~$1$ are freely absorbed);
  \item no frictions preventing the intertemporal rent transfer (e.g., no limited liability of the date-1 monitor, no intervening default that would interrupt the relationship).
\end{itemize}
When these conditions fail---for example, when monitoring capital is scarce, when monitors are risk-averse, or when the relationship can be interrupted by the borrower's exit or the monitor's failure---the intertemporal dissipation breaks down, and the incumbent's rent becomes a net welfare cost.
The literature on relationship banking explores these richer environments at length.

%====================================================================
\section{Further Reading}
\label{sec:mon-further}
%====================================================================

The monitoring framework of Part~A is due to \citet{HolmstromTirole1997}, who introduce the distinction between monitored and unmonitored borrowers and derive the balance-sheet thresholds of Proposition~\ref{prop:mon-regimes}.
The empirical counterpart of the three-regime partition has been extensively studied: \citet{James1987}, \citet{LummerMcConnell1989}, and \citet{CantilloWright2000} document the pattern of firms' choices between bank lending and direct market access, with strong support for the theoretical predictions.
\citet{BergerMillerPetersenRajanStein2005} extend the analysis to the internal organization of banks, showing that smaller banks specialize in the soft-information-intensive monitoring that dominates intermediate-quality lending.

The optimal ownership-structure result of Section~\ref{sec:mon-optimal} (Proposition~\ref{prop:mon-optimal})---that the monitor should not hold all external shares---is developed in detail by \citet{PaganoRoell1998}, who analyze the ownership-structure implications of monitoring in the context of the decision to go public.
Their analysis connects the moral-hazard foundations of intermediated finance to observable patterns in ownership concentration.

The learning-by-lending analysis of Part~B builds on a large literature on relationship banking and switching costs.
The winner's-curse mechanism underlying the incumbent's rent is a classical theme in the adverse-selection literature, paralleling the mechanism that drives pricing in sealed-bid auctions with common values.
The intertemporal dissipation result (Proposition~\ref{prop:mon-dissipation}) provides a clean microfoundation for the ``front-loading'' of relationship costs that banking empiricists have documented: young firms pay high rates for their first bank loan, with rates declining as the bank relationship matures, until---in some applications---the relationship locks in excess rents later in the firm's life cycle.

\citet{Tirole2006} (Chapter~9) provides the standard textbook exposition and extends the analysis in several directions not covered here, including passive versus active monitoring, the cost structure of different types of monitors, and the interaction between monitoring and capital structure choice.

The dynamic dimension of the monitoring framework has been substantially extended by the modern literature on continuous-time contracting.
\citet{DeMarzoSannikov2006} provide the canonical continuous-time analysis of optimal security design and dynamic capital structure in a moral-hazard environment, implementing the optimal contract through a combination of credit line, long-term debt, and equity with agent-controlled payout policy.
\citet{Sannikov2008} develops the associated continuous-time principal--agent problem for a risk-averse agent, providing a sharp characterization of the optimal contract via stochastic calculus.
\citet{BiaisMariottiRochetVilleneuve2010} extend the framework to an environment with large risks and limited liability, showing how the optimal contract manages the trade-off between risk-sharing and incentive provision dynamically.
\citet{PiskorskiWesterfield2016} integrate endogenous monitoring with dynamic moral hazard, showing how monitoring technology interacts with the termination threat and compensation dynamics---a direct dynamic extension of the static monitoring framework developed in Part~A of this chapter.
Collectively, these papers provide the modern analytical machinery for studying the dynamics of monitoring, relationship lending, and their financial-contracting implications.

%====================================================================
\chapter{Control Rights and Corporate Governance}
\label{ch:control}
%====================================================================

Chapter~\ref{ch:monitoring} introduced active monitoring as a mechanism for mitigating moral hazard in settings where the entrepreneur's effort is privately chosen.
The monitor acquires information about the firm's internal operations and uses that information to restrict the set of available alternatives, reducing the scope for opportunism.
The mechanism was purely informational: the monitor's power came from knowing which actions the entrepreneur was about to take.

This chapter turns to a different but complementary mechanism: the explicit allocation of \emph{decision rights} between the entrepreneur and the investor.
Corporate contracts typically assign the right to make specific kinds of decisions---hiring and firing key personnel, selling off divisions, altering the firm's strategic direction, and so on---to either the manager-entrepreneur or to outside creditors and shareholders, often contingent on observable events such as financial performance.
These \emph{control rights} are a distinct dimension of the corporate contract from the monetary incentives (cash-flow rights) that have been the focus of Parts~II through~IV.

The central insight of the control-rights literature, due to \citet{AghionBolton1992} and built on the incomplete-contracts framework of \citet{HartMoore1990}, is that the transfer of control rights to investors can \emph{increase pledgeable income} by inducing investor-preferred decisions that generate more cash flow (even if they reduce the entrepreneur's private benefits).
The central trade-off is that the allocation of control that maximizes pledgeable income is generally not the allocation that maximizes social NPV: investors' preferences over decisions differ systematically from the entrepreneur's, and the choice of control-rights allocation reflects a balance between the two.

The model of this chapter---based on \citet{Tirole2006} Chapter~10---captures this trade-off in its starkest form.
An interim action is available that would raise the probability of success at a private cost to the entrepreneur that exceeds the public gain in expected output.
The action is \emph{first-best suboptimal} because its cost exceeds its gross benefit.
But if control is allocated to investors, they will nonetheless choose the action---they bear none of the cost but capture part of the gain---and the resulting pledgeable income is strictly greater than it would be under entrepreneur control.
For entrepreneurs with sufficiently weak balance sheets, relinquishing control to investors is the \emph{second-best} optimal choice: the first-best suboptimal action is preferred because it expands the range of projects that can be financed.

The analysis yields a sharp three-regime partition of entrepreneurs by balance-sheet strength, paralleling the partition of Chapter~\ref{ch:monitoring} but with control rights rather than monitoring as the mechanism.
Strong balance sheets retain control and finance at first-best; intermediate balance sheets relinquish control and finance at a distorted second-best; weak balance sheets cannot finance at all.
The result provides a microfoundation for the widespread empirical observation that corporate governance structures---board composition, voting rights, covenants---systematically correlate with firm financial conditions: financially weaker firms are subject to tighter investor oversight.

The chapter corresponds to Chapter~10 of \citet{Tirole2006}.
The broader theoretical framework is developed in \citet{Hart1995} and surveyed in \citet{Hart2001}.

%====================================================================
\section{Setup: Moral Hazard with an Interim Action}
\label{sec:control-setup}
%====================================================================

\subsection{The baseline framework}

We adopt the fixed-investment model of Chapter~\ref{ch:net-worth}.
A risk-neutral entrepreneur with initial wealth $A < I$ needs to raise $I - A$ from outside investors.
The project yields a verifiable income $R > 0$ in success and zero in failure.

The moral-hazard stage is standard: the entrepreneur privately chooses to \emph{behave} or \emph{misbehave}.
Under behaving, the probability of success is $p = p_h$ and the entrepreneur receives no private benefit; under misbehaving, the probability is $p = p_\ell < p_h$ and the entrepreneur receives private benefit $B > 0$.
We write $\Delta p \coloneqq p_h - p_\ell$.

\subsection{The interim action}

The novelty of this chapter is the introduction of an \emph{interim action}, taken at a date after the initial financing but before the moral-hazard stage.
If the action is taken, it raises the success probability \emph{uniformly} by $\tau > 0$: the probabilities under behaving and misbehaving become $p_h + \tau$ and $p_\ell + \tau$, respectively.
Taking the action imposes a private, non-monetary cost $\gamma > 0$ on the entrepreneur.

The interpretation of the action varies across applications:
\begin{itemize}
  \item switching from a preferred strategy to a more routine but more profitable one;
  \item severing a long-standing relationship with a collaborator or supplier;
  \item firing workers in a restructuring;
  \item divesting a division that the entrepreneur personally values.
\end{itemize}
In each case, the action improves expected financial performance (raising the success probability) but imposes a subjective cost on the entrepreneur---the loss of a preferred strategy, relationship, or activity.

\subsection{Incentive compatibility}

A key structural observation is that the interim action does \emph{not} alter the incentive-compatibility constraint for effort.
Whether or not the action is taken, the wedge between the two success probabilities is the same:
\begin{equation*}
  (p_h + \tau) - (p_\ell + \tau) = p_h - p_\ell = \Delta p.
\end{equation*}
Therefore, the entrepreneur's incentive-compatibility constraint remains
\begin{equation}\label{eq:control-IC}
  (\Delta p) R_b \geq B,
\end{equation}
where $R_b$ is the entrepreneur's stake in success (zero in failure).
This separability means that the interim-action decision and the effort-incentive problem can be analyzed independently; in particular, the timing of the action relative to the moral-hazard stage is irrelevant.

%====================================================================
\section{The Allocation of Control Rights}
\label{sec:control-allocation}
%====================================================================

The central contractual question of the chapter is: who decides whether the interim action is taken?
We assume that the action itself cannot be contracted upon at the initial financing stage---it is, in the language of \citet{HartMoore1990}, \emph{non-contractible}.
However, the two parties \emph{can} contract on who holds the right to decide.

\subsection{The first-best suboptimality assumption}

Throughout the chapter, we maintain the following assumption about the interim action:

\begin{assumption}[First-best suboptimality]\label{ass:control-FBsub}
The interim action reduces aggregate welfare:
\begin{equation}\label{eq:control-FBsub}
  \tau R < \gamma.
\end{equation}
\end{assumption}

Assumption~\ref{ass:control-FBsub} says that the private cost $\gamma$ exceeds the gross expected benefit $\tau R$.
From a first-best perspective, the action should \emph{not} be taken: it destroys value.
The economic interest of the analysis lies in the fact that, despite its first-best suboptimality, the action may nonetheless be chosen in equilibrium under investor control---and this choice may be second-best optimal for the entrepreneur's financing problem.

\subsection{Investor control}

Suppose the control right over the interim action is allocated to the investors.
Investors care only about their share of the verifiable output---they receive $R - R_b$ in success and zero in failure.
Their decision rule is to take the action if and only if it strictly raises their expected payoff:
\begin{equation*}
  \tau (R - R_b) > 0,
\end{equation*}
which holds for any feasible $R_b < R$.
Investors \emph{always} take the action.
They bear none of the private cost $\gamma$ and capture part of the gross benefit $\tau R$; the action is strictly in their interest.

The pledgeable income under investor control is computed with the enhanced success probabilities.
Using the IC-binding stake $R_b = B/\Delta p$:
\begin{equation}\label{eq:control-P-investor}
  \mathcal{P}^{\text{c-inv}} = (p_h + \tau)\left(R - \frac{B}{\Delta p}\right).
\end{equation}
The enhanced probability $p_h + \tau$ multiplies the per-success pledgeable contribution, raising the total pledgeable income by $\tau(R - B/\Delta p)$ relative to the no-action case.

Under competitive lending, the lender's breakeven binds, and the entrepreneur captures the project's NPV \emph{net} of the interim-action cost:
\begin{equation}\label{eq:control-Pi-investor}
  \Pi_b^{\text{c-inv}} = \NPV^{\text{c-inv}} = (p_h + \tau) R - I - \gamma.
\end{equation}

\subsection{Entrepreneur control}

Suppose instead the control right is allocated to the entrepreneur.
The entrepreneur's decision rule is to take the action if and only if
\begin{equation*}
  \tau R_b \geq \gamma.
\end{equation*}
Since $R_b \leq R$ (the entrepreneur's stake cannot exceed the total output) and Assumption~\ref{ass:control-FBsub} gives $\tau R < \gamma$, we have
\begin{equation*}
  \tau R_b \leq \tau R < \gamma,
\end{equation*}
so the entrepreneur \emph{never} takes the action.

The economic force is transparent: the entrepreneur bears the entire private cost $\gamma$ but captures only a fraction (determined by $R_b/R < 1$) of the gross benefit $\tau R$.
Since even the full benefit is smaller than the cost (Assumption~\ref{ass:control-FBsub}), it is a fortiori unprofitable for the entrepreneur to act when he internalizes only a fraction.

The pledgeable income under entrepreneur control is therefore the standard no-action expression:
\begin{equation}\label{eq:control-P-entrepreneur}
  \mathcal{P}^{\text{c-ent}} = p_h\left(R - \frac{B}{\Delta p}\right) = \rho_0.
\end{equation}
Under competitive lending, the entrepreneur captures the first-best NPV:
\begin{equation}\label{eq:control-Pi-entrepreneur}
  \Pi_b^{\text{c-ent}} = \NPV^{\text{c-ent}} = p_h R - I.
\end{equation}

\subsection{The trade-off}

Comparing \eqref{eq:control-Pi-investor} and \eqref{eq:control-Pi-entrepreneur}:
\begin{equation}\label{eq:control-NPV-diff}
  \NPV^{\text{c-ent}} - \NPV^{\text{c-inv}} = \bigl\{p_h R - I\bigr\} - \bigl\{(p_h + \tau) R - I - \gamma\bigr\} = \gamma - \tau R > 0,
\end{equation}
where the final inequality uses Assumption~\ref{ass:control-FBsub}.
\emph{Entrepreneur control delivers strictly higher NPV.}

Comparing pledgeable incomes \eqref{eq:control-P-investor} and \eqref{eq:control-P-entrepreneur}:
\begin{equation}\label{eq:control-P-diff}
  \mathcal{P}^{\text{c-inv}} - \mathcal{P}^{\text{c-ent}} = \tau\left(R - \frac{B}{\Delta p}\right) > 0
\end{equation}
(assuming $R > B/\Delta p$, which must hold for any financing to be feasible).
\emph{Investor control delivers strictly higher pledgeable income.}

This is the central trade-off of the chapter.
The two criteria---social surplus (NPV) and lender breakeven (pledgeable income)---rank the two allocations in opposite directions.
Entrepreneur control is first-best but has lower pledgeable capacity; investor control has higher pledgeable capacity but is first-best suboptimal.

%====================================================================
\section{The Three Balance-Sheet Regimes}
\label{sec:control-regimes}
%====================================================================

The resolution of the trade-off depends on the entrepreneur's balance-sheet strength.
An entrepreneur with sufficient net worth can secure financing without relinquishing control; an entrepreneur with weaker net worth may need to relinquish control to satisfy the lender's breakeven; an entrepreneur with very weak net worth cannot finance even with investor control.

\subsection{Two thresholds}

Define the \emph{high threshold}
\begin{equation}\label{eq:control-Aupper}
  \overline{A} \coloneqq I - p_h\left(R - \frac{B}{\Delta p}\right) = I - \mathcal{P}^{\text{c-ent}}.
\end{equation}
This is the minimum wealth under which entrepreneur-control financing is feasible: the pledgeable income $\mathcal{P}^{\text{c-ent}}$ must cover the outside financing need $I - A$.

Define the \emph{low threshold}
\begin{equation}\label{eq:control-Alower}
  \underline{A} \coloneqq I - (p_h + \tau)\left(R - \frac{B}{\Delta p}\right) = I - \mathcal{P}^{\text{c-inv}}.
\end{equation}
This is the minimum wealth under which investor-control financing is feasible.
By \eqref{eq:control-P-diff}, $\mathcal{P}^{\text{c-inv}} > \mathcal{P}^{\text{c-ent}}$, so
\begin{equation*}
  \underline{A} < \overline{A}.
\end{equation*}

\subsection{Three regimes}

\begin{proposition}[Balance-sheet regimes]\label{prop:control-regimes}
Under Assumption~\ref{ass:control-FBsub}, the financing outcome for an entrepreneur with wealth $A$ is:
\begin{enumerate}
  \item[\textup{(i)}] \emph{Strong balance sheet} ($A \geq \overline{A}$): financing with entrepreneur control.
  The interim action is not taken, and the entrepreneur captures the first-best NPV, $\Pi_b = p_h R - I$.
  \item[\textup{(ii)}] \emph{Intermediate balance sheet} ($\underline{A} \leq A < \overline{A}$): financing with transfer of control to investors.
  The interim action is taken (first-best suboptimal), and the entrepreneur captures the distorted NPV, $\Pi_b = (p_h + \tau) R - I - \gamma$.
  \item[\textup{(iii)}] \emph{Weak balance sheet} ($A < \underline{A}$): no financing.
  The entrepreneur's payoff is zero.
\end{enumerate}
\end{proposition}

\begin{proof}
\emph{Case (i).}
If $A \geq \overline{A}$, the pledgeable income $\mathcal{P}^{\text{c-ent}}$ covers $I - A$, so entrepreneur-control financing is feasible.
The entrepreneur strictly prefers entrepreneur control to investor control by \eqref{eq:control-NPV-diff}.
Hence the equilibrium is entrepreneur control.

\emph{Case (ii).}
If $\underline{A} \leq A < \overline{A}$, entrepreneur-control financing is infeasible (the pledgeable income does not cover $I - A$).
Investor-control financing is feasible (by definition of $\underline{A}$) and delivers strictly positive NPV:
\begin{equation*}
  (p_h + \tau) R - I - \gamma > 0,
\end{equation*}
where positivity follows from the feasibility threshold (if $(p_h + \tau)R - I - \gamma \leq 0$, even at $A = \underline{A}$ the entrepreneur's surplus would be non-positive, making the project uninteresting).
The entrepreneur prefers investor-control financing to no financing, so the equilibrium is investor control.

\emph{Case (iii).}
If $A < \underline{A}$, even investor-control financing is infeasible, so no contract is signed.
\end{proof}

\subsection{Economic interpretation}

Proposition~\ref{prop:control-regimes} delivers the central structural insight of the chapter.
The second-best outcome for an entrepreneur with intermediate wealth is to \emph{voluntarily} give up control---even though this triggers a first-best suboptimal decision---in exchange for access to financing that would otherwise be foreclosed.
The first-best suboptimal decision is \emph{second-best optimal} once the credit-market imperfection is taken into account.

The result provides a natural microfoundation for several well-documented empirical patterns in corporate governance:

\begin{itemize}
  \item \emph{Firms with weaker balance sheets are subject to tighter covenants.}
  Loan agreements with highly leveraged or low-net-worth firms typically include extensive \emph{negative covenants}---restrictions on the firm's actions that effectively transfer decision rights to the lender in exchange for improved access to credit.
  Firms with stronger balance sheets borrow on less restrictive terms.
  \item \emph{Venture-capital contracts feature extensive staged-control transfers.}
  Early-stage ventures with little collateral and high informational asymmetry typically cede substantial control rights to their venture investors, including board seats, veto rights, and vesting schedules.
  As the firm matures and builds financial strength, control rights typically revert toward the founders.
  \item \emph{Financial distress triggers control transfers.}
  When a firm approaches bankruptcy, creditors gain enhanced control rights---through debt covenants, chapter-11 proceedings, or informal workouts.
  The transfer of control is precisely the mechanism that sustains financing in the distress phase.
\end{itemize}

In each case, the observed structure is consistent with the logic of Proposition~\ref{prop:control-regimes}: \emph{weaker balance sheets require more investor-favorable decision rights}, and the allocation of control systematically correlates with financial condition.

%====================================================================
\section{Further Reading}
\label{sec:control-further}
%====================================================================

The control-rights approach to financial contracting was introduced by \citet{AghionBolton1992}, who developed the framework in a richer continuous-decision environment and demonstrated that state-contingent transfers of control can implement outcomes close to the first-best when information about the firm's state is publicly available.
The simpler binary-action model of this chapter, following the exposition of \citet{Tirole2006} (Chapter~10), isolates the central economic logic in its starkest form.

The broader theoretical foundations of the control-rights approach rest on the \emph{incomplete-contracts} literature, developed primarily by \citet{HartMoore1990} in their analysis of property rights and the nature of the firm.
The key insight of that literature is that when some decisions cannot be contracted upon ex ante---because they are too complex, too state-contingent, or too costly to verify---the allocation of residual decision rights becomes a first-order economic variable.
\citet{Hart1995} provides a comprehensive textbook treatment of incomplete contracts in the corporate context; \citet{Hart2001} offers a more compact survey of the literature's development through the early 2000s.

The connection between control rights and debt covenants has been explored extensively in the theoretical and empirical corporate-finance literature.
Debt covenants can be understood as \emph{state-contingent} transfers of control from entrepreneurs to creditors: a technical default triggers the covenant, and the creditor gains rights to intervene in the firm's operations.
The empirical literature has documented rich cross-sectional variation in covenant structures and has connected this variation to firm-level measures of informational opacity, balance-sheet strength, and industry characteristics.

A related line of research examines the choice between debt and equity as alternative forms of control allocation.
Debt typically transfers control only in distress states (when covenants or payment defaults are triggered), while equity provides ongoing monitoring and oversight through board membership and voting rights.
The optimal mix of the two depends on the distribution of the firm's states and the specific decisions that must be allocated.
This broader question connects the control-rights theory of this chapter to the classical capital-structure literature of Parts~II through~IV.

%====================================================================
\backmatter
\bibliographystyle{kluwer}
\bibliography{references}
%====================================================================

\end{document}
